\documentclass[a4paper,11pt]{article}
\pdfoutput=1
\usepackage{jheppub}   % JHEP class file (SISSA Medialab, v1.1227), copied into this directory
\usepackage[T1]{fontenc}
\usepackage{amsmath,amssymb,amsfonts,mathtools}
\usepackage{braket}
\usepackage{booktabs}
\usepackage{array}
\allowdisplaybreaks[2]
\usepackage{xcolor}
\usepackage{tikz}
\usetikzlibrary{calc}
\usetikzlibrary{decorations.pathmorphing}
\usetikzlibrary{decorations.markings}
\usetikzlibrary{decorations.text}
\usetikzlibrary{decorations.pathreplacing}
\usepackage{tikzit}
\input{stiles.tikzstyles}

% ---------------------------------------------------------------------
% Notation (identical to arXiv:2607.18373 and to the calculation files)
% ---------------------------------------------------------------------
\def\mf#1{\mathbb{#1}}          % coefficients of Omega: \mf{N},\mf{A},\mf{B},\mf{C}
\def\v#1{\mathbf{#1}}
\newcommand{\Kv}{\mathcal{K}}    % Weizsaecker-Williams field  K^i(x)= x^i/x^2
\newcommand{\MS}{\overline{\rm MS}}
\newcommand{\Tr}{{\rm Tr}}
\newcommand{\NLOfirst}{\mathrm{NLO}}
\newcommand{\Lam}{\Lambda}
% upper cut-off of the soft modes is denoted \vee, as in 2607.18373 and 1610.03453 (\vee = \Lambda e^{\delta Y})

\title{Single inclusive gluon production at mid rapidity at next-to-leading order in the Color Glass Condensate}

\author{AUTHORS}
\affiliation{AFFILIATIONS}
% \emailAdd{}

\abstract{We compute the single inclusive gluon spectrum in the collision of a dilute projectile with a dense target at next-to-leading order in the Color Glass Condensate framework, for a gluon produced at mid rapidity, in pure Yang--Mills theory. The calculation is performed in light-cone perturbation theory: the projectile is described by its soft gluon light-cone wave function, generated from the valence colour charge density $\rho$ by a unitary dressing operator $\Omega$, and the target by eikonal adjoint Wilson lines. The coefficient of $\Omega$ that creates one gluon at order $g^3$ is computed here; the others were obtained previously. The valence charges are treated as non-commuting operators throughout. We organise the $O(g^4)$ cross section into interference terms with one gluon in the final state and into contributions with two gluons in the final state, one of which is integrated over. All terms enhanced by the logarithm of the longitudinal cut-off ratio have the structure required by high-energy factorisation: the logarithm splits at the longitudinal momentum $k^+$ of the produced gluon, and the coefficient of $\log(\vee/k^+)$ is exactly the leading-order BFKL evolution of the projectile colour charge correlator, while the coefficient of $\log(k^+/\Lambda)$ is exactly the JIMWLK evolution of the target Wilson-line operator of the leading-order cross section. The terms quartic and (in the symmetrised sense) cubic in $\rho$ that accompany these logarithms cancel. The remaining contributions are a collinear term, which is the final-state DGLAP fragmentation of the produced gluon with the $\overline{\rm MS}$ splitting function, a running-coupling term, which sets the scale of the leading-order coupling, and finite terms. We account for all poles of dimensional regularisation. We give the finite terms row by row, after the cancellations that are exact, with every product of three charges written as its symmetrised product plus two-charge contact terms, for an arbitrary projectile and, averaged over the projectile, in colour-singlet form; the colour structures are written through differences of Wilson lines and the longitudinal integrals with plus prescriptions in a momentum fraction. The large-distance tails of individual terms cancel in pairs of virtual and real terms, without any condition on the target at large distance. The terms cubic in $\rho$ do not vanish after the average over the projectile: their symmetrised part, odd in the transverse momentum, is proportional to the symmetrised ($d^{abc}$) three-point function of the projectile, for which no model is assumed, and to C-odd target structures, while the commutators of the charges leave contact terms of the same order as the quadratic terms. Energy denominators that vanish inside the integration region are defined as principal values. In the dilute limit the real-emission part reproduces exactly the two-gluon production cross section in quasi-multi-Regge kinematics. The points that remain open are listed.
}

\begin{document}
\maketitle
\flushbottom

% ======================================================================
\section{Introduction}\label{sec:intro}
% ======================================================================

At high energy the wave function of a hadron or nucleus contains a large number of soft gluons, and in a sufficiently dense system their density saturates~\cite{Gribov:1983ivg,Mueller:1985wy}. A description of this regime from first principles was proposed by McLerran and Venugopalan~\cite{McLerran:1993ni,McLerran:1993ka}: the fast (valence) partons of a large nucleus are treated as static sources of colour charge, with a density $\rho^a(\v x)$ in the transverse plane; the soft gluons are the classical Yang--Mills field radiated by these sources; and observables are obtained by averaging over the sources with a Gaussian weight. This picture has developed into the effective theory known as the Color Glass Condensate (CGC)~\cite{Iancu:2003xm,Weigert:2005us,Gelis:2010nm,Kovchegov:2012mbw}. The dependence of observables on the rapidity interval over which the soft modes are integrated out is governed by the BFKL equation~\cite{Kuraev:1977fs,Balitsky:1978ic} in the dilute regime and by the non-linear Balitsky--Kovchegov (BK) and JIMWLK equations~\cite{Balitsky:1995ub,Kovchegov:1999yj,Kovchegov:1999ua,Jalilian-Marian:1997qno,Jalilian-Marian:1997jx,Jalilian-Marian:1997ubg,Kovner:2000pt,Weigert:2000gi,Iancu:2000hn,Iancu:2001ad,Ferreiro:2001qy} in the dense regime, which replace the Gaussian weight by an evolving distribution of the sources.

The simplest observable that is sensitive both to the projectile and to the target wave functions is the single inclusive gluon spectrum in the collision of a dilute projectile with a dense target. At leading order (LO), for a gluon produced at mid rapidity, i.e.\ far in rapidity from both the projectile and the target, the spectrum is given by the formula of Kovchegov and Tuchin~\cite{Kovchegov:2001sc,Kovchegov:1998bi}: the gluon is emitted from the valence charges of the projectile in the Weizs\"acker--Williams field, and the emission amplitude, together with the colour rotation of the emitting charge, scatters eikonally on the target. The cross section is a product of the projectile colour charge correlator and of a target operator built from two adjoint Wilson lines; in the limit of a dilute target it reduces to the $k_T$-factorised form with the Lipatov vertex~\cite{Gribov:1983ivg,Kharzeev:2003wz,Blaizot:2004wu}. This is a result of the classical theory: it is quadratic in the charges of the projectile, and after the average over a colour-neutral projectile it does not depend on whether these charges are numbers or operators. Its leading logarithmic (LL) rapidity evolution was obtained by Kovner and Lublinsky~\cite{Kovner:2006wr} (see also~\cite{Kovner:2006ge}) and, within the Schwinger--Keldysh formulation, by Gelis, Lappi and Venugopalan~\cite{Gelis:2008rw}: the projectile is to be evolved from its own rapidity down to the rapidity of the produced gluon, and the target from the rapidity of the produced gluon to its own rapidity.

Beyond leading order the situation is different, and the picture of classical fields of static, commuting colour charges with a Gaussian weight, which was designed for the leading-order problem and solves it, is not by itself a complete description of gluon production. Three elements have to be added to it. First, the colour charges of the valence partons are operators, $[\rho^a(\v x),\rho^b(\v y)]=if^{abc}\rho^c(\v x)\delta^{(2)}(\v x-\v y)$~\cite{Kovner:2005pe,Kovner:2005nq,Kovner:2005jc,Kovner:2005en}. At the order $g^4$ of the next-to-leading order (NLO) correction the amplitudes contain products of three and four charges, and a commutator of two charges is of the same order as the two-charge terms that carry the evolution and the NLO coefficient. The order of the charges therefore matters: the commutators leave contact terms, and the totally symmetrised product of three charges brings in the three-point function of the projectile with the symmetric structure constants $d^{abc}$, which is odd under charge conjugation and absent for a Gaussian weight. Second, the logarithms of an NLO calculation are of several kinds: those of the rapidity evolution of the projectile and of the target, the collinear logarithms of the fragmentation of the produced gluon, which are governed by the DGLAP equations~\cite{Gribov:1972ri,Altarelli:1977zs,Dokshitzer:1977sg}, and the ultraviolet logarithms that make the coupling run. In the leading-order framework they are supplied by separate considerations; in a complete description all of them, with their scales and their overlaps, have to emerge from one calculation. Third, the terms of relative order $\alpha_s$ that are not enhanced by any logarithm, which are needed for any quantitative use of the leading-order formula, have not been available for this observable.

Many of the ingredients have become available over the last years. The evolution equations themselves are known at NLO: the NLO BK equation~\cite{Balitsky:2007feb} and the NLO JIMWLK Hamiltonian~\cite{Balitsky:2013fea,Kovner:2013ona,Kovner:2014lca,Grabovsky:2013mba}, together with the running-coupling corrections to the LO kernels~\cite{Balitsky:2006wa,Kovchegov:2006vj,Kovchegov:2006wf}. In particular, Lublinsky and Mulian~\cite{Lublinsky:2016meo} computed the soft gluon light-cone wave function of a dilute projectile to order $g^3$ and derived the NLO JIMWLK Hamiltonian from it. For particle production, the NLO corrections to single inclusive hadron production at forward rapidity were computed in the hybrid formalism~\cite{Dumitru:2005gt,Altinoluk:2011qy} by Chirilli, Xiao and Yuan~\cite{Chirilli:2011km,Chirilli:2012jd}, and the structure of the rapidity subtraction and of the kinematic constraints in this framework was analysed in~\cite{Stasto:2013cha,Altinoluk:2014eka,Ducloue:2016shw}. Running-coupling corrections to inclusive gluon production at mid rapidity were obtained by Horowitz and Kovchegov~\cite{Horowitz:2010yg}, and the relation between $\beta_0$-dependent terms of NLO calculations and the running of the coupling was clarified in~\cite{Kovner:2023vsy}. The general structure of the NLO correction to inclusive spectra in the CGC was discussed by Gelis and Laidet~\cite{Gelis:2012ri}. Corrections of higher order in the projectile colour charge (``first saturation corrections'' in the projectile) were computed in the classical approximation in~\cite{Chirilli:2015tea,Li:2021zmf,Li:2021ntm}, and their contribution to single inclusive production, odd in $\v k$, was analysed in~\cite{McLerran:2016snu,Kovchegov:2018jun}. A complete NLO calculation of the single inclusive gluon spectrum at mid rapidity has, to our knowledge, not been presented.

The purpose of this paper is to supply the missing pieces for this observable. In a previous paper~\cite{Radhakrishnan:2026soft} the soft gluon wave function of a dilute projectile and the associated unitary evolution (dressing) operator $\Omega$ were constructed up to order $g^2$, in pure Yang--Mills theory, within light-cone perturbation theory (LCPT)~\cite{Lepage:1980fj,Brodsky:1997de}. The operator $\Omega$ was written as a normal-ordered series in soft gluon creation and annihilation operators with coefficients that are functionals of the non-commuting valence charge density $\rho$, and the coefficients were fixed by matching to LCPT and by unitarity. The NLO cross section requires in addition the coefficient of $\Omega$ that creates one gluon at order $g^3$; its ingredients, the $O(g^3)$ one-gluon components of the light-cone wave function, were computed by Lublinsky and Mulian~\cite{Lublinsky:2016meo}, and we recompute them here with the longitudinal cut-offs and the operator ordering required by the cross section. The observable is the expectation value of the number operator of dressed gluons in the scattered, dressed projectile state~\cite{Kovner:2006wr}, so that the full calculation reduces to an algebraic evaluation of products of the coefficients of $\Omega$ and of Wilson lines, in which the order of the charges is kept, followed by longitudinal and transverse integrations. No average over the projectile is taken until the end, and no model for it is assumed.

Our main results are the following.
\begin{itemize}
\item The $O(g^4)$ cross section naturally decomposes into three groups of terms: interference of the LO amplitude with the $O(g^3)$ one-gluon amplitude (Group~I); the squares of the $O(g^2)$ two-gluon amplitudes, with one of the gluons integrated over (Group~II); and the interference terms between different two-gluon amplitudes (Group~III). We list all products of coefficients that contribute, with their longitudinal integration ranges and their order in $\rho$ (Section~\ref{sec:observable}).
\item All terms enhanced by $\log(\vee/\Lambda)$, where $\Lambda<p^+<\vee$ is the longitudinal momentum range of the soft modes, are the LL evolution of the LO result. The logarithm splits at the longitudinal momentum $k^+$ of the produced gluon into $\log(\vee/k^+)+\log(k^+/\Lambda)$; the coefficient of $\log(\vee/k^+)$ is exactly the BFKL evolution of the projectile charge correlator inserted into the LO formula, and the coefficient of $\log(k^+/\Lambda)$ is exactly the JIMWLK Hamiltonian acting on the target operator of the LO formula. We give both terms for an arbitrary (ordered) projectile correlator and in colour-singlet form. The contributions quartic in $\rho$, all of which carry $\log(\vee/\Lambda)$, cancel identically for commuting charges, and the totally symmetrised cubic terms cancel as well; with non-commuting charges, their reordering produces the two-$\rho$ terms that complete the evolution (Section~\ref{sec:rapidity}).
\item The remaining terms are a collinear term, associated with the fragmentation of the produced gluon and governed by the gluon--gluon splitting function, a running-coupling term, and finite terms. We give an account of all poles $1/\epsilon$: the ultraviolet poles of the virtual corrections to the produced gluon cancel among themselves, the coupling is renormalised by the normalisation of the dressed produced gluon, and the collinear poles combine into the full $\MS$ splitting function $\mathcal P_{gg}$ and are absorbed into the fragmentation function. The running-coupling term is the $\frac{\beta_0}{2}\delta(1-\xi)$ part of the collinear term; it sets the scale of the LO coupling. The scheme is conventional dimensional regularisation with the $\MS$ coupling and the standard $\MS$ number-density fragmentation function, with which standard one-hadron-inclusive coefficient functions are reproduced (Section~\ref{sec:dglap_rc}).
\item We give the finite terms Group by Group and Row by Row, after the cancellations that are exact, in two forms: for an arbitrary projectile, with every product of three charges written as its totally symmetrised product plus two-charge contact terms, and averaged over the projectile, in colour-singlet form. All colour structures are written, as at leading order, through differences of Wilson lines, and all longitudinal integrals with plus prescriptions in a momentum fraction. After the average the terms cubic in $\rho$ do not vanish: their symmetrised part, odd in $\v k$, is proportional to the symmetrised ($d^{abc}$) three-point function of the projectile, which is kept general, and to C-odd structures of the target, while the commutators of the ordered charges leave contact terms, even in $\v k$, which are present for any projectile. The large-distance tails of individual terms cancel in pairs of a virtual and a real term, without any condition on the target at large distance; we state which terms have to be integrated together. A compact summary of the result is given in Section~\ref{sec:result_summary} (Section~\ref{sec:finite}, Appendices~\ref{app:finite} and~\ref{app:tails}).
\item In the dilute--dilute limit the real emission reproduces exactly the two-gluon production cross section in quasi-multi-Regge kinematics~\cite{Fadin:1998py,Bartels:2006hg}; for a point-like projectile the coefficients of the rapidity logarithms in the virtual correction are the gluon Regge trajectories. The comparison of the finite part of the virtual correction with the one-loop central-emission vertex, which requires the real emission with matched regulators, remains open, as do the other points listed in Section~\ref{sec:openpoints} (Section~\ref{sec:checks}).
\end{itemize}
The result can be written schematically as
\begin{align}\label{eq:schematic}
\frac{d^3N}{d^3k} &= \frac{d^3N_{\rm LO}}{d^3k}
+ \log\frac{\vee}{k^+}\,\frac{d^3N_{\rm BFKL}}{d^3k}
+ \log\frac{k^+}{\Lambda}\,\frac{d^3N_{\rm JIMWLK}}{d^3k}\notag\\
&\quad+ \frac{d^3N_{\rm DGLAP}}{d^3k}
+ \frac{d^3N_{\rm RC}}{d^3k}
+ \frac{d^3N_{\rm finite}}{d^3k}\,,
\end{align}
where every term on the right-hand side except the first is of relative order $\alpha_s$. In this formula the evolution of the projectile, the evolution of the target, the collinear evolution of the final state, the running of the coupling and the finite terms, with the charges of the projectile treated as operators throughout, come out of a single calculation. In this sense the present work completes the NLO description of the single inclusive gluon spectrum at mid rapidity within the eikonal, dilute-projectile approximation and in pure Yang--Mills theory, up to the open points of Section~\ref{sec:openpoints}; it does not address the dense--dense problem, corrections beyond the eikonal approximation, or quarks.

The paper is organised as follows. In Section~\ref{sec:setup} we summarise the setup: light-cone conventions, the Hamiltonian of the soft modes in the eikonal approximation, LCPT with operator-valued matrix elements, and the expansion of $\Omega$ with its coefficients in coordinate space. In Section~\ref{sec:observable} we define the observable, rederive the LO cross section and organise the NLO contributions. Section~\ref{sec:rapidity} is devoted to the rapidity logarithms and their identification with the BFKL and JIMWLK evolution. The collinear and running-coupling terms are discussed in Section~\ref{sec:dglap_rc}, the finite terms in Section~\ref{sec:finite}, and limits and comparisons with the literature in Section~\ref{sec:checks}. We conclude in Section~\ref{sec:conclusions}, where the open points are listed. The appendices collect the Hamiltonian and its matrix elements (Appendix~\ref{app:hamiltonian}), the light-cone wave functions from which the coefficients of $\Omega$ are obtained (Appendix~\ref{app:wavefunctions}), the unitarity relations among the coefficients (Appendix~\ref{app:unitarity}), the Fourier transforms (Appendix~\ref{app:fourier}), the prescriptions and distributions used in the longitudinal and transverse integrals (Appendix~\ref{app:distributions}), the finite terms for an arbitrary projectile (Appendix~\ref{app:finite}) and the analysis of their large-distance behaviour (Appendix~\ref{app:tails}).

% ======================================================================
\section{Setup: soft gluon wave function and the evolution operator}\label{sec:setup}
% ======================================================================

In this section we summarise the framework of~\cite{Radhakrishnan:2026soft}, on which the present calculation is built, and collect the coefficients of the evolution operator $\Omega$ that enter the single inclusive cross section. The coefficients of order $g$ and $g^2$ were derived in~\cite{Radhakrishnan:2026soft}; the coefficient of order $g^3$ that creates one gluon, $\mf A^{(3)}_2$, is obtained here. The derivations of the light-cone wave functions from which the coefficients are obtained are summarised in Appendix~\ref{app:wavefunctions}.

\subsection{Light-cone conventions}\label{sec:conventions}

We use light-cone coordinates $x^\pm=(x^0\pm x^3)/\sqrt2$ and denote transverse vectors by $\v x=(x^1,x^2)$; where no confusion is possible the transverse vectors are written without boldface, and $\int_{x}\equiv\int d^2x$. The projectile moves in the $+$ direction, $x^+$ plays the role of time, and $p^+>0$ is the longitudinal momentum of a physical particle. The light-cone energy of an on-shell gluon is $p^-=\v p^2/2p^+$, and we write
\begin{align}
&E_g(p)=\frac{\v p^2}{2p^+}\,,\qquad E_{gg}(p,q)=\frac{\v p^2}{2p^+}+\frac{\v q^2}{2q^+}\,,\notag\\ &d^3p\equiv dp^+d^2p\,,\qquad \delta^{(3)}(p-q)\equiv\delta(p^+-q^+)\,\delta^{(2)}(\v p-\v q)\,.
\end{align}
We work in the light-cone gauge $A^+=0$, in which the transverse components $A_i^a$, $i=1,2$, are the dynamical fields. The gluon creation and annihilation operators satisfy
\begin{align}\label{eq:aadagger}
\big[a^a_i(k),a^{\dagger b}_j(p)\big]=(2\pi)^3\,\delta^{ab}\,\delta_{ij}\,\delta^{(3)}(k-p)\,.
\end{align}
For the Fourier transforms in the transverse plane we use the conventions
\begin{align}\label{eq:FTconventions}
&a(k^+,\v k)=\int d^2z\,e^{-i\v k\cdot\v z}\,a(k^+,\v z)\,,\qquad
a^\dagger(k^+,\v k)=\int d^2z\,e^{i\v k\cdot\v z}\,a^\dagger(k^+,\v z)\,,\notag\\
&f(\v k)=\int d^2x\,e^{i\v k\cdot\v x}\,f(\v x)\,,
\end{align}
so that $\big[a^a_i(k^+,\v z),a^{\dagger b}_j(p^+,\v z')\big]=2\pi\,\delta^{ab}\delta_{ij}\,\delta(k^+-p^+)\,\delta^{(2)}(\v z-\v z')$. One- and two-gluon states are normalised as
\begin{align}\label{eq:states}
\ket{g^a_i(k^+,\v z)}=\frac{a^{\dagger a}_i(k^+,\v z)}{(2\pi)^{1/2}}\ket{0}\,,\qquad
\ket{g^a_i(k^+,\v z)\,g^b_j(p^+,\v z')}=\frac{a^{\dagger a}_i(k^+,\v z)\,a^{\dagger b}_j(p^+,\v z')}{2\pi}\ket{0}\,,
\end{align}
with $\braket{g^b_j(p^+,\v z')|g^a_i(k^+,\v z)}=\delta^{ab}\delta_{ij}\,\delta(k^+-p^+)\,\delta^{(2)}(\v z-\v z')$; in momentum space $\ket{g^a_i(k)}=(2\pi)^{-3/2}a^{\dagger a}_i(k)\ket0$ (Appendix~\ref{app:states}).

Throughout the paper we use the Weizs\"acker--Williams field of a point charge and the transverse kernels built from it,
\begin{align}\label{eq:kernels}
\Kv^i(\v x)\equiv\frac{x^i}{\v x^2}\,,\qquad
K(\v x,\v y,\v z)\equiv\Kv(\v x-\v z)\cdot\Kv(\v y-\v z)=\frac{(\v x-\v z)\cdot(\v y-\v z)}{(\v x-\v z)^2(\v y-\v z)^2}\,,
\end{align}
so that $K(\v x,\v x,\v z)=1/(\v x-\v z)^2$ and $K(\v x,\v x,\v z)-2K(\v x,\v y,\v z)+K(\v y,\v y,\v z)=(\v x-\v y)^2/[(\v x-\v z)^2(\v y-\v z)^2]$ is the dipole kernel.

\subsection{Soft and valence modes; the eikonal Hamiltonian}\label{sec:hamiltonian}

Following the CGC logic~\cite{McLerran:1993ni,McLerran:1993ka,Jalilian-Marian:1997jx,Kovner:2005en}, the gluon modes of the projectile are divided according to their longitudinal momentum. The valence modes have $k^+>\vee$, the soft modes have
\begin{align}
\Lambda<k^+<\vee\,,\qquad \vee\equiv\Lambda\,e^{\delta Y}\,,
\end{align}
i.e.\ $\vee$ is the upper and $\Lambda$ the lower cut-off on the plus momentum of the soft modes, and the modes with $k^+<\Lambda$ do not belong to the projectile wave function; their effect is part of the target. The rapidity interval $\delta Y=\log(\vee/\Lambda)$ is the interval over which the soft modes are integrated; the notation is that of~\cite{Lublinsky:2016meo,Radhakrishnan:2026soft}. Since the valence modes are much faster than the soft ones, they act on the soft modes only through their colour charge density
\begin{align}\label{eq:rho}
\rho^a(\v x)\,,\qquad \big[\rho^a(\v x),\rho^b(\v y)\big]=i f^{abc}\,\rho^c(\v x)\,\delta^{(2)}(\v x-\v y)\,,\qquad
\big[\rho^a(\v k),\rho^b(\v p)\big]=i f^{abc}\,\rho^c(\v k+\v p)\,,
\end{align}
which is an operator on the valence Hilbert space; its microscopic definition in terms of the valence gluon operators is given in~\cite{Radhakrishnan:2026soft}. We keep the charges non-commuting throughout. This is necessary: products of $\rho$'s appear in the cross section in a definite order, determined by the time ordering of the light-cone perturbation theory, and the commutators~\eqref{eq:rho} reduce the number of $\rho$'s by one. These reductions produce some of the logarithmic terms of Section~\ref{sec:rapidity}.

After the decomposition of the gauge field into valence and soft components, all terms suppressed by powers of $k^+_{\rm soft}/k^+_{\rm valence}$ are dropped (eikonal approximation). The light-cone Yang--Mills Hamiltonian ($N_f=0$) then takes the form~\cite{Lublinsky:2016meo,Radhakrishnan:2026soft}
\begin{align}\label{eq:Hdecomp}
&H=H_0+H_g+H_{ggg}+H_{gg\text{-inst}}+H_{gggg}+\dots,\notag\\
&H_0=\int_\Lambda^\vee\frac{dk^+}{2\pi}\int\frac{d^2k}{(2\pi)^2}\,\frac{\v k^2}{2k^+}\,a^{\dagger a}_i(k^+,\v k)\,a^a_i(k^+,\v k)\,,
\end{align}
where the dots stand for terms which do not contribute at the order considered. The coupling of the soft gluons to the valence charges is
\begin{align}\label{eq:Hg}
H_g=\int_\Lambda^\vee\frac{dk^+}{2\pi}\int\frac{d^2k}{(2\pi)^2}\,\frac{g\,k^i}{\sqrt2\,|k^+|^{3/2}}\Big[a^{\dagger a}_i(k^+,\v k)\,\rho^a(-\v k)+a^a_i(k^+,\v k)\,\rho^a(\v k)\Big]\,.
\end{align}
It describes the emission and absorption of a soft gluon by the valence charges. $H_{ggg}$ is the triple-gluon vertex among the soft modes, $H_{gg\text{-inst}}$ the instantaneous (Coulomb-like) interaction between two soft gluons and the valence charge density, and $H_{gggg}$ the four-gluon interaction of the soft modes. Their explicit forms and the matrix elements used in this paper are listed in Appendix~\ref{app:hamiltonian}.

\subsection{Light-cone perturbation theory with operator-valued matrix elements}\label{sec:lcpt}

The soft gluon wave function of the projectile is obtained in time-independent (old-fashioned) perturbation theory in $H-H_0$ on the light cone~\cite{Lepage:1980fj,Brodsky:1997de}. Since the matrix elements of $H_g$ and $H_{gg\text{-inst}}$ are operators on the valence space, the perturbative series is written with the matrix elements in the order of the light-cone time ordering: the vertex acting later stands to the left. With $E_0=0$ for the soft vacuum, the first terms of the expansion of the dressed state in terms of free Fock states are
\begin{align}\label{eq:lcpt}
\ket{\Psi}=\ &\ket0-\sum_n\ket n\frac{V_{n0}}{E_n}+\sum_{n,m}\ket n\frac{V_{nm}V_{m0}}{E_nE_m}
-\sum_{n,m,l}\ket n\frac{V_{nm}V_{ml}V_{l0}}{E_nE_mE_l}
+\sum_{n,m}\ket n\frac{V_{n0}\,V_{0m}V_{m0}}{E_n^2E_m}
\notag\\
&-\frac12\sum_m\frac{V_{0m}V_{m0}}{E_m^2}\ket0
+\frac12\sum_{n,m}\ket n\frac{V_{n0}\,V_{0m}V_{m0}}{E_nE_m^2}+\dots,
\end{align}
where $V_{nm}=\braket{n|H-H_0|m}$ and the sums run over Fock states of the soft sector. The first line is the Rayleigh--Schr\"odinger series with intermediate normalisation; the second line implements the canonical normalisation of the state and is needed at orders $g^2$ (vacuum term) and $g^3$ (one-gluon term). In the last term of each line the single-emission matrix element $V_{n0}$ stands to the left of the norm factor $V_{0m}V_{m0}$; this ordering follows from the perturbative construction of the unitary operator that diagonalises $H$ with operator-valued couplings, and we have verified it against an explicit construction of $\Omega$ with finite-dimensional matrices representing the charges. Energy denominators that vanish inside an integration region are understood as principal values (Appendix~\ref{app:distributions}); this is the standing-wave dressing, and the out-state dressing of the observed gluon differs by the finite term $\Delta N$ of Section~\ref{sec:fin_pv}.

\subsection{The evolution operator}\label{sec:omega}

The dressed (physical) states of the projectile are obtained from the free Fock states by a unitary operator $\Omega$, $\Omega^\dagger\Omega=1$. Since the soft modes are treated in the eikonal approximation, $\Omega$ can be written as a normal-ordered series in the soft gluon operators, with coefficients that are functionals of $\rho$~\cite{Radhakrishnan:2026soft}:
\begin{align}\label{eq:OmegaExp}
\Omega = \ &\mf{N}
+\int\frac{d^3p}{(2\pi)^3}\Big[\mf{A}^b_{1j}(p)\,a^b_j(p)+\mf{A}^b_{2j}(p)\,a^{\dagger b}_j(p)\Big]
\notag\\&
+\int\frac{d^3p}{(2\pi)^3}\frac{d^3q}{(2\pi)^3}\Big[\mf{B}^{cb}_{1kj}(p,q)\,a^b_j(p)a^c_k(q)+\mf{B}^{cb}_{2kj}(p,q)\,a^{\dagger b}_j(p)a^{\dagger c}_k(q)
\notag\\&\hspace{3.4cm}+\mf{B}^{cb}_{3kj}(p,q)\,a^{\dagger c}_k(q)a^b_j(p)\Big]
\notag\\&
+\int\frac{d^3p}{(2\pi)^3}\frac{d^3q}{(2\pi)^3}\Big[\mf{C}^{dcb}_{1lkj}(p,q)\,a^{\dagger d}_l(q)a^{\dagger c}_k(p-q)a^b_j(p)
\notag\\&\hspace{3.4cm}+\mf{C}^{dcb}_{2lkj}(p,q)\,a^{\dagger b}_j(p)a^c_k(p-q)a^d_l(q)\Big]+\dots
\end{align}
The dots denote terms with more gluon operators, which do not contribute to the single inclusive cross section at $O(g^4)$. The coefficients have definite orders in $g$: $\mf N=1+\mf N_2$ with $\mf N_2=O(g^2)$; $\mf A_{1,2}=\mf A^{(1)}_{1,2}+\mf A^{(3)}_{1,2}$ with $\mf A^{(n)}=O(g^n)$; $\mf B_{1,2,3}=O(g^2)$; $\mf C_{1,2}=O(g)$. The coefficients that multiply only creation operators, $\mf A_2$ and $\mf B_2$, together with $\mf N$, are fixed by the action of $\Omega$ on the soft vacuum, i.e.\ by the light-cone wave function of the projectile; $\mf B_3$ and $\mf C_1$ are fixed by the action of $\Omega$ on a one-gluon state. The remaining coefficients follow from unitarity, imposed order by order in $g$ (Appendix~\ref{app:unitarity}):
\begin{align}\label{eq:unitarity_short}
&\mf A^{(1)\dagger b}_{1j}(p)=-\mf A^{(1)b}_{2j}(p)\,,\qquad
\mf C^{\dagger dcb}_{2lkj}(p,q)=-\mf C^{dcb}_{1lkj}(p,q)\,,\notag\\
&\mf N=1-\frac12\int\frac{d^3p}{(2\pi)^3}\,\mf A^{(1)b}_{1j}(p)\,\mf A^{(1)\dagger b}_{1j}(p)\,,
\end{align}
together with the relations for $\mf A^{(3)}_1$ and $\mf B_1$ given in Eqs.~\eqref{eq:A31unit} and~\eqref{eq:B1unit}. In the relation for $\mf A^{(3)}_1$ the operator character of the coefficients matters: $\mf N_2$ does not commute with $\mf A^{(1)}$, since $[\rho^a(\v x)\rho^a(\v y),\rho^b(\v w)]\neq0$ for $\v x\ne\v y$.

In transverse coordinate space, after the unitarity relations~\eqref{eq:unitarity_short} are used, the part of $\Omega$ that contributes is
\begin{align}\label{eq:OmegaX}
\Omega = \ &\mf N
+\int\frac{dp^+}{2\pi}\int_x\Big[\mf A^{(1)b}_{1j}(p^+,\v x)\,a^b_j(p^+,\v x)-\mf A^{\dagger(1)b}_{1j}(p^+,\v x)\,a^{\dagger b}_j(p^+,\v x)
\notag\\&\hspace{2.6cm}
+\mf A^{(3)b}_{1j}(p^+,\v x)\,a^b_j(p^+,\v x)+\mf A^{(3)b}_{2j}(p^+,-\v x)\,a^{\dagger b}_j(p^+,\v x)\Big]
\notag\\&
+\int\frac{dp^+}{2\pi}\frac{dq^+}{2\pi}\int_{x,y}\Big[\mf B^{cb}_{1kj}(p^+,\v x;q^+,\v y)\,a^b_j(p^+,\v x)a^c_k(q^+,\v y)
\notag\\&\hspace{3.2cm}
+\mf B^{cb}_{2kj}(p^+,-\v x;q^+,-\v y)\,a^{\dagger b}_j(p^+,\v x)a^{\dagger c}_k(q^+,\v y)
\notag\\&\hspace{3.2cm}
+\mf B^{cb}_{3kj}(p^+,\v x;q^+,-\v y)\,a^{\dagger c}_k(q^+,\v y)a^b_j(p^+,\v x)\Big]
\notag\\&
+\int\frac{dp^+}{2\pi}\frac{dq^+}{2\pi}\int_{x,y,z}\Big[\mf C^{dcb}_{1lkj}(p^+,\v z-\v y;q^+,\v y-\v x)
\notag\\&\hspace{3.2cm}\times a^{\dagger d}_l(q^+,\v x)a^{\dagger c}_k(p^+-q^+,\v y)a^b_j(p^+,\v z)
\notag\\&\hspace{3.2cm}
-\mf C^{\dagger dcb}_{1lkj}(p^+,\v z-\v y;q^+,\v y-\v x)
\notag\\&\hspace{3.2cm}\times a^{\dagger b}_j(p^+,\v z)a^c_k(p^+-q^+,\v y)a^d_l(q^+,\v x)\Big]+\dots
\end{align}
The transverse arguments of the coefficients are fixed by the Fourier conventions~\eqref{eq:FTconventions}.

\subsection{The coefficients}\label{sec:coefficients}

We now list the coefficients in transverse coordinate space. They are obtained by comparing $\Omega\ket0$ and $\Omega\ket{g^a_i(k^+,\v z)}$ with the LCPT wave functions of Appendix~\ref{app:wavefunctions}, sector by sector in Fock space; the factors of $2\pi$ come from the normalisation~\eqref{eq:states}. The momentum-space expressions and the intermediate steps are given in Appendix~\ref{app:wavefunctions} (for the coefficients of order $g$ and $g^2$ also in~\cite{Radhakrishnan:2026soft}).

\paragraph{Normalisation.} The coefficient of the identity is fixed by the norm of the LO wave function,
\begin{align}\label{eq:N2}
\mf N=1+\mf N_2\,,\qquad
\mf N_2&=-\frac{g^2}{8\pi^3}\,\log\frac{\vee}{\Lambda}\int_{x,y,z}K(\v x,\v y,\v z)\,\rho^a(\v x)\rho^a(\v y)
\notag\\&=-\frac{g^2}{8\pi^3}\,\log\frac{\vee}{\Lambda}\int d^2p\,\frac{\rho^a(-\v p)\rho^a(\v p)}{\v p^2}\,.
\end{align}
$\mf N_2$ is the (virtual) probability of not emitting a soft gluon anywhere in the interval $\Lambda<p^+<\vee$, and it carries the full logarithm $\log(\vee/\Lambda)$.

\paragraph{Single-gluon emission at $O(g)$.} The Weizs\"acker--Williams amplitude for the emission of a soft gluon at $\v z$ is
\begin{align}\label{eq:A1}
\mf A^{(1)a}_{2i}(k^+,-\v z)=-\mf A^{(1)\dagger a}_{1i}(k^+,\v z)=\frac{ig}{\sqrt2\,\pi\sqrt{k^+}}\int_x\Kv^i(\v x-\v z)\,\rho^a(\v x)\,.
\end{align}
It corresponds to the LO wave function $\ket{\Psi^{\rm LO}}=\frac{ig}{2\pi^{3/2}}\int_\Lambda^\vee\frac{dk^+}{\sqrt{k^+}}\int_{x,z}\Kv^i(\v x-\v z)\,\rho^a(\v x)\ket{g^a_i(k^+,\v z)}$, Eq.~\eqref{eq:PsiLO}. In what follows we use the shorthand $\mf A\equiv\mf A^{(1)}_2$ when no confusion can arise, and a bar denotes the rotation of all charges by the target, $\bar{\mf A}^a(\v x)=\mf A^a(\v x)\big|_{\rho(\v u)\to U(\v u)\rho(\v u)}$, see Section~\ref{sec:observable}.

\paragraph{Single-gluon emission at $O(g^3)$.} The coefficient $\mf A^{(3)}_2$ contains terms with one, two and three charges. In momentum space the one-$\rho$ part reads, for $\Lambda\to0$ wherever this is possible,
\begin{align}\label{eq:A3onerho}
\Big[\mf A^{(3)a}_{2i}(k^+,\v k)\Big]_{\rho}=\ &\frac{g^3N_c\,\rho^a(-\v k)\,k^i}{8\sqrt{2k^+}\,\pi^2\,\v k^2}\Bigg\{\Big[-\frac2\epsilon+\log\frac{\v k^2}{\mu^2_{\MS}}\Big]
\notag\\&\hspace{1.5cm}\times\Big[\frac{11}{3}-2\log\frac{k^+}{\Lambda}+\log\Big(1-\frac{k^+}{\vee}\Big)\Big]
\notag\\&\hspace{1.5cm}
+\frac{\pi^2}{3}-\frac{67}{9}+\log^2\Big(1-\frac{k^+}{\vee}\Big)+\log\Big(1-\frac{k^+}{\vee}\Big)\log\frac{\vee}{k^+}\Bigg\},
\end{align}
where the transverse integrals are performed in $d=2-\epsilon$ dimensions in the $\MS$ scheme. It collects the one-loop self-energy of the emitted gluon, the vertex corrections in which the emitted gluon splits and recombines, and the one-$\rho$ (commutator) parts of the diagrams with two valence vertices (Figures~\ref{fig:vac_g3_onerho} and~\ref{fig:vac_g3_tworho}). The pole multiplies the combination $\frac{11}{3}-2\log(k^+/\Lambda)$, whose first term is $\beta_0/N_c$ for pure glue, $\beta_0=\frac{11}{3}N_c$. In coordinate space,
\begin{align}\label{eq:A3onerhoX}
\Big[\mf A^{(3)a}_{2i}(k^+,-\v z)\Big]_{\rho}=\ &-\int_x\frac{ig^3N_c\,\rho^a(\v x)\,\Kv^i(\v x-\v z)}{16\sqrt2\,\pi^3\sqrt{k^+}}\Bigg\{\Big[-\frac2\epsilon-2\gamma_E-\log\frac{(\v x-\v z)^2\mu^2_{\MS}}{4}\Big]
\notag\\&\hspace{1.5cm}\times\Big[\frac{11}{3}-2\log\frac{k^+}{\Lambda}+\log\Big(1-\frac{k^+}{\vee}\Big)\Big]
\notag\\&\hspace{1.5cm}
-\frac{67}{9}+\frac{\pi^2}{3}+\log^2\Big(1-\frac{k^+}{\vee}\Big)+\log\Big(1-\frac{k^+}{\vee}\Big)\log\frac{\vee}{k^+}\Bigg\}.
\end{align}
The two-$\rho$ part is antisymmetric in colour, $\propto f^{abc}\{\rho^b,\rho^c\}$, and is conveniently written in a mixed representation,
\begin{align}\label{eq:A3tworho}
\Big[\mf A^{(3)a}_{2i}(k^+,-\v z)\Big]_{\rho\rho}=\ &\int d^2p\,d^2k\int_{x,y}e^{i\v k\cdot(\v z-\v y)}e^{i\v p\cdot(\v y-\v x)}\,\frac{ig^3f^{abc}\{\rho^b(\v x),\rho^c(\v y)\}}{32\sqrt{2k^+}\,\pi^5\,\v p^2\,\v k^2}
\notag\\&\times\Bigg[\log\frac{\Lambda}{k^+}\,\frac{\v k^2p^i-(\v k\cdot\v p)\,k^i}{(\v k-\v p)^2}
-2k^i\log\frac{\vee}{\Lambda}
-2k^i\log\frac{k^+\v p^2}{\vee\,\v k^2+k^+\v p^2}
\notag\\&\quad
+\frac{(2\v p^2-2\v k\cdot\v p+\v k^2)\,k^i+\v k^2p^i}{(\v k-\v p)^2}\log\Big[1+\Big(\frac{\vee}{k^+}-1\Big)\frac{\v k^2}{(\v k-\v p)^2}\Big]
\notag\\&
-\frac{\v p^2k^i+\v k^2p^i}{(\v k-\v p)^2}\log\frac{\vee-k^+}{\Lambda}-k^i\log\frac{\vee}{k^+}\Bigg],
\end{align}
where, inside the integral, $\v k$ and $\v p$ are integration variables. The three-$\rho$ part originates from the third-order terms of Eq.~\eqref{eq:lcpt} in which a gluon is emitted and reabsorbed by the valence charges before or after the emission of the observed gluon (Figure~\ref{fig:vac_g3_threerho}). The terms proportional to $1/\Lambda$ and $1/\vee$ cancel exactly, and the remainder is
\begin{align}\label{eq:A3threerho}
\Big[\mf A^{(3)a}_{2i}(k^+,-\v z)\Big]_{\rho\rho\rho}&=-\log\frac{\vee}{\Lambda}\int_{w,x,y,u}\frac{ig^3\,\rho^a(\v w)\,\rho^b(\v x)\,\rho^b(\v y)\,K(\v x,\v y,\v u)\,\Kv^i(\v w-\v z)}{8\sqrt2\,\pi^4\sqrt{k^+}}
\notag\\&=\mf A^{(1)a}_{2i}(k^+,-\v z)\,\mf N_2\,.
\end{align}
The operator ordering in Eq.~\eqref{eq:A3threerho}, with $\rho^a(\v w)$ to the left of the norm factor, is dictated by Eq.~\eqref{eq:lcpt}. It is important in what follows: since $\mf N_2\mf A^{(1)}\neq\mf A^{(1)}\mf N_2$, the opposite ordering would differ by $-[\mf N_2,\mf A^{(1)}]$, which is a two-$\rho$ term equal and opposite to the term $-2k^i\log(\vee/\Lambda)$ of Eq.~\eqref{eq:A3tworho}.

The one-$\rho$, two-$\rho$ and three-$\rho$ parts of the wave function from which $\mf A^{(3)}_2$ is obtained agree with the NLO light-cone wave function of Lublinsky and Mulian~\cite{Lublinsky:2016meo} (eqs.~3.97--3.102 there), including the coefficient $11/3$, the constant $67/9-\pi^2/3$, the dependence on the longitudinal cut-offs and the three-$\rho$ term; the two-$\rho$ parts agree after the antisymmetrisation enforced by $f^{abc}\{\rho^b,\rho^c\}$. The coordinate-space forms~\eqref{eq:A3onerhoX}--\eqref{eq:A3threerho} go beyond the published material.

\paragraph{Two-gluon emission.} The coefficient $\mf B_2$ gives the amplitude for two soft gluons, with longitudinal momenta $k^+$ and $p^+$, at transverse positions $\v w$ and $\v z$:
\begin{align}\label{eq:B2}
\mf B^{bc}_{2ji}(k^+,-\v w;p^+,-\v z)=\ &\int_x\frac{ig^2f^{cba}\,\rho^a(\v x)\,\sqrt{p^+k^+}}{4\pi^2\big[p^+(\v x-\v z)^2+k^+(\v x-\v w)^2\big](p^++k^+)}
\notag\\&\times\Bigg[\frac{\delta^{ij}\big[(\v x-\v z)^2-(\v x-\v w)^2\big]}{2(\v z-\v w)^2}
\notag\\&\quad
+\frac{p^++k^+}{k^+}\Big(\frac{(x-z)^j(z-w)^i}{(\v z-\v w)^2}-\frac{(x-z)^j(x-w)^i}{2(\v x-\v w)^2}\Big)
\notag\\&\quad
+\frac{p^++k^+}{p^+}\Big(\frac{(x-w)^i(z-w)^j}{(\v z-\v w)^2}+\frac{(x-z)^j(x-w)^i}{2(\v x-\v z)^2}\Big)\Bigg]
\notag\\&\quad\times\Theta(\vee-k^+-p^+)
\notag\\&
-\int_{x,y}\frac{g^2\,\{\rho^c(\v y),\rho^b(\v x)\}\,\Kv^i(\v y-\v w)\,\Kv^j(\v x-\v z)}{8\pi^2\sqrt{p^+k^+}}
\notag\\&\quad\times\Theta(\vee-k^+)\,\Theta(\vee-p^+)\,,
\end{align}
with $k^+,p^+>\Lambda$. The first term, linear in $\rho$, arises from the emission of a gluon of momentum $k^++p^+$ which splits into two (triple-gluon vertex), from the instantaneous emission of the pair, and from the commutator part of the emission of two gluons by two charges (Figure~\ref{fig:vac_gg}); the parent gluon is a soft mode, hence $k^++p^+<\vee$. The second term is the independent emission of the two gluons by two valence charges; each gluon is separately a soft mode, and there is no condition on $k^++p^+$.
% CHOICE FOR THE AUTHORS: the support of the commutator part of the two-emission diagram (contained in the first term) is taken as k^+ + p^+ < vee, as for the splitting diagrams; by the kinematics of that diagram it could also be the full square Lambda<k^+,p^+<vee.

\paragraph{Scattering of a soft gluon on the valence charges.} The coefficient $\mf B_3$ describes the transition of a soft gluon of momentum $k^+$ at $\v z$ into a soft gluon of momentum $p^+$ at $\v u$ through its interaction with the valence charges: instantaneous interaction, emission of a second gluon followed by absorption of the first, and absorption of the gluon followed by re-emission (Figure~\ref{fig:g_g}). For $p^+>k^+$ the valence charges transfer longitudinal momentum to the gluon, for $p^+<k^+$ the gluon transfers momentum to the valence charges. With $\Theta_>\equiv\Theta(p^+-k^+-\Lambda)$ and $\Theta_<\equiv\Theta(k^+-\Lambda-p^+)$,
\begin{align}\label{eq:B3}
\mf B^{ba}_{3ji}(k^+,\v z;p^+,-\v u)=\ &\frac{ig^2f^{abc}}{(2\pi)^2}\int_x\frac{\sqrt{p^+k^+}\,(p^++k^+)}{(k^+-p^+)^2}\,\frac{\rho^c(\v x)\,\delta_{ij}}{p^+(\v x-\v u)^2-k^+(\v x-\v z)^2}
\notag\\&\qquad\times\big[\Theta_>+\Theta_<\big]
\notag\\&
+\frac{ig^2f^{abc}}{2\pi^2k^+}\int_x\frac{\sqrt{k^+p^+}}{k^+-p^+}
\notag\\&\qquad\times\frac{\big(k^+(x-z)^m-p^+(x-u)^m\big)\big((x-z)^l-(x-u)^l\big)}{k^+(\v x-\v z)^2-p^+(\v x-\v u)^2}\,\rho^c(\v x)
\notag\\&\qquad\times\Big(\delta_{il}\delta_{jm}-\frac{k^+}{p^+}\delta_{jl}\delta_{im}-\frac{k^+}{k^+-p^+}\delta_{ij}\delta_{lm}\Big)
\notag\\&\qquad\times\Bigg[-\frac{\Theta_>}{(\v u-\v z)^2}-\frac{p^+k^+\,\Theta_<}{\big(k^+(\v x-\v z)-p^+(\v x-\v u)\big)^2}\Bigg]
\notag\\&
-\frac{g^2}{4\pi^2\sqrt{p^+k^+}}\int_{x,y}\Kv^i(\v x-\v z)\,\Kv^j(\v y-\v u)\,\{\rho^a(\v x),\rho^b(\v y)\}
\notag\\&
-\frac{ig^2f^{abc}}{4\pi^2\sqrt{p^+k^+}}\int_x\Kv^i(\v x-\v z)\,\Kv^j(\v x-\v u)
\notag\\&\qquad\times\frac{k^+(\v x-\v z)^2+p^+(\v x-\v u)^2}{p^+(\v x-\v u)^2-k^+(\v x-\v z)^2}\,\rho^c(\v x)\,.
\end{align}
The double pole $1/(k^+-p^+)^2$ at the boundary $|p^+-k^+|=\Lambda$ cancels between the instantaneous and the non-instantaneous terms of the first two lines, so that no term proportional to $k^+/\Lambda$ arises. The vanishing denominators $p^+(\v x-\v u)^2=k^+(\v x-\v z)^2$ are principal values; for the out-state prescription see Section~\ref{sec:fin_pv}.

\paragraph{Gluon splitting.} The coefficient $\mf C_1$ is independent of $\rho$; it is the light-cone amplitude for the splitting of a soft gluon $(k^+,\v z)$ into two soft gluons of momenta $p^+$ and $k^+-p^+$ (Figure~\ref{fig:g_gg}):
\begin{align}\label{eq:C1}
\mf C^{bca}_{1jki}(k^+,\v z;p^+,\v x)=\ &-\frac{gf^{abc}}{2\pi\sqrt{2k^+}}\,\frac{\sqrt{p^+(k^+-p^+)}}{k^+}\,\delta^{(2)}\Big(\v z+\frac{p^+}{k^+}\v x\Big)\,\Kv^m(\v x)
\notag\\&\times\Big[\delta_{im}\delta_{jk}-\frac{k^+}{p^+}\delta_{jm}\delta_{ik}-\frac{k^+}{k^+-p^+}\delta_{km}\delta_{ij}\Big]\Theta(k^+-p^+-\Lambda)\,.
\end{align}
The $\delta$-function expresses the conservation of the transverse centre of mass $p^+\v x_1+(k^+-p^+)\v x_2=k^+\v x_0$ in the splitting.

\paragraph{Coefficients fixed by unitarity.} $\mf A^{(3)}_1$ and $\mf B_1$, which multiply annihilation operators, are not independent; they are given in terms of the coefficients above by Eqs.~\eqref{eq:A31unit} and~\eqref{eq:B1unit}. In the cross section they appear only through these relations.

% ======================================================================
\section{The observable and its expansion}\label{sec:observable}
% ======================================================================

\subsection{Single inclusive gluon spectrum}\label{sec:observable_def}

The projectile, before the collision, is described by its valence state $\ket{\Psi_{\rm v}}$ times the dressing of the soft modes, $\Omega\ket{0}$, where $\ket0$ is the vacuum of the soft modes. The target is described by the eikonal scattering matrix $S$. At high energy the interaction of a soft gluon and of the valence charges with the target is a colour rotation by the adjoint Wilson line $U(\v x)$ at fixed transverse position:
\begin{align}\label{eq:Srotation}
S^\dagger\rho^a(\v x)S&=U^{ab}(\v x)\,\rho^b(\v x)\,,\notag\\
S^\dagger a^a_i(k^+,\v x)S&=U^{ab}(\v x)\,a^b_i(k^+,\v x)\,,\qquad
S^\dagger a^{\dagger a}_i(k^+,\v x)S=U^{ab}(\v x)\,a^{\dagger b}_i(k^+,\v x)\,.
\end{align}
After the collision the state is $S\,\Omega\ket0$. A gluon observed in the final state is a dressed (physical) gluon; the number of dressed gluons is counted by the bare number operator in the basis rotated back by $\Omega^\dagger$~\cite{Kovner:2006wr}. The single inclusive spectrum is therefore
\begin{align}\label{eq:observable}
\frac{d^3N}{d^3k}=\frac{1}{(2\pi)^3}\int_{w,w'}e^{-i\v k\cdot(\v w-\v w')}\,
\Big\langle\bra0\,\Omega^\dagger S^\dagger\Omega\;a^{\dagger a}_i(k^+,\v w')\,a^a_i(k^+,\v w)\;\Omega^\dagger S\,\Omega\ket0\Big\rangle\,,
\end{align}
where $\langle\dots\rangle$ denotes the average over the valence state, i.e.\ over the projectile charge configurations, taken at the end. The factor $(2\pi)^{-3}$ and the phase follow from $a^{\dagger a}_i(k)a^a_i(k)=\int_{w,w'}e^{-i\v k\cdot(\v w-\v w')}a^{\dagger a}_i(k^+,\v w')a^a_i(k^+,\v w)$ and the normalisation~\eqref{eq:aadagger}. The amplitude stands at $\v w$ and the conjugate amplitude at $\v w'$; we use this convention, $e^{-i\v k\cdot(\v w-\v w')}$ with the conjugated factor at the primed points, throughout the paper.

Inserting $SS^\dagger=1$ and defining $\bar\Omega\equiv S^\dagger\Omega S$, which is $\Omega$ with $\rho\to U\rho$ and $a\to Ua$ according to Eq.~\eqref{eq:Srotation}, we obtain
\begin{align}
\frac{d^3N}{d^3k}=\frac{1}{(2\pi)^3}\int_{w,w'}e^{-i\v k\cdot(\v w-\v w')}\,\Big\langle\bra0\,\Omega^\dagger\bar\Omega\;U^{ab'}(\v w')a^{\dagger b'}_i(k^+,\v w')\,U^{ab}(\v w)a^b_i(k^+,\v w)\;\bar\Omega^\dagger\Omega\ket0\Big\rangle\,.
\end{align}
A bar on any coefficient of $\Omega$ denotes, here and below, the replacement $\rho^a(\v u)\to U^{ab}(\v u)\rho^b(\v u)$ of every charge in it. Between $a(k^+,\v w)$ and $\bar\Omega^\dagger\Omega\ket0$ we insert a complete set of soft states. The vacuum does not contribute, and states with three or more gluons contribute at order $g^6$. Hence
\begin{align}\label{eq:N_amplitudes}
\frac{d^3N}{d^3k}=\frac{1}{(2\pi)^3}\Big\langle\mathcal A^{a\dagger}_i(k)\,\mathcal A^a_i(k)+\int dq^+\!\int_{v}\mathcal A^{ab\dagger}_{ij}(k;q^+,\v v)\,\mathcal A^{ab}_{ij}(k;q^+,\v v)\Big\rangle\,,
\end{align}
with the one-gluon and two-gluon amplitudes
\begin{align}\label{eq:amplitudes_def}
\mathcal A^a_i(k)&=\sqrt{2\pi}\int_{w}e^{-i\v k\cdot\v w}\,U^{ab}(\v w)\,\bra{g^b_i(k^+,\v w)}\bar\Omega^\dagger\Omega\ket0\,,\notag\\
\mathcal A^{ab}_{ij}(k;q^+,\v v)&=\sqrt{2\pi}\int_{w}e^{-i\v k\cdot\v w}\,U^{ac}(\v w)\,\bra{g^b_j(q^+,\v v)\,g^c_i(k^+,\v w)}\bar\Omega^\dagger\Omega\ket0\,.
\end{align}
The colour index $b$ of the second gluon in $\mathcal A^{ab}_{ij}$ is that of the bare state; since the second gluon is integrated over, its rotation by $U(\v v)$ drops out of the sum. The operators $\bar\Omega^\dagger\Omega$ are evaluated with the expansion~\eqref{eq:OmegaX}: $\Omega$ acting on $\ket0$ involves only $\mf N$, $\mf A_2$ and $\mf B_2$; $\bar\Omega^\dagger$ then acts on states with up to two gluons. Note that $\bar\Omega^\dagger\Omega=1$ for $U=1$, so that every amplitude vanishes in the absence of the target.

From here on $\mf A\equiv\mf A^{(1)}_2$ and $\mf A^\dagger\equiv\mf A^{(1)\dagger}_2$; the unitarity relation $\mf A^{(1)\dagger}_1(p^+,\v x)=-\mf A(p^+,-\v x)$ is used to eliminate $\mf A^{(1)}_1$ wherever possible, and $\mf A^{(3)}\equiv\mf A^{(3)}_2$.

\subsection{Leading order}\label{sec:LO}

At order $g$ only $\mf A$ contributes:
\begin{align}\label{eq:ALO}
\mathcal A^{a}_{{\rm LO}\,i}(k)&=\int_{w}e^{-i\v k\cdot\v w}\Big[U^{ab}(\v w)\,\mf A^b_i(k^+,-\v w)-\bar{\mf A}^a_i(k^+,-\v w)\Big]
\notag\\&=\frac{ig}{\sqrt2\,\pi\sqrt{k^+}}\int_{w,x}e^{-i\v k\cdot\v w}\,\Kv^i(\v x-\v w)\big[U(\v w)-U(\v x)\big]^{ab}\rho^b(\v x)\,.
\end{align}
The first term is the emission of the gluon before the collision, followed by the scattering of the gluon; the second is the emission after the collision by the rotated charge. The LO spectrum is
\begin{align}\label{eq:LO}
\frac{d^3N_{\rm LO}}{d^3k}=\ &\frac{1}{(2\pi)^3}\,\frac{g^2}{2\pi^2}\,\frac{1}{k^+}\int_{x,x',w,w'}e^{-i\v k\cdot(\v w-\v w')}\,
\Kv(\v x-\v w)\cdot\Kv(\v x'-\v w')
\notag\\&\times
\Big\langle\big[U(\v x')-U(\v w')\big]^{ab'}\rho^{b'}(\v x')\,\big[U(\v x)-U(\v w)\big]^{ab}\rho^{b}(\v x)\Big\rangle\,.
\end{align}
Note that $\frac{1}{(2\pi)^3}\frac{g^2}{2\pi^2}=\frac{\alpha_s}{4\pi^4}$.

\paragraph{Colour-singlet projectile.} For a projectile that is a colour singlet, or more generally whose charge correlator is colour diagonal,
\begin{align}\label{eq:singlet}
\big\langle\rho^a(\v x)\rho^b(\v y)\big\rangle=\frac{\delta^{ab}}{N_c^2-1}\,n(\v x,\v y)\,,\qquad n(\v x,\v y)=\big\langle\rho^f(\v x)\rho^f(\v y)\big\rangle\,,
\end{align}
and the LO spectrum becomes
\begin{align}\label{eq:LOsinglet}
\frac{d^3N_{\rm LO}}{d^3k}=\ &\frac{\alpha_s}{4\pi^4(N_c^2-1)}\,\frac{1}{k^+}\int_{x,y,w,w'}e^{-i\v k\cdot(\v w-\v w')}\,\Kv(\v x-\v w)\cdot\Kv(\v y-\v w')
\notag\\&\times\,n(\v x,\v y)\,
\mathcal O_{\rm LO}(\v x,\v y;\v w,\v w')\,,
\end{align}
with the target operator
\begin{align}\label{eq:OLO}
\mathcal O_{\rm LO}(\v x,\v y;\v w,\v w')={\rm Tr}\Big\{\big[U^\dagger(\v x)-U^\dagger(\v w)\big]\big[U(\v y)-U(\v w')\big]\Big\}\,.
\end{align}
Since $U$ is real, ${\rm Tr}\{A^\dagger B\}={\rm Tr}\{B^\dagger A\}$ for adjoint matrices, and for a symmetric $n$ the integrand is invariant under $(\v x,\v w)\leftrightarrow(\v y,\v w')$; in the colour-singlet form the sign of the exponent is therefore immaterial. Eq.~\eqref{eq:LOsinglet} is the formula of Kovchegov and Tuchin~\cite{Kovchegov:2001sc}: for a single quark, $n\to C_F$ times a $\delta$-function, it reproduces their normalisation $\alpha_sC_F/(4\pi^4)$ exactly. It is also the operator that appears in the LL analysis of Kovner and Lublinsky~\cite{Kovner:2006wr,Kovner:2006ge}.

\paragraph{Momentum space and $k_T$ factorisation.} With $f(\v k)=\int_xe^{i\v k\cdot\v x}f(\v x)$ and Eq.~\eqref{eq:FT_K} of Appendix~\ref{app:fourier}, the LO amplitude takes the form
\begin{align}\label{eq:ALOmom}
\mathcal A^{a}_{{\rm LO}\,i}(k)=\frac{\sqrt2\,g}{\sqrt{k^+}}\int\frac{d^2q}{(2\pi)^2}\,\Big[\frac{k^i}{\v k^2}-\frac{q^i}{\v q^2}\Big]\,U^{ab}(\v q-\v k)\,\rho^b(-\v q)\,,\qquad
U^{ab}(\v l)=\int_xe^{i\v l\cdot\v x}U^{ab}(\v x)\,,
\end{align}
so that the gluon is produced with the Lipatov-type vertex $\v k/\v k^2-\v q/\v q^2$, where $\v q$ is the momentum supplied by the projectile and $\v k-\v q$ the momentum transferred by the target. The term $U^{ab}(\v l)\propto\delta^{ab}\delta^{(2)}(\v l)$ of a trivial target does not contribute, since the vertex vanishes at $\v q=\v k$. For a dilute target, $U^{ab}(\v x)=\delta^{ab}+f^{acb}\alpha^c(\v x)+O(\alpha^2)$, the amplitude is linear in $\alpha$, and with $\big|\frac{\v k}{\v k^2}-\frac{\v q}{\v q^2}\big|^2=\frac{(\v k-\v q)^2}{\v k^2\v q^2}$ the spectrum reduces to the $k_T$-factorised form~\cite{Gribov:1983ivg,Kovchegov:2001sc,Kharzeev:2003wz,Blaizot:2004wu}. For a projectile and a target that are homogeneous in the transverse plane, with $\langle\rho^{b'}(\v q')\rho^b(-\v q)\rangle=\delta^{bb'}(2\pi)^2\delta^{(2)}(\v q-\v q')P(\v q)$ and $\langle\alpha^{c}(\v l)\alpha^{d}(-\v l')\rangle=\delta^{cd}(2\pi)^2\delta^{(2)}(\v l-\v l')T(\v l)$,
\begin{align}\label{eq:kTfact}
\frac{d^3N_{\rm LO}}{d^3k}\simeq\frac{\alpha_s N_c(N_c^2-1)}{\pi^2}\,\frac{S_\perp}{k^+}\int\frac{d^2q}{(2\pi)^2}\,\frac{(\v k-\v q)^2}{\v k^2\,\v q^2}\,P(\v q)\,T(\v k-\v q)\,,
\end{align}
where $S_\perp$ is the transverse area.

\subsection{Next-to-leading order: amplitudes}\label{sec:NLOamplitudes}

At order $g^3$ the one-gluon amplitude receives five types of contributions. Collecting the terms of $\bar\Omega^\dagger\Omega$ that connect $\ket0$ with $\bra{g}$ at this order we find
\begin{align}\label{eq:A3amp}
\mathcal A^a_{(3)\,i}(k)=\int_w e^{-i\v k\cdot\v w}\Big\{\mathcal A_{N}+\mathcal A_{A^{(3)}}+\mathcal A_{B_3}+\mathcal A_{AB_2}+\mathcal A_{CB_2}\Big\}^a_i\,,
\end{align}
with
\begin{align}
\mathcal A_{N}&=-\bar{\mf A}^a_i(k^+,-\v w)\,\mf N_2+U^{ab}(\v w)\,\bar{\mf N}_2\,\mf A^b_i(k^+,-\v w)-\big[\bar{\mf N}_2,\bar{\mf A}^a_i(k^+,-\v w)\big]\,,
\notag\\
\mathcal A_{A^{(3)}}&=U^{ab}(\v w)\,\mf A^{(3)b}_i(k^+,-\v w)-\bar{\mf A}^{(3)a}_i(k^+,-\v w)\,,
\notag\\
\mathcal A_{B_3}&=\int\frac{dp^+}{2\pi}\int_x\bar{\mf B}^{\dagger ba}_{3ji}(k^+,\v w;p^+,-\v x)\Big[U^{bc}(\v x)\,\mf A^c_j(p^+,-\v x)-\bar{\mf A}^b_j(p^+,-\v x)\Big]\,,
\notag\\
\mathcal A_{AB_2}&=2\int\frac{dp^+}{2\pi}\int_x\Big[U^{ad}(\v w)U^{bc}(\v x)\,\bar{\mf A}^{\dagger b}_j(p^+,-\v x)\,\mf B^{cd}_{2ji}(k^+,-\v w;p^+,-\v x)
\notag\\&\hspace{2.2cm}
-\bar{\mf A}^{\dagger b}_j(p^+,-\v x)\,\bar{\mf B}^{ba}_{2ji}(k^+,-\v w;p^+,-\v x)\Big]\,,
\notag\\
\mathcal A_{CB_2}&=2\int\frac{dp^+}{2\pi}\int_{x,z}\Big[\bar{\mf C}^{\dagger dca}_{1jki}(k^+,\v w-\v z;p^+,\v z-\v x)\,U^{db}(\v x)U^{ce}(\v z)\,\mf B^{be}_{2jk}(k^+-p^+,-\v z;p^+,-\v x)
\notag\\&\hspace{2.2cm}
-\bar{\mf C}^{\dagger bca}_{1jki}(k^+,\v w-\v z;p^+,\v z-\v x)\,\bar{\mf B}^{bc}_{2jk}(k^+-p^+,-\v z;p^+,-\v x)\Big]\,.
\end{align}
The term $\mathcal A_N$ contains the normalisation: the term $\bar{\mf A}^{(3)}_1$ of $\bar\Omega^\dagger$ contributes, through the unitarity relation~\eqref{eq:A31unit}, the commutator $-[\bar{\mf N}_2,\bar{\mf A}]$. $\mathcal A_{A^{(3)}}$ is the $O(g^3)$ one-gluon amplitude of the wave function, before and after the collision. $\mathcal A_{B_3}$ describes the scattering of an already emitted gluon of momentum $p^+$ into the observed gluon by its interaction with the valence charges after the collision ($\bar{\mf B}_3^\dagger$); $\mathcal A_{AB_2}$, the absorption, after the collision, of one of the two gluons of the two-gluon component of the wave function; $\mathcal A_{CB_2}$, the recombination of the two gluons of that component into the observed one. In each term the two contributions correspond to an interaction of the soft gluons with the target before and after the respective valence interaction; their difference vanishes for $U=1$.

The terms of $\mathcal A_N$ are not all of the same order in $\rho$. With the ordering of the three-$\rho$ part of $\mf A^{(3)}$ derived in Section~\ref{sec:coefficients}, $[\mf A^{(3)}]_{\rho\rho\rho}=\mf A\,\mf N_2$, the commutator in $\mathcal A_N$ cancels against the corresponding part of $\mathcal A_{A^{(3)}}$, and the four-$\rho$ (normalisation) terms of the amplitude add up to
\begin{align}\label{eq:Nterms}
-\bar{\mf A}\,\mf N_2+U\,\bar{\mf N}_2\,\mf A+U\,\mf A\,\mf N_2-\bar{\mf N}_2\,\bar{\mf A}\,,
\end{align}
which vanishes for $U=1$, as it must. With the opposite ordering, $\mf N_2\mf A$, a spurious two-$\rho$ term $\propto\log(\vee/\Lambda)$ would remain.

The two-gluon amplitude at order $g^2$ is
\begin{align}\label{eq:A2gamp}
\mathcal A^{ab}_{ij}(k;q^+,\v v)=\frac{2}{\sqrt{2\pi}}\int_w e^{-i\v k\cdot\v w}\Big\{\mathcal A_B+\mathcal A^{(1)}_{AA}+\mathcal A^{(2)}_{AA}-\mathcal A_{CA}\Big\}^{ab}_{ij}\,,
\end{align}
with
\begin{align}
\mathcal A_B&=U^{ac}(\v w)\,\mf B^{bc}_{2ji}(k^+,-\v w;q^+,-\v v)-U^{\dagger bc}(\v v)\,\bar{\mf B}^{ca}_{2ji}(k^+,-\v w;q^+,-\v v)\,,
\notag\\
\mathcal A^{(1)}_{AA}&=\tfrac12\Big[U^{\dagger bc}(\v v)\,\bar{\mf A}^c_j(q^+,-\v v)\,\bar{\mf A}^a_i(k^+,-\v w)-U^{ac}(\v w)\,\bar{\mf A}^d_j(q^+,-\v v)\,U^{\dagger bd}(\v v)\,\mf A^c_i(k^+,-\v w)\Big]\,,
\notag\\
\mathcal A^{(2)}_{AA}&=\tfrac12\Big[U^{\dagger bc}(\v v)\,\bar{\mf A}^a_i(k^+,-\v w)\,\bar{\mf A}^c_j(q^+,-\v v)-\bar{\mf A}^a_i(k^+,-\v w)\,\mf A^b_j(q^+,-\v v)\Big]\,,
\notag\\
\mathcal A_{CA}&=\int_z\Big[\bar{\mf C}^{dac}_{1jik}(q^++k^+,\v z-\v w;q^+,\v w-\v v)\,U^{\dagger bd}(\v v)\,U^{ce}(\v z)\,\mf A^e_k(q^++k^+,-\v z)
\notag\\&\hspace{1.2cm}
-U^{\dagger bd}(\v v)\,\bar{\mf C}^{dac}_{1jik}(q^++k^+,\v z-\v w;q^+,\v w-\v v)\,\bar{\mf A}^c_k(q^++k^+,-\v z)\Big]\,.
\end{align}
$\mathcal A_B$ is the two-gluon component of the wave function scattered on the target; $\mathcal A^{(1)}_{AA}$ and $\mathcal A^{(2)}_{AA}$ arise from the product of two single emissions, one before and one after the collision, in the two operator orderings generated by $\bar\Omega^\dagger$; and $\mathcal A_{CA}$ is the splitting, after the collision, of a gluon of momentum $q^++k^+$ emitted before the collision into the observed gluon and a second gluon of momentum $q^+$.

\subsection{Next-to-leading order: Groups and rows}\label{sec:groups}

The $O(g^4)$ cross section is
\begin{align}\label{eq:NLOdecomp}
\frac{d^3N}{d^3k}=\frac{d^3N_{\rm LO}}{d^3k}+\frac{d^3N^{\rm(I)}}{d^3k}+\frac{d^3N^{\rm(II)}}{d^3k}+\frac{d^3N^{\rm(III)}}{d^3k}+O(g^6)\,,
\end{align}
where
\begin{align}\label{eq:Groups}
\frac{d^3N^{\rm(I)}}{d^3k}&=\frac{1}{(2\pi)^3}\Big\langle\mathcal A^{a\dagger}_{{\rm LO}\,i}(k)\,\mathcal A^a_{(3)\,i}(k)\Big\rangle+{\rm c.c.}\,,\notag\\
\frac{d^3N^{\rm(II)}}{d^3k}&=\frac{1}{(2\pi)^3}\int dq^+\!\int_v\Big\langle\big|\mathcal A_B\big|^2+\big|\mathcal A^{(1)}_{AA}\big|^2+\big|\mathcal A^{(2)}_{AA}\big|^2+\big|\mathcal A_{CA}\big|^2\Big\rangle\,,\notag\\
\frac{d^3N^{\rm(III)}}{d^3k}&=\frac{1}{(2\pi)^3}\int dq^+\!\int_v\Big\langle\mathcal A^{(1)\dagger}_{AA}\mathcal A_B+\mathcal A^{(2)\dagger}_{AA}\mathcal A_B-\mathcal A^{\dagger}_{CA}\mathcal A_B
\notag\\&\hspace{3.3cm}+\mathcal A^{(2)\dagger}_{AA}\mathcal A^{(1)}_{AA}-\mathcal A^{\dagger}_{CA}\mathcal A^{(1)}_{AA}-\mathcal A^{\dagger}_{CA}\mathcal A^{(2)}_{AA}\Big\rangle+{\rm c.c.}
\end{align}
Here $|\mathcal A|^2\equiv\mathcal A^{ab\dagger}_{ij}\mathcal A^{ab}_{ij}$, the sign of $\mathcal A_{CA}$ is that of Eq.~\eqref{eq:A2gamp}, and the momentum of the integrated gluon is renamed $q^+\to p^+$ below. Group~I is the interference of the LO amplitude with the $O(g^3)$ one-gluon amplitude; Group~II consists of the squares of the four components of the two-gluon amplitude; Group~III of their interferences. Groups~I and~III are written with their complex conjugates; Group~II is Hermitian by itself. We call the individual products of coefficients \emph{rows}. They are listed in Table~\ref{tab:rows}, together with the range of the longitudinal integration, the number of charges and the parts of the final result to which they contribute.

\begin{table}[t]
\centering
\renewcommand{\arraystretch}{1.2}
\footnotesize
\begin{tabular}{@{}l p{4.0cm} p{2.9cm} l p{3.4cm}@{}}
\toprule
Row & product (Eqs.~\eqref{eq:A3amp},~\eqref{eq:A2gamp}) & range of $p^+$ & $\#\rho$ & contributes to\\
\midrule
I.I & LO$^\dagger\times$ normalisation terms~\eqref{eq:Nterms} & --- & 4 & 4$\rho$ ($\log\frac{\vee}{\Lambda}$)\\
I.II & LO$^\dagger\times[U\mf A^{(3)}-\bar{\mf A}^{(3)}]$ & --- & 2--4 & 4$\rho$, BFKL, JIMWLK, 3$\rho$, finite\\
I.III & LO$^\dagger\times\bar{\mf B}_3^\dagger\mf A$ & $\Lambda<p^+<\vee$, $|p^+-k^+|>\Lambda$ & 3, 4 & BFKL, 3$\rho$, finite\\
I.IV & LO$^\dagger\times\bar{\mf A}^\dagger\mf B_2$ & $\Lambda<p^+<\vee-k^+$ $[\vee]$ & 3, 4 & 4$\rho$, BFKL, 3$\rho$; finite part $=-$III.I\\
I.V & LO$^\dagger\times\bar{\mf C}_1^\dagger\mf B_2$ & $\Lambda<p^+<k^+-\Lambda$ & 2, 3 & JIMWLK, 3$\rho$, finite\\
\midrule
II.I & $|\mathcal A_B|^2$ & $\Lambda<p^+<\vee-k^+$ $[\vee]$ & 2--4 & 4$\rho$, BFKL, JIMWLK, 3$\rho$, finite\\
II.II & $|\mathcal A^{(1)}_{AA}|^2$ & $\Lambda<p^+<\vee$ & 4 & 4$\rho$\\
II.III & $|\mathcal A^{(2)}_{AA}|^2$ & $\Lambda<p^+<\vee$ & 4 & 4$\rho$\\
II.IV & $|\mathcal A_{CA}|^2$ & $\Lambda<p^+<\vee-k^+$ & 2 & BFKL, JIMWLK, DGLAP, RC, finite\\
\midrule
III.I & $\mathcal A^{(1)\dagger}_{AA}\mathcal A_B$ & $\Lambda<p^+<\vee-k^+$ $[\vee]$ & 3, 4 & 4$\rho$, BFKL, 3$\rho$; finite part $=-$I.IV\\
III.II & $\mathcal A^{(2)\dagger}_{AA}\mathcal A_B$ & $\Lambda<p^+<\vee-k^+$ $[\vee]$ & 3, 4 & 4$\rho$, BFKL, 3$\rho$, finite\\
III.III & $-\mathcal A^{\dagger}_{CA}\mathcal A_B$ & $\Lambda<p^+<\vee-k^+$ & 2, 3 & BFKL, JIMWLK, 3$\rho$, finite\\
III.IV & $\mathcal A^{(2)\dagger}_{AA}\mathcal A^{(1)}_{AA}$ & $\Lambda<p^+<\vee$ & 4 & 4$\rho$\\
III.V & $-\mathcal A^{\dagger}_{CA}\mathcal A^{(1)}_{AA}$ & $\Lambda<p^+<\vee-k^+$ & 3 & BFKL, 3$\rho$, finite\\
III.VI & $-\mathcal A^{\dagger}_{CA}\mathcal A^{(2)}_{AA}$ & $\Lambda<p^+<\vee-k^+$ & 3 & BFKL, 3$\rho$, finite\\
\bottomrule
\end{tabular}
\caption{The rows of the $O(g^4)$ cross section. ``LO$^\dagger$'' is the conjugate LO amplitude, Eq.~\eqref{eq:ALO}; $p^+$ is the longitudinal momentum of the second (integrated or virtual) gluon. Where two ranges apply, the first refers to the part linear in $\rho$ of $\mf B_2$ (its parent gluon carries $p^++k^+$) and the bracket $[\vee]$ to its $\{\rho,\rho\}$ part, $\Lambda<p^+<\vee$. ``$\#\rho$'' is the number of charges before reordering; ``4$\rho$'' and ``3$\rho$'' denote the terms quartic and cubic in $\rho$ multiplying a rapidity logarithm, which cancel (Section~\ref{sec:rapidity}). Row I.II contains in addition the ultraviolet $\beta_0$ terms of the vertex correction, which cancel against Row I.V (Section~\ref{sec:dglap_rc}).}
\label{tab:rows}
\end{table}

The longitudinal ranges in Table~\ref{tab:rows} follow from the kinematics of the coefficients. All soft gluons satisfy $\Lambda<p^+<\vee$. Whenever a soft gluon of momentum $p^++k^+$ appears as an intermediate state (the parent of $\mf C_1$, the part of $\mf B_2$ linear in $\rho$), the condition $p^++k^+<\vee$ applies. The $\{\rho,\rho\}$ part of $\mf B_2$ describes two independent emissions and has support on the full square $\Lambda<k^+,p^+<\vee$; this is required for the cancellation of the four-$\rho$ terms (Section~\ref{sec:fourrho}). In $\mf B_3$ the momentum difference of the incoming and outgoing gluon is carried by the valence charges and satisfies $|p^+-k^+|>\Lambda$.

The evaluation of each row consists of the colour algebra of the products of coefficients, with the charges kept in their order, the longitudinal integration over $p^+$, and the transverse integrations that can be done in closed form. Three kinds of terms emerge:
\begin{enumerate}
\item terms proportional to $\log(\vee/k^+)$ or $\log(k^+/\Lambda)$, quadratic, cubic or quartic in $\rho$, which arise from the regions $p^+\gg k^+$ and $p^+\ll k^+$ of the longitudinal integrals. They are discussed in Section~\ref{sec:rapidity};
\item terms with a collinear singularity, from the region in which the integrated gluon of Row~II.IV is collinear with the observed one, and terms proportional to $\beta_0$ (Section~\ref{sec:dglap_rc});
\item the remaining finite terms (Section~\ref{sec:finite}).
\end{enumerate}
Two rows require a remark. The non-logarithmic parts of Rows~I.IV and III.I cancel identically. Rows~II.II, II.III and III.IV are products of four coefficients $\mf A$, each linear in $\rho$; they contain only four-$\rho$ terms.

% ======================================================================
\section{Rapidity logarithms: BFKL evolution of the projectile and JIMWLK evolution of the target}\label{sec:rapidity}
% ======================================================================

In this section we collect all terms of the $O(g^4)$ cross section that are enhanced by a logarithm of the longitudinal cut-offs and show that they are the leading-logarithmic evolution of the LO spectrum: the BFKL evolution of the projectile charge correlator over the rapidity interval between the projectile and the produced gluon, and the JIMWLK evolution of the target operator over the interval between the produced gluon and the target.

\subsection{Origin of the logarithms}\label{sec:loglong}

Rapidity logarithms arise from longitudinal integrals of the form $\int dp^+/p^+$ over the momentum of the second gluon, real (integrated over in Groups~II and~III) or virtual (inside the coefficients $\mf N_2$, $\mf A^{(3)}$, $\mf B_3$ and in the interference terms of Group~I). In every row the integrand behaves as $1/p^+$ in two regions: for $k^+\ll p^+<\vee$, where the second gluon is faster than the observed one and belongs to the projectile, and for $\Lambda<p^+\ll k^+$, where it is slower and is resolved by the observed gluon as part of the target. The first region produces $\log(\vee/k^+)$, the second $\log(k^+/\Lambda)$. Terms that do not depend on $k^+$ at all, in particular the normalisation $\mf N_2$ and the products of four single-emission amplitudes, carry the full logarithm
\begin{align}\label{eq:logsplit}
\log\frac{\vee}{\Lambda}=\log\frac{\vee}{k^+}+\log\frac{k^+}{\Lambda}\,,
\end{align}
which we split at the longitudinal momentum $k^+$ of the produced gluon. With this split every logarithmic term is assigned either to the projectile or to the target side. Terms proportional to $\log(1-k^+/\vee)$, which arise from the upper limit $\vee-k^+$ of the integrals with a soft parent gluon, vanish for $\vee\to\infty$ and are dropped.

By the order of the charges, the logarithmic terms are quartic, cubic or quadratic in $\rho$ (Table~\ref{tab:rows}). We discuss them in this order, because the reordering of the non-commuting charges feeds terms with more charges into terms with fewer.

\subsection{Terms quartic in \texorpdfstring{$\rho$}{rho}}\label{sec:fourrho}

Terms with four charges arise in Rows I.I, I.II (through the three-$\rho$ part $\mf A\mf N_2$ of $\mf A^{(3)}$), I.IV, II.I, II.II, II.III, III.I, III.II and III.IV, from the products of four single-emission coefficients $\mf A$, of $\mf A$ with $\mf N_2$, and of $\mf A$ with the $\{\rho,\rho\}$ part of $\mf B_2$. In all of them the longitudinal integral runs over the full interval $\Lambda<p^+<\vee$ and gives $\log(\vee/\Lambda)$, split according to Eq.~\eqref{eq:logsplit}.

The totally symmetrised four-$\rho$ part of their sum vanishes identically, i.e.\ the four-$\rho$ terms cancel for commuting charges. This cancellation has a simple origin. For commuting charges the part of $\Omega$ that is linear in $\rho$ exponentiates: to the order considered, $\Omega$ acting on the vacuum is the coherent state
\begin{align}\label{eq:coherent}
\Omega_{\rm coh}\ket0=\exp\Big\{\int\frac{d^3p}{(2\pi)^3}\Big[\mf A^b_j(p)\,a^{\dagger b}_j(p)-\mf A^{b\dagger}_j(p)\,a^b_j(p)\Big]\Big\}\ket0\,,
\end{align}
whose expansion to second order produces $\mf N_2$ and the $\{\rho,\rho\}$ part of $\mf B_2$, $\mf B_2\supset\frac12\mf A\mf A$. Since $\bar\Omega_{\rm coh}=S^\dagger\Omega_{\rm coh}S$ is again a displacement operator, $\bar\Omega^\dagger_{\rm coh}\Omega_{\rm coh}$ is, up to a phase, a displacement operator, and the expectation value of the number operator in the corresponding coherent state is the squared modulus of the displacement. The displacement is the LO amplitude~\eqref{eq:ALO}. Hence for commuting charges the spectrum is exactly the LO result, with no correction of order $g^4\rho^4$. Two features of the coefficients are necessary for this. First, the $\{\rho,\rho\}$ part of $\mf B_2$ must have support on the full square $\Lambda<k^+,p^+<\vee$, as follows from the kinematics of two independent emissions; with the support $k^++p^+<\vee$ the four-$\rho$ terms of the different rows would carry different logarithms and would not cancel. Second, the three-$\rho$ part of $\mf A^{(3)}$ must be ordered as $\mf A\mf N_2$ (Section~\ref{sec:coefficients}), so that the normalisation terms combine into Eq.~\eqref{eq:Nterms}. We have verified the cancellation by an explicit term-by-term evaluation, with the charges represented by $SU(2)$ and $SU(3)$ matrices and the order of all operators kept.

For non-commuting charges the four-$\rho$ terms do not vanish as operators: rewriting each product in terms of totally symmetrised products leaves remainders with commutators. For two and three operators
\begin{align}\label{eq:symmetrise}
\rho_1\rho_2&=\tfrac12\{\rho_1,\rho_2\}+\tfrac12[\rho_1,\rho_2]\,,\notag\\
\rho_1\rho_2\rho_3&=S(\rho_1\rho_2\rho_3)+\tfrac12\Big(S\big([\rho_1,\rho_2]\rho_3\big)+S\big([\rho_1,\rho_3]\rho_2\big)+S\big([\rho_2,\rho_3]\rho_1\big)\Big)\notag\\&\quad+\tfrac16\Big(\big[[\rho_1,\rho_2],\rho_3\big]-\big[[\rho_2,\rho_3],\rho_1\big]\Big)\,,
\end{align}
where $S$ denotes the totally symmetrised product, $S(XY)=\frac12\{X,Y\}$, and similarly for four operators. By Eq.~\eqref{eq:rho} each commutator reduces the number of charges by one. The four-$\rho$ terms therefore leave a three-$\rho$ remainder, multiplied by $\log(\vee/k^+)$ and by $\log(k^+/\Lambda)$, and a two-$\rho$ remainder, which is part of the BFKL and JIMWLK terms below.

\subsection{Terms cubic in \texorpdfstring{$\rho$}{rho}}\label{sec:threerho}

Logarithmic terms with three charges arise directly in Rows I.II--I.V, II.I, III.I--III.III, III.V and III.VI, from the parts of $\mf B_2$, $\mf B_3$ and $\mf A^{(3)}$ that are linear or quadratic in $\rho$ multiplied by single emissions, and in addition from the reordering of the four-$\rho$ terms. They are of the form $if^{abc}\times$(real kernel)$\times$(Wilson lines)$\times\rho\rho\rho$, they change sign together with the sign of the exponent $e^{-i\v k\cdot(\v w-\v w')}$, and they are of the type of the ``first saturation corrections'' in the projectile.

For each of the two logarithms, the totally symmetrised three-$\rho$ part of the sum of all these terms vanishes identically, for an arbitrary function of $\v w$ and $\v w'$ in place of the phase. We have organised the sum by transverse kernel and by the number of Wilson lines (two, three and four) and verified the cancellation in each class separately, analytically, at the operator level with explicit $SU(2)$ matrices for the charges on a transverse lattice, so that contact terms are included, and for the totally symmetrised colour tensors with random $SU(3)$ and $SU(4)$ Wilson lines (to $2\times10^{-15}$). Since the colour tensor itself vanishes, the cancellation holds for any symmetrised three-point function of the projectile, in particular for the $d^{abc}$ (odderon) one of Eq.~\eqref{eq:A2d}, whose projection we checked separately: no rapidity logarithm multiplies the three-point function. That no logarithm can multiply $\sigma$ follows from charge conjugation. C acts as $\rho^at^a\to-(\rho^at^a)^T$, i.e.\ $\rho^a\to-\eta_a\rho^a$ with $(t^a)^T=\eta_at^a$, on the soft gluons as $A^a\to-\eta_aA^a$ and on the target as $V\to V^*$, $U^{ab}\to\eta_a\eta_bU^{ab}$; since $\eta_a\eta_b\eta_cf^{abc}=-f^{abc}$ and $\eta_a\eta_b\eta_cd^{abc}=d^{abc}$, it preserves Eq.~\eqref{eq:rho} and all couplings, and the spectrum is invariant under the joint transformation. Under C, $n$ is even and $\sigma$ odd, and the Wilson-line traces with $T^e$ are even and those with $D^h$ odd (Section~\ref{sec:fin_singlet}). By the factorisation of Section~\ref{sec:proofsketch}, the coefficient of $\log(\vee/k^+)$ is the LO spectrum with the projectile correlator replaced by its one-step evolution, which does not involve the target: a three-charge term can enter it only through a $c$-number invariant tensor, $d^{abc}$ or $f^{abc}$, and C-invariance of the evolution admits only $f^{abc}$, whose contraction with Eq.~\eqref{eq:A2d} vanishes; the target enters through the C-even $\mathcal O_{\rm LO}$. The coefficient of $\log(k^+/\Lambda)$ is $-H_{\rm JIMWLK}$ acting on the target operator of the LO spectrum, in which the projectile enters only through $n$. Hence neither logarithm can multiply $\sigma$; a dependence on $\sigma$ requires a C-odd target structure and is non-logarithmic (Section~\ref{sec:fin_rows}). C alone does not forbid terms $\sigma\times$(C-odd target structure), which are exactly the symmetrised three-$\rho$ finite terms; the factorised form of the logarithms is needed in addition. The explicit cancellation confirms this argument and is stronger: individual terms have $d^{abc}$ projections of order one, with C-odd target factors, which C alone would allow, and they cancel in the sum; the cancellation of the whole symmetrised tensor also removes an $f^{abc}g_f$ part, which is C-even and not forbidden by C. As operators the three-$\rho$ terms leave, by Eq.~\eqref{eq:symmetrise}, two-$\rho$ and one-$\rho$ remainders. The two-$\rho$ remainders are part of the BFKL and JIMWLK terms. The one-$\rho$ remainders are contact terms $\propto f^{abc}\rho^c(\v x)\delta^{(2)}(\v x-\v y)$, which do not contribute in the colour-singlet form.

The non-logarithmic three-$\rho$ terms do not cancel. They belong to the finite terms (Section~\ref{sec:finite}).

\subsection{The BFKL term}\label{sec:BFKL}

Collecting all two-$\rho$ terms proportional to $\log(\vee/k^+)$, the genuine ones of the rows of Table~\ref{tab:rows} (I.II--I.IV, II.I, II.IV, III.I--III.III, III.V, III.VI) and the remainders of the reordered three- and four-$\rho$ terms, we find
\begin{align}\label{eq:BFKLrho}
\log\frac{\vee}{k^+}\,\frac{d^{3}N_{\rm BFKL}}{d^{3}k} & = \frac{1}{(2\pi)^{3}}\,\frac{g^{4}}{16\pi^{5}}\,\frac{1}{k^{+}}\,\log\Big[\frac{\vee}{k^{+}}\Big]\notag\\&\times\int_{x,x',z,w,w'}e^{-ik\cdot (w-w')}\Bigg[ \notag\\
& +\mathcal{K}(x-w)\cdot\mathcal{K}(x'-w')\,\frac{1}{(x-z)^{2}}\Big[\notag\\
& -N_{c}\,U^{ab}(w')\,U^{ac}(w)\,\rho^{b}(x')\,\rho^{c}(x)\notag\\
&\quad +N_{c}\,U^{ab}(w')\,U^{ac}(x)\,\rho^{b}(x')\,\rho^{c}(x)\notag\\
& +N_{c}\,U^{ab}(x')\,U^{ac}(w)\,\rho^{b}(x')\,\rho^{c}(x) \notag\\
&\quad -N_{c}\,U^{ab}(x')\,U^{ac}(x)\,\rho^{b}(x')\,\rho^{c}(x)\Big]\notag\\
& +\mathcal{K}(x-w)\cdot\mathcal{K}(x'-w')\,\frac{1}{(x'-z)^{2}}\Big[\notag\\
& -N_{c}\,U^{ab}(w')\,U^{ac}(w)\,\rho^{b}(x')\,\rho^{c}(x)\notag\\
&\quad +N_{c}\,U^{ab}(w')\,U^{ac}(x)\,\rho^{b}(x')\,\rho^{c}(x)\notag\\
& +N_{c}\,U^{ab}(x')\,U^{ac}(w)\,\rho^{b}(x')\,\rho^{c}(x) \notag\\
&\quad -N_{c}\,U^{ab}(x')\,U^{ac}(x)\,\rho^{b}(x')\,\rho^{c}(x)\Big]\notag\\
& +\mathcal{K}(x-w)\cdot\mathcal{K}(x'-w')\,\mathcal{K}(x-z)\cdot\mathcal{K}(x'-z)\Big[\notag\\
& +2\,f^{abc}\,f^{dbe}\,U^{ha}(w')\,U^{hd}(w)\,\rho^{c}(x')\,\rho^{e}(x)\notag\\
&\quad -2\,f^{abc}\,f^{dbe}\,U^{ha}(w')\,U^{hd}(x)\,\rho^{c}(x')\,\rho^{e}(x)\notag\\
& -2\,f^{abc}\,f^{dbe}\,U^{ha}(x')\,U^{hd}(w)\,\rho^{c}(x')\,\rho^{e}(x) \notag\\
&\quad +2\,f^{abc}\,f^{dbe}\,U^{ha}(x')\,U^{hd}(x)\,\rho^{c}(x')\,\rho^{e}(x)\Big]\notag\\
& +\mathcal{K}(w-z)\cdot\mathcal{K}(w'-z)\,\mathcal{K}(x-z)\cdot\mathcal{K}(x'-z)\Big[\notag\\
& +2\,N_{c}\,\rho^{a}(x)\,\rho^{a}(x')\notag\\
&\quad -\,f^{abc}\,f^{dce}\,U^{ha}(w')\,U^{hd}(w)\,\rho^{b}(x)\,\rho^{e}(x')\notag\\
& -\,f^{abc}\,f^{dce}\,U^{ha}(w')\,U^{hd}(w)\,\rho^{b}(x')\,\rho^{e}(x) \notag\\
& +\,f^{abc}\,f^{dce}\,U^{ha}(w')\,U^{hd}(z)\,\rho^{b}(x)\,\rho^{e}(x')\notag\\
&\quad +\,f^{abc}\,f^{dce}\,U^{ha}(w')\,U^{hd}(z)\,\rho^{b}(x')\,\rho^{e}(x)\notag\\
& +\,f^{abc}\,f^{dce}\,U^{ha}(z)\,U^{hd}(w)\,\rho^{b}(x)\,\rho^{e}(x') \notag\\
&\quad +\,f^{abc}\,f^{dce}\,U^{ha}(z)\,U^{hd}(w)\,\rho^{b}(x')\,\rho^{e}(x)\Big]\notag\\
& +\mathcal{K}(x'-w')\cdot\mathcal{K}(w-z)\,\frac{1}{(x-z)^{2}}\Big[\notag\\
& -N_{c}\,U^{ab}(w')\,U^{ac}(w)\,\rho^{b}(x')\,\rho^{c}(x) \notag\\
& +N_{c}\,U^{ab}(w')\,U^{ac}(z)\,\rho^{b}(x')\,\rho^{c}(x)\notag\\
&\quad +N_{c}\,U^{ab}(x')\,U^{ac}(w)\,\rho^{b}(x')\,\rho^{c}(x)\notag\\
& -N_{c}\,U^{ab}(x')\,U^{ac}(z)\,\rho^{b}(x')\,\rho^{c}(x)\Big]\notag\\
& +\mathcal{K}(x'-w')\cdot\mathcal{K}(w-z)\,\mathcal{K}(x-z)\cdot\mathcal{K}(x'-z)\Big[\notag\\
& +2\,f^{abc}\,f^{cde}\,U^{hd}(w')\,U^{ha}(w)\,\rho^{b}(x)\,\rho^{e}(x')\notag\\
&\quad -2\,f^{abc}\,f^{cde}\,U^{hd}(w')\,U^{ha}(z)\,\rho^{b}(x)\,\rho^{e}(x')\notag\\
& -2\,f^{abc}\,f^{cde}\,U^{hd}(x')\,U^{ha}(w)\,\rho^{b}(x)\,\rho^{e}(x') \notag\\
&\quad +2\,f^{abc}\,f^{cde}\,U^{hd}(x')\,U^{ha}(z)\,\rho^{b}(x)\,\rho^{e}(x')\Big]\notag\\
& +\mathcal{K}(x-w)\cdot\mathcal{K}(w'-z)\,\frac{1}{(x'-z)^{2}}\Big[\notag\\
& -N_{c}\,U^{ab}(w')\,U^{ac}(w)\,\rho^{c}(x)\,\rho^{b}(x')\notag\\
&\quad +N_{c}\,U^{ab}(w')\,U^{ac}(x)\,\rho^{c}(x)\,\rho^{b}(x')\notag\\
& +N_{c}\,U^{ab}(z)\,U^{ac}(w)\,\rho^{c}(x)\,\rho^{b}(x') \notag\\
&\quad -N_{c}\,U^{ab}(z)\,U^{ac}(x)\,\rho^{c}(x)\,\rho^{b}(x')\Big]\notag\\
& +\mathcal{K}(x-w)\cdot\mathcal{K}(w'-z)\,\mathcal{K}(x-z)\cdot\mathcal{K}(x'-z)\Big[\notag\\
& +2\,f^{abc}\,f^{cde}\,U^{ha}(w')\,U^{hd}(w)\,\rho^{b}(x')\,\rho^{e}(x)\notag\\&\quad -2\,f^{abc}\,f^{cde}\,U^{ha}(w')\,U^{hd}(x)\,\rho^{b}(x')\,\rho^{e}(x)\notag\\
& -2\,f^{abc}\,f^{cde}\,U^{ha}(z)\,U^{hd}(w)\,\rho^{b}(x')\,\rho^{e}(x) \notag\\
&\quad +2\,f^{abc}\,f^{cde}\,U^{ha}(z)\,U^{hd}(x)\,\rho^{b}(x')\,\rho^{e}(x)\Big]\Bigg]
\end{align}
where the average over the projectile is implied. With non-commuting charges the projectile evolution produces the symmetrised product $\frac12\{\rho^b(\v x'),\rho^c(\v x)\}$ in the last four blocks; the fixed order written there differs from it by contact terms $\propto f^{abc}\rho^c(\v x)\delta^{(2)}(\v x-\v x')$, which do not contribute in the colour-singlet form. The prefactor is $\frac{1}{(2\pi)^3}\frac{g^4}{16\pi^5}=\frac{\alpha_s^2}{8\pi^6}$.

For a colour-singlet projectile, Eq.~\eqref{eq:singlet}, Eq.~\eqref{eq:BFKLrho} becomes
\begin{align}\label{eq:BFKLsinglet}
\log\frac{\vee}{k^+}\,\frac{d^{3}N_{\rm BFKL}}{d^{3}k} & = \frac{1}{(2\pi)^{3}}\,\frac{g^{4}}{16\pi^{5}}\,\frac{1}{k^{+}}\,\frac{1}{N^{2}_{c}-1}\,\log\Big[\frac{\vee}{k^{+}}\Big]\,\int_{x,y,z,w,w'}e^{-ik\cdot (w-w')}\notag\\&\times\langle\rho^{f}(x)\rho^{f}(y) \rangle\,\Bigg[ \notag\\ & \mathcal K(x-w)\cdot\mathcal K(y-w')\,K(x,y,z)\Big[2 N_{c}\mathcal O_{\rm LO}(x,y;w,w')\Big]\notag\\
&+\mathcal K(x-w)\cdot\mathcal K(y-w')\,K(x,x,z)\Big[-N_{c}\mathcal O_{\rm LO}(x,y;w,w')\Big]\notag\\
&+\mathcal K(x-w)\cdot\mathcal K(y-w')\,K(y,y,z)\Big[-N_{c}\mathcal O_{\rm LO}(x,y;w,w')\Big]\notag\\
&+K(w,w',z)\,K(x,y,z)\Big[2 N_{c}\mathcal O_{\rm LO}(z,z;w,w')\Big]\notag\\
&+\mathcal K(x-w')\cdot\mathcal K(w-z)\,K(x,y,z)\Big[2 N_{c}\mathcal O_{\rm LO}(z,x;w,w')\Big]\notag\\
&+\mathcal K(x-w')\cdot\mathcal K(w-z)\,K(y,y,z)\Big[-N_{c}\mathcal O_{\rm LO}(z,x;w,w')\Big]\notag\\
&+\mathcal K(x-w)\cdot\mathcal K(w'-z)\,K(x,y,z)\Big[2 N_{c}\mathcal O_{\rm LO}(x,z;w,w')\Big]\notag\\
&+\mathcal K(x-w)\cdot\mathcal K(w'-z)\,K(y,y,z)\Big[-N_{c}\mathcal O_{\rm LO}(x,z;w,w')\Big]\Bigg]
\end{align}
with $K(\v x,\v y,\v z)$ of Eq.~\eqref{eq:kernels} and $\mathcal O_{\rm LO}$ of Eq.~\eqref{eq:OLO}. As for LO, the sign of the exponent is immaterial in this form.

Eq.~\eqref{eq:BFKLsinglet} is the LO spectrum~\eqref{eq:LOsinglet} with the charge correlator replaced by its one-step BFKL evolution over the rapidity interval $\log(\vee/k^+)$:
\begin{align}\label{eq:BFKLasLO}
\log\frac{\vee}{k^+}\,\frac{d^3N_{\rm BFKL}}{d^3k}=\log\frac{\vee}{k^+}\;\frac{d^3N_{\rm LO}}{d^3k}\bigg|_{n(\v x,\v y)\to\frac{d}{dY}n(\v x,\v y)}\,,
\end{align}
where the evolution of the projectile charge correlator in the dilute regime is~\cite{Kovner:2006wr,Kovner:2006ge}
\begin{align}\label{eq:BFKLrhorho}
\frac{d}{dY}\big\langle\rho^f(\v x)\rho^f(\v y)\big\rangle=-\frac{\alpha_sN_c}{2\pi^2}\int_z\Big[&K_D(\v x,\v y,\v z)\,\big\langle\rho^f(\v x)\rho^f(\v y)\big\rangle
\notag\\&-2\,\delta^{(2)}(\v x-\v y)\int_{z'}K(\v z,\v z',\v x)\,\big\langle\rho^f(\v z)\rho^f(\v z')\big\rangle
\notag\\&
+\big(2K(\v y,\v z,\v x)-K(\v z,\v z,\v x)\big)\big\langle\rho^f(\v y)\rho^f(\v z)\big\rangle
\notag\\&+\big(2K(\v x,\v z,\v y)-K(\v z,\v z,\v y)\big)\big\langle\rho^f(\v x)\rho^f(\v z)\big\rangle\Big]\,,
\end{align}
with the dipole kernel $K_D(\v x,\v y,\v z)=K(\v x,\v x,\v z)-2K(\v x,\v y,\v z)+K(\v y,\v y,\v z)$. The first term of Eq.~\eqref{eq:BFKLrhorho} gives the first three lines of Eq.~\eqref{eq:BFKLsinglet} (virtual terms, the produced gluon is emitted by the original charges); the second term, in which the charge correlator is concentrated at the position $\v z$ of the new gluon, gives the fourth line (the produced gluon is emitted by the new gluon); and the last two terms give the remaining four lines (the produced gluon is emitted by the new gluon in the amplitude and by an original charge in the conjugate amplitude, or vice versa). In terms of the colour-singlet correlator of the dipole picture, Eq.~\eqref{eq:BFKLrhorho} is the linearised BK (dipole BFKL) equation; the check for a single colour dipole is described in Section~\ref{sec:checks}.

\subsection{The JIMWLK term}\label{sec:JIMWLK}

The two-$\rho$ terms proportional to $\log(k^+/\Lambda)$, the genuine ones of Rows I.II, I.V, II.I, II.IV and III.III and the remainders of the reordered three- and four-$\rho$ terms, add up to
\begin{align}\label{eq:JIMWLKrho}
\log\frac{k^+}{\Lambda}\,\frac{d^{3}N_{\rm JIMWLK}}{d^{3}k} & = \frac{1}{(2\pi)^{3}}\,\frac{g^{4}}{16\pi^{5}}\,\frac{1}{k^{+}}\,\log{\left(\frac{k^{+}}{\Lambda}\right)}\,\int_{x,x',z,w,w'}e^{-ik\cdot(w-w')}\,\Bigg[
\notag \\ & \mathcal{K}(x-w)\cdot\mathcal{K}(x'-w')\,\frac{1}{(x-z)^{2}}\Big[\notag\\
&+2\,N_{c}\,U^{ab}(w')\,U^{ac}(x)\,\rho^{b}(x')\,\rho^{c}(x) -2\,N_{c}\,U^{ab}(x')\,U^{ac}(x)\,\rho^{b}(x')\,\rho^{c}(x)\notag\\
& +2\,f^{abc}\,f^{deh}\,U^{da}(z)\,U^{el}(w')\,U^{hb}(x)\,\rho^{l}(x')\,\rho^{c}(x) \notag\\
&\quad -2\,f^{abc}\,f^{deh}\,U^{da}(z)\,U^{el}(x')\,U^{hb}(x)\,\rho^{l}(x')\,\rho^{c}(x)\Big]\notag\\
& +\mathcal{K}(x-w)\cdot\mathcal{K}(x'-w')\,\frac{1}{(x'-z)^{2}}\Big[\notag\\
& +2\,N_{c}\,U^{ab}(x')\,U^{ac}(w)\,\rho^{b}(x')\,\rho^{c}(x)\notag\\
&\quad -2\,N_{c}\,U^{ab}(x')\,U^{ac}(x)\,\rho^{b}(x')\,\rho^{c}(x)\notag\\
& +2\,f^{abc}\,f^{deh}\,U^{da}(z)\,U^{el}(w)\,U^{hb}(x')\,\rho^{c}(x')\,\rho^{l}(x) \notag\\
&\quad -2\,f^{abc}\,f^{deh}\,U^{da}(z)\,U^{el}(x)\,U^{hb}(x')\,\rho^{c}(x')\,\rho^{l}(x)\Big]\notag\\
& +\mathcal{K}(x-w)\cdot\mathcal{K}(x'-w')\,\mathcal{K}(x-z)\cdot\mathcal{K}(x'-z)\Big[\notag\\
& +2\,N_{c}\,U^{ab}(x')\,U^{ac}(x)\,\rho^{b}(x')\,\rho^{c}(x)\notag\\
&\quad +2\,f^{abc}\,f^{ade}\,U^{hb}(x')\,U^{hd}(x)\,\rho^{c}(x')\,\rho^{e}(x)\notag\\
& -2\,f^{abc}\,f^{deh}\,U^{da}(z)\,U^{eb}(x)\,U^{hl}(x')\,\rho^{l}(x')\,\rho^{c}(x) \notag\\
&\quad -2\,f^{abc}\,f^{deh}\,U^{da}(z)\,U^{eb}(x')\,U^{hl}(x)\,\rho^{c}(x')\,\rho^{l}(x)\Big]\notag\\
& +\mathcal{K}(x-w)\cdot\mathcal{K}(x'-w')\,\frac{1}{(w-z)^{2}}\Big[\notag\\
& -2\,N_{c}\,U^{ab}(w')\,U^{ac}(w)\,\rho^{b}(x')\,\rho^{c}(x)\notag\\
&\quad +2\,N_{c}\,U^{ab}(x')\,U^{ac}(w)\,\rho^{b}(x')\,\rho^{c}(x)\notag\\
& -2\,f^{abc}\,f^{deh}\,U^{da}(z)\,U^{el}(w')\,U^{hb}(w)\,\rho^{l}(x')\,\rho^{c}(x) \notag\\
&\quad +2\,f^{abc}\,f^{deh}\,U^{da}(z)\,U^{el}(x')\,U^{hb}(w)\,\rho^{l}(x')\,\rho^{c}(x)\Big]\notag\\
& +\mathcal{K}(x-w)\cdot\mathcal{K}(x'-w')\,\frac{1}{(w'-z)^{2}}\Big[\notag\\
& -2\,N_{c}\,U^{ab}(w')\,U^{ac}(w)\,\rho^{b}(x')\,\rho^{c}(x)\notag\\
&\quad +2\,N_{c}\,U^{ab}(w')\,U^{ac}(x)\,\rho^{b}(x')\,\rho^{c}(x)\notag\\
& -2\,f^{abc}\,f^{deh}\,U^{da}(z)\,U^{el}(w)\,U^{hb}(w')\,\rho^{c}(x')\,\rho^{l}(x) \notag\\
&\quad +2\,f^{abc}\,f^{deh}\,U^{da}(z)\,U^{el}(x)\,U^{hb}(w')\,\rho^{c}(x')\,\rho^{l}(x)\Big]\notag\\
& +\mathcal{K}(x-w)\cdot\mathcal{K}(x'-w')\,\mathcal{K}(w-z)\cdot\mathcal{K}(w'-z)\Big[\notag\\
& +2\,N_{c}\,U^{ab}(w')\,U^{ac}(w)\,\rho^{b}(x')\,\rho^{c}(x)\notag\\
&\quad +2\,f^{abc}\,f^{ade}\,U^{hb}(w')\,U^{hd}(w)\,\rho^{c}(x')\,\rho^{e}(x)\notag\\
& -2\,f^{abc}\,f^{deh}\,U^{da}(z)\,U^{eb}(w)\,U^{hl}(w')\,\rho^{l}(x')\,\rho^{c}(x) \notag\\
&\quad -2\,f^{abc}\,f^{deh}\,U^{da}(z)\,U^{eb}(w')\,U^{hl}(w)\,\rho^{c}(x')\,\rho^{l}(x)\Big]\notag\\
& +\mathcal{K}(x-w)\cdot\mathcal{K}(x'-w')\,\mathcal{K}(x-z)\cdot\mathcal{K}(w'-z)\Big[\notag\\
& -2\,N_{c}\,U^{ab}(w')\,U^{ac}(x)\,\rho^{b}(x')\,\rho^{c}(x)\notag\\
& -2\,f^{abc}\,f^{ade}\,U^{hb}(w')\,U^{hd}(x)\,\rho^{c}(x')\,\rho^{e}(x) \notag\\
&\quad +2\,f^{abc}\,f^{deh}\,U^{da}(z)\,U^{eb}(x)\,U^{hl}(w')\,\rho^{l}(x')\,\rho^{c}(x)\notag\\
& +2\,f^{abc}\,f^{deh}\,U^{da}(z)\,U^{eb}(w')\,U^{hl}(x)\,\rho^{c}(x')\,\rho^{l}(x)\Big]\notag\\
& +\mathcal{K}(x-w)\cdot\mathcal{K}(x'-w')\,\mathcal{K}(x'-z)\cdot\mathcal{K}(w-z)\Big[\notag\\
& -2\,N_{c}\,U^{ab}(x')\,U^{ac}(w)\,\rho^{b}(x')\,\rho^{c}(x)\notag\\&\quad -2\,f^{abc}\,f^{ade}\,U^{hb}(x')\,U^{hd}(w)\,\rho^{c}(x')\,\rho^{e}(x)\notag\\
& +2\,f^{abc}\,f^{deh}\,U^{da}(z)\,U^{eb}(w)\,U^{hl}(x')\,\rho^{l}(x')\,\rho^{c}(x) \notag\\
&\quad +2\,f^{abc}\,f^{deh}\,U^{da}(z)\,U^{eb}(x')\,U^{hl}(w)\,\rho^{c}(x')\,\rho^{l}(x)\Big]\Bigg]
\end{align}
For a colour-singlet projectile this becomes
\begin{align}\label{eq:JIMWLKsinglet}
    \log\frac{k^+}{\Lambda}\,\frac{d^{3}N_{\rm JIMWLK}}{d^{3}k} & = \frac{1}{(2\pi)^{3}}\,\frac{g^{4}}{16\pi^{5}}\,\frac{1}{k^{+}}\,\frac{1}{N^{2}_{c}-1}\,\log{\left(\frac{k^{+}}{\Lambda}\right)}\,\int_{x,y,z,w,w'}e^{-ik\cdot(w-w')}\notag\\&\times\langle\rho^{f}(x)\rho^{f}(y) \rangle\,\Bigg[\notag \\ & \mathcal K(x-w)\cdot\mathcal K(y-w')\,K(w,w',z)\Big[4 N_{c}{\rm Tr}\{U^\dagger(w)U(w')\}\notag\\
&-2{\rm Tr}\{T^eU^\dagger(w)U(w')T^eU^\dagger(w')U(z)\}\notag\\
&\quad-2{\rm Tr}\{T^eU^\dagger(w)U(w')T^eU^\dagger(z)U(w)\}\Big]\notag\\
&+\mathcal K(x-w)\cdot\mathcal K(y-w')\,K(x,w',z)\Big[-4 N_{c}{\rm Tr}\{U^\dagger(x)U(w')\}\notag\\
&+2{\rm Tr}\{T^eU^\dagger(w')U(x)T^eU^\dagger(x)U(z)\}\notag\\
&\quad+2{\rm Tr}\{T^eU^\dagger(w')U(x)T^eU^\dagger(z)U(w')\}\Big]\notag\\
&+\mathcal K(x-w)\cdot\mathcal K(y-w')\,K(x,y,z)\Big[4 N_{c}{\rm Tr}\{U^\dagger(x)U(y)\}\notag\\
&-2{\rm Tr}\{T^eU^\dagger(x)U(y)T^eU^\dagger(y)U(z)\}\notag\\
&\quad-2{\rm Tr}\{T^eU^\dagger(x)U(y)T^eU^\dagger(z)U(x)\}\Big]\notag\\
&+\mathcal K(x-w)\cdot\mathcal K(y-w')\,K(y,w,z)\Big[-4 N_{c}{\rm Tr}\{U^\dagger(w)U(y)\}\notag\\
&+2{\rm Tr}\{T^eU^\dagger(w)U(y)T^eU^\dagger(y)U(z)\}\notag\\
&\quad+2{\rm Tr}\{T^eU^\dagger(w)U(y)T^eU^\dagger(z)U(w)\}\Big]\notag\\
&+\mathcal K(x-w)\cdot\mathcal K(y-w')\,K(w',w',z)\Big[\notag\\
&\quad-2 N_{c}{\rm Tr}\{U^\dagger(w)U(w')\}+2 N_{c}{\rm Tr}\{U^\dagger(x)U(w')\}\notag\\
&-{\rm Tr}\{T^eU^\dagger(w')U(x)T^eU^\dagger(w')U(z)\}\notag\\
&\quad-{\rm Tr}\{T^eU^\dagger(w')U(x)T^eU^\dagger(z)U(w')\}\notag\\
&+{\rm Tr}\{T^eU^\dagger(w)U(w')T^eU^\dagger(w')U(z)\}\notag\\
&\quad+{\rm Tr}\{T^eU^\dagger(w)U(w')T^eU^\dagger(z)U(w')\}\notag\\
&+{\rm Tr}\{T^eU^\dagger(w')U(x)\}\,{\rm Tr}\{T^eU^\dagger(w')U(z)\}\notag\\
&\quad+{\rm Tr}\{T^eU^\dagger(w)U(w')\}\,{\rm Tr}\{T^eU^\dagger(w')U(z)\}\Big]\notag\\
&+\mathcal K(x-w)\cdot\mathcal K(y-w')\,K(w,w,z)\Big[\notag\\
&\quad-2 N_{c}{\rm Tr}\{U^\dagger(w)U(w')\}+2 N_{c}{\rm Tr}\{U^\dagger(w)U(y)\}\notag\\
&+{\rm Tr}\{T^eU^\dagger(w)U(w')T^eU^\dagger(w)U(z)\}\notag\\
&\quad+{\rm Tr}\{T^eU^\dagger(w)U(w')T^eU^\dagger(z)U(w)\}\notag\\
&-{\rm Tr}\{T^eU^\dagger(w)U(y)T^eU^\dagger(w)U(z)\}\notag\\
&\quad-{\rm Tr}\{T^eU^\dagger(w)U(y)T^eU^\dagger(z)U(w)\}\notag\\
&-{\rm Tr}\{T^eU^\dagger(w)U(w')\}\,{\rm Tr}\{T^eU^\dagger(w)U(z)\}\notag\\
&\quad+{\rm Tr}\{T^eU^\dagger(w)U(y)\}\,{\rm Tr}\{T^eU^\dagger(w)U(z)\}\Big]\notag\\
&+\mathcal K(x-w)\cdot\mathcal K(y-w')\,K(x,x,z)\Big[\notag\\
&\quad2 N_{c}{\rm Tr}\{U^\dagger(x)U(w')\}-2 N_{c}{\rm Tr}\{U^\dagger(x)U(y)\}\notag\\
&-{\rm Tr}\{T^eU^\dagger(w')U(x)T^eU^\dagger(x)U(z)\}\notag\\
&\quad-{\rm Tr}\{T^eU^\dagger(w')U(x)T^eU^\dagger(z)U(x)\}\notag\\
&+{\rm Tr}\{T^eU^\dagger(x)U(y)T^eU^\dagger(x)U(z)\}\notag\\
&\quad+{\rm Tr}\{T^eU^\dagger(x)U(y)T^eU^\dagger(z)U(x)\}\notag\\
&-{\rm Tr}\{T^eU^\dagger(w')U(x)\}\,{\rm Tr}\{T^eU^\dagger(x)U(z)\}\notag\\
&\quad-{\rm Tr}\{T^eU^\dagger(x)U(y)\}\,{\rm Tr}\{T^eU^\dagger(x)U(z)\}\Big]\notag\\
&+\mathcal K(x-w)\cdot\mathcal K(y-w')\,K(y,y,z)\Big[\notag\\
&\quad2 N_{c}{\rm Tr}\{U^\dagger(w)U(y)\}-2 N_{c}{\rm Tr}\{U^\dagger(x)U(y)\}\notag\\
&-{\rm Tr}\{T^eU^\dagger(w)U(y)T^eU^\dagger(y)U(z)\}\notag\\
&\quad-{\rm Tr}\{T^eU^\dagger(w)U(y)T^eU^\dagger(z)U(y)\}\notag\\
&+{\rm Tr}\{T^eU^\dagger(x)U(y)T^eU^\dagger(y)U(z)\}\notag\\
&\quad+{\rm Tr}\{T^eU^\dagger(x)U(y)T^eU^\dagger(z)U(y)\}\notag\\
&-{\rm Tr}\{T^eU^\dagger(w)U(y)\}\,{\rm Tr}\{T^eU^\dagger(y)U(z)\}\notag\\
&\quad+{\rm Tr}\{T^eU^\dagger(x)U(y)\}\,{\rm Tr}\{T^eU^\dagger(y)U(z)\}\Big]\Bigg]
\end{align} 
where the products of traces with $T^e$ are the adjoint rewriting of the $ffUUU$ contractions of Eq.~\eqref{eq:JIMWLKrho}.

Eq.~\eqref{eq:JIMWLKrho} is the JIMWLK evolution of the target operator of the LO spectrum over the rapidity interval $\log(k^+/\Lambda)$. With the JIMWLK Hamiltonian in the normalisation of~\cite{Lublinsky:2016meo},
\begin{align}\label{eq:HJIMWLK}
H_{\rm JIMWLK}=\frac{\alpha_s}{2\pi^2}\int_{x,y,z}K(\v x,\v y,\v z)\Big[J_L^e(\v x)J_L^e(\v y)+J_R^e(\v x)J_R^e(\v y)-2\,c^{ef}(\v z)\,J_L^e(\v x)J_R^f(\v y)\Big]\,,
\end{align}
where $J_L^e(\v x)$ and $J_R^e(\v x)$ act on a Wilson line at $\v x$ as $J_L^eU=T^eU$ and $J_R^eU=UT^e$, $c^{ef}(\v z)$ is defined by $U^\dagger(\v z)T^eU(\v z)=c^{ef}(\v z)T^f$, and $\frac{d}{dY}=-H_{\rm JIMWLK}$, we find
\begin{align}\label{eq:JIMWLKasLO}
\log\frac{k^+}{\Lambda}\,\frac{d^3N_{\rm JIMWLK}}{d^3k}=-\log\frac{k^+}{\Lambda}\;H_{\rm JIMWLK}\,\frac{d^3N_{\rm LO}}{d^3k}\,,
\end{align}
where $H_{\rm JIMWLK}$ acts on the Wilson lines of Eq.~\eqref{eq:LO} only. The prefactors match since $\frac{g^4}{16\pi^5}\big/\frac{g^2}{2\pi^2}=\frac{\alpha_s}{2\pi^2}$. Because $\mathcal O_{\rm LO}$ is bilinear in $U$, $H_{\rm JIMWLK}$ acting on it can be evaluated exactly, for an arbitrary (not necessarily colour-singlet or symmetric) ordered correlator $\langle\rho^{b'}(\v x')\rho^b(\v x)\rangle$; Eq.~\eqref{eq:JIMWLKrho} agrees with the result pointwise in all transverse coordinates, and Eq.~\eqref{eq:JIMWLKsinglet} with its colour-singlet projection (Section~\ref{sec:checks}). For a single adjoint dipole the action of $H_{\rm JIMWLK}$ reduces to the BK form, $-H_{\rm JIMWLK}{\rm Tr}[U^\dagger(\v a)U(\v b)]=\frac{\alpha_s}{\pi^2}\int_zK_D(\v a,\v b,\v z)\big[c^{ef}(\v z){\rm Tr}(U^\dagger(\v a)T^eU(\v b)T^f)-N_c{\rm Tr}(U^\dagger(\v a)U(\v b))\big]$.

\subsection{Why the logarithms are the evolution of the LO spectrum}\label{sec:proofsketch}

The identifications~\eqref{eq:BFKLasLO} and~\eqref{eq:JIMWLKasLO} are the expansion to first order of the leading-logarithmic factorisation of the single inclusive spectrum~\cite{Kovner:2006wr,Kovner:2006ge,Gelis:2008rw}: the spectrum of a gluon at rapidity $y$ is the LO expression evaluated with the projectile evolved from its own rapidity $Y_P$ down to $y$ and with the target evolved from $y$ to its own rapidity,
\begin{align}\label{eq:LLfact}
\frac{d^3N}{d^3k}\bigg|_{\rm LL}=\Big\langle\,\frac{d^3N_{\rm LO}}{d^3k}\,\Big\rangle_{P,\,Y_P-y;\;T,\,y-Y_T}\,.
\end{align}
The structure of the calculation shows why this holds, and we sketch the argument. Consider first the region $p^+\gg k^+$. Here the second gluon is much faster than the observed one, and its emission and absorption amplitudes factorise from those of the observed gluon: in all coefficients the dependence on $k^+$ disappears at leading power in $k^+/p^+$, and the observed gluon is emitted by the total colour charge of the valence charges and of the fast gluon, $\rho^a(\v x)\to\rho^a(\v x)+$ (charge density of the gluon at $\v z$). The sum over the rows is then the LO spectrum evaluated in the projectile state that contains one more soft gluon, minus the probability that it is not there (the virtual terms with $\mf N_2$ and the self-energy parts of $\mf A^{(3)}$), i.e.\ the LO spectrum with the projectile correlator evolved by one step. For a dilute projectile this step is linear in $\langle\rho\rho\rangle$, Eq.~\eqref{eq:BFKLrhorho}; terms with three and four charges, which would describe the non-linear (KLWMIJ) evolution of a dense projectile, must and do cancel at this order in the symmetrised sense (Sections~\ref{sec:fourrho},~\ref{sec:threerho}). In the region $p^+\ll k^+$ the second gluon is much slower than the observed one. It is emitted coherently by the valence charges and by the observed gluon, and at leading power its emission does not change the momentum of the observed gluon. The rows then combine into the change of the scattering amplitudes of the valence charges and of the observed gluon due to one soft gluon that scatters together with them, i.e.\ into one step of the JIMWLK evolution of the Wilson lines of the LO amplitude and of its conjugate, including the virtual terms. Both statements are verified by the explicit results~\eqref{eq:BFKLrho}--\eqref{eq:JIMWLKsinglet}.

Two features of the calculation enter this argument in an essential way. First, the logarithms are split at the rapidity of the observed gluon, Eq.~\eqref{eq:logsplit}; the four-$\rho$ terms and the normalisation, which carry $\log(\vee/\Lambda)$, are distributed between the two evolutions accordingly. Second, the operator ordering of the charges matters: it is the commutator terms of the reordering that produce part of the two-$\rho$ structures of both evolutions.

\subsection{Summary of the logarithmic terms}\label{sec:logsummary}

The rapidity-enhanced part of the $O(g^4)$ cross section is
\begin{align}\label{eq:logsummary}
\frac{d^3N}{d^3k}\bigg|_{\log}=\log\frac{\vee}{k^+}\,\frac{d^3N_{\rm BFKL}}{d^3k}+\log\frac{k^+}{\Lambda}\,\frac{d^3N_{\rm JIMWLK}}{d^3k}\,,
\end{align}
with the BFKL term~\eqref{eq:BFKLrho} and the JIMWLK term~\eqref{eq:JIMWLKrho}, and no term with three or four charges multiplies a rapidity logarithm. Individual rows contain also double logarithms, $\log^2(k^+/\Lambda)$, from the region in which the second gluon is both soft and far from the observed one in the transverse plane. They come from the recoil of the observed gluon in the rows with $\mf B_3$ and $\mf C_1$ and cancel in the sum over rows (Section~\ref{sec:fin_soft}).
% Source of the ln^2 cancellation: limits_checked/Cancellation_of_flagged_terms.tex (numerical, SU(2), SU(3)), presented in sec_finite.

% ======================================================================
\section{Collinear (DGLAP) and running-coupling terms}\label{sec:dglap_rc}
% ======================================================================
% Sources: paper/notes/fix_poles_dglap.tex (P3, P4, P5, Row II.IV tail), notes/lit_dglap_rc.tex,
% limits_checked/Finite_terms_final.tex (Group II Row IV), limits_checked/Divergent_terms_checked.tex.

In this section we discuss the terms that arise when the integrated gluon of Row~II.IV is collinear with the observed one, and we give the account of all poles $1/\epsilon$ of the calculation. Although, following the convention of the rest of the paper, no pole is written in the final result, the poles determine what the collinear and the running-coupling terms are. We use the notation
\begin{align}\label{eq:dg_not}
\mathcal O(y,y')&\equiv\Kv(x-y)\cdot\Kv(x'-y')\,\big[U(y')-U(x')\big]^{ce'}\big[U(y)-U(x)\big]^{ce}\rho^{e'}(x')\,\rho^{e}(x)\,,\notag\\
\widehat L(q)&\equiv\int_{x,y,x',y'}e^{-iq\cdot(y-y')}\,\mathcal O(y,y')\,,\notag\\
\frac{d^3N_{\rm LO}}{d^3k}&=\mathcal N_0\,\widehat L(k)\,,\qquad
\mathcal N_0\equiv\frac{1}{(2\pi)^3}\frac{g^2}{2\pi^2}\frac{1}{k^+}\,,\notag\\
\mathcal C&\equiv\frac{\alpha_s}{2\pi}\,\mathcal N_0=\frac{1}{(2\pi)^3}\frac{g^4}{16\pi^4}\frac{1}{k^+}\,,
\end{align}
so that Eq.~\eqref{eq:LO} is $\mathcal N_0\widehat L(k)$, and
\begin{align}\label{eq:dg_B}
\mathcal B(r)\equiv2\gamma_E+\log\frac{r^2\mu^2}{4}=\log\frac{r^2\mu^2}{c_0^2}\,,\qquad c_0=2e^{-\gamma_E}\,,\qquad
\mathfrak a\equiv\frac{\alpha_sN_c}{\pi}\,,\qquad \frac{\alpha_s\beta_0}{4\pi}=\frac{11}{12}\,\mathfrak a\,,
\end{align}
with $\mu\equiv\mu_{\MS}$, $\beta_0=11N_c/3$, $L\equiv\log(k^+/\Lambda)$ and $\ell_\vee\equiv\log(1-k^+/\vee)$. The phase is $e^{-i\v k\cdot(\text{unprimed}-\text{primed})}$ throughout (Section~\ref{sec:fin_conv}); in the collinear terms it is $e^{-i\v k\cdot(\v y-\v y')/\xi}$.

\subsection{Origin: final-state splitting of the LO gluon}\label{sec:dg_origin}

Row~II.IV is the square of the amplitude $\mathcal A_{CA}$ of Eq.~\eqref{eq:A2gamp}: a gluon of momentum $p^++k^+$ produced at LO splits after the collision, through $\bar{\mf C}_1$, into the observed gluon (at $\v w$ in the amplitude, $\v w'$ in the conjugate amplitude) and a second gluon at $\v z$, which is integrated over. With $\xi\equiv k^+/(p^++k^+)$ the momentum fraction of the observed gluon in its parent, the parent sits at $\v y=\v z+\xi(\v w-\v z)$ and $\v y'=\v z+\xi(\v w'-\v z)$; the Wilson line $U(\v z)$ cancels by unitarity, and the parent carries the Wilson-line structure of the LO amplitude. The integral over the daughter separation $\v r=\v w-\v z$ at fixed $\v u=\v w'-\v w$ diverges only at large $|\v r|$, i.e.\ at small relative transverse momentum of the two daughters: it is a collinear (infrared) divergence. Its coefficient is the real splitting function
\begin{align}\label{eq:Pggplus}
\mathcal P_{gg+}(\xi)=2N_c\Big[\xi(1-\xi)+\frac{\xi}{1-\xi}+\frac{1-\xi}{\xi}\Big]\,,
\end{align}
without a $\delta(1-\xi)$ term, as appropriate for a real emission. The two end points are the rapidity logarithms of the row: $\xi\to1$ (the integrated gluon soft) gives its JIMWLK term, $\xi\to0$ (the observed gluon soft with respect to its parent) its BFKL term (Section~\ref{sec:rapidity}). It is convenient to write the LO structure of a parent gluon of transverse momentum $\v k/\xi$ as $\Pi(\xi)\equiv\xi^{-2}\widehat L(\v k/\xi)$, and to define
\begin{align}\label{eq:dg_X}
\mathcal X\equiv\int_{x,x',z}\mathcal O(z,z)=\int_{x,x',z}&\Kv(x-z)\cdot\Kv(x'-z)\notag\\&\times\big[U(z)-U(x')\big]^{ce'}\big[U(z)-U(x)\big]^{ce}\rho^{e'}(x')\,\rho^e(x)\,,
\end{align}
the structure that multiplies the BFKL logarithm of this row, $+2N_c\,\mathcal C\,\frac{4\pi}{\v k^2}\log\frac{\vee}{k^+}\,\mathcal X$.

\subsection{Account of the poles}\label{sec:dg_poles}

We regularise the transverse loop and collinear integrals in $d=2-\epsilon$ dimensions. A scaleless integral is $\int d^dp/\v p^2=\pi\big(\frac{2}{\epsilon_{\rm UV}}-\frac{2}{\epsilon_{\rm IR}}\big)$; with a single $\epsilon$, $\epsilon_{\rm UV}=\epsilon_{\rm IR}$, it vanishes, which is the convention used elsewhere in the paper. We give both bookkeepings; they lead to the same finite result. Four sources of poles have to be considered.

\paragraph{Scheme.} We use conventional dimensional regularisation (CDR): transverse momenta and gluon polarisations are continued to $d_\perp=2-\epsilon$ dimensions wherever a polarisation sum multiplies a pole, i.e.\ in Rows~I.II, I.V and II.IV; in the finite terms, which are locally integrable, $d_\perp=2$. The coupling is the $\MS$ coupling and the fragmentation function the $\MS$ number density (pure-pole counterterms with the four-dimensional $\mathcal P_{gg}$, Section~\ref{sec:dg_msbar}), both in this scheme; Eq.~\eqref{eq:A3onerho} and the constant $\frac{67}{9}-\frac{\pi^2}{3}$ refer to it. In the dilute limit the two-gluon amplitude of the real rows reproduces, for any number $d_\perp$ of transverse polarisations, the $(D-2)$ dependence of the CDR two-gluon production cross section of~\cite{Fadin:1998py,Bartels:2006hg} (Section~\ref{sec:check_dilute}). With two-dimensional polarisations and $d$-dimensional momenta (a dimensional-reduction-like scheme) the reduction $\tilde p^l\tilde p^i\to\delta^{li}\tilde p^2/d$ gives a factor $2/d$ instead of $d_\perp/d=1$: the constant of Row~I.II would be $\frac{64}{9}-\frac{\pi^2}{3}$ (a shift $-\frac{\alpha_sN_c}{6\pi}\times$LO), that of Row~I.V would change by the opposite amount, as it must by Eq.~\eqref{eq:dg_SEid} below, Row~II.IV by $-\mathcal C\int_0^1d\xi\,2N_c\,\xi(1-\xi)\,\Pi(\xi)$, and the coupling and the fragmentation function would have to be converted, $\alpha_s^{\rm DR}=\alpha_s^{\MS}\big(1+\frac{\alpha_sN_c}{12\pi}\big)$ and $D^{\rm DR}=D^{\MS}\otimes\big[\delta(1-\xi)+\frac{\alpha_s}{2\pi}\big(2N_c\xi(1-\xi)-\frac{N_c}{6}\delta(1-\xi)\big)\big]$; the coefficient $\frac{N_c}{6}$ of the $\delta$ term is fixed independently by the momentum sum rule and equals the coupling conversion. With these conversions the result is unchanged; the shift of order $N_c/6$ cancels in the virtual part and appears only in the conversion of $\alpha_s$ and $D$. (The sign of the Row~II.IV term was obtained from the general $d_\perp/d$ rule and from this closure, not by re-deriving the row.)

\paragraph{The vertex correction (Row I.II).} The one-$\rho$ part of $\mf A^{(3)}$, Eq.~\eqref{eq:A3onerho}, carries $\big(\frac{11}{3}-2L+\ell_\vee\big)\big[-\frac{2}{\epsilon}+\log\frac{\v k^2}{\mu^2}\big]$ times $\mf A^{(1)}$; in coordinate space $\log\frac{\v k^2}{\mu^2}\frac{k^i}{\v k^2}\to-\mathcal B(\v x-\v w)\frac{(x-w)^i}{(\v x-\v w)^2}$, Eq.~\eqref{eq:FT_log}. In the cross section, with its complex conjugate, the $\frac{11}{3}$ part gives
\begin{align}\label{eq:dg_II}
\frac{d^3N^{\rm I.II}_{11/3}}{d^3k}=\frac{\alpha_s\beta_0}{4\pi}\,\mathcal N_0\int_{x,w,x',w'}e^{-ik\cdot(w-w')}\Big[\frac{4}{\epsilon}+\mathcal B(x-w)+\mathcal B(x'-w')\Big]\,\mathcal O(w,w')\,.
\end{align}
This is the virtual term $\frac{\alpha_s\beta_0}{4\pi}\big[\frac4\epsilon+L(\v x-\v y)+L(\v x'-\v y')\big]\times$LO expected from the self-energy of the emitted gluon. The pole comes from the region of large loop momentum of the $gg$ loop, $\int d^2\tilde p\,\tilde p^2/(\tilde p^2+\xi(1-\xi)\v k^2)$, which is finite at small $\tilde p$: it is ultraviolet.

\paragraph{The recombination term (Row I.V).} The part of $\bar{\mf C}_1^\dagger\mf B_2$ in which the two gluons of $\mf B_2$ are at one point (``$N_{2b\,\rm UV}$'') has the longitudinal coefficient $\frac{1}{k^+}\int_\Lambda^{k^+-\Lambda}dp^+[\dots]=2L-\frac{11}{6}$ and the bracket $\frac2\epsilon+\mathcal B(\v y-\v x)$, the pole coming from $\v z\to\v x$, again ultraviolet (the $\v z$ integral has no large-distance divergence). After the relabelling $\v x\to\v w$, $\v y\to\v x$, which turns the phase into $e^{-i\v k\cdot(\v w-\v w')}$ and leaves the order of the charges unchanged, and with $\frac{g^4N_c}{16\pi^4}\frac{11}{6}=\frac{g^2}{2\pi^2}\frac{\alpha_s\beta_0}{4\pi}$, its $-\frac{11}{6}$ part is
\begin{align}\label{eq:dg_V}
\frac{d^3N^{\rm I.V}_{2b\,\rm UV}}{d^3k}\bigg|_{-\frac{11}{6}}&=-\frac{\alpha_s\beta_0}{4\pi}\,\mathcal N_0\int_{x,w,x',w'}e^{-ik\cdot(w-w')}\Big[\frac{4}{\epsilon}+\mathcal B(x-w)+\mathcal B(x'-w')\Big]\,\mathcal O(w,w')\notag\\&=-\frac{d^3N^{\rm I.II}_{11/3}}{d^3k}\,,
\end{align}
pole, $2\gamma_E$ and logarithm, provided the same $\mu$ is used. The $-\frac{11}{6}$ of $N_{2b\,\rm UV}$ is the dimensionally regularised $\v z$ integral of a subtraction term; the same term, unintegrated and with the opposite sign, remains in the finite part of the row (in $N_{2b\,\rm rest}$), where it is the term with $\frac{11}{6}$ of Eq.~\eqref{eq:S-I.V-2o}, and makes the point $\v z\to\v x$ integrable; its $\v z$ integral alone diverges. The cancellation against Row~I.II concerns the integrated subtraction only. The cancellation of Eqs.~\eqref{eq:dg_II} and~\eqref{eq:dg_V} is a cancellation of an ultraviolet pole against an ultraviolet pole. Its physical origin is that a point-like $gg$ fluctuation of the produced gluon does not resolve the shock wave, $f^{adc}f^{fbe}U^{db}(\v x)U^{ce}(\v z)\to N_cU^{af}(\v x)$ at $\v z=\v x$. The observable counts dressed gluons, $\Omega^\dagger S^\dagger\Omega\,a^\dagger a\,\Omega^\dagger S\Omega$, and where $S$ is locally constant $\Omega^\dagger S\Omega=S$: the ultraviolet part of the self-energy before the shock wave (Row~I.II) is undone by the same loop straddling it (Row~I.V). 

\paragraph{Result and grouping.} The cancellation of the poles and of the dependence on $\mu$ between Rows~I.II and~I.V is a result, for any grouping of the finite terms. Which finite terms accompany it depends on the grouping. Let $a$ be the light-cone energy of the emitted gluon, $D$ the relative energy of the gluon pair into which it fluctuates and $b=a+D$ the energy of the pair. The self-energy part of $\mf A^{(3)}$ in Row~I.II has the energy denominators $1/(ab)$; the part of Row~I.V in which the pair produced by the triple-gluon vertex (together with the instantaneous term, which cancels the longitudinal part of the polarisation sum) is recombined by $\bar{\mf C}_1^\dagger$ without resolving the target has $1/(bD)$, with the same vertices and the same longitudinal supports $\Lambda<p^+<k^+-\Lambda$; and
\begin{align}\label{eq:dg_SEid}
\frac{1}{ab}+\frac{1}{bD}=\frac{1}{aD}\,,
\end{align}
whose transverse integral is scaleless and vanishes in dimensional regularisation, for each momentum fraction $\xi$, in $d$ dimensions. Hence, if the colour of the pair in Row~I.V is evaluated at its transverse centre of mass $\xi\v x+(1-\xi)\v z$, the self-energy part of Row~I.II, with its pole, its logarithm $\mathcal B$, its soft end points and its constant $\frac{67}{9}-\frac{\pi^2}{3}$, is cancelled identically by the unresolved part of Row~I.V, and the $\beta_0$-type coefficient survives only in the region in which the pair resolves the target. (The commutator parts of $\mf A^{(3)}$ and of $\mf B_2$ are not covered by Eq.~\eqref{eq:dg_SEid}.) In the form used in this paper the colour of the pair is evaluated at the position $\v x$ of one of its gluons; then no term $\frac{\alpha_s\beta_0}{4\pi}[\mathcal B(\v x-\v w)+\mathcal B(\v x'-\v w')]\times$LO remains explicitly, and the region in which the pair resolves the target is in the finite terms of Row~I.V ($N_{2a}$, $N_{2b\,\rm rest}$), which carry the same coefficient $2L-\frac{11}{6}$ multiplying the small-pair kernel. For an emitter distance $|\v x-\v w|$ larger than the inverse saturation scale of the target they contain, by power counting of the kernel, a logarithm of $|\v x-\v w|Q_s$ with this coefficient (not evaluated with a target model). The statement that the emission coupling does not run with the distance $\v x-\v w$ between the emitter and the produced gluon is therefore a statement about this grouping, valid as long as the pair does not resolve the target within this distance. In NLO JIMWLK, where the soft gluon is not measured, the $\beta_0$ terms with the emitter distance survive~\cite{Kovner:2023vsy}. The identity~\eqref{eq:dg_SEid} was confirmed independently in momentum space, in the dilute limit with a point-like projectile, where the unresolved part of Row~I.V cancels the whole self-energy part of Row~I.II $\xi$ by $\xi$, to $10^{-13}$ (Section~\ref{sec:check_dilute}).

\paragraph{The normalisation of the dressed observed gluon.} The single-gluon normalisation contains the scaleless integral $\braket{\Psi_{gg}|\Psi_{gg}}=\frac{g^2N_c}{8\pi^3}\big(2L-\frac{11}{6}\big)\int\frac{d^2\tilde p}{\tilde p^2}$, which is set to zero by the convention above. In the cross section it enters as
\begin{align}\label{eq:dg_norm}
\frac{d^3N_{\rm norm}}{d^3k}=-\frac{g^2N_c}{8\pi^3}\Big(2L-\frac{11}{6}\Big)\int\frac{d^d\tilde p}{\tilde p^2}\;\frac{d^3N_{\rm LO}}{d^3k}
=\Big[\frac{\alpha_s\beta_0}{4\pi}-\mathfrak aL\Big]\Big(\frac{2}{\epsilon_{\rm UV}}-\frac{2}{\epsilon_{\rm IR}}\Big)\frac{d^3N_{\rm LO}}{d^3k}\,.
\end{align}
Its ultraviolet half, $\frac{\alpha_sN_c}{4\pi}\big(\frac{11}{3}-4L\big)\frac{2}{\epsilon_{\rm UV}}$, is the light-cone-gauge $Z_3-1$ of the produced gluon; its infrared half ($\tilde p\to0$, collinear pair) is the virtual partner of the real splitting of Row~II.IV.

\paragraph{The real splitting (Row II.IV).} After the extraction of the end points, the collinear poles of Row~II.IV are: in its JIMWLK term $-(L+\ell_\vee)\,\mathfrak a\,\frac2\epsilon\,\frac{d^3N_{\rm LO}}{d^3k}$; in its DGLAP term
\begin{align}\label{eq:dg_IVpole}
-\mathcal C\,\frac2\epsilon\int_0^1d\xi\,\mathcal P_{gg+}(\xi)\,\Pi(\xi)=-\mathcal C\,\frac2\epsilon\,\big[\mathcal P_{gg}\otimes\Pi\big]+\frac{\alpha_s\beta_0}{4\pi}\,\frac2\epsilon\,\frac{d^3N_{\rm LO}}{d^3k}\,,
\end{align}
where $\mathcal P_{gg}=\mathcal P_{gg+}+\frac{\beta_0}{2}\delta(1-\xi)$ is the full $\MS$ splitting function at $N_f=0$, with the plus prescriptions of Eq.~\eqref{eq:plusxi}, and $\mathcal C\frac{\beta_0}{2}\Pi(1)=\frac{\alpha_s\beta_0}{4\pi}\mathcal N_0\widehat L(k)$. The subtraction at $\xi\to0$ has no pole, since the coefficient of the pole vanishes at $\xi=0$ for $\v k\neq0$. The second term of Eq.~\eqref{eq:dg_IVpole} is the pole of what we call the running-coupling term below: it is collinear, not ultraviolet.

\paragraph{Counterterms.} The coupling is renormalised in the $\MS$ scheme, $g_0^2=\mu^\epsilon g^2(\mu)\big[1-\frac{\alpha_s\beta_0}{4\pi}\frac2\epsilon\big]$; in the LO spectrum this gives $-\frac{\alpha_s\beta_0}{4\pi}\frac{2}{\epsilon_{\rm UV}}\times$LO. The charges $\rho$ and the Wilson lines $U$ are not renormalised: the source is classical and $U$ is the given target field. The collinear pole of the observed gluon is absorbed into an $\MS$ fragmentation function with the full $\mathcal P_{gg}$, which gives the counterterm $+\mathcal C\frac{2}{\epsilon_{\rm IR}}[\mathcal P_{gg}\otimes\Pi]$, with the $d$-dimensional Jacobian of Section~\ref{sec:dg_msbar}.

\begin{table}[t]
\centering
\renewcommand{\arraystretch}{1.25}
\footnotesize
\begin{tabular}{@{}p{3.1cm}p{2.6cm}p{0.6cm}p{2.4cm}p{4.0cm}@{}}
\toprule
pole & region & type & coefficient & cancelled by / absorbed into\\
\midrule
Row I.II, $\mf A^{(3)}$ & $gg$ loop, $\tilde p\to\infty$ & UV & $\frac{11}{6}-L+\frac{\ell_\vee}{2}$ & $\frac{11}{6}$: Row I.V; $-L$: JIMWLK sector\\
Row I.V, $N_{2b\,\rm UV}$ & pair at one point, $\v z\to\v x$ & UV & $-\frac{11}{6}+2L$ & $-\frac{11}{6}$: Row I.II; $2L$: JIMWLK sector\\
normalisation, UV half & $\tilde p\to\infty$ & UV & $\frac{11}{12}-L$ & $\frac{11}{12}$: coupling counterterm; $-L$: JIMWLK sector\\
coupling counterterm & $g_0\to g(\mu)$ & UV & $-\frac{11}{12}$ & UV half of the normalisation\\
\midrule
normalisation, IR half & collinear pair, $\tilde p\to0$ & IR & $-\frac{11}{12}+L$ & $-\frac{11}{12}$: $\delta(1-\xi)$ pole of Row II.IV; $L$: Row II.IV, $\xi\to1$\\
Row II.IV, $\xi\to1$ (JIMWLK) & large $\v r$ & IR & $-L-\ell_\vee$ & $-L$: normalisation, IR half; $-\ell_\vee$: left\\
Row II.IV, real $\mathcal P_{gg+}$ & large $\v r$ & IR & $-[\mathcal P_{gg}\otimes\Pi]+\frac{11}{12}$ & $-[\mathcal P_{gg}\otimes\Pi]$: FF counterterm; $\frac{11}{12}$: normalisation, IR half\\
Row II.IV, $\xi\to0$ (BFKL) & --- & --- & $0$ & ---\\
FF counterterm, full $\mathcal P_{gg}$ & $D_g\to D_g(\mu)$ & IR & $+[\mathcal P_{gg}\otimes\Pi]$ & real $\mathcal P_{gg+}$ and virtual $\frac{\beta_0}{2}\delta(1-\xi)$\\
\midrule
sum & & & UV: $\frac{\ell_\vee}{2}$;\ IR: $-\ell_\vee$ & $\to0$ for $\vee\to\infty$\\
\bottomrule
\end{tabular}
\caption{The poles of the calculation. Coefficients in units of $\mathfrak a\,\frac2\epsilon\,\frac{d^3N_{\rm LO}}{d^3k}$, and of $\mathcal C\frac2\epsilon[\,\cdot\otimes\Pi]$ for the collinear (DGLAP) entries; complex conjugates included for Group~I; $\mathfrak a=\alpha_sN_c/\pi$, so that $\frac{11}{12}\mathfrak a=\frac{\alpha_s\beta_0}{4\pi}$. The upper block contains the ultraviolet poles, the lower block the infrared (collinear) ones.}
\label{tab:poles}
\end{table}

\paragraph{The balance.} The poles are collected in Table~\ref{tab:poles}; the coefficients were obtained in exact rational arithmetic. In the single-$\epsilon$ bookkeeping, in which the normalisation~\eqref{eq:dg_norm} vanishes, the sum is Row~I.II $+$ Row~I.V $+$ Row~II.IV $+$ coupling counterterm $+$ fragmentation counterterm $=\big(\frac{11}{6}-L+\frac{\ell_\vee}{2}\big)+\big(-\frac{11}{6}+2L\big)+\big(-L-\ell_\vee+\frac{11}{12}-[\mathcal P_{gg}\otimes\Pi]\big)-\frac{11}{12}+[\mathcal P_{gg}\otimes\Pi]=-\frac{\ell_\vee}{2}$, the same total as in the two-$\epsilon$ bookkeeping. The poles proportional to $L$ cancel within the JIMWLK sector in both bookkeepings. The remainder $-\frac{\ell_\vee}{2}\frac2\epsilon$ is a pole multiplied by $\ell_\vee=\log(1-k^+/\vee)=-k^+/\vee+O(k^{+2}/\vee^2)$, a power correction in the ratio of the soft and the valence longitudinal momenta. It is not cancelled, and we do not claim that it is: it is outside the eikonal approximation. $\ell_\vee=-\int_0^{k^+/\vee}d\xi/(1-\xi)$ is the part of the vertex correction (and, with the opposite sign of the pole, of the real splitting) in which the parent of the observed gluon is faster than $\vee$, i.e.\ a valence mode, and which is not enhanced at small momentum fraction (the enhanced part $2/\xi$ of the same integrand gives the $\log(\vee/k^+)$ that cancels in the pole coefficient). For this part the description of the valence modes by a static source is not valid, and its coefficient is not defined by the effective theory: for a smooth separation $f(p^+/\vee)$ of soft and valence modes ($f(0)=1$) $\ell_\vee$ becomes $-\kappa_fk^+/\vee$ with $\kappa_f=\int_0^\infty dv\,[1-f(v)]/v^2$ ($\kappa_f=1,\sqrt\pi,\pi/2$ for a sharp cut-off, $e^{-u^2}$ and $1/(1+u^2)$), while $\log(\vee/k^+)$ only shifts by a universal constant (Section~\ref{sec:fin_inst}). We therefore define the eikonal approximation as the leading term of the expansion in $k^+/\vee$ and $\Lambda/k^+$, with the logarithms kept, taken at the level of the longitudinal integrands at fixed transverse momenta and at fixed $\epsilon$, i.e.\ before the transverse integrals are done and the regulator is removed. Since the pole factor of Row~I.II is a product of a longitudinal and a transverse factor, $\big(\frac{11}{3}-2L+\ell_\vee\big)\big[-\frac2\epsilon+\log\frac{\v k^2}{\mu^2}\big]$, this expansion is uniform in the transverse loop momentum; at fixed $\epsilon$ it is equivalent to $\vee\to\infty$ before $\epsilon\to0$. A pole multiplying a power correction would have to be renormalised, at that power, by non-eikonal interactions of the valence sector that are not part of the Hamiltonian~\eqref{eq:Hdecomp}; at leading power it is absent. The same holds for the terms $O(\Lambda/k^+)$ on the target side.

The account answers the questions raised by the virtual $\frac{11}{3}$ terms. (i) The virtual term of the vertex correction is present, Eq.~\eqref{eq:dg_II}, and is cancelled, ultraviolet against ultraviolet, by Eq.~\eqref{eq:dg_V}; hence it does not appear in the final result. (ii) The pole of the term $\frac{\beta_0}{2}\delta(1-\xi)$ as it arises in Row~II.IV is collinear. What renormalises the coupling is the ultraviolet half of the normalisation of the dressed produced gluon, $Z_3$ in light-cone gauge, consistent with the Ward identity $Z_g=Z_3^{-1/2}$ ($g_0A_0=gA$, so that the coupling of the classical source to the field needs no counterterm of its own): after the cancellation of Rows~I.II and~I.V the only ultraviolet pole left is that of $Z_3$, $+\frac{\alpha_s\beta_0}{4\pi}\frac{2}{\epsilon_{\rm UV}}$, which is what $g_0^2\to g^2(\mu)$ removes. (iii) The infrared half of the same normalisation is the virtual $\frac{\beta_0}{2}\delta(1-\xi)$ (and $-2N_cL\,\delta(1-\xi)$) of the final-state splitting; real $\mathcal P_{gg+}$ and virtual term together are the full $\mathcal P_{gg}$, which is absorbed by the fragmentation counterterm. No infrared pole is left. In the single-$\epsilon$ bookkeeping the coupling counterterm cancels the collinear pole of the $\delta(1-\xi)$ term formally; this is legitimate only because the normalisation integral is scaleless.

\paragraph{Remark (pure Yang--Mills; identification of the $\beta_0$ terms).} The calculation and the result are for pure Yang--Mills theory, $N_f=0$, where $\beta_0/(2N_c)=\frac{11}{6}$ coincides with integrals of the splitting function. Which of the numbers $\frac{11}{6}$ are $\beta_0/(2N_c)$ can be decided by the BLM-type device of adding a quark loop, without computing the $N_f$ terms of the result, which we do not give. With the vertices normalised as here (the paper's triple-gluon matrix element is, exactly, minus the covariant vertex contracted with light-cone polarisation vectors, and the quark vertex is taken in the same normalisation), the angular-averaged squared vertices satisfy $\frac{n_fT_R\sum|V_{q\bar q}|^2}{\frac12N_c\sum|V_{gg}|^2}=\frac{n_fT_R}{N_c}\frac{\xi^2+(1-\xi)^2}{K_{gg}(\xi)}$, $K_{gg}=\xi(1-\xi)-2+\frac{1}{\xi(1-\xi)}=\mathcal P_{gg+}/(2N_c)$, for every $\xi$: a $q\bar q$ loop adds to the kernel $K_{gg}$ of the self-energy the kernel $\frac{n_fT_R}{N_c}[\xi^2+(1-\xi)^2]$ (with its $d_\perp$-dimensional spin trace), and the corresponding $q\bar q$ pair crosses the shock wave in Row~I.V, with colour factor $T_R[U(\v y)-U(\v x)]$ for a point-like pair. Hence the $\frac{11}{3}$ of Row~I.II becomes $\beta_0/N_c$, $\beta_0=\frac{11}{3}N_c-\frac43T_Rn_f$, the $-\frac{11}{6}$ of $N_{2b\,\rm UV}$ becomes $-\beta_0/(2N_c)$, the $\frac{11}{6}$ of the subtraction term in $\mathfrak K^{\rm V}_{2b}$ becomes $\beta_0/(2N_c)$, the $\frac{11}{12}$ of the normalisation, of the coupling counterterm and of the RC term become $\beta_0/(4N_c)$, and the $\frac{\beta_0}{2}$ of $\delta(1-\xi)$ becomes the full $\frac{\beta_0}{2}$: all of these are integrals of the complete one-loop gluon self-energy kernel. In contrast, the coefficient $-\frac{11}{6}$ of the large-distance tail, Eq.~\eqref{eq:dg_tailtot}, is the double-plus integral of the real $\mathcal P_{gg+}$ and stays $-\frac{11}{6}$ for any $N_f$ (the $\delta(1-\xi)$ parts of the DGLAP and RC terms cancel in the tail for any $\beta_0$). The constant of Row~I.II would become $\frac{67}{9}-\frac{\pi^2}{3}-\frac{20}{9}\frac{n_fT_R}{N_c}=\frac{2}{N_c}K_{\rm CMW}$, $K_{\rm CMW}=N_c\big(\frac{67}{18}-\frac{\pi^2}{6}\big)-\frac{10}{9}T_Rn_f$. There is no real $N_f$ term for a measured gluon at this order. The cancellation of the ultraviolet poles between Rows~I.II and I.V, and the pole balance of Table~\ref{tab:poles}, hold for any $\beta_0$.

\subsection{The collinear term in the \texorpdfstring{$\MS$}{MS-bar} scheme}\label{sec:dg_msbar}

\paragraph{The bare collinear term.} In the variables of Row~II.IV, $\v r=\v w-\v z$, $\v u=\v w'-\v w$ and $\v Y=\frac{\v y+\v y'}{2}$, with $\v y'-\v y=\xi\v u$, the Jacobian of $(\v z,\v w,\v w')\to(\v Y,\v r,\v u)$ is one in any dimension, and the $\v r$ integral gives $-\pi\Gamma(1-\frac\epsilon2)\big(\frac{\mu^2\v u^2e^{\gamma_E}}{4}\big)^{\epsilon/2}\big[\frac{\delta^{mm'}}{\epsilon}+\frac{u^mu^{m'}}{\v u^2}\big]$. If $\v u$ is continued to $d$ dimensions as well and the separation of the parents, $\v s=\v y'-\v y=\xi\v u$, is used as variable,
\begin{align}\label{eq:dg_resc}
\int d^du\,(\v u^2)^{\epsilon/2}\,e^{i\v k\cdot\v u}\,f(\xi\v u)=\xi^{-d}\xi^{-\epsilon}\int d^ds\,(\v s^2)^{\epsilon/2}\,e^{i\frac{\v k}{\xi}\cdot\v s}\,f(\v s)=\xi^{-2}\int d^ds\,(\v s^2)^{\epsilon/2}\,e^{i\frac{\v k}{\xi}\cdot\v s}\,f(\v s)\,.
\end{align}
The bare collinear term is therefore $\xi^{-2}\big[\frac2\epsilon+\mathcal B(\v y'-\v y)\big]_d$ times the LO structure of the parent, without $\log\xi^2$ in the logarithm; a $\log\xi^2$ arises if only $\v r$ is continued to $d$ dimensions, and the two sides of Eq.~\eqref{eq:dg_resc} then differ by $\xi^{-\epsilon}$.

\paragraph{Definition of the fragmentation function and the counterterm.} What survives in the subtracted result depends on the power of $\xi$ that multiplies the counterterm, i.e.\ on the definition of the fragmentation function. We define the fragmentation function as an $\MS$ number density in the momentum fraction (operator definition, momentum sum rule $\sum_h\int dz\,z\,D_h(z)=1$ in $d$ dimensions). Collinear kinematics $\v P=z\v p$ in all $d$ transverse components then gives
\begin{align}
\frac{dN_h}{dP^+d^dP}=\int\frac{dz}{z^{1+d}}\,D(z)\,\frac{dN_g}{dk^+d^dk}\bigg|_{k=P/z}\,,
\end{align}
and the $\MS$ renormalisation $D_0=\big[\delta(1-\xi)+\frac{\alpha_s}{2\pi}\frac2\epsilon\mathcal P_{gg}(\xi)+\dots\big]\otimes D(\mu)$ (pure pole, Mellin convolution) gives the partonic counterterm $+\mathcal C\frac2\epsilon\int d\xi\,\mathcal P_{gg}(\xi)\,\xi^{-d}\,\widehat L_d(\v k/\xi)$, with $\xi^{-d}$ rather than $\xi^{-2}$. This is the counterterm used throughout. It is the standard $\MS$ convention: for a cross section differential in all components of the momentum of the observed hadron the $\MS$ collinear counterterm carries $dz/z^{2-2\epsilon}$ in $4-2\epsilon$ dimensions, i.e.\ $\xi^{-d}$ here, Eq.~(6.17) of~\cite{Catani:1996vz}. With it, and with the same pure-pole renormalisation, the one-hadron-inclusive gluon coefficient functions of $e^+e^-\to gX$~\cite{Altarelli:1979kv,Kretzer:2000yf} and, for $g\to g$ in pure Yang--Mills, of $H\to gX$~\cite{Almasy:2011eq} are reproduced exactly (we computed both, with the $d$-dimensional real-emission matrix elements, phase space and counterterm, in exact arithmetic). With $\xi^{-2}$ in place of $\xi^{-d}$ both come out wrong by $-\frac{\alpha_s}{2\pi}P(z)\log z^2$ times the Born term, the structure of $\Delta$ in Eq.~\eqref{eq:dg_Delta} below. These checks are angle-integrated; for the transverse-momentum Jacobian they rely on the universality of the fragmentation function and on $d$-dimensional collinear kinematics, which is the derivation of~\cite{Catani:1996vz}. With $\xi^{-d}=\xi^{-2}(1+\epsilon\log\xi)$ (CT: counterterm),
\begin{align}\label{eq:dg_two}
\text{bare}-\text{CT}(\xi^{-d})&\to\xi^{-2}\!\int e^{-i\frac{k}{\xi}\cdot(y-y')}\,\log\frac{(y'-y)^2\mu^2}{c_0^2\,\xi^2}\,\mathcal O(y,y')\,,\notag\\
\text{bare}-\text{CT}(\xi^{-2})&\to\xi^{-2}\!\int e^{-i\frac{k}{\xi}\cdot(y-y')}\,\log\frac{(y'-y)^2\mu^2}{c_0^2}\,\mathcal O(y,y')\,.
\end{align}
The first line is the collinear term for the fragmentation function defined above, and is the form used in this paper: the argument of its logarithm is the transverse separation $(\v y'-\v y)/\xi$ of the observed gluon in the amplitude and in the conjugate amplitude, the variable conjugate to $\v k$. The second line is the form of Chirilli, Xiao and Yuan (CXY)~\cite{Chirilli:2011km,Chirilli:2012jd} and of the $\MS$ form of~\cite{Altinoluk:2014eka}, with the separation of the parent. Both statements were verified numerically in $d=2\pm0.002$ with $d$-dimensional Fourier--Bessel transforms of a test function, to $O(\epsilon^2)$. Hence the $\log\xi^2$ cancels in the bare result, but not in the subtracted one, since with our definition the counterterm carries the $d$-dimensional Jacobian. CXY obtain the collinear pole of the final-state splitting from a $d$-dimensional integral over the momentum of the parent at fixed observed momentum, which yields exactly $\xi^{-2}$ times the parent spectrum, and absorb it by a pure-pole renormalisation of $D_{h/g}$ in a hadronic formula written with $dz/z^2$~\cite{Chirilli:2012jd}: their counterterm carries $\xi^{-2}$. The appendix of~\cite{Altinoluk:2014eka} writes the hadronic formula with the $d$-dimensional Jacobian $\zeta^{-(D-2)}$, but defines the renormalised fragmentation function with the measure $d\xi/(1-\xi)^{5-D}$, which is not the pure-pole renormalisation of the number density and amounts to the same choice as CXY. In the cut-off scheme of the main text of~\cite{Altinoluk:2014eka} our form is obtained with $\mu^2_{\MS}=c_0^2\mu^2_{\rm cut}\approx1.26\,\mu^2_{\rm cut}$.

\paragraph{Remark: the CXY convention.} With $\mathcal G_0(\xi)\equiv\xi^{-2}\int_{x,y,x',y'}e^{-i\v k\cdot(\v y-\v y')/\xi}\,\mathcal O(y,y')=\Pi(\xi)$, the two forms of the collinear term differ by
\begin{align}\label{eq:dg_Delta}
\Delta\equiv\frac{d^3N_{\rm DGLAP}}{d^3k}-\frac{d^3N^{\rm CXY}_{\rm DGLAP}}{d^3k}=+\,\mathcal C\int_0^1d\xi\,\mathcal P_{gg}(\xi)\,\log\xi^2\;\mathcal G_0(\xi)\,.
\end{align}
The $\delta(1-\xi)$ part of $\Delta$ vanishes, $\Delta$ is integrable at $\xi\to0$ ($\mathcal G_0(\xi)\to0$), and it does not depend on $\mu$. Equivalently $D^{\rm CXY}=\big[\delta(1-\xi)+\frac{\alpha_s}{2\pi}\mathcal P_{gg}(\xi)\log\xi^2\big]\otimes D$. The results of CXY and of~\cite{Altinoluk:2014eka} are meant to be used with standard $\MS$ fragmentation functions, which obey the momentum sum rule and for which the $\MS$ coefficient is the first line of Eq.~\eqref{eq:dg_two}; their form of the collinear term is the coefficient appropriate for the functions $D^{\rm CXY}$, i.e.\ it misses $\Delta$ at $O(\alpha_s)$ when used with standard fragmentation functions. Since $\int_0^1d\xi\,\xi\,\mathcal P_{gg}(\xi)\log\xi^2=4N_c\big(\frac{65}{144}-\zeta_2\big)=\frac{N_c(65-24\pi^2)}{36}\neq0$ ($-14.3225$ for $N_c=3$; with quarks the combination $\int_0^1d\xi\,\xi\log\xi^2[\mathcal P_{gg}+2N_f\mathcal P_{qg}]=\frac{N_c(65-24\pi^2)}{36}-\frac{11N_f}{36}$ would enter), the functions $D^{\rm CXY}$ do not satisfy the momentum sum rule at $O(\alpha_s)$; equivalently, the energy sums of the $e^+e^-$ and $H\to gX$ coefficient functions computed with $\xi^{-2}$ are not those of the total rates. $\Delta$ is given only so that the CXY convention can be recovered; all collinear terms of this paper refer to the first line of Eq.~\eqref{eq:dg_two}. A comparison with the NLO coefficient of hadroproduction, $pp\to gX$, and a collinear (large-$\v k$) limit of the hybrid formula of~\cite{Chirilli:2012jd} have not been carried out; the bare final-state collinear terms of~\cite{Chirilli:2012jd} and of this paper are the same $d$-dimensional integral, so that such a comparison would re-test only the counterterm, which is fixed above.

\subsection{The DGLAP and running-coupling terms}\label{sec:dg_final}

We give the DGLAP and running-coupling terms in two exactly equivalent forms: in coordinate space, Eqs.~\eqref{eq:DGLAP}--\eqref{eq:LORC}, and in a local momentum-space form, Eqs.~\eqref{eq:dg_local}--\eqref{eq:dg_R}; their relation is Eq.~\eqref{eq:dg_equiv}. Both refer to the $\MS$ number-density fragmentation function of Section~\ref{sec:dg_msbar}.

\paragraph{Coordinate-space form.} With one scale $\mu$ and the phase $e^{-i\v k\cdot(\v y-\v y')/\xi}$, define
\begin{align}\label{eq:dg_F}
\mathcal F(\xi)\equiv\frac{1}{\xi^2}\int_{x,y,x',y'}e^{-ik\cdot(y-y')\frac1\xi}\,\log\frac{(y'-y)^2\mu^2}{c_0^2\,\xi^2}\;\mathcal O(y,y')\,,\qquad
\mathcal F(0)\equiv\lim_{\xi\to0}\mathcal F(\xi)=-\frac{4\pi}{\v k^2}\,\mathcal X\,,
\end{align}
where the integral over $\v y-\v y'$ is done before the integral over $\xi$ ($\mathcal F(0)$ is the same in the CXY form). The DGLAP term is
\begin{align}\label{eq:DGLAP}
\frac{d^3N_{\rm DGLAP}}{d^3k}&=-\mathcal C\int_0^1d\xi\,\widetilde{\mathcal P}_{gg}(\xi)\,\mathcal F(\xi)\,,\notag\\
\widetilde{\mathcal P}_{gg}(\xi)&=2N_c\Big[\xi(1-\xi)+\frac{\xi}{(1-\xi)_+}+\Big[\frac{1-\xi}{\xi}\Big]_+\Big]+\frac{\beta_0}{2}\,\delta(1-\xi)\,,\notag\\
\frac{d^3N_{\rm DGLAP}}{d^3k}&=-2N_c\,\mathcal C\int_0^1d\xi\Big\{\xi(1-\xi)\,\mathcal F(\xi)+\frac{\xi\mathcal F(\xi)-\mathcal F(1)}{1-\xi}+\frac{1-\xi}{\xi}\big[\mathcal F(\xi)-\mathcal F(0)\big]\Big\}\notag\\&\quad-\mathcal C\,\frac{\beta_0}{2}\,\mathcal F(1)\,,
\end{align}
with the prescriptions~\eqref{eq:plusxi}. $\widetilde{\mathcal P}_{gg}$ is the $\MS$ splitting function $\mathcal P_{gg}$ with, in addition, the subtraction at $\xi\to0$, whose subtracted term belongs to the BFKL term of the row. With this reading the left-over of that subtraction,
\begin{align}\label{eq:remnant}
\frac{d^3N_{\rm rem}}{d^3k}=-2N_c\,\mathcal C\,\frac{4\pi}{\v k^2}\,\mathcal X=-\frac{1}{(2\pi)^3}\,\frac{g^4N_c}{2\pi^3}\,\frac{1}{\v k^2\,k^+}\,\mathcal X\,,
\end{align}
the term $-\frac{1}{k^+}\times$(BFKL structure), is counted among the finite terms of Row~II.IV (Section~\ref{sec:finite}). It arises because $\int_{k^+/\vee}^1d\xi\frac{1-\xi}{\xi}=\log\frac{\vee}{k^+}-1+\frac{k^+}{\vee}$. The sum of Eqs.~\eqref{eq:DGLAP} and~\eqref{eq:remnant} can be written with the plus prescription at $\xi\to0$ acting on $(1-\xi)\mathcal F(\xi)$,
\begin{align}\label{eq:DGLAPdoubleplus}
\frac{d^3N_{\rm DGLAP}}{d^3k}+\frac{d^3N_{\rm rem}}{d^3k}&=-\mathcal C\int_0^1d\xi\,\Big\{\mathcal P_{gg}(\xi)\,\mathcal F(\xi)-\frac{2N_c}{\xi}\,\mathcal F(0)\Big\}\notag\\
&=-2N_c\,\mathcal C\int_0^1\frac{d\xi}{\xi_+(1-\xi)_+}\,N(\xi)\,\mathcal F(\xi)-\mathcal C\,\frac{\beta_0}{2}\,\mathcal F(1)\,,
\end{align}
where $\mathcal P_{gg}$ has no subtraction at $\xi=0$ and the last form is the ``double plus'' prescription of~\cite{Kovner:2023vsy} (Appendix~\ref{app:distributions}); the two subtractions are there for the same reason as in~\cite{Kovner:2023vsy}: both poles are part of the leading-logarithmic evolution, here BFKL at $\xi\to0$ and JIMWLK at $\xi\to1$. Eqs.~\eqref{eq:DGLAP}--\eqref{eq:remnant} and Eq.~\eqref{eq:DGLAPdoubleplus} are the same quantity; only one of the two forms is to be used. The dependence on the scale is $\partial N_{\rm DGLAP}/\partial\log\mu^2=-\mathcal C\int_0^1d\xi\,\mathcal P_{gg}(\xi)\,\Pi(\xi)$, the standard DGLAP evolution of $D_g(\mu)$ including its $\frac{\beta_0}{2}\delta(1-\xi)$; the subtraction at $\xi=0$ does not affect it, since $\Pi(0)=0$.

The running-coupling term is the $\delta(1-\xi)$ part of Eq.~\eqref{eq:DGLAP} with the opposite sign,
\begin{align}\label{eq:RC}
\frac{d^3N_{\rm RC}}{d^3k}=+\,\mathcal C\,\frac{\beta_0}{2}\,\mathcal F(1)=+\frac{\alpha_s\beta_0}{4\pi}\,\mathcal N_0\int_{x,y,x',y'}e^{-ik\cdot(y-y')}\,\log\frac{(y'-y)^2\mu^2}{c_0^2}\;\mathcal O(y,y')\,.
\end{align}
It is added and subtracted: $N_{\rm DGLAP}+N_{\rm RC}$ does not depend on $\beta_0$, since the real emission produces no $\delta(1-\xi)$ and the virtual $\frac{11}{3}$ terms cancel between Rows~I.II and~I.V. Its content is the dependence on the scale: $\partial N_{\rm RC}/\partial\log\mu^2=+\frac{\alpha_s\beta_0}{4\pi}\times$LO compensates the running of $g^2(\mu)$ in the LO term, and the $\delta(1-\xi)$ part of $N_{\rm DGLAP}$ the corresponding part of the evolution of $D_g(\mu)$. Together with the LO term,
\begin{align}\label{eq:LORC}
\frac{d^3N_{\rm LO}}{d^3k}+\frac{d^3N_{\rm RC}}{d^3k}&=\frac{1}{(2\pi)^3}\,\frac{g^2(\mu)}{2\pi^2}\,\frac{1}{k^+}\notag\\&\times\int_{x,y,x',y'}e^{-ik\cdot(y-y')}\Big[1+\frac{\alpha_s(\mu)\beta_0}{4\pi}\log\frac{(y'-y)^2\mu^2}{c_0^2}\Big]\mathcal O(y,y')\,,
\end{align}
and to this order the bracket is $\alpha_s(c_0^2/\v r^2)/\alpha_s(\mu^2)$, $\v r=\v y'-\v y$. Since $N_{\rm RC}$ is minus the $\delta(1-\xi)$ term of $N_{\rm DGLAP}$, the scale $c_0^2/\v r^2$ is the choice $\mu=c_0/r$ in the $\delta(1-\xi)$ part of the collinear logarithm, the choice made by CXY for the factorisation scale~\cite{Chirilli:2012jd}; Eq.~\eqref{eq:LORC} holds only together with the $\delta(1-\xi)$ term kept in $\widetilde{\mathcal P}_{gg}$ at the same $\mu$.

\paragraph{Separate renormalisation and factorisation scales.} The add-and-subtract of the $\delta(1-\xi)$ term identifies the renormalisation scale with the factorisation scale; the statement that the residual dependence on $\mu$ comes only from the $\delta(1-\xi)$ term holds only for $\mu_R=\mu_F$. With $g_0^2=\mu_R^\epsilon g^2(\mu_R)\big[1-\frac{\alpha_s\beta_0}{4\pi}\frac2\epsilon\big]$ and $D_0=\big[\delta(1-\xi)+\frac{\alpha_s}{2\pi}\frac2\epsilon\mathcal P_{gg}\big]\otimes D(\mu_F)$, the two halves of the scaleless normalisation~\eqref{eq:dg_norm} are cancelled separately, the ultraviolet half by the coupling counterterm at $\mu_R$ and the infrared half by the $\delta(1-\xi)$ part of the fragmentation counterterm at $\mu_F$, and
\begin{align}\label{eq:dg_muRmuF}
\frac{d^3N_{\rm LO}}{d^3k}+\frac{d^3N_{\rm RC}}{d^3k}&=\frac{1}{(2\pi)^3}\,\frac{g^2(\mu_R)}{2\pi^2}\,\frac{1}{k^+}\notag\\&\quad\times\int_{x,y,x',y'}e^{-ik\cdot(y-y')}\Big[1+\frac{\alpha_s\beta_0}{4\pi}\log\frac{(y'-y)^2\mu_R^2}{c_0^2}\Big]\mathcal O(y,y')\,,\notag\\
\frac{d^3N_{\rm DGLAP}}{d^3k}&=-\mathcal C\int_0^1d\xi\,\widetilde{\mathcal P}_{gg}(\xi)\,\mathcal F(\xi)\Big|_{\mu\to\mu_F}\,,
\end{align}
including the $\frac{\beta_0}{2}\delta(1-\xi)$ term at $\mu_F$. Their $\beta_0$ part is the running-coupling term at $\mu_R$ minus the $\delta(1-\xi)$ term at $\mu_F$, $N_{\rm RC}(\mu_R)-\mathcal C\frac{\beta_0}{2}\mathcal F(1)|_{\mu_F}$, which is $+\frac{\alpha_s\beta_0}{4\pi}\log\frac{\mu_R^2}{\mu_F^2}\times$LO, and the result is independent of $\mu_R$ and $\mu_F$ to this order. In the local form, $N^{\rm loc}_{\rm RC}=-\frac{\alpha_s\beta_0}{4\pi}\log\frac{\v k^2}{\mu_R^2}\,\mathcal N_0\widehat L(k)$ and $N^{\rm loc}_{\rm DGLAP}=\mathcal C\log\frac{\v k^2}{\mu_F^2}\int_0^1d\xi\,\mathcal P_{gg}(\xi)\,\Pi(\xi)$, while $\mathcal R$ and $N_{\mathcal R}$ are unchanged.

In the classification of~\cite{Kovner:2023vsy}, where of the two $\beta_0$'s of the NLO JIMWLK kernel one is the ultraviolet running of the coupling and the other a DGLAP-type resolution logarithm, both are here tied to the produced gluon: the ultraviolet one is the $Z_3$ pole of the scaleless normalisation~\eqref{eq:dg_norm}, removed by the coupling counterterm and not visible in the finite result; the infrared one is the $\frac{\beta_0}{2}\delta(1-\xi)$ of $\mathcal P_{gg}$, the only $\beta_0$ that appears explicitly. They are equal in magnitude because the normalisation integral is scaleless. This agrees in structure with Kovchegov and Weigert~\cite{Kovchegov:2007vf} and Horowitz and Kovchegov~\cite{Horowitz:2010yg}, in whose calculation (quark--quark scattering, BLM resummation of quark bubbles) the observable includes collinear pairs whose invariant mass, i.e.\ the virtuality of the produced-gluon line, is below a fixed cut-off $\Lambda_{\rm coll}$, and the coupling of the emission vertex of the produced gluon is $\alpha_s(\Lambda^2_{\rm coll}e^{-5/3})$ (Eqs.~(3.10)--(3.16) of~\cite{Horowitz:2010yg}); $\Lambda_{\rm coll}$ does not depend on the exchanged momenta and plays the role of the factorisation scale of the fragmentation function~\cite{Kovchegov:2007vf}, while the remaining couplings run with $\v q^2$, $(\v k-\v q)^2$ and a complex scale $Q(\v q,\v k-\v q)$, Eqs.~(3.31), (3.32) of~\cite{Horowitz:2010yg}. (Equation numbers of~\cite{Horowitz:2010yg} refer to the arXiv version.)

\paragraph{Local (momentum-space) form.} With $\log\frac{\v r^2\mu^2}{c_0^2}=-\frac1\pi\int d^2l\,\frac{e^{-i\v l\cdot\v r}-\theta(\mu-|\v l|)}{\v l^2}$ one has exactly
\begin{align}\label{eq:dg_loc}
\mathcal F(\xi)=-\log\frac{\v k^2}{\mu^2}\,\Pi(\xi)+\frac{1}{\xi^2}\,\mathcal R\Big(\frac{\v k}{\xi}\Big)\,,\qquad
\mathcal R(\v q)\equiv-\frac1\pi\int d^2l\,\frac{\widehat L(\v q+\v l)-\theta(|\v q|-|\v l|)\,\widehat L(\v q)}{\v l^2}\,,
\end{align}
with $\lim_{\xi\to0}\xi^{-2}\mathcal R(\v k/\xi)=\mathcal F(0)$; $\mathcal R$ does not depend on $\mu$. Inserted into Eqs.~\eqref{eq:DGLAP}--\eqref{eq:RC}, it splits each term into a local part and a non-local, $\mu$-independent part. The local DGLAP and running-coupling terms are
\begin{align}\label{eq:dg_local}
\frac{d^3N^{\rm loc}_{\rm DGLAP}}{d^3k}&\equiv\mathcal C\,\log\frac{\v k^2}{\mu^2}\int_0^1d\xi\,\mathcal P_{gg}(\xi)\,\Pi(\xi)\,,\qquad
\frac{d^3N^{\rm loc}_{\rm RC}}{d^3k}\equiv-\frac{\alpha_s\beta_0}{4\pi}\,\log\frac{\v k^2}{\mu^2}\,\mathcal N_0\,\widehat L(k)\,,\notag\\
\frac{d^3N_{\rm LO}}{d^3k}+\frac{d^3N^{\rm loc}_{\rm RC}}{d^3k}&=\mathcal N_0\,\widehat L(k)\Big[1-\frac{\alpha_s\beta_0}{4\pi}\log\frac{\v k^2}{\mu^2}\Big]\,,
\end{align}
and the coordinate-space terms are
\begin{align}\label{eq:dg_locsplit}
\frac{d^3N_{\rm DGLAP}}{d^3k}+\frac{d^3N_{\rm rem}}{d^3k}&=\frac{d^3N^{\rm loc}_{\rm DGLAP}}{d^3k}-\mathcal C\int_0^1d\xi\Big\{\mathcal P_{gg}(\xi)\,\frac{\mathcal R(\v k/\xi)}{\xi^2}-\frac{2N_c}{\xi}\,\mathcal F(0)\Big\}\,,\notag\\
\frac{d^3N_{\rm RC}}{d^3k}&=\frac{d^3N^{\rm loc}_{\rm RC}}{d^3k}+\mathcal C\,\frac{\beta_0}{2}\,\mathcal R(\v k)\,.
\end{align}
The local terms carry the whole dependence on $\mu$ and need no subtraction at $\xi=0$ ($\Pi(\xi)\to0$); to this order $N_{\rm LO}+N^{\rm loc}_{\rm RC}$ is the LO spectrum with the coupling at the measured momentum, $\alpha_s(\v k^2)/\alpha_s(\mu^2)$. The non-local remainders of the two lines of Eq.~\eqref{eq:dg_locsplit}, whose $\delta(1-\xi)$ parts $\mp\mathcal C\frac{\beta_0}{2}\mathcal R(\v k)$ cancel, add up to
\begin{align}\label{eq:dg_R}
\frac{d^3N_{\mathcal R}}{d^3k}=-\mathcal C\int_0^1d\xi\Big\{\mathcal P_{gg+}(\xi)\,\frac{\mathcal R(\v k/\xi)}{\xi^2}-\frac{2N_c}{\xi}\,\mathcal F(0)\Big\}\,,
\end{align}
which belongs to the finite terms of Row~II.IV. It contains the BFKL ($\xi\to0$) subtraction and the region of a soft parent discussed next. The two forms are related by
\begin{align}\label{eq:dg_equiv}
\frac{d^3N_{\rm DGLAP}}{d^3k}+\frac{d^3N_{\rm rem}}{d^3k}+\frac{d^3N_{\rm RC}}{d^3k}=\frac{d^3N^{\rm loc}_{\rm DGLAP}}{d^3k}+\frac{d^3N^{\rm loc}_{\rm RC}}{d^3k}+\frac{d^3N_{\mathcal R}}{d^3k}\,,
\end{align}
an identity; either side may be used, with the same finite terms of the other rows.

\paragraph{The two forms at large distance.} The forms differ in where the large-distance (soft-parent) region of Row~II.IV sits (Section~\ref{sec:dg_tail}). In the coordinate-space form the DGLAP and running-coupling terms (and the remnant) are each logarithmically divergent at large position of the parent gluon. Their sum with the finite terms of the second part of Row~II.I and of Row~III.III is not: the tail cancels term by term between the three rows (Table~\ref{tab:tails}), which therefore have to be integrated together over the position of the gluon pair. In the local form $N^{\rm loc}_{\rm DGLAP}$ and $N^{\rm loc}_{\rm RC}$ have no large-distance tail; the tail of the coordinate-space DGLAP, running-coupling and remnant terms together, $-\frac{11}{6}\,\mathcal U$ in the units of Table~\ref{tab:tails}, is carried by the remainder $N_{\mathcal R}$, which is then integrated together with the second part of Row~II.I and with Row~III.III. In both forms the $\cos2\varphi$ part of the tail of Row~II.IV is in the tensor term of its finite part.

\subsection{The large-distance tail of Row II.IV}\label{sec:dg_tail}

At large parent position $\v Y$ the LO structure tends to $G_\infty(\v Y)/\v Y^2$, $G_\infty=\int_{x,x'}[U(\v Y)-U(\v x')]^{ce'}[U(\v Y)-U(\v x)]^{ce}\rho^{e'}(\v x')\rho^e(\v x)$, independently of the separation $\v s=\v y'-\v y$. The logarithm then survives the Fourier transform, $\int d^2s\,e^{-i\v q\cdot\v s}\log\v s^2=-4\pi/\v q^2$, and $\mathcal F(\xi)\to-\frac{4\pi}{\v k^2}G_\infty/\v Y^2$ for every $\xi$. In units of $\mathcal C\frac{4\pi}{\v k^2}\int\frac{d^2Y}{\v Y^2}G_\infty$ the tails of the coordinate-space terms are: DGLAP $+\frac{N_c}{6}$, RC $-\frac{11N_c}{6}$, remnant $-2N_c$, and the traceless tensor term of the finite part of the row, Eq.~\eqref{eq:H-II.IVrem}, $+\frac{N_c}{3}\cos2\varphi$, $\varphi$ the angle between $\v k$ and $\v Y$; in total
\begin{align}\label{eq:dg_tailtot}
2N_c\,\mathcal C\,\frac{4\pi}{\v k^2}\int\frac{d^2Y}{\v Y^2}\Big[-\frac{11}{6}+\frac16\cos2\varphi\Big]\,G_\infty(\v Y)\,.
\end{align}
The coefficient $-\frac{11}{6}=\int_0^1\frac{d\xi}{\xi_+(1-\xi)_+}N(\xi)$ is the double-plus integral of the real splitting function. The RC term and the $\delta(1-\xi)$ part of the DGLAP term cancel each other also in the tail, so the tail is not related to the running coupling, nor to the pole ($\mathcal G_0$ has no tail for $\v k\neq0$). $G_\infty\neq0$ also for a colour-neutral projectile and in the colour-singlet projection.

\paragraph{Origin.} In momentum space the tail is the region $\v k/\xi+\v l\to0$ of $\mathcal R$: a parent much softer than the measured gluon, $|\v q|\ll|\v k|$, which splits into two hard gluons. Since the LO spectrum behaves as $\widehat L(\v q)\propto1/\v q^2$ at small $\v q$ (Weizs\"acker--Williams field of the colour charge $\int[U_\infty-U(\v x)]\rho(\v x)$, non-zero also for a neutral projectile), $\int d^2q\,\widehat L(\v q)/\v k^2$ diverges logarithmically. In a model with an infrared cut-off $m$ the coordinate-space form grows like $-\frac{2}{\v k^2}\log\frac1m$, while the local form $-\log\frac{\v k^2}{\mu^2}\widehat L(\v k)$ is independent of $m$. The DGLAP and RC terms as written in coordinate space are thus separately logarithmically divergent at large distance; the local terms~\eqref{eq:dg_local} are not, and the divergence resides in $N_{\mathcal R}$.

\paragraph{Cancellation.} The two-gluon final state of Group~II and~III contains the amplitude $\mathcal A_B-\mathcal A_{CA}$ (Eq.~\eqref{eq:A2gamp}), where $\mathcal A_B$ is the two-gluon component of the wave function scattered on the target, through the single-$\rho$ part of $\mf B_2$, and $\mathcal A_{CA}$ the LO gluon split by $\bar{\mf C}_1$ after the shock wave. $|\mathcal A_B|^2$ is the second part of Row~II.I, $|\mathcal A_{CA}|^2$ is Row~II.IV, and $-\mathcal A^\dagger_{CA}\mathcal A_B$ with its complex conjugate is Row~III.III. In LCPT, with $E_1$ the energy of the one-gluon (parent) state and $E_2$ that of the pair, the two amplitudes are $b=\frac{V_{\rho g}V_{ggg}}{E_1E_2}$ and $c\,A=\frac{V_{\rho g}V_{ggg}}{E_1(E_2-E_1)}$, so that
\begin{align}\label{eq:dg_bcA}
b-c\,A=-\frac{V_{\rho g}V_{ggg}}{E_2(E_2-E_1)}\,:
\end{align}
the parent pole $1/E_1\propto q^+/\v q^2$ cancels. With the colour factors $[U(\v w)U(\v z)-U(\v x)]$ of $\mathcal A_B$ (in the reduced form) and $[U(\v y)-U(\v x)]$ of $\mathcal A_{CA}$, the amplitude is $[U(\v w)U(\v z)-U(\v y)]\,b+[U(\v y)-U(\v x)]\,(b-cA)$. We verified this numerically with the explicit momentum-space forms of $\mf B_2$ and of $\mf C_1\mf A^{(1)}$, Eqs.~\eqref{eq:B2mom} and~\eqref{eq:C1mom}: at small parent momentum the two amplitudes grow like $1/|\v q|$ while their difference tends to a constant, and the coefficients of $1/\v q^2$ in $|\mathcal A_B|^2$, $|\mathcal A_{CA}|^2$ and $2\,{\rm Re}\,\mathcal A_{CA}^\dagger\mathcal A_B$ are equal, i.e.\ they enter with the weights $1,1,-2$. Consequently, the large-distance (soft-parent) tail of Row~II.IV is cancelled by the large-distance region of the second part of Row~II.I and of Row~III.III, in which both final-state gluons are far from the projectile together ($\v w,\v w',\v z\to\infty$ at fixed $\v w-\v z$). The cancellation is pointwise in the momentum fraction and in the relative momentum of the pair, and hence holds separately in the BFKL, JIMWLK and finite sectors, including the $\cos2\varphi$ part, provided the three rows are split with the same end-point subtractions. It requires $U(\v w)U(\v z)f=fU(\v y)$ for the pair at large distance, i.e.\ a Wilson line that varies slowly there, as for a target of finite size. For a colour-neutral projectile no tail survives even without this condition, since the remaining term $[U(\v w)U(\v z)-U(\v y)]\,b$ multiplies $\int_x\Kv(\v x-\v Y)\rho(\v x)$, which falls like $1/\v Y^2$. The cancellation does not require a change of the split into DGLAP and RC terms, but the coordinate-space terms~\eqref{eq:DGLAP}--\eqref{eq:RC} are separately divergent at large distance: with the local form~\eqref{eq:dg_local} DGLAP and RC terms are finite, and the remainder~\eqref{eq:dg_R} has to be integrated together with the finite terms of the second part of Row~II.I and of Row~III.III, under one integral over the position of the pair. 
\paragraph{Term-by-term check.} We have verified the cancellation on the finite terms of the three rows. In the joint limit the three rows have the same tail density in $t=p^+/k^+$,
\begin{align}\label{eq:dg_taildens}
\rho(t)=\frac{(t^2+t+1)^2+t^2\cos2\varphi}{t(1+t)^4}\,,
\end{align}
with colour factors $+N_c\Gamma$, $+N_c\Gamma$ and $-N_c\Gamma$ (the last twice, with its complex conjugate), so that the sum vanishes at every $t$. Integrated, in units of $\mathcal U\equiv2N_c\,\mathcal C\frac{4\pi}{\v k^2}\int\frac{d^2Y}{\v Y^2}G_\infty$, the tails of the rows as written are given in Table~\ref{tab:tails}.
\begin{table}[h]\centering
\begin{tabular}{lc}\hline
row & tail$/\,\mathcal U$\\\hline
II.IV: DGLAP $\frac1{12}$, RC $-\frac{11}{12}$, remnant $-1$, tensor $\frac16\cos2\varphi$ & $-\frac{11}{6}+\frac16\cos2\varphi$\\
II.I, second part & $-\frac{11}{6}+\frac16\cos2\varphi$\\
III.III, second part, with c.c. & $+\frac{11}{3}-\frac13\cos2\varphi$\\
II.I and III.III, first parts & $0$\\\hline
sum & $0$\\\hline
\end{tabular}
\caption{Large-distance (soft-parent) tails of the finite terms in the joint limit $\v w,\v w',\v z\to\infty$ at fixed $\v w-\v z$.}\label{tab:tails}
\end{table}
In the first parts the measured gluon on one side is tied to the projectile by the Fourier transform, and there is no tail. In the BFKL and JIMWLK sectors each row carries $(1+1-2)\times$ the same logarithm; the piece moved between sectors by the recoil rewriting of Row~III.III is the companion of its logarithmic term and together with it reproduces the unsubtracted term. The rule $\int d^2z/\v z^2=0$ plays no role here, since the poles it removes multiply $\delta^2(\v k)$. The colour factors were checked with random $SU(2)$ and $SU(3)$ Wilson lines: for Wilson lines that are uniform at large distance the cancellation is exact (to $10^{-14}$), and for a colour-neutral projectile only the terms bilinear in $U(\v x)$, $U(\v x')$ survive, with factors $+N_c,+N_c,-N_c$ independent of the far Wilson lines, so that the tail cancels for any target. For a coloured projectile and a target that is not uniform at large distance a remainder proportional to the total charge survives; it is the non-singlet term of the JIMWLK formula, not a defect of the finite terms. The region in which only part of the pair is far does not contribute at this order (Section~\ref{sec:fin_integrable}, condition (C10)). Not verified are the interchange of the $p^+$ integral with the joint limit in Row~III.III (checked numerically for Row~II.I only) and a full transverse integration with Wilson lines.

% ======================================================================
\begingroup\sloppy
\section{Finite terms}\label{sec:finite}
% ======================================================================
% The equations of the rows (files rows/*.tex) are generated mechanically by scripts/gen_rows.py from
% Sec. 8 of nlo_review/final_explicit/Final_terms_explicit_2_fixed.tex; do not edit them by hand.

The finite terms are, row by row, the limit $\Lambda\to0$, $\vee\to\infty$ of
\begin{align}\label{eq:fin_def}
&\big[\text{row with the cut-offs}\big]\notag\\&\quad-\log\frac{\vee}{k^+}\times\big[\text{its term in the BFKL list}\big]-\log\frac{k^+}{\Lambda}\times\big[\text{its term in the JIMWLK list}\big]\,,
\end{align}
where the lists are those whose sums are Eqs.~\eqref{eq:BFKLrho} and~\eqref{eq:JIMWLKrho}. The collinear poles are treated as in Section~\ref{sec:dglap_rc}, and the DGLAP and running-coupling terms of Row~II.IV are those of Section~\ref{sec:dg_final}. Finite terms arise in Rows I.II, I.III, I.V, II.I, II.IV, III.II, III.III, III.V and III.VI. The rows with four charges only (I.I, II.II, II.III, III.IV) contribute to the evolution terms only, and the finite parts of Rows~I.IV and~III.I cancel identically.

In this section we give the finite terms Group by Group and Row by Row, in two forms. In the first form the cancellations that are exact have been carried out and every product of three charges has been written as its totally symmetrised product plus commutator terms, which contain two charges; this form holds for an arbitrary projectile and is given in Appendix~\ref{app:finite}. In the second form the result is averaged over the projectile and written in terms of colour-singlet combinations of Wilson lines, Section~\ref{sec:fin_rows}. In both forms the Wilson lines appear, as at leading order, through differences $U-U$, and the longitudinal integrals are written with plus prescriptions in a momentum fraction $\xi$. We begin with a summary of the complete result.

\subsection{Summary of the result}\label{sec:result_summary}

The NLO single inclusive gluon spectrum, Eq.~\eqref{eq:schematic}, is
\begin{align}\label{eq:result}
\frac{d^3N}{d^3k}&=\Big[\frac{d^3N_{\rm LO}}{d^3k}+\frac{d^3N_{\rm RC}}{d^3k}\Big]+\log\frac{\vee}{k^+}\,\frac{d^3N_{\rm BFKL}}{d^3k}+\log\frac{k^+}{\Lambda}\,\frac{d^3N_{\rm JIMWLK}}{d^3k}\notag\\
&\quad+\frac{d^3N_{\rm DGLAP}}{d^3k}+\sum_{\rm rows}\Big[\frac{d^3N^{\rm row}_{\rm sym}}{d^3k}+\frac{d^3N^{\rm row}_{\rm con}}{d^3k}+\frac{d^3N^{\rm row}_{\rm quad}}{d^3k}\Big]\,,
\end{align}
in pure Yang--Mills theory ($N_f=0$), in the $\MS$ scheme with conventional dimensional regularisation (Section~\ref{sec:dg_poles}), with the $\MS$ number-density fragmentation function (Section~\ref{sec:dg_msbar}), at leading power in $k^+/\vee$ and $\Lambda/k^+$ (eikonal approximation, Section~\ref{sec:dg_poles}). The finite terms of a row consist of up to three parts: the symmetrised three-$\rho$ part (``sym''), the two-$\rho$ contact terms that the symmetrisation of the three charges leaves behind (``con''), and the genuine two-$\rho$ terms (``quad''). The pieces are:
\begin{center}
\renewcommand{\arraystretch}{1.18}
\footnotesize\setlength{\tabcolsep}{3pt}
\begin{tabular}{@{}>{\raggedright\arraybackslash}p{0.17\textwidth}>{\raggedright\arraybackslash}p{0.27\textwidth}>{\raggedright\arraybackslash}p{0.27\textwidth}>{\raggedright\arraybackslash}p{0.24\textwidth}@{}}
\toprule
term & ordered charges & averaged, colour singlet & content\\
\midrule
LO $+$ RC & \eqref{eq:LO}, \eqref{eq:LORC}; local form \eqref{eq:dg_local} & \eqref{eq:LOsinglet}, \eqref{eq:sLOrepl} & LO with $\alpha_s(c_0^2/\v r^2)$, resp.\ $\alpha_s(\v k^2)$; separate $\mu_R$, $\mu_F$: \eqref{eq:dg_muRmuF}\\
BFKL & \eqref{eq:BFKLrho} & \eqref{eq:BFKLsinglet} & one LO BFKL step of the projectile correlator, \eqref{eq:BFKLasLO}\\
JIMWLK & \eqref{eq:JIMWLKrho} & \eqref{eq:JIMWLKsinglet} & $-H_{\rm JIMWLK}$ on the LO target operator, \eqref{eq:JIMWLKasLO}\\
DGLAP & \eqref{eq:DGLAP}; local form \eqref{eq:dg_local}, \eqref{eq:dg_R} & \eqref{eq:sLOrepl} & final-state $\MS$ fragmentation with $\widetilde{\mathcal P}_{gg}$\\
\midrule
Rows I.I, II.II, II.III, III.IV & --- & --- & evolution terms only\\
Rows I.IV $+$ III.I & --- & --- & finite parts cancel identically\\
Row I.II & \eqref{eq:H-I.II-3}, \eqref{eq:H-I.II-2}, \eqref{eq:H-I.II-2o} & \eqref{eq:S-I.II-3}, \eqref{eq:S-I.II-2}, \eqref{eq:S-I.II-2o} & vertex correction, $\mf A^{(3)}$\\
Row I.III & \eqref{eq:H-I.III-3}, \eqref{eq:H-I.III-2} & \eqref{eq:S-I.III-3}, \eqref{eq:S-I.III-2} & rescattering, $\mf B_3$\\
Row I.V & \eqref{eq:H-I.V-3}, \eqref{eq:H-I.V-2}, \eqref{eq:H-I.V-2o} & \eqref{eq:S-I.V-3}, \eqref{eq:S-I.V-2}, \eqref{eq:S-I.V-2o} & recombination, $\mf C_1^\dagger\mf B_2$\\
Row II.I & \eqref{eq:H-II.I-3}, \eqref{eq:H-II.I-2}, \eqref{eq:H-II.I-2a}--\eqref{eq:H-II.I-2e} & \eqref{eq:S-II.I-3}, \eqref{eq:S-II.I-2}, \eqref{eq:S-II.I-2a}--\eqref{eq:S-II.I-2e} & $|\mathcal A_B|^2$\\
Row II.IV & \eqref{eq:H-II.IVrem} & \eqref{eq:S-II.IVrem} & $|\mathcal A_{CA}|^2$: tensor term and remnant\\
Row III.II & \eqref{eq:H-III.II-3}, \eqref{eq:H-III.II-2} & \eqref{eq:S-III.II-3}, \eqref{eq:S-III.II-2} & interference\\
Row III.III & \eqref{eq:H-III.III-3}, \eqref{eq:H-III.III-2}, \eqref{eq:H-III.III-2o} & \eqref{eq:S-III.III-3}, \eqref{eq:S-III.III-2}, \eqref{eq:S-III.III-2o} & $-2{\rm Re}\,\mathcal A_{CA}^\dagger\mathcal A_B$\\
Row III.V & \eqref{eq:H-III.V-3}, \eqref{eq:H-III.V-2} & \eqref{eq:S-III.V-3}, \eqref{eq:S-III.V-2} & interference\\
Row III.VI & \eqref{eq:H-III.VI-3}, \eqref{eq:H-III.VI-2} & \eqref{eq:S-III.VI-3}, \eqref{eq:S-III.VI-2} & interference\\
\bottomrule
\end{tabular}
\end{center}
For each row the equations are listed in the order sym, con, quad. The first two parts are present in the rows that are cubic in the charges (we also call them the first, or cubic, part of a row, and the genuine two-$\rho$ terms its second part); Rows~I.II, I.V, II.I and~III.III have in addition, and Row~II.IV has only, genuine two-$\rho$ terms. After the average over the projectile (Section~\ref{sec:fin_singlet}) the symmetrised part is proportional to the symmetrised three-point function $\sigma$ of the projectile, Eq.~\eqref{eq:A2d}, and is odd in $\v k$, while the contact and the genuine two-$\rho$ terms are proportional to the two-point function $n$ and are even in $\v k$. The prescriptions that define the finite terms are part of the result: the vanishing energy denominators of Row~I.III are principal values, and certain terms, whose large-distance tails cancel only in combination, have to be integrated together (Section~\ref{sec:fin_together}). The finite terms are given in coordinate space; no numerical evaluation of them, and no plot, is given in this paper.

\subsection{Conventions}\label{sec:fin_conv}

\paragraph{Phase and notation.} All terms are written with the phase $e^{-i\v k\cdot(\text{unprimed}-\text{primed})}$, with the conjugated factor at the primed points and to the left, e.g.\ $e^{-i\v k\cdot(\v w-\v w')}$. Where the position of the observed gluon has been integrated out, the phase contains a recoil factor: $e^{-i\v k\cdot(\v z-\v w)\xi}$ in Row~I.III, $e^{-i\v k\cdot(\v z-\v w')}e^{-i\v k\cdot(\v x-\v z)\xi}$ in Row~I.V, $e^{-i\v k\cdot(\v w-\v y')}e^{-i\v k\cdot(\v z-\v y')\xi}$ in Rows~III.III, III.V and~III.VI, and $e^{-i\v k\cdot(\v y-\v y')/\xi}$ in Row~II.IV and in the DGLAP term. The convention matters: a change of convention in a row amounts to $\v k\to-\v k$ there, which leaves the part of a term that is even in $\v k$ unchanged and flips the sign of its odd part, i.e.\ of the symmetrised three-$\rho$ part (Section~\ref{sec:fin_singlet}); mixing conventions between rows would give wrong relative signs of the three-$\rho$ terms. The complex conjugate is implied for Groups~I and~III and is written as ``$+{\rm c.c.}$''; Group~II has none. The charges $\rho$ are operators in the order written. In the formulas of this section and of Appendix~\ref{app:finite} transverse vectors are not printed in bold, $\int_x\equiv\int d^2x$, $\Kv(x)=x/x^2$ and $K(x,y,z)=\Kv(x-z)\cdot\Kv(y-z)$ as in Eq.~\eqref{eq:kernels}, and $\{\dots\}_{\rm sym}$ denotes the average of a product of charges over all orderings. The adjoint Wilson lines $U^{ab}$ are real, $U^\dagger=U^T$, and
\begin{align}\label{eq:TD}
(T^e)_{bc}=-if^{ebc}\,,\qquad (D^e)_{bc}=d^{ebc}\,,\qquad{\rm Tr}\,T^aT^b=N_c\,\delta^{ab}\,.
\end{align}

\paragraph{Longitudinal momentum fraction.} The longitudinal integrals are written in terms of a momentum fraction $\xi$, with two definitions. In all rows except Row~II.IV,
\begin{align}\label{eq:xi_def}
\xi=\frac{p^+}{k^+}\,,
\end{align}
where $p^+$ is the longitudinal momentum of the gluon that is integrated over: $0\le\xi\le1$ in Group~I, where this gluon is part of the observed one, and $0\le\xi<\infty$ in Group~III, where the upper limit $\vee-k^+$ has been sent to infinity. In Row~II.IV, as in the DGLAP term of Section~\ref{sec:dg_final},
\begin{align}\label{eq:xi_defIV}
\xi=\frac{k^+}{p^++k^+}
\end{align}
is the fraction of the longitudinal momentum of the parent gluon that is carried by the observed gluon; then $p^+=k^+\frac{1-\xi}{\xi}$, $dp^+=-k^+d\xi/\xi^2$, and the range $0\le p^+\le\vee-k^+$ is $k^+/\vee\le\xi\le1$. Every row has the overall factor $1/k^+$ of the leading-order spectrum.

\paragraph{Plus prescriptions.} The plus prescriptions act on the whole $\xi$-dependent function $f$ of the term in which they stand, i.e.\ on the recoil phase and on the energy denominators:
\begin{align}\label{eq:plus_xi}
\int_0^1d\xi\,\Big[\frac{1}{\xi}\Big]_+f(\xi)&=\int_0^1\frac{d\xi}{\xi}\,\big[f(\xi)-f(0)\big]\,,\notag\\
\int_0^1d\xi\,\Big[\frac{1}{1-\xi}\Big]_+f(\xi)&=\int_0^1\frac{d\xi}{1-\xi}\,\big[f(\xi)-f(1)\big]\,,\notag\\
\int_0^\infty d\xi\,\Big[\frac{1}{\xi\,(1+\xi)}\Big]_+f(\xi)&=\int_0^\infty\frac{d\xi}{\xi\,(1+\xi)}\,\big[f(\xi)-f(0)\big]\,.
\end{align}
Where a plus prescription at $\xi\to0$ arises from the subtraction of Eq.~\eqref{eq:fin_def}, the subtracted term is the term of the row in the JIMWLK list.

\paragraph{Instantaneous terms of Group III.} The terms of the kernel $\mathfrak K^{(y')}_{\rm III}$ and the kernel $\mathfrak K^{(z)}_{\rm III}$ below that carry no factor depending on $\xi$ other than the phase are defined as distributions,
\begin{align}\label{eq:inst_xi}
\int_0^\infty d\xi\,e^{-i\v k\cdot(\v z-\v y')\,\xi}=\pi\,\delta(s)-i\,\mathcal P\frac1s\,,\qquad s=\v k\cdot(\v z-\v y')\,,
\end{align}
which act on the transverse kernels, with the two instantaneous terms of a row (those with $\mathfrak K^{(y')}_{\rm III}$ and with $\mathfrak K^{(z)}_{\rm III}$) under one $\v y'$ integral, done before the $\xi$ integral (Appendix~\ref{app:tails}).

\paragraph{Limits and their order.} The limits $\Lambda\to0$ and $\vee\to\infty$ are taken after the logarithmic (BFKL and JIMWLK) terms have been taken out, Eq.~\eqref{eq:fin_def}; terms of order $\Lambda/k^+$ and $k^+/\vee$ are dropped. The limit $\vee\to\infty$ is taken before $\epsilon\to0$; this order is part of the eikonal approximation (Section~\ref{sec:dg_poles}). Row~II.IV is written with its lower limit $\xi=k^+/\vee$; the limit $\vee\to\infty$ of this term exists once the integral over $\v y-\v y'$ has been done (see below Eq.~\eqref{eq:S-II.IVrem}).

\paragraph{Vanishing energy denominators.} The energy denominators of $\mf B_3$, which enter Row~I.III and its complex conjugate, vanish inside the region of integration. They are defined as principal values, as in stationary perturbation theory (Appendix~\ref{app:distributions}); the finite terms of this paper refer to this prescription throughout. In the kernels of Row~I.III as printed the poles of the separate terms have cancelled, and all logarithms have positive arguments. An $i\varepsilon$ prescription appropriate to an asymptotic out-state would add to Row~I.III a finite term that is not included here.

\subsection{Ordered charges: symmetrised products and contact terms}\label{sec:fin_sym}

The rows that are cubic in the charges contain ordered products $\rho\rho\rho$, $\rho\{\rho,\rho\}$ or $\{\rho,\rho\}\rho$; the order is the one in which the coefficients of $\bar\Omega^\dagger\Omega$ and of the conjugate amplitude act, with the conjugate amplitude to the left, and is not a choice. With the commutator of the charges, Eq.~\eqref{eq:rho}, an ordered product of three charges is exactly the totally symmetrised product, plus single commutators, which are products of \emph{two} charges at coincident transverse positions, plus double commutators, which contain one charge. For any $c$-number colour tensor $C^{abc}$ multiplying the three charges (the Wilson lines are inside $C^{abc}$),
\begin{align}\label{eq:symord}
&C^{abc}\rho^{a}(p_{1})\rho^{b}(p_{2})\rho^{c}(p_{3})=C^{abc}\{\rho^{a}(p_{1})\rho^{b}(p_{2})\rho^{c}(p_{3})\}_{\rm sym}\notag\\
&\quad +\frac{i}{2}C^{abc}\Big[f^{abr}\delta^{(2)}(p_{1}-p_{2})\{\rho^{r}(p_{1})\rho^{c}(p_{3})\}_{\rm sym} +f^{bcr}\delta^{(2)}(p_{2}-p_{3})\{\rho^{a}(p_{1})\rho^{r}(p_{2})\}_{\rm sym}\notag\\
&\quad\quad +f^{acr}\delta^{(2)}(p_{1}-p_{3})\{\rho^{r}(p_{1})\rho^{b}(p_{2})\}_{\rm sym}\Big] -\frac{1}{6}C^{abc}\big[f^{abr}f^{rcs}\notag\\
&\quad\quad -f^{bcr}f^{ras}\big]\delta^{(2)}(p_{1}-p_{2})\, \delta^{(2)}(p_{1}-p_{3})\, \rho^{s}(p_{1})\,,
\end{align}

\begin{align}\label{eq:symr12}
&C^{abc}\rho^{a}(p_{1})\{\rho^{b}(p_{2}),\rho^{c}(p_{3})\}=2C^{abc}\{\rho^{a}(p_{1})\rho^{b}(p_{2})\rho^{c}(p_{3})\}_{\rm sym}\notag\\
&\quad +i\, C^{abc}\Big[f^{abr}\delta^{(2)}(p_{1}-p_{2})\{\rho^{r}(p_{1})\rho^{c}(p_{3})\}_{\rm sym} +f^{acr}\delta^{(2)}(p_{1}-p_{3})\{\rho^{r}(p_{1})\rho^{b}(p_{2})\}_{\rm sym}\Big]\notag\\
&\quad -\frac{1}{6}C^{abc}\big[f^{abr}f^{rcs} +f^{acr}f^{rbs}\big]\delta^{(2)}(p_{1}-p_{2})\, \delta^{(2)}(p_{1}-p_{3})\, \rho^{s}(p_{1})\,,
\end{align}

\begin{align}\label{eq:symr21}
&C^{abc}\{\rho^{a}(p_{1}),\rho^{b}(p_{2})\}\rho^{c}(p_{3})=2C^{abc}\{\rho^{a}(p_{1})\rho^{b}(p_{2})\rho^{c}(p_{3})\}_{\rm sym}\notag\\
&\quad +i\, C^{abc}\Big[f^{bcr}\delta^{(2)}(p_{2}-p_{3})\{\rho^{a}(p_{1})\rho^{r}(p_{2})\}_{\rm sym} +f^{acr}\delta^{(2)}(p_{1}-p_{3})\{\rho^{r}(p_{1})\rho^{b}(p_{2})\}_{\rm sym}\Big]\notag\\
&\quad +\frac{1}{6}C^{abc}\big[f^{bcr}f^{ras} +f^{acr}f^{rbs}\big]\delta^{(2)}(p_{1}-p_{2})\, \delta^{(2)}(p_{1}-p_{3})\, \rho^{s}(p_{1})\,.
\end{align}

Accordingly every cubic row is written as two equations: its \emph{symmetrised three-$\rho$ term}, and its \emph{two-$\rho$ terms from the symmetrisation} (contact terms), in which the $\delta$-functions put two of the three charges, and with them the corresponding Wilson lines of the colour tensor, at the same point. The one-$\rho$ terms, the last lines of Eqs.~\eqref{eq:symord}--\eqref{eq:symr21}, are not written in the rows; they vanish after the average over a projectile that is invariant under global colour rotations, $\langle\rho^s\rangle=0$, so that the averaged form of Section~\ref{sec:fin_rows} is complete, while as operator identities the rows of Appendix~\ref{app:finite} equal the ordered rows only together with them. The symmetrised term is the same operator for every ordering of the three charges; the ordering enters only through the contact terms. The contact terms are of the same order in the charges as the genuine two-$\rho$ terms, are present for any projectile, and are not removed by colour neutrality.

\subsection{Average over the projectile and colour-singlet form}\label{sec:fin_singlet}

\paragraph{Assumptions.} We assume
\begin{itemize}
\item[(A1)] \emph{Global colour invariance:} the projectile density matrix commutes with the total charge $Q^a=\int_x\rho^a(\v x)$. Then $\langle\rho^a(\v x)\rangle=0$, and the two-point function is colour diagonal,
\begin{align}\label{eq:A1n}
\big\langle\{\rho^a(\v x)\rho^b(\v y)\}_{\rm sym}\big\rangle=\frac{\delta^{ab}}{N_c^2-1}\,n(\v x,\v y)\,,\qquad n(\v x,\v y)\equiv\big\langle\rho^f(\v x)\rho^f(\v y)\big\rangle\,,
\end{align}
as in Eq.~\eqref{eq:singlet}, with $n(\v x,\v y)=n(\v y,\v x)$ real; the symmetry follows from $\langle[\rho^a(\v x),\rho^b(\v y)]\rangle=0$.
\item[(A2)] \emph{Symmetrised three-point function of $d$-type:}
\begin{align}\label{eq:A2d}
\big\langle\{\rho^a(\v x)\rho^b(\v y)\rho^c(\v z)\}_{\rm sym}\big\rangle&=d^{abc}\,\frac{N_c}{(N_c^2-4)(N_c^2-1)}\,\sigma(\v x,\v y,\v z)\,,\notag\\
\sigma(\v x,\v y,\v z)&\equiv\big\langle d^{efg}\{\rho^e(\v x)\rho^f(\v y)\rho^g(\v z)\}_{\rm sym}\big\rangle\,.
\end{align}
\end{itemize}
The normalisation follows from $d^{abc}d^{abc}=(N_c^2-4)(N_c^2-1)/N_c$; $\sigma$ is real and totally symmetric in its arguments. By (A1) the symmetrised three-point function is of the form $d^{abc}g_d+f^{abc}g_f$ with $g_d$ totally symmetric and $g_f$ totally antisymmetric in the three points; (A2) states that $g_f=0$ and gives $g_d$ the normalisation of Eq.~\eqref{eq:A2d}. The function $\sigma$ is odd under charge conjugation of the projectile: it is the three-point function of the cubic Casimir density, the odderon of the projectile. No model for $\sigma$ is assumed: $\sigma$ is a general real, totally symmetric function of three points, and all results below hold for any such $\sigma$. For $SU(2)$, $d^{abc}=0$ and the terms with $\sigma$ are absent. Nothing else is assumed; in particular no assumption on the four-point function is needed, since the rows with four charges contribute only to the evolution terms.

\paragraph{Quadratic and cubic terms.} For a genuine two-$\rho$ term, $\langle C^{ab}\rho^a(p)\rho^b(q)\rangle=C^{aa}\,n(p,q)/(N_c^2-1)$. For a cubic term, Eqs.~\eqref{eq:symord}--\eqref{eq:symr21} with (A1) and (A2) give, with $\Phi_{ij}\equiv C^{abc}f^{abc}$ evaluated with the Wilson lines at $p_i=p_j$,
\begin{align}\label{eq:cubicavg}
&\langle C^{abc}\rho^a(p_1)\rho^b(p_2)\rho^c(p_3)\rangle=\frac{N_c}{(N_c^2-4)(N_c^2-1)}\,C^{abc}d^{abc}\,\sigma(p_1,p_2,p_3)\notag\\&\quad+\frac{i}{2(N_c^2-1)}\Big[\Phi_{12}\,\delta^{(2)}(p_1-p_2)\,n(p_1,p_3)+\Phi_{23}\,\delta^{(2)}(p_2-p_3)\,n(p_1,p_2)\notag\\&\qquad\qquad-\Phi_{13}\,\delta^{(2)}(p_1-p_3)\,n(p_1,p_2)\Big]\,,\notag\\
&\langle C^{abc}\rho^a(p_1)\{\rho^b(p_2),\rho^c(p_3)\}\rangle=\frac{2N_c}{(N_c^2-4)(N_c^2-1)}\,C^{abc}d^{abc}\,\sigma(p_1,p_2,p_3)\notag\\&\quad+\frac{i}{N_c^2-1}\Big[\Phi_{12}\,\delta^{(2)}(p_1-p_2)\,n(p_1,p_3)-\Phi_{13}\,\delta^{(2)}(p_1-p_3)\,n(p_1,p_2)\Big]\,,\notag\\
&\langle C^{abc}\{\rho^a(p_1),\rho^b(p_2)\}\rho^c(p_3)\rangle=\frac{2N_c}{(N_c^2-4)(N_c^2-1)}\,C^{abc}d^{abc}\,\sigma(p_1,p_2,p_3)\notag\\&\quad+\frac{i}{N_c^2-1}\Big[\Phi_{23}\,\delta^{(2)}(p_2-p_3)\,n(p_1,p_2)-\Phi_{13}\,\delta^{(2)}(p_1-p_3)\,n(p_1,p_2)\Big]\,.
\end{align}
Every cubic row therefore has two parts after the average, and both are part of the result: the symmetrised three-$\rho$ part, proportional to $\sigma$, with the factor $t\,i$ in the equations below, $t=1$ for an ordered product and $t=2$ for a product with an anticommutator, and the contact part, two-point terms $\propto\delta^{(2)}(p_i-p_j)\,n$ with the coefficients $(\pm\frac{i}{2})$, resp.\ $(\pm i)$, of Eq.~\eqref{eq:cubicavg}. Since charge conjugation of the projectile reverses $\sigma$ and leaves $n$ unchanged, the two parts are the parts of the cubic terms that are odd and even under charge conjugation of the projectile. Section~\ref{sec:check_cubic} shows the split on a single quark and on a dipole, computed exactly in charge space.

\paragraph{Parity in $\v k$.} Write a term as $A(\v k)\langle O\rangle$, with $A$ the $c$-number factor and $O$ the colour operator. All kernels depend on $\v k$ through the phases, including the recoil phases under the $\xi$ integrals, and through the distribution~\eqref{eq:inst_xi}, which turn into their complex conjugates under $\v k\to-\v k$; the genuine two-$\rho$ rows have real prefactors and the cubic rows an explicit $i$ (in Row~I.II the $i$ is in the Fourier transform of the odd function $R^i$). Hence $A(-\v k)=A(\v k)^*$ for the genuine two-$\rho$ and $A(-\v k)=-A(\v k)^*$ for the cubic terms; for Row~II.I, which has no complex conjugate, the same holds for the sum of its two terms. With $\langle O\rangle=\mathsf S+i\mathsf C$, $\mathsf S$ the symmetrised and $\mathsf C$ the contact part, both real, a cubic term plus its complex conjugate is $\mathsf S[A(\v k)-A(-\v k)]+i\mathsf C[A(\v k)+A(-\v k)]$. The symmetrised part of the cubic terms is odd under $\v k\to-\v k$, their contact part is even, and the genuine two-$\rho$ terms are even. The part of the finite terms that is odd in $\v k$ is therefore exactly the symmetrised three-$\rho$ part, which is proportional to the projectile three-point function $\sigma$ and to C-odd structures of the target (see below). This refines the statement that the cubic, ``first saturation'' terms are odd in $\v k$~\cite{McLerran:2016snu,Kovchegov:2018jun}: it holds for their symmetrised part, configuration by configuration, but not for the commutator remainders of the ordered products.

\paragraph{Colour structures.} In all rows the colour tensors are written as sums of products of \emph{differences} of Wilson lines, as in the LO structure $\big[U^{ce'}(y')-U^{ce'}(x')\big]\big[U^{ce}(y)-U^{ce}(x)\big]$. The rewriting uses only the invariance of the structure constants, $f^{gde}U^{db}(p)U^{ec}(p)=U^{gh}(p)f^{hbc}$, the orthogonality $U^{mh}(p)U^{kh}(p)=\delta^{mk}$ and $f^{acd}f^{bcd}=N_c\delta^{ab}$, and is exact for any $N_c$; for instance,
\begin{align}\label{eq:factex}
&f^{hbc}U^{gh}(w)-f^{gde}U^{db}(x)U^{ec}(y)\notag\\&\quad=-f^{gde}\Big\{\big[U^{db}(x)-U^{db}(w)\big]U^{ec}(y)+U^{db}(w)\big[U^{ec}(y)-U^{ec}(w)\big]\Big\}\,.
\end{align}
In the averaged form the contractions $C^{abc}d^{abc}$ of the symmetrised part and $C^{abc}f^{abc}$ of the contact part are written as traces in the adjoint representation. For matrices $B$, $C$ and an orthogonal matrix $O$,
\begin{align}\label{eq:traceid}
f^{ijk}d^{lmn}\,O^{il}B^{jm}C^{kn}&=i\,{\rm Tr}\big\{T^hO^\dagger C\,D^hB^\dagger O\big\}\,,\notag\\
f^{ijk}f^{lmn}\,O^{il}B^{jm}C^{kn}&={\rm Tr}\big\{T^eO^\dagger C\,T^eB^\dagger O\big\}\,,
\end{align}
where $B$ and $C$ are the differences of Wilson lines and $O$ is a Wilson line or a product of Wilson lines, e.g.\ ${\rm Tr}\big\{T^hU^\dagger(x)\big(U(z)-U(y)\big)D^h\big(U^\dagger(x')-U^\dagger(w')\big)U(x)\big\}$. The traces with $D^h$ are imaginary; $i$ times such a trace is real and is a C-odd (odderon) structure of the target: it changes sign when the fundamental Wilson lines, $U^{ab}=2\,{\rm tr}(t^aVt^bV^\dagger)$, are replaced by their complex conjugates, it starts at third order in the target field for a weak target, and it vanishes for $SU(2)$. The traces with $T^e$ and the traces of two differences are C-even.

The factorisation into differences is complete in the following sense.
\begin{enumerate}
\item A three-$\rho$ colour tensor $C^{abc}$ has three adjoint indices, one for each charge. In every term exactly two of the three colour legs carry a difference of Wilson lines; the third leg carries one Wilson line $U(p)$ (or a product $U^\dagger(z)U(\dots)$, or the unit matrix). This third Wilson line cannot be turned into a difference: writing $U(p)=[U(p)-U(q)]+U(q)$ adds a term that does not vanish. The three-$\rho$ structures therefore have the form $[U-U]\times[U-U]\times U\times f$ and, after the average, are traces with two differences and one Wilson line, the matrix $O$ of Eq.~\eqref{eq:traceid}.
\item In Row~II.I, in Rows~III.II and~III.III (genuine two-$\rho$ terms) and in the contact terms of Rows~II.I and~III.II the factorisation produces products $U^\dagger(z)\,[U(z)-U(y)]$ and $[U^\dagger(w)-U^\dagger(z)]\,U(x')$, i.e.\ a difference times a Wilson line in one leg; this is the closest exact form, for the same reason.
\item In the three-$\rho$ part of Row~I.II the two traces of the factorised form add up to a shorter expression with one difference, since the term $f^{hbc}U^{gh}(w)$ of the tensor~\eqref{eq:factex} drops out of $C^{abc}d^{abc}$ ($f^{hbc}d^{abc}=0$); the factorised form is kept for uniformity, and both are exact.
\item The genuine two-$\rho$ structures and all contact structures are sums of products of two differences, with $f$, $N_c$ and one Wilson line; nothing is left unfactorised.
\end{enumerate}
Every term of every row therefore vanishes for $U=1$, and for any Wilson line that does not depend on the transverse position.

\paragraph{The LO structure.} In Row~II.IV, and in the DGLAP and running-coupling terms~\eqref{eq:DGLAP}, \eqref{eq:RC} and in Eq.~\eqref{eq:LORC}, the colour factor is that of the LO spectrum, and the average is the replacement
\begin{align}\label{eq:sLOrepl}
&\big[U^{ce'}(y')-U^{ce'}(x')\big]\big[U^{ce}(y)-U^{ce}(x)\big]\rho^{e'}(x')\rho^e(x)\notag\\&\to\;\frac{n(x',x)}{N_c^2-1}\,{\rm Tr}\big\{\big(U^\dagger(y')-U^\dagger(x')\big)\big(U(y)-U(x)\big)\big\}\,.
\end{align}
In the labelling of the BFKL term ($x'\to y$, $y\to w$, $y'\to w'$) the trace is $\mathcal O_{\rm LO}$ of Eq.~\eqref{eq:OLO}.

\subsection{Exact cancellations, and terms that must be integrated together}\label{sec:fin_together}

\paragraph{Exact cancellations.} Two cancellations among the finite terms are exact, i.e.\ operator identities at the level of the integrand, and have been carried out: the finite part of Row~I.IV cancels that of Row~III.I, and the terms proportional to $\frac{11}{3}$ of Rows~I.II and~I.V, which carry the ultraviolet poles of the virtual corrections, cancel each other (Section~\ref{sec:dg_poles}). All other finite terms remain and are listed below.

\paragraph{Terms that must be integrated together.} Several of the remaining terms have a tail $\propto1/\v z^2$ at large transverse distance $\v z$ of the integrated gluon, which alone would make the $\v z$ integral diverge. These tails cancel between terms of different rows, but only the tails cancel, not the terms; none of them is dropped, and the terms concerned have to be combined under one integral. They are:
\begin{itemize}
\item \emph{Pair A:} the term of Row~I.III with the difference $\big[U(w)-U(y)\big]$, i.e.\ the term with the kernel $\mathfrak K^{(w)}_{\rm I.III}$, with Row~III.V.
\item \emph{Pair B:} the terms of the cubic part of Row~I.V with $\big[\frac1\xi\big]_+$, kernel $\mathfrak K_{\rm I.V}$, with the term of the cubic part of Row~III.III with $\big[\frac{1}{\xi(1+\xi)}\big]_+$, kernel $\mathfrak K^{(y')}_{\rm III}$.
\item \emph{Pair C:} the terms of the genuine two-$\rho$ part of Row~I.V with $\big[\frac1\xi\big]_+$ and $\big[\frac{1}{1-\xi}\big]_+$, kernel $\mathfrak K^{\rm quad}_{\rm I.V}$, with the term of the genuine two-$\rho$ part of Row~III.III with $\big[\frac{1}{\xi(1+\xi)}\big]_+$.
\item \emph{Pair D:} the part $\int_0^1d\xi\big[\frac1\xi\big]_+e^{-ik\cdot(z-w)\xi}$ of the term of Row~I.III with $\big[U(z)-U(y)\big]$, i.e.\ the last term of the kernel $\mathfrak K^{(z)}_{\rm I.III}$, with Row~III.VI.
\end{itemize}
Each pair stands under one $\v z$ integral, with $\v y'$ renamed $\v w'$ in the Group~III member (and $\v x\leftrightarrow\v y$ in pair~D). The members of a pair are a virtual term (Group~I, the integrated gluon is absorbed by the observed gluon in the amplitude, recoil phase with $\v w$) and a real term (Group~III, the integrated gluon is emitted by the observed gluon in the conjugate amplitude, recoil phase with $\v w'$) with the same coefficient, the same ordered colour operator and the same transverse kernel up to $\v w\to\v w'$; the cancellation of their tails holds for every value of $U(\v z)$ and requires no condition on the target at large distance. In Row~I.III, moreover, the remaining terms of the kernel $\mathfrak K^{(z)}_{\rm I.III}$, i.e.\ its first three groups of terms and the term with $-\frac12\log\frac{(z-w)^2}{(x-w)^2}-\int_0^1d\xi\big[\frac1\xi\big]_+e^{-ik\cdot(z-w)\xi}$, are integrable only with the $\v w$ integral, the Fourier transform at the measured momentum, done before the $\v z$ integral. Finally, Row~II.IV, the genuine two-$\rho$ part of Row~II.I and the genuine two-$\rho$ part of Row~III.III have to be integrated together over the position of the pair of final-state gluons: in the region in which both gluons are far from the projectile their tails cancel with the weights $1:1:-2$ (Section~\ref{sec:dg_tail}). The origin of these statements, the closed forms of the longitudinal integrals involved and the checks are collected in Appendix~\ref{app:tails}.

\paragraph{Row I.II: the momentum-space kernel.} The cubic part of Row~I.II is given as a Fourier integral of a momentum-space kernel, $\mathfrak K_{\rm I.II}$ of Eq.~\eqref{eq:K-rm_I_II}. This kernel is the antisymmetrised one, $\frac12\big[R^i(\v q,\v p)/\v p^2-R^i(\v q,\v q-\v p)/(\v q-\v p)^2\big]$: each of its two parts alone behaves as $\log(1/r)/r^2$ at $\v p\to\v q$ ($r=|\v p-\v q|$) and at $\v p\to0$ ($r=|\v p|$), which is not integrable over $d^2p$, while in the antisymmetrised combination these terms cancel and the kernel is $O(\log(1/r)/r)$, which is integrable. The antisymmetrised form is equivalent to the form with $R^i(\v q,\v p)/\v p^2$ alone, since the colour structure of the row is antisymmetric under $\v x\leftrightarrow\v y$, which under the phases $e^{iq\cdot(w-y)}e^{ip\cdot(y-x)}$ is $\v p\to\v q-\v p$; the antisymmetrised form makes explicit the value that the scaleless rule $\int d^2p/\v p^2=0$ of dimensional regularisation assigns to these points.

\paragraph{Row I.V: the collinear point.} In the genuine two-$\rho$ part of Row~I.V the kernel $\mathfrak K^{\rm quad}_{\rm I.V}$ is singular at the collinear point $\v z\to\v x$. There the term is finite only together with its $\frac{11}{6}$ line, the second term of Eq.~\eqref{eq:S-I.V-2o}, and with a rotation-invariant treatment of the point $\v z=\v x$, i.e.\ with the angular average around this point taken first: the sum of the two behaves as a pure second harmonic in the angle of $\v z-\v x$, divided by $(\v z-\v x)^2$. The rotation-invariant treatment is the one consistent with the $\frac{11}{6}$ term, which comes from a dimensionally regularised, rotation-invariant integral.

\subsection{The finite terms in colour-singlet form}\label{sec:fin_rows}

We now list the finite terms, averaged over the projectile according to Eqs.~\eqref{eq:A1n}--\eqref{eq:cubicavg}. The corresponding expressions before the average are given in Appendix~\ref{app:finite}; the two lists have the same prefactors, kernels and phases. Blocks of transverse and longitudinal kernels that occur in several equations are given names, $\mathfrak K$ with the row as a subscript. A kernel depends on the transverse integration variables of the equation in which it stands, on $\v k$ and, where it stands under a $\xi$ integral, on $\xi$; a plus prescription inside a kernel acts on everything in the term that depends on $\xi$, including a recoil phase written in front of the kernel, according to Eq.~\eqref{eq:plus_xi}. In the contact terms the $\delta$-functions are to be used to set the charge positions equal in the kernel as well.

\subsubsection{Group I}

\paragraph{Row I.II.} With the kernel
\begin{align}\label{eq:K-rm_I_II}
&\mathfrak{K}_{\rm I.II}\equiv \int d^{2}p\, d^{2}q\;e^{iq\cdot(w-y)}\, e^{ip\cdot(y-x)}\, \frac{(x'-w')^{i}}{(x'-w')^{2}\,\mathbf{q}^{2}}\, \frac{1}{2}\bigg[\frac{R^{i}(\mathbf{q},\mathbf{p})}{\mathbf{p}^{2}} -\frac{R^{i}(\mathbf{q},\mathbf{q}-\mathbf{p})}{(\mathbf{q}-\mathbf{p})^{2}}\bigg]
\end{align}

where
\begin{align}\label{eq:Rdef}
&R^{i}(\mathbf{q},\mathbf{p})=-2q^{i}\log\left(\frac{\mathbf{p}^{2}}{\mathbf{q}^{2}}\right) +\frac{(2\mathbf{p}^{2}-2\,\mathbf{q}\cdot\mathbf{p}+\mathbf{q}^{2})\,q^{i}+\mathbf{q}^{2}p^{i}}{(\mathbf{q}-\mathbf{p})^{2}}\, \log\left[\frac{\mathbf{q}^{2}}{(\mathbf{q}-\mathbf{p})^{2}}\right]\,.
\end{align}

the symmetrised three-$\rho$ part is
\begin{align}\label{eq:S-I.II-3}
&\Big\langle\frac{d^{3}N^{\rm I.II}_{{\rm sym}}}{d^{3}k}\Big\rangle=-\frac{1}{(2\pi)^{3}}\, \frac{g^{4}}{64\pi^{6}}\, \frac{1}{k^{+}} \int_{w,w'}e^{-ik\cdot(w-w')} \int_{x,y,x'} \times \mathfrak{K}_{\rm I.II}\notag\\
&\quad \times \frac{N_{c}}{(N_{c}^{2}-4)(N_{c}^{2}-1)}\, \Big(+2i\Big)\notag\\
&\quad \times \sigma(x',x,y)\Big[-{\rm Tr}\Big\{T^{h}U^{\dagger}(y)\big(U(x)-U(w)\big)D^{h}\big(U^{\dagger}(x')-U^{\dagger}(w')\big)U(y)\Big\}\notag\\
&\quad\quad -{\rm Tr}\Big\{T^{h}U^{\dagger}(w)\big(U(x')-U(w')\big)D^{h}\big(U^{\dagger}(y)-U^{\dagger}(w)\big)U(w)\Big\}\Big] +{\rm c.c.}
\end{align}

the contact part is
\begin{align}\label{eq:S-I.II-2}
&\Big\langle\frac{d^{3}N^{\rm I.II}_{{\rm con}}}{d^{3}k}\Big\rangle=-\frac{1}{(2\pi)^{3}}\, \frac{g^{4}}{64\pi^{6}}\, \frac{1}{k^{+}} \int_{w,w'}e^{-ik\cdot(w-w')} \int_{x,y,x'} \times \mathfrak{K}_{\rm I.II}\notag\\
&\quad \times \frac{1}{N_{c}^{2}-1}\bigg[\Big(+i\Big)\delta^{(2)}(x'-x)\notag\\
&\quad\quad \times n(x,y)\Big[-{\rm Tr}\Big\{T^{e}U^{\dagger}(y)\big(U(x)-U(w)\big)T^{e}\big(U^{\dagger}(x)-U^{\dagger}(w')\big)U(y)\Big\}\notag\\
&\quad\quad\quad -{\rm Tr}\Big\{T^{e}U^{\dagger}(w)\big(U(y)-U(w)\big)T^{e}\big(U^{\dagger}(x)-U^{\dagger}(w')\big)U(w)\Big\}\Big]\notag\\
&\quad\quad +\Big(-i\Big)\delta^{(2)}(x'-y)\notag\\
&\quad\quad \times n(y,x)\Big[-{\rm Tr}\Big\{T^{e}U^{\dagger}(y)\big(U(x)-U(w)\big)T^{e}\big(U^{\dagger}(y)-U^{\dagger}(w')\big)U(y)\Big\}\notag\\
&\quad\quad\quad -{\rm Tr}\Big\{T^{e}U^{\dagger}(w)\big(U(y)-U(w)\big)T^{e}\big(U^{\dagger}(y)-U^{\dagger}(w')\big)U(w)\Big\}\Big]\bigg] +{\rm c.c.}
\end{align}

and the genuine two-$\rho$ term, which is proportional to the LO structure, is
\begin{align}\label{eq:S-I.II-2o}
&\Big\langle\frac{d^{3}N^{\rm I.II}_{{\rm quad}}}{d^{3}k}\Big\rangle=+\frac{1}{(2\pi)^{3}}\, \frac{g^{4}\,N_{c}}{32\pi^{4}}\, \frac{1}{k^{+}}\Big(-\frac{67}{9}+\frac{\pi^{2}}{3}\Big)\notag\\
&\quad \times \int_{w,w',x,x'} e^{-ik\cdot(w-w')}\;\Kv(x'-w')\cdot\Kv(x-w)\notag\\
&\quad \times\frac{n(x',x)}{N_{c}^{2}-1}\Big[{\rm Tr}\Big\{\big(U^{\dagger}(x')-U^{\dagger}(w')\big)\big(U(w)-U(x)\big)\Big\}\Big] +{\rm c.c.}
\end{align}

The constant $-\frac{67}{9}+\frac{\pi^2}{3}$ is the one left in this row by the cancellation of the $\frac{11}{3}$ terms against Row~I.V; it is a property of this grouping of the terms and not a constant of the result (Section~\ref{sec:lit_constant}).

\paragraph{Row I.III.} The kernels are
\begin{align}\label{eq:K-z_rm_I_III}
&\mathfrak{K}^{(z)}_{\rm I.III}\equiv \bigg[\frac{\Kv(x'-w')\cdot\Kv(y-z)}{2\,(z-w)^{2}} +\Kv(x'-w')\cdot\Kv(z-w)\, \Kv(y-z)\cdot\Kv(x-z)\notag\\
&\quad\quad -\Kv(x'-w')\cdot\Kv(x-w)\, \Kv(y-z)\cdot\Kv(z-w)\notag\\
&\quad\quad -\Kv(x'-w')\cdot\Kv(x-w)\, \Kv(y-z)\cdot\Kv(x-z)\bigg]\log\frac{(x-z)^{2}}{(x-w)^{2}}\notag\\
&\quad +\frac{\arccos\left(\frac{(x-w)^{2}-(x-w)\cdot (x-z)}{\sqrt{(x-w)^{2}\,(z-w)^{2}}}\right)}{\sqrt{(x-w)^{2}\,(x-z)^{2}-\left[(x-w)\cdot (x-z)\right]^{2}}\,(z-w)^{2}}\, \bigg[\Kv(x'-w')\cdot(z-w)\notag\\
&\quad\quad \times \Kv(y-z)\cdot(x-w) -\frac{(x-w)\cdot(x-z)}{(x-z)^{2}}\, \Kv(x'-w')\cdot(z-w)\, \Kv(y-z)\cdot(x-z)\notag\\
&\quad\quad +\Kv(x'-w')\cdot(x-z)\, \Kv(y-z)\cdot(z-w) -\frac{(x-w)\cdot(x-z)}{(x-w)^{2}}\, \Kv(x'-w')\cdot(x-w)\notag\\
&\quad\quad \times \Kv(y-z)\cdot(z-w)\bigg] -\frac{1}{2\,(z-w)^{2}}\log\frac{(z-w)^{2}}{(x-w)^{2}}\, \Kv(x'-w')\cdot(z-w)\notag\\
&\quad \times \Kv(y-z)\cdot\Kv(x-z) +\bigg[-\frac{1}{2}\log\frac{(z-w)^{2}}{(x-w)^{2}} -\notag\\
&\quad\quad \times \int_{0}^{1}d\xi\Big[\frac{1}{\xi}\Big]_{+}e^{-ik\cdot(z-w)\,\xi}\bigg]\Kv(y-z)\cdot\Kv(w-z)\, \Kv(x'-w')\cdot\Kv(x-w)\notag\\
&\quad + \int_{0}^{1}d\xi\Big[\frac{1}{\xi}\Big]_{+}e^{-ik\cdot(z-w)\,\xi}\, \Kv(y-z)\cdot\Kv(w-z)\, \Kv(x'-w')\cdot\Kv(x-w)
\end{align}

and
\begin{align}\label{eq:K-w_rm_I_III}
&\mathfrak{K}^{(w)}_{\rm I.III}\equiv - \int_{0}^{1}d\xi\Big[\frac{1}{\xi}\Big]_{+}e^{-ik\cdot(z-w)\,\xi}\, \Kv(x-z)\cdot\Kv(w-z)\, \Kv(x'-w')\cdot\Kv(y-w)\,.
\end{align}

The last two terms of $\mathfrak K^{(z)}_{\rm I.III}$ are written separately, although the integrals over $\xi$ cancel between them, because they are integrated differently: the last term is the Group~I member of pair~D, while the term before it belongs with the rest of the kernel (Section~\ref{sec:fin_together}). The symmetrised three-$\rho$ part of the row is
\begin{align}\label{eq:S-I.III-3}
&\Big\langle\frac{d^{3}N^{\rm I.III}_{{\rm sym}}}{d^{3}k}\Big\rangle=\frac{1}{(2\pi)^{3}}\, \frac{ig^{4}}{8\pi^{5}}\, \frac{1}{k^{+}} \int_{w,w'}e^{-ik\cdot(w-w')} \int_{x,y,z,x'}\Bigg[ \mathfrak{K}^{(z)}_{\rm I.III}\notag\\
&\quad\quad \times \frac{N_{c}}{(N_{c}^{2}-4)(N_{c}^{2}-1)}\, \Big(+i\Big)\notag\\
&\quad\quad \times \sigma(x',x,y)\Big[{\rm Tr}\Big\{T^{h}U^{\dagger}(x)\big(U(z)-U(y)\big)D^{h}\big(U^{\dagger}(x')-U^{\dagger}(w')\big)U(x)\Big\}\Big] + \mathfrak{K}^{(w)}_{\rm I.III}\notag\\
&\quad\quad \times \frac{N_{c}}{(N_{c}^{2}-4)(N_{c}^{2}-1)}\, \Big(+i\Big)\notag\\
&\quad\quad \times \sigma(x',x,y)\Big[{\rm Tr}\Big\{T^{h}U^{\dagger}(x)\big(U(w)-U(y)\big)D^{h}\big(U^{\dagger}(x')-U^{\dagger}(w')\big)U(x)\Big\}\Big]\Bigg] +{\rm c.c.}
\end{align}

and the contact part is
\begin{align}\label{eq:S-I.III-2}
&\Big\langle\frac{d^{3}N^{\rm I.III}_{{\rm con}}}{d^{3}k}\Big\rangle=\frac{1}{(2\pi)^{3}}\, \frac{ig^{4}}{8\pi^{5}}\, \frac{1}{k^{+}} \int_{w,w'}e^{-ik\cdot(w-w')} \int_{x,y,z,x'}\Bigg[ \mathfrak{K}^{(z)}_{\rm I.III}\notag\\
&\quad\quad \times \frac{1}{N_{c}^{2}-1}\bigg[\Big(+\frac{1}{2}i\Big)\delta^{(2)}(x'-x)\notag\\
&\quad\quad\quad \times n(x,y)\Big[-{\rm Tr}\Big\{T^{e}U^{\dagger}(x)\big(U(z)-U(y)\big)T^{e}\big(U^{\dagger}(x)-U^{\dagger}(w')\big)U(x)\Big\}\Big]\notag\\
&\quad\quad\quad +\Big(+\frac{1}{2}i\Big)\delta^{(2)}(y-x)\notag\\
&\quad\quad\quad \times n(x,x')\Big[-{\rm Tr}\Big\{T^{e}U^{\dagger}(x)\big(U(z)-U(x)\big)T^{e}\big(U^{\dagger}(x')-U^{\dagger}(w')\big)U(x)\Big\}\Big]\notag\\
&\quad\quad\quad +\Big(-\frac{1}{2}i\Big)\delta^{(2)}(x'-y)\notag\\
&\quad\quad\quad \times n(y,x)\Big[-{\rm Tr}\Big\{T^{e}U^{\dagger}(x)\big(U(z)-U(y)\big)T^{e}\big(U^{\dagger}(y)-U^{\dagger}(w')\big)U(x)\Big\}\Big]\bigg] + \mathfrak{K}^{(w)}_{\rm I.III}\notag\\
&\quad\quad \times \frac{1}{N_{c}^{2}-1}\bigg[\Big(+\frac{1}{2}i\Big)\delta^{(2)}(x'-x)\notag\\
&\quad\quad\quad \times n(x,y)\Big[-{\rm Tr}\Big\{T^{e}U^{\dagger}(x)\big(U(w)-U(y)\big)T^{e}\big(U^{\dagger}(x)-U^{\dagger}(w')\big)U(x)\Big\}\Big]\notag\\
&\quad\quad\quad +\Big(+\frac{1}{2}i\Big)\delta^{(2)}(y-x)\notag\\
&\quad\quad\quad \times n(x,x')\Big[-{\rm Tr}\Big\{T^{e}U^{\dagger}(x)\big(U(w)-U(x)\big)T^{e}\big(U^{\dagger}(x')-U^{\dagger}(w')\big)U(x)\Big\}\Big]\notag\\
&\quad\quad\quad +\Big(-\frac{1}{2}i\Big)\delta^{(2)}(x'-y)\notag\\
&\quad\quad\quad \times n(y,x)\Big[-{\rm Tr}\Big\{T^{e}U^{\dagger}(x)\big(U(w)-U(y)\big)T^{e}\big(U^{\dagger}(y)-U^{\dagger}(w')\big)U(x)\Big\}\Big]\bigg]\Bigg] +{\rm c.c.}
\end{align}

The row has no genuine two-$\rho$ term.

\paragraph{Row I.V.} The kernel of the cubic part is
\begin{align}\label{eq:K-rm_I_V}
&\mathfrak{K}_{\rm I.V}\equiv e^{-ik\cdot(x-z)\,\xi}\, \Kv(x'-w')\cdot\Kv(z-x)\, \Kv(y-z)\cdot\Kv(y'-x)\notag\\
&\quad -\Big[\frac{1}{\xi}\Big]_{+}e^{-ik\cdot(x-z)\,\xi}\, \Kv(x'-w')\cdot\Kv(y-z)\, \Kv(z-x)\cdot\Kv(y'-x)\notag\\
&\quad -\Big[\frac{1}{\xi}\Big]_{+}e^{-ik\cdot(x-z)}e^{-ik\cdot(z-x)\,\xi}\, \Kv(x'-w')\cdot\Kv(y'-x)\, \Kv(z-x)\cdot\Kv(y-z)
\end{align}

The symmetrised three-$\rho$ part is
\begin{align}\label{eq:S-I.V-3}
&\Big\langle\frac{d^{3}N^{\rm I.V}_{{\rm sym}}}{d^{3}k}\Big\rangle=\frac{1}{(2\pi)^{3}}\, \frac{ig^{4}}{16\pi^{4}}\, \frac{1}{k^{+}} \int_{0}^{1}\frac{d\xi}{2\pi} \int_{x,z,y,y',x',w'}e^{-ik\cdot(z-w')} \times \mathfrak{K}_{\rm I.V}\notag\\
&\quad \times \frac{N_{c}}{(N_{c}^{2}-4)(N_{c}^{2}-1)}\, \Big(+2i\Big)\notag\\
&\quad \times \sigma(x',y,y')\Big[{\rm Tr}\Big\{T^{h}U^{\dagger}(z)\big(U(x)-U(y')\big)D^{h}\big(U^{\dagger}(x')-U^{\dagger}(w')\big)U(z)\Big\}\notag\\
&\quad\quad +{\rm Tr}\Big\{T^{h}U^{\dagger}(y')\big(U(x')-U(w')\big)D^{h}\big(U^{\dagger}(z)-U^{\dagger}(y)\big)U(y')\Big\}\Big] +{\rm c.c.}
\end{align}

and the contact part is
\begin{align}\label{eq:S-I.V-2}
&\Big\langle\frac{d^{3}N^{\rm I.V}_{{\rm con}}}{d^{3}k}\Big\rangle=\frac{1}{(2\pi)^{3}}\, \frac{ig^{4}}{16\pi^{4}}\, \frac{1}{k^{+}} \int_{0}^{1}\frac{d\xi}{2\pi} \int_{x,z,y,y',x',w'}e^{-ik\cdot(z-w')} \times \mathfrak{K}_{\rm I.V}\notag\\
&\quad \times \frac{1}{N_{c}^{2}-1}\bigg[\Big(+i\Big)\delta^{(2)}(x'-y)\notag\\
&\quad\quad \times n(y,y')\Big[-{\rm Tr}\Big\{T^{e}U^{\dagger}(z)\big(U(x)-U(y')\big)T^{e}\big(U^{\dagger}(y)-U^{\dagger}(w')\big)U(z)\Big\}\notag\\
&\quad\quad\quad -{\rm Tr}\Big\{T^{e}U^{\dagger}(y')\big(U(z)-U(y)\big)T^{e}\big(U^{\dagger}(y)-U^{\dagger}(w')\big)U(y')\Big\}\Big]\notag\\
&\quad\quad +\Big(-i\Big)\delta^{(2)}(x'-y')\notag\\
&\quad\quad \times n(y',y)\Big[-{\rm Tr}\Big\{T^{e}U^{\dagger}(z)\big(U(x)-U(y')\big)T^{e}\big(U^{\dagger}(y')-U^{\dagger}(w')\big)U(z)\Big\}\notag\\
&\quad\quad\quad -{\rm Tr}\Big\{T^{e}U^{\dagger}(y')\big(U(z)-U(y)\big)T^{e}\big(U^{\dagger}(y')-U^{\dagger}(w')\big)U(y')\Big\}\Big]\bigg] +{\rm c.c.}
\end{align}

The genuine two-$\rho$ part has the kernel
\begin{align}\label{eq:K-rm_quad_rm_I_V}
&\mathfrak{K}^{\rm quad}_{\rm I.V}\equiv \frac{e^{-ik\cdot(x-z)\,\xi}}{\xi\,(y-x)^{2}+(1-\xi)\,(y-z)^{2}}\bigg[\xi(1-\xi)\, \frac{(y-x)^{2}-(y-z)^{2}}{(x-z)^{2}}\, \Kv(x'-w')\cdot\Kv(z-x)\notag\\
&\quad\quad +(1-\xi)\, \Kv(x'-w')\cdot\Kv(z-x)\Big[(y-z)\cdot\Kv(x-z) +\frac{1}{2}(y-z)\cdot\Kv(y-x)\notag\\
&\quad\quad\quad -\frac{(y-x)^{2}-(y-z)^{2}}{2\,(x-z)^{2}}\Big] +\xi\, \Kv(x'-w')\cdot\Kv(z-x)\Big[(y-x)\cdot\Kv(x-z)\notag\\
&\quad\quad\quad -\frac{1}{2}(y-x)\cdot\Kv(y-z) -\frac{(y-x)^{2}-(y-z)^{2}}{2\,(x-z)^{2}}\Big]\notag\\
&\quad\quad +\frac{(y-x)^{2}}{(y-z)^{2}}\Big[-\frac{(y-z)\cdot\Kv(x'-w')}{(x-z)^{2}} +\frac{1}{2}\, (y-z)\cdot\Kv(x'-w')\, \Kv(z-x)\cdot\Kv(y-x)\Big]\notag\\
&\quad\quad +\frac{(y-z)^{2}}{(y-x)^{2}}\Big[-\frac{(y-x)\cdot\Kv(x'-w')}{(x-z)^{2}} -\frac{1}{2}\, (y-x)\cdot\Kv(x'-w')\, \Kv(z-x)\cdot\Kv(y-z)\Big]\notag\\
&\quad\quad -(y-x)\cdot\Kv(z-x)\, \Kv(x-z)\cdot\Kv(x'-w')\notag\\
&\quad\quad +\frac{1}{2}\, (y-x)\cdot\Kv(z-x)\, \Kv(y-z)\cdot\Kv(x'-w')\notag\\
&\quad\quad -(y-z)\cdot\Kv(z-x)\, \Kv(x-z)\cdot\Kv(x'-w') -\frac{1}{2}\, \Kv(y-x)\cdot\Kv(x'-w')\notag\\
&\quad\quad \times (y-z)\cdot\Kv(z-x)\bigg] -\Big[\frac{1}{\xi}\Big]_{+}\frac{e^{-ik\cdot(x-z)\,\xi}}{(y-z)^{2}}\Big[-\frac{(y-z)\cdot\Kv(x'-w')}{(x-z)^{2}}\notag\\
&\quad\quad +\frac{1}{2}\, (y-z)\cdot\Kv(x'-w')\, \Kv(z-x)\cdot\Kv(y-x)\Big]\notag\\
&\quad -\Big[\frac{1}{1-\xi}\Big]_{+}\frac{e^{-ik\cdot(x-z)\,\xi}}{(y-x)^{2}}\Big[-\frac{(y-x)\cdot\Kv(x'-w')}{(x-z)^{2}}\notag\\
&\quad\quad -\frac{1}{2}\, (y-x)\cdot\Kv(x'-w')\, \Kv(z-x)\cdot\Kv(y-z)\Big]
\end{align}

and reads
\begin{align}\label{eq:S-I.V-2o}
&\Big\langle\frac{d^{3}N^{\rm I.V}_{{\rm quad}}}{d^{3}k}\Big\rangle=\frac{1}{(2\pi)^{3}}\, \frac{g^{4}}{8\pi^{4}}\, \frac{1}{k^{+}} \int_{0}^{1}\frac{d\xi}{2\pi} \int_{x,y,z,x',w'}e^{-ik\cdot(z-w')} \times \mathfrak{K}^{\rm quad}_{\rm I.V}\notag\\
&\quad \times\frac{n(x',y)}{N_{c}^{2}-1}\Big[-{\rm Tr}\Big\{T^{e}U^{\dagger}(x)\big(U(x')-U(w')\big)T^{e}\big(U^{\dagger}(z)-U^{\dagger}(x)\big)U(x)\Big\}\notag\\
&\quad\quad +N_{c}{\rm Tr}\Big\{\big(U^{\dagger}(x')-U^{\dagger}(w')\big)\big(U(y)-U(x)\big)\Big\}\Big]\notag\\
&\quad +\frac{1}{(2\pi)^{3}}\, \frac{g^{4}N_{c}}{16\pi^{5}}\, \frac{1}{k^{+}} \int_{x,y,z,x',w'}e^{-ik\cdot(z-w')}\, e^{-ik\cdot(x-z)}\, \frac{11}{6}\, \Kv(x'-w')\cdot\Kv(y-x)\notag\\
&\quad \times \frac{(y-x)\cdot\Kv(y-z)}{(x-z)^{2}} \times\frac{n(x',y)}{N_{c}^{2}-1}\Big[{\rm Tr}\Big\{\big(U^{\dagger}(x')-U^{\dagger}(w')\big)\big(U(y)-U(x)\big)\Big\}\Big]\, +{\rm c.c.}\, ,
\end{align}

The second term, with $\frac{11}{6}$, is the one required at the collinear point $\v z\to\v x$ (Section~\ref{sec:fin_together}).

\subsubsection{Group II}

\paragraph{Row II.I.} The kernels of the cubic part are
\begin{align}\label{eq:K-1_rm_II_I}
&\mathfrak{K}^{(1)}_{\rm II.I}\equiv -\bigg[\frac{\Kv(x'-w')\cdot\Kv(y-z)}{2\,(z-w)^{2}} +\Kv(x'-w')\cdot\Kv(z-w)\, \Kv(y-z)\cdot\Kv(x-z)\notag\\
&\quad\quad -\Kv(x'-w')\cdot\Kv(x-w)\, \Kv(y-z)\cdot\Kv(x-z)\notag\\
&\quad\quad -\Kv(x'-w')\cdot\Kv(x-w)\, \Kv(y-z)\cdot\Kv(z-w)\bigg]\log\frac{(x-z)^{2}}{(x-w)^{2}}
\end{align}

and
\begin{align}\label{eq:K-2_rm_II_I}
&\mathfrak{K}^{(2)}_{\rm II.I}\equiv \bigg[\frac{\Kv(x'-w)\cdot\Kv(y-z)}{2\,(z-w')^{2}} +\Kv(x'-w)\cdot\Kv(z-w')\, \Kv(y-z)\cdot\Kv(x-z)\notag\\
&\quad\quad -\Kv(x'-w)\cdot\Kv(x-w')\, \Kv(y-z)\cdot\Kv(x-z)\notag\\
&\quad\quad -\Kv(x'-w)\cdot\Kv(x-w')\, \Kv(y-z)\cdot\Kv(z-w')\bigg]\log\frac{(x-z)^{2}}{(x-w')^{2}}\,.
\end{align}

which is $-\mathfrak K^{(1)}_{\rm II.I}$ with $\v w\leftrightarrow\v w'$. The symmetrised three-$\rho$ part is
\begin{align}\label{eq:S-II.I-3}
&\Big\langle\frac{d^{3}N^{\rm II.I}_{{\rm sym}}}{d^{3}k}\Big\rangle=\frac{1}{(2\pi)^{3}}\, \frac{ig^{4}}{16\pi^{5}}\, \frac{1}{k^{+}} \int_{x,y,z,w,x',w'}e^{-ik\cdot(w-w')}\Bigg[ \mathfrak{K}^{(1)}_{\rm II.I}\notag\\
&\quad\quad \times \frac{N_{c}}{(N_{c}^{2}-4)(N_{c}^{2}-1)}\, \Big(+2i\Big)\notag\\
&\quad\quad \times \sigma(x',y,x)\Big[-{\rm Tr}\Big\{T^{h}U^{\dagger}(z)\big(U(w')-U(x')\big)D^{h}\big(U^{\dagger}(x)-U^{\dagger}(z)\big)U(z)\Big\}\notag\\
&\quad\quad\quad +{\rm Tr}\Big\{T^{h}U^{\dagger}(z)\big(U(z)-U(y)\big)D^{h}U^{\dagger}(x')\big(U(w)-U(z)\big)\Big\}\notag\\
&\quad\quad\quad -{\rm Tr}\Big\{T^{h}U^{\dagger}(x')\big(U(x)-U(z)\big)D^{h}\big(U^{\dagger}(z)-U^{\dagger}(y)\big)U(x')\Big\}\Big] + \mathfrak{K}^{(2)}_{\rm II.I}\notag\\
&\quad\quad \times \frac{N_{c}}{(N_{c}^{2}-4)(N_{c}^{2}-1)}\, \Big(+2i\Big)\notag\\
&\quad\quad \times \sigma(x,x',y)\Big[-{\rm Tr}\Big\{T^{h}U^{\dagger}(z)\big(U(w)-U(x')\big)D^{h}\big(U^{\dagger}(x)-U^{\dagger}(z)\big)U(z)\Big\}\notag\\
&\quad\quad\quad +{\rm Tr}\Big\{T^{h}U^{\dagger}(z)\big(U(z)-U(y)\big)D^{h}U^{\dagger}(x')\big(U(w')-U(z)\big)\Big\}\notag\\
&\quad\quad\quad -{\rm Tr}\Big\{T^{h}U^{\dagger}(x')\big(U(x)-U(z)\big)D^{h}\big(U^{\dagger}(z)-U^{\dagger}(y)\big)U(x')\Big\}\Big]\Bigg]
\end{align}

and the contact part is
\begin{align}\label{eq:S-II.I-2}
&\Big\langle\frac{d^{3}N^{\rm II.I}_{{\rm con}}}{d^{3}k}\Big\rangle=\frac{1}{(2\pi)^{3}}\, \frac{ig^{4}}{16\pi^{5}}\, \frac{1}{k^{+}} \int_{x,y,z,w,x',w'}e^{-ik\cdot(w-w')}\Bigg[ \mathfrak{K}^{(1)}_{\rm II.I}\notag\\
&\quad\quad \times \frac{1}{N_{c}^{2}-1}\bigg[\Big(+i\Big)\delta^{(2)}(y-x)\, n(x,x')\Big[N_{c}{\rm Tr}\Big\{\big(U^{\dagger}(w')-U^{\dagger}(x')\big)\big(U(w)-U(z)\big)\Big\}\notag\\
&\quad\quad\quad\quad -{\rm Tr}\Big\{T^{e}U^{\dagger}(z)\big(U(x)-U(z)\big)T^{e}\big(U^{\dagger}(w')-U^{\dagger}(x')\big)U(z)\Big\}\notag\\
&\quad\quad\quad\quad +{\rm Tr}\Big\{T^{e}U^{\dagger}(z)\big(U(z)-U(x)\big)T^{e}U^{\dagger}(x')\big(U(w)-U(z)\big)\Big\}\notag\\
&\quad\quad\quad\quad -{\rm Tr}\Big\{T^{e}U^{\dagger}(x')\big(U(x)-U(z)\big)T^{e}\big(U^{\dagger}(z)-U^{\dagger}(x)\big)U(x')\Big\}\Big]\notag\\
&\quad\quad\quad +\Big(-i\Big)\delta^{(2)}(x'-x)\, n(x,y)\Big[N_{c}{\rm Tr}\Big\{\big(U^{\dagger}(w')-U^{\dagger}(x)\big)\big(U(w)-U(z)\big)\Big\}\notag\\
&\quad\quad\quad\quad -{\rm Tr}\Big\{T^{e}U^{\dagger}(z)\big(U(x)-U(z)\big)T^{e}\big(U^{\dagger}(w')-U^{\dagger}(x)\big)U(z)\Big\}\notag\\
&\quad\quad\quad\quad +{\rm Tr}\Big\{T^{e}U^{\dagger}(z)\big(U(z)-U(y)\big)T^{e}U^{\dagger}(x)\big(U(w)-U(z)\big)\Big\}\notag\\
&\quad\quad\quad\quad -{\rm Tr}\Big\{T^{e}U^{\dagger}(x)\big(U(x)-U(z)\big)T^{e}\big(U^{\dagger}(z)-U^{\dagger}(y)\big)U(x)\Big\}\Big]\bigg] + \mathfrak{K}^{(2)}_{\rm II.I}\notag\\
&\quad\quad \times \frac{1}{N_{c}^{2}-1}\bigg[\Big(+i\Big)\delta^{(2)}(x'-x)\, n(x,y)\Big[N_{c}{\rm Tr}\Big\{\big(U^{\dagger}(w)-U^{\dagger}(x)\big)\big(U(w')-U(z)\big)\Big\}\notag\\
&\quad\quad\quad\quad -{\rm Tr}\Big\{T^{e}U^{\dagger}(z)\big(U(x)-U(z)\big)T^{e}\big(U^{\dagger}(w)-U^{\dagger}(x)\big)U(z)\Big\}\notag\\
&\quad\quad\quad\quad +{\rm Tr}\Big\{T^{e}U^{\dagger}(z)\big(U(z)-U(y)\big)T^{e}U^{\dagger}(x)\big(U(w')-U(z)\big)\Big\}\notag\\
&\quad\quad\quad\quad -{\rm Tr}\Big\{T^{e}U^{\dagger}(x)\big(U(x)-U(z)\big)T^{e}\big(U^{\dagger}(z)-U^{\dagger}(y)\big)U(x)\Big\}\Big]\notag\\
&\quad\quad\quad +\Big(-i\Big)\delta^{(2)}(y-x)\, n(x,x')\Big[N_{c}{\rm Tr}\Big\{\big(U^{\dagger}(w)-U^{\dagger}(x')\big)\big(U(w')-U(z)\big)\Big\}\notag\\
&\quad\quad\quad\quad -{\rm Tr}\Big\{T^{e}U^{\dagger}(z)\big(U(x)-U(z)\big)T^{e}\big(U^{\dagger}(w)-U^{\dagger}(x')\big)U(z)\Big\}\notag\\
&\quad\quad\quad\quad +{\rm Tr}\Big\{T^{e}U^{\dagger}(z)\big(U(z)-U(x)\big)T^{e}U^{\dagger}(x')\big(U(w')-U(z)\big)\Big\}\notag\\
&\quad\quad\quad\quad -{\rm Tr}\Big\{T^{e}U^{\dagger}(x')\big(U(x)-U(z)\big)T^{e}\big(U^{\dagger}(z)-U^{\dagger}(x)\big)U(x')\Big\}\Big]\bigg]\Bigg]
\end{align}

Group~II has no complex conjugate. The genuine two-$\rho$ part consists of five terms with the same colour factor. With the kernels
\begin{align}\label{eq:K-rm_quad_a_rm_II_I}
&\mathfrak{K}^{{\rm quad},a}_{\rm II.I}\equiv \frac{1}{\big[(x'-w')^{2}\,(x-z)^{2}-(x-w)^{2}\,(x'-z)^{2}\big]}\, \log\left[\frac{(x-z)^{2}}{(x'-z)^{2}}\right]\notag\\
&\quad \times\Bigg[(x-w)\cdot\Kv(z-w')\, (x'-z)\cdot\Kv(z-w) -\frac{1}{2}\, (x-w)\cdot\Kv(x'-w')\notag\\
&\quad\quad \times \Kv(z-w)\cdot(x'-z) +\frac{1}{2}\, \Kv(x-z)\cdot(x'-z)\, (x-w)\cdot\Kv(z-w')\notag\\
&\quad\quad -\frac{1}{4}\, (x'-z)\cdot\Kv(x-z)\, \Kv(x'-w')\cdot(x-w)\notag\\
&\quad\quad +(x'-w')\cdot\Kv(z-w)\, (x-z)\cdot\Kv(z-w') +\frac{1}{2}\, (x'-w')\cdot\Kv(z-w)\notag\\
&\quad\quad \times (x-z)\cdot\Kv(x'-z) -\frac{1}{2}\, (x-z)\cdot\Kv(z-w')\, (x'-w')\cdot\Kv(x-w)\notag\\
&\quad\quad -\frac{1}{4}\, (x'-w')\cdot\Kv(x-w)\, (x-z)\cdot\Kv(x'-z)\Bigg]\,.
\end{align}

\begin{align}\label{eq:K-rm_quad_b_rm_II_I}
&\mathfrak{K}^{{\rm quad},b}_{\rm II.I}\equiv \frac{1}{\big[(x'-w')^{2}\,(x-z)^{2}-(x-w)^{2}\,(x'-z)^{2}\big]}\, \log\left[\frac{(x-w)^{2}}{(x'-w')^{2}}\right]\notag\\
&\quad \times\Bigg[(x-w)\cdot\Kv(z-w')\, (x'-z)\cdot\Kv(z-w) -\frac{1}{2}\, (x-w)\cdot\Kv(x'-w')\notag\\
&\quad\quad \times \Kv(z-w)\cdot(x'-z) +\frac{1}{2}\, \Kv(x-z)\cdot(x'-z)\, (x-w)\cdot\Kv(z-w')\notag\\
&\quad\quad -\frac{1}{4}\, (x'-z)\cdot\Kv(x-z)\, \Kv(x'-w')\cdot(x-w)\notag\\
&\quad\quad +(x'-w')\cdot\Kv(z-w)\, (x-z)\cdot\Kv(z-w') +\frac{1}{2}\, (x'-w')\cdot\Kv(z-w)\notag\\
&\quad\quad \times (x-z)\cdot\Kv(x'-z) -\frac{1}{2}\, (x-z)\cdot\Kv(z-w')\, (x'-w')\cdot\Kv(x-w)\notag\\
&\quad\quad -\frac{1}{4}\, (x'-w')\cdot\Kv(x-w)\, (x-z)\cdot\Kv(x'-z)\Bigg]\,.
\end{align}

\begin{align}\label{eq:K-rm_quad_c_rm_II_I}
&\mathfrak{K}^{{\rm quad},c}_{\rm II.I}\equiv \frac{1}{\big[(x'-w')^{2}\,(x-z)^{2}-(x-w)^{2}\,(x'-z)^{2}\big]}\, \log\left[\frac{(x-z)^{2}}{(x-w)^{2}}\right]\notag\\
&\quad \times\Bigg[-\frac{(x-z)^{2}}{(x-w)^{2}}\bigg[\frac{(x'-w')\cdot(x-w)\,(z-w)\cdot(z-w')}{(z-w)^{2}\,(z-w')^{2}}\notag\\
&\quad\quad\quad +\frac{(x'-w')\cdot(x-w)\,(x'-z)\cdot(z-w)}{2\,(z-w)^{2}\,(x'-z)^{2}} +\frac{(x-w)\cdot(x'-w')\,(x-z)\cdot(z-w')}{2\,(x-z)^{2}\,(z-w')^{2}}\notag\\
&\quad\quad\quad +\frac{(x-z)\cdot(x'-z)\,(x'-w')\cdot(x-w)}{4\,(x-z)^{2}\,(x'-z)^{2}}\bigg]\notag\\
&\quad\quad -\frac{(x-w)^{2}}{(x-z)^{2}}\bigg[\frac{(x'-z)\cdot(x-z)\,(z-w)\cdot(z-w')}{(z-w)^{2}\,(z-w')^{2}}\notag\\
&\quad\quad\quad -\frac{(x'-z)\cdot(x-z)\,(x'-w')\cdot(z-w)}{2\,(z-w)^{2}\,(x'-w')^{2}} -\frac{(x-w)\cdot(z-w')\,(x-z)\cdot(x'-z)}{2\,(x-w)^{2}\,(z-w')^{2}}\notag\\
&\quad\quad\quad +\frac{(x-z)\cdot(x'-z)\,(x'-w')\cdot(x-w)}{4\,(x-w)^{2}\,(x'-w')^{2}}\bigg]\notag\\
&\quad\quad +\frac{1}{(x-z)^{2}-(x-w)^{2}}\bigg\{(x-z)^{2}\bigg[\frac{(x-w)\cdot(z-w)\,\big[(x'-z)^{2}-(x'-w')^{2}\big]}{2\,(z-w)^{2}\,(z-w')^{2}}\notag\\
&\quad\quad\quad\quad +\frac{(x-z)\cdot(x-w)\,\big[(x'-z)^{2}-(x'-w')^{2}\big]}{4\,(x-z)^{2}\,(z-w')^{2}}\notag\\
&\quad\quad\quad\quad +\frac{(x'-w')\cdot(z-w')\,\big[(x-z)^{2}-(x-w)^{2}\big]}{2\,(z-w)^{2}\,(z-w')^{2}}\notag\\
&\quad\quad\quad\quad +\frac{(x'-z)\cdot(x'-w')\,\big[(x-z)^{2}-(x-w)^{2}\big]}{4\,(z-w)^{2}\,(x'-z)^{2}}\bigg]\notag\\
&\quad\quad\quad -(x-w)^{2}\bigg[\frac{(x-z)\cdot(z-w)\,\big[(x'-z)^{2}-(x'-w')^{2}\big]}{2\,(z-w)^{2}\,(z-w')^{2}}\notag\\
&\quad\quad\quad\quad -\frac{(x-z)\cdot(x-w)\,\big[(x'-z)^{2}-(x'-w')^{2}\big]}{4\,(x-w)^{2}\,(z-w')^{2}}\notag\\
&\quad\quad\quad\quad +\frac{(x'-z)\cdot(z-w')\,\big[(x-z)^{2}-(x-w)^{2}\big]}{2\,(z-w)^{2}\,(z-w')^{2}}\notag\\
&\quad\quad\quad\quad -\frac{(x'-z)\cdot(x'-w')\,\big[(x-z)^{2}-(x-w)^{2}\big]}{4\,(z-w)^{2}\,(x'-w')^{2}}\bigg]\bigg\}\notag\\
&\quad\quad -\frac{1}{2}\, \frac{(x-z)^{2}\,(x-w)^{2}\,\big[(x'-z)^{2}-(x'-w')^{2}\big]}{(z-w)^{2}\,(z-w')^{2}\,\big[(x-z)^{2}-(x-w)^{2}\big]}\Bigg]
\end{align}

\begin{align}\label{eq:K-rm_quad_d_rm_II_I}
&\mathfrak{K}^{{\rm quad},d}_{\rm II.I}\equiv \frac{1}{\big[(x'-w')^{2}\,(x-z)^{2}-(x-w)^{2}\,(x'-z)^{2}\big]}\, \log\left[\frac{(x'-z)^{2}}{(x'-w')^{2}}\right]\notag\\
&\quad \times\Bigg[\frac{(x'-z)^{2}}{(x'-w')^{2}}\bigg[(x'-w')\cdot(x-w)\, \Kv(z-w)\cdot\Kv(z-w')\notag\\
&\quad\quad\quad +\frac{1}{2}\, (x'-w')\cdot(x-w)\, \Kv(x'-z)\cdot\Kv(z-w)\notag\\
&\quad\quad\quad +\frac{1}{2}\, (x-w)\cdot(x'-w')\, \Kv(x-z)\cdot\Kv(z-w')\notag\\
&\quad\quad\quad +\frac{1}{4}\, \Kv(x-z)\cdot\Kv(x'-z)\, (x'-w')\cdot(x-w)\bigg]\notag\\
&\quad\quad +\frac{(x'-w')^{2}}{(x'-z)^{2}}\bigg[(x'-z)\cdot(x-z)\, \Kv(z-w)\cdot\Kv(z-w')\notag\\
&\quad\quad\quad -\frac{1}{2}\, (x'-z)\cdot(x-z)\, \Kv(x'-w')\cdot\Kv(z-w)\notag\\
&\quad\quad\quad -\frac{1}{2}\, \Kv(x-w)\cdot\Kv(z-w')\, (x-z)\cdot(x'-z)\notag\\
&\quad\quad\quad +\frac{1}{4}\, (x-z)\cdot(x'-z)\, \Kv(x'-w')\cdot\Kv(x-w)\bigg]\notag\\
&\quad\quad +\frac{1}{(x'-w')^{2}-(x'-z)^{2}}\bigg\{(x'-z)^{2}\bigg[\frac{(x-w)\cdot\Kv(z-w)\,\big[(x'-z)^{2}-(x'-w')^{2}\big]}{2\,(z-w')^{2}}\notag\\
&\quad\quad\quad\quad +\frac{(x-w)\cdot\Kv(x-z)\,\big[(x'-z)^{2}-(x'-w')^{2}\big]}{4\,(z-w')^{2}}\notag\\
&\quad\quad\quad\quad +\frac{(x'-w')\cdot\Kv(z-w')\,\big[(x-z)^{2}-(x-w)^{2}\big]}{2\,(z-w)^{2}}\notag\\
&\quad\quad\quad\quad +\frac{(x'-w')\cdot\Kv(x'-z)\,\big[(x-z)^{2}-(x-w)^{2}\big]}{4\,(z-w)^{2}}\bigg]\notag\\
&\quad\quad\quad -(x'-w')^{2}\bigg[\frac{(x-z)\cdot\Kv(z-w)\,\big[(x'-z)^{2}-(x'-w')^{2}\big]}{2\,(z-w')^{2}}\notag\\
&\quad\quad\quad\quad -\frac{(x-z)\cdot\Kv(x-w)\,\big[(x'-z)^{2}-(x'-w')^{2}\big]}{4\,(z-w')^{2}}\notag\\
&\quad\quad\quad\quad +\frac{(x'-z)\cdot\Kv(z-w')\,\big[(x-z)^{2}-(x-w)^{2}\big]}{2\,(z-w)^{2}}\notag\\
&\quad\quad\quad\quad -\frac{(x'-z)\cdot\Kv(x'-w')\,\big[(x-z)^{2}-(x-w)^{2}\big]}{4\,(z-w)^{2}}\bigg]\bigg\}\notag\\
&\quad\quad +\frac{1}{2}\, \frac{(x'-w')^{2}\,(x'-z)^{2}\,\big[(x-z)^{2}-(x-w)^{2}\big]}{(z-w)^{2}\,(z-w')^{2}\,\big[(x'-z)^{2}-(x'-w')^{2}\big]}\Bigg]
\end{align}

the five terms are
\begingroup\small
\begin{align}\label{eq:S-II.I-2a}
&\Big\langle\frac{d^{3}N^{\rm II.I}_{{\rm quad},a}}{d^{3}k}\Big\rangle=\frac{1}{(2\pi)^{3}}\, \frac{g^{4}}{8\pi^{5}}\, \frac{1}{k^{+}} \int_{z,w,w',x,x'}e^{-ik\cdot(w-w')}\, \mathfrak{K}^{{\rm quad},a}_{\rm II.I} \times\frac{n(x',x)}{N_{c}^{2}-1}\Big[\notag\\
&\quad\quad N_{c}{\rm Tr}\Big\{\big(U^{\dagger}(w')-U^{\dagger}(z)\big)\big(U(w)-U(z)\big)\Big\} +N_{c}{\rm Tr}\Big\{\big(U^{\dagger}(x')-U^{\dagger}(z)\big)\big(U(x)-U(z)\big)\Big\}\notag\\
&\quad\quad -{\rm Tr}\Big\{T^{e}U^{\dagger}(z)\big(U(x)-U(z)\big)T^{e}\big(U^{\dagger}(w')-U^{\dagger}(z)\big)U(z)\Big\}\notag\\
&\quad\quad -{\rm Tr}\Big\{T^{e}U^{\dagger}(z)\big(U(x')-U(z)\big)T^{e}\big(U^{\dagger}(w)-U^{\dagger}(z)\big)U(z)\Big\}\Big]
\end{align}
\endgroup

\begin{align}\label{eq:S-II.I-2b}
&\Big\langle\frac{d^{3}N^{\rm II.I}_{{\rm quad},b}}{d^{3}k}\Big\rangle=-\frac{1}{(2\pi)^{3}}\, \frac{g^{4}}{8\pi^{5}}\, \frac{1}{k^{+}} \int_{z,w,w',x,x'}e^{-ik\cdot(w-w')}\, \mathfrak{K}^{{\rm quad},b}_{\rm II.I}\notag\\
&\quad \times\frac{n(x',x)}{N_{c}^{2}-1}\Big[N_{c}{\rm Tr}\Big\{\big(U^{\dagger}(w')-U^{\dagger}(z)\big)\big(U(w)-U(z)\big)\Big\}\notag\\
&\quad\quad +N_{c}{\rm Tr}\Big\{\big(U^{\dagger}(x')-U^{\dagger}(z)\big)\big(U(x)-U(z)\big)\Big\}\notag\\
&\quad\quad -{\rm Tr}\Big\{T^{e}U^{\dagger}(z)\big(U(x)-U(z)\big)T^{e}\big(U^{\dagger}(w')-U^{\dagger}(z)\big)U(z)\Big\}\notag\\
&\quad\quad -{\rm Tr}\Big\{T^{e}U^{\dagger}(z)\big(U(x')-U(z)\big)T^{e}\big(U^{\dagger}(w)-U^{\dagger}(z)\big)U(z)\Big\}\Big]
\end{align}

\begin{align}\label{eq:S-II.I-2c}
&\Big\langle\frac{d^{3}N^{\rm II.I}_{{\rm quad},c}}{d^{3}k}\Big\rangle=\frac{1}{(2\pi)^{3}}\, \frac{g^{4}}{8\pi^{5}}\, \frac{1}{k^{+}} \int_{z,w,w',x,x'}e^{-ik\cdot(w-w')}\, \mathfrak{K}^{{\rm quad},c}_{\rm II.I}\notag\\
&\quad \times\frac{n(x',x)}{N_{c}^{2}-1}\Big[N_{c}{\rm Tr}\Big\{\big(U^{\dagger}(w')-U^{\dagger}(z)\big)\big(U(w)-U(z)\big)\Big\}\notag\\
&\quad\quad +N_{c}{\rm Tr}\Big\{\big(U^{\dagger}(x')-U^{\dagger}(z)\big)\big(U(x)-U(z)\big)\Big\}\notag\\
&\quad\quad -{\rm Tr}\Big\{T^{e}U^{\dagger}(z)\big(U(x)-U(z)\big)T^{e}\big(U^{\dagger}(w')-U^{\dagger}(z)\big)U(z)\Big\}\notag\\
&\quad\quad -{\rm Tr}\Big\{T^{e}U^{\dagger}(z)\big(U(x')-U(z)\big)T^{e}\big(U^{\dagger}(w)-U^{\dagger}(z)\big)U(z)\Big\}\Big]
\end{align}

\begin{align}\label{eq:S-II.I-2d}
&\Big\langle\frac{d^{3}N^{\rm II.I}_{{\rm quad},d}}{d^{3}k}\Big\rangle=\frac{1}{(2\pi)^{3}}\, \frac{g^{4}}{8\pi^{5}}\, \frac{1}{k^{+}} \int_{z,w,w',x,x'}e^{-ik\cdot(w-w')}\, \mathfrak{K}^{{\rm quad},d}_{\rm II.I}\notag\\
&\quad \times\frac{n(x',x)}{N_{c}^{2}-1}\Big[N_{c}{\rm Tr}\Big\{\big(U^{\dagger}(w')-U^{\dagger}(z)\big)\big(U(w)-U(z)\big)\Big\}\notag\\
&\quad\quad +N_{c}{\rm Tr}\Big\{\big(U^{\dagger}(x')-U^{\dagger}(z)\big)\big(U(x)-U(z)\big)\Big\}\notag\\
&\quad\quad -{\rm Tr}\Big\{T^{e}U^{\dagger}(z)\big(U(x)-U(z)\big)T^{e}\big(U^{\dagger}(w')-U^{\dagger}(z)\big)U(z)\Big\}\notag\\
&\quad\quad -{\rm Tr}\Big\{T^{e}U^{\dagger}(z)\big(U(x')-U(z)\big)T^{e}\big(U^{\dagger}(w)-U^{\dagger}(z)\big)U(z)\Big\}\Big]
\end{align}

\begin{align}\label{eq:S-II.I-2e}
&\Big\langle\frac{d^{3}N^{\rm II.I}_{{\rm quad},e}}{d^{3}k}\Big\rangle=-\frac{1}{(2\pi)^{3}}\, \frac{g^{4}}{16\pi^{5}}\, \frac{1}{k^{+}} \int_{z,w,w',x,x'}e^{-ik\cdot(w-w')}\, \frac{1}{(z-w)^{2}\,(z-w')^{2}}\notag\\
&\quad \times\frac{n(x',x)}{N_{c}^{2}-1}\Big[N_{c}{\rm Tr}\Big\{\big(U^{\dagger}(w')-U^{\dagger}(z)\big)\big(U(w)-U(z)\big)\Big\}\notag\\
&\quad\quad +N_{c}{\rm Tr}\Big\{\big(U^{\dagger}(x')-U^{\dagger}(z)\big)\big(U(x)-U(z)\big)\Big\}\notag\\
&\quad\quad -{\rm Tr}\Big\{T^{e}U^{\dagger}(z)\big(U(x)-U(z)\big)T^{e}\big(U^{\dagger}(w')-U^{\dagger}(z)\big)U(z)\Big\}\notag\\
&\quad\quad -{\rm Tr}\Big\{T^{e}U^{\dagger}(z)\big(U(x')-U(z)\big)T^{e}\big(U^{\dagger}(w)-U^{\dagger}(z)\big)U(z)\Big\}\Big]
\end{align}

The poles of the kernels $\mathfrak K^{{\rm quad},a}_{\rm II.I},\dots,\mathfrak K^{{\rm quad},d}_{\rm II.I}$ on the surface $(x'-w')^2(x-z)^2=(x-w)^2(x'-z)^2$ cancel in the sums of the terms $a$ and $b$ and of the terms $c$ and $d$.

\paragraph{Row II.IV.} With $\xi$ of Eq.~\eqref{eq:xi_defIV},
\begin{align}\label{eq:S-II.IVrem}
&\Big\langle\frac{d^{3}N^{\rm II.IV}_{{\rm fin}}}{d^{3}k}\Big\rangle=+\frac{1}{(2\pi)^{3}}\, \frac{g^{4}\,N_{c}}{4\pi^{3}}\, \frac{1}{k^{+}} \int_{k^{+}/\vee}^{1}\frac{d\xi}{2\pi}\, \frac{1-\xi}{\xi}\notag\\
&\quad \times \int_{x,y,x',y'}e^{-ik\cdot(y-y')\frac{1}{\xi}}\Big[\Kv(x-y)\cdot\Kv(x'-y') -2\, \Kv(x-y)\cdot(y'-y)\notag\\
&\quad\quad \times \Kv(x'-y')\cdot\Kv(y'-y)\Big] \times\frac{n(x',x)}{N_{c}^{2}-1}\, {\rm Tr}\Big\{\big(U^{\dagger}(y')-U^{\dagger}(x')\big)\big(U(y)-U(x)\big)\Big\}\notag\\
&\quad -\frac{1}{(2\pi)^{3}}\, \frac{g^{4}\,N_{c}}{2\pi^{3}}\, \frac{1}{k^{2}}\, \frac{1}{k^{+}} \int_{x,x',z}\Kv(x-z)\cdot\Kv(x'-z)\frac{n(x',x)}{N_{c}^{2}-1}\notag\\
&\quad \times {\rm Tr}\Big\{\big(U^{\dagger}(z)-U^{\dagger}(x')\big)\big(U(z)-U(x)\big)\Big\}
\end{align}

The first term is the finite term of the real splitting that is left beside the DGLAP term~\eqref{eq:DGLAP}; its tensor structure averages to zero over the direction of $\v y'-\v y$. After the integral over $\v y-\v y'$ its integrand is $O(\xi)$ at $\xi\to0$, so that the limit $\vee\to\infty$ exists; at fixed coordinates the $\xi$ integral converges conditionally and gives the same result. The second term is the remnant~\eqref{eq:remnant} of the subtraction at $\xi\to0$ in $\widetilde{\mathcal P}_{gg}$.

\subsubsection{Group III}

\paragraph{Row III.II.} The kernel is
\begin{align}\label{eq:K-rm_III_II}
&\mathfrak{K}_{\rm III.II}\equiv \bigg[\frac{\Kv(x'-w')\cdot\Kv(y-z)}{2\,(z-w)^{2}} +\Kv(x'-w')\cdot\Kv(z-w)\, \Kv(y-z)\cdot\Kv(x-z)\notag\\
&\quad\quad -\Kv(x'-w')\cdot\Kv(x-w)\, \Kv(y-z)\cdot\Kv(x-z)\notag\\
&\quad\quad -\Kv(x'-w')\cdot\Kv(x-w)\, \Kv(y-z)\cdot\Kv(z-w)\bigg]\log\frac{(x-z)^{2}}{(x-w)^{2}}\,.
\end{align}

so that $\mathfrak K^{(1)}_{\rm II.I}=-\mathfrak K_{\rm III.II}$. The symmetrised three-$\rho$ part is
\begin{align}\label{eq:S-III.II-3}
&\Big\langle\frac{d^{3}N^{\rm III.II}_{{\rm sym}}}{d^{3}k}\Big\rangle=-\frac{1}{(2\pi)^{3}}\, \frac{ig^{4}}{8\pi^{5}}\, \frac{1}{k^{+}} \int_{x,y,x',z,w,w'}e^{-ik\cdot(w-w')} \times \mathfrak{K}_{\rm III.II}\notag\\
&\quad \times \frac{N_{c}}{(N_{c}^{2}-4)(N_{c}^{2}-1)}\, \Big(+i\Big)\notag\\
&\quad \times \sigma(y,x',x)\Big[-{\rm Tr}\Big\{T^{h}U^{\dagger}(z)\big(U(z)-U(y)\big)D^{h}U^{\dagger}(x')\big(U(w)-U(z)\big)\Big\}\notag\\
&\quad\quad +{\rm Tr}\Big\{T^{h}U^{\dagger}(x')\big(U(x)-U(z)\big)D^{h}\big(U^{\dagger}(z)-U^{\dagger}(y)\big)U(x')\Big\}\Big] +{\rm c.c.}
\end{align}

and the contact part is
\begin{align}\label{eq:S-III.II-2}
&\Big\langle\frac{d^{3}N^{\rm III.II}_{{\rm con}}}{d^{3}k}\Big\rangle=-\frac{1}{(2\pi)^{3}}\, \frac{ig^{4}}{8\pi^{5}}\, \frac{1}{k^{+}} \int_{x,y,x',z,w,w'}e^{-ik\cdot(w-w')} \times \mathfrak{K}_{\rm III.II}\notag\\
&\quad \times \frac{1}{N_{c}^{2}-1}\bigg[\Big(+\frac{1}{2}i\Big)\delta^{(2)}(x'-y)\notag\\
&\quad\quad \times n(y,x)\Big[{\rm Tr}\Big\{T^{e}U^{\dagger}(z)\big(U(z)-U(y)\big)T^{e}U^{\dagger}(y)\big(U(w)-U(z)\big)\Big\}\notag\\
&\quad\quad\quad -{\rm Tr}\Big\{T^{e}U^{\dagger}(y)\big(U(x)-U(z)\big)T^{e}\big(U^{\dagger}(z)-U^{\dagger}(y)\big)U(y)\Big\}\Big]\notag\\
&\quad\quad +\Big(+\frac{1}{2}i\Big)\delta^{(2)}(x'-x)\, n(x,y)\Big[{\rm Tr}\Big\{T^{e}U^{\dagger}(z)\big(U(z)-U(y)\big)T^{e}U^{\dagger}(x)\big(U(w)-U(z)\big)\Big\}\notag\\
&\quad\quad\quad -{\rm Tr}\Big\{T^{e}U^{\dagger}(x)\big(U(x)-U(z)\big)T^{e}\big(U^{\dagger}(z)-U^{\dagger}(y)\big)U(x)\Big\}\Big]\notag\\
&\quad\quad +\Big(-\frac{1}{2}i\Big)\delta^{(2)}(y-x)\, n(x,x')\Big[{\rm Tr}\Big\{T^{e}U^{\dagger}(z)\big(U(z)-U(x)\big)T^{e}U^{\dagger}(x')\big(U(w)-U(z)\big)\Big\}\notag\\
&\quad\quad\quad -{\rm Tr}\Big\{T^{e}U^{\dagger}(x')\big(U(x)-U(z)\big)T^{e}\big(U^{\dagger}(z)-U^{\dagger}(x)\big)U(x')\Big\}\Big]\bigg] +{\rm c.c.}
\end{align}


\paragraph{Row III.III.} Rows~III.III, III.V and~III.VI have the same kernels in their cubic parts,
\begin{align}\label{eq:K-y_rm_III}
&\mathfrak{K}^{(y')}_{\rm III}\equiv \frac{1}{1+\xi}\Big[\Kv(y'-z)\cdot\Kv(x'-y')\, \Kv(y-w)\cdot\Kv(x-z)\notag\\
&\quad\quad -\Kv(y'-z)\cdot\Kv(x-z)\, \Kv(x'-y')\cdot\Kv(y-w)\Big]\notag\\
&\quad -\Big[\frac{1}{\xi\,(1+\xi)}\Big]_{+}\Kv(y'-z)\cdot\Kv(x-z)\, \Kv(x'-y')\cdot\Kv(y-w)\notag\\
&\quad -\Kv(y'-z)\cdot\Kv(y-w)\, \Kv(x-z)\cdot\Kv(x'-y')
\end{align}

and
\begin{align}\label{eq:K-z_rm_III}
&\mathfrak{K}^{(z)}_{\rm III}\equiv \Kv(y'-z)\cdot\Kv(y-w)\, \Kv(x-z)\cdot\Kv(x'-z)
\end{align}

The last term of $\mathfrak K^{(y')}_{\rm III}$ and the kernel $\mathfrak K^{(z)}_{\rm III}$ are the instantaneous terms, defined by Eq.~\eqref{eq:inst_xi}. The symmetrised three-$\rho$ part of Row~III.III is
\begin{align}\label{eq:S-III.III-3}
&\Big\langle\frac{d^{3}N^{\rm III.III}_{{\rm sym}}}{d^{3}k}\Big\rangle=\frac{1}{(2\pi)^{3}}\, \frac{ig^{4}}{8\pi^{4}}\, \frac{1}{k^{+}} \int_{0}^{\infty}\frac{d\xi}{2\pi} \int_{w,x',y',z,y,x}e^{-ik\cdot(w-y')}e^{-ik\cdot(z-y')\,\xi}\Bigg[ \mathfrak{K}^{(y')}_{\rm III}\notag\\
&\quad\quad \times \frac{N_{c}}{(N_{c}^{2}-4)(N_{c}^{2}-1)}\, \Big(+2i\Big)\notag\\
&\quad\quad \times \sigma(x',y,x)\Big[{\rm Tr}\Big\{T^{h}U^{\dagger}(w)\big(U(z)-U(x)\big)D^{h}\big(U^{\dagger}(y')-U^{\dagger}(x')\big)U(w)\Big\}\notag\\
&\quad\quad\quad +{\rm Tr}\Big\{T^{h}U^{\dagger}(x)\big(U(y')-U(x')\big)D^{h}\big(U^{\dagger}(w)-U^{\dagger}(y)\big)U(x)\Big\}\Big] + \mathfrak{K}^{(z)}_{\rm III}\notag\\
&\quad\quad \times \frac{N_{c}}{(N_{c}^{2}-4)(N_{c}^{2}-1)}\, \Big(+2i\Big)\notag\\
&\quad\quad \times \sigma(x',y,x)\Big[{\rm Tr}\Big\{T^{h}U^{\dagger}(w)\big(U(z)-U(x)\big)D^{h}\big(U^{\dagger}(z)-U^{\dagger}(x')\big)U(w)\Big\}\notag\\
&\quad\quad\quad +{\rm Tr}\Big\{T^{h}U^{\dagger}(x)\big(U(z)-U(x')\big)D^{h}\big(U^{\dagger}(w)-U^{\dagger}(y)\big)U(x)\Big\}\Big]\Bigg] +{\rm c.c.}
\end{align}

and the contact part is
\begin{align}\label{eq:S-III.III-2}
&\Big\langle\frac{d^{3}N^{\rm III.III}_{{\rm con}}}{d^{3}k}\Big\rangle=\frac{1}{(2\pi)^{3}}\, \frac{ig^{4}}{8\pi^{4}}\, \frac{1}{k^{+}} \int_{0}^{\infty}\frac{d\xi}{2\pi} \int_{w,x',y',z,y,x}e^{-ik\cdot(w-y')}e^{-ik\cdot(z-y')\,\xi}\Bigg[ \mathfrak{K}^{(y')}_{\rm III}\notag\\
&\quad\quad \times \frac{1}{N_{c}^{2}-1}\bigg[\Big(+i\Big)\delta^{(2)}(x'-y)\notag\\
&\quad\quad\quad \times n(y,x)\Big[-{\rm Tr}\Big\{T^{e}U^{\dagger}(w)\big(U(z)-U(x)\big)T^{e}\big(U^{\dagger}(y')-U^{\dagger}(y)\big)U(w)\Big\}\notag\\
&\quad\quad\quad\quad -{\rm Tr}\Big\{T^{e}U^{\dagger}(x)\big(U(w)-U(y)\big)T^{e}\big(U^{\dagger}(y')-U^{\dagger}(y)\big)U(x)\Big\}\Big]\notag\\
&\quad\quad\quad +\Big(-i\Big)\delta^{(2)}(x'-x)\notag\\
&\quad\quad\quad \times n(x,y)\Big[-{\rm Tr}\Big\{T^{e}U^{\dagger}(w)\big(U(z)-U(x)\big)T^{e}\big(U^{\dagger}(y')-U^{\dagger}(x)\big)U(w)\Big\}\notag\\
&\quad\quad\quad\quad -{\rm Tr}\Big\{T^{e}U^{\dagger}(x)\big(U(w)-U(y)\big)T^{e}\big(U^{\dagger}(y')-U^{\dagger}(x)\big)U(x)\Big\}\Big]\bigg] + \mathfrak{K}^{(z)}_{\rm III}\notag\\
&\quad\quad \times \frac{1}{N_{c}^{2}-1}\bigg[\Big(+i\Big)\delta^{(2)}(x'-y)\notag\\
&\quad\quad\quad \times n(y,x)\Big[-{\rm Tr}\Big\{T^{e}U^{\dagger}(w)\big(U(z)-U(x)\big)T^{e}\big(U^{\dagger}(z)-U^{\dagger}(y)\big)U(w)\Big\}\notag\\
&\quad\quad\quad\quad -{\rm Tr}\Big\{T^{e}U^{\dagger}(x)\big(U(w)-U(y)\big)T^{e}\big(U^{\dagger}(z)-U^{\dagger}(y)\big)U(x)\Big\}\Big]\notag\\
&\quad\quad\quad +\Big(-i\Big)\delta^{(2)}(x'-x)\notag\\
&\quad\quad\quad \times n(x,y)\Big[-{\rm Tr}\Big\{T^{e}U^{\dagger}(w)\big(U(z)-U(x)\big)T^{e}\big(U^{\dagger}(z)-U^{\dagger}(x)\big)U(w)\Big\}\notag\\
&\quad\quad\quad\quad -{\rm Tr}\Big\{T^{e}U^{\dagger}(x)\big(U(w)-U(y)\big)T^{e}\big(U^{\dagger}(z)-U^{\dagger}(x)\big)U(x)\Big\}\Big]\bigg]\Bigg] +{\rm c.c.}
\end{align}

The genuine two-$\rho$ part is
\begingroup\small
\begin{align}\label{eq:S-III.III-2o}
&\Big\langle\frac{d^{3}N^{\rm III.III}_{{\rm quad}}}{d^{3}k}\Big\rangle=\frac{1}{(2\pi)^{3}}\, \frac{g^{4}}{4\pi^{4}}\, \frac{1}{k^{+}} \int_{0}^{\infty}\frac{d\xi}{2\pi} \int_{w,x',y',z,x}e^{-ik\cdot(w-y')}e^{-ik\cdot(z-y')\,\xi}\, \frac{1}{\xi\,(x-z)^{2}+(x-w)^{2}}\notag\\
&\quad \times\Bigg\{\Kv(y'-z)\cdot\Kv(x'-y')\bigg[\frac{\xi\big[(x-z)^{2}-(x-w)^{2}\big]}{(z-w)^{2}\,(1+\xi)^{2}}\notag\\
&\quad\quad\quad +\frac{\xi}{1+\xi}\Big((x-z)\cdot\Kv(z-w)-\frac{1}{2}(x-z)\cdot\Kv(x-w)-\frac{(x-z)^{2}-(x-w)^{2}}{2\,(z-w)^{2}}\Big)\notag\\
&\quad\quad\quad +\frac{1}{1+\xi}\Big((x-w)\cdot\Kv(z-w)+\frac{1}{2}(x-w)\cdot\Kv(x-z)-\frac{(x-z)^{2}-(x-w)^{2}}{2\,(z-w)^{2}}\Big)\bigg]\notag\\
&\quad\quad +\frac{(x-z)^{2}-(x-w)^{2}}{(1+\xi)\,(x-w)^{2}}\Big[(x-w)\cdot\Kv(x'-y')\, \Kv(z-w)\cdot\Kv(y'-z)\notag\\
&\quad\quad\quad +\frac{1}{2}\, \Kv(x-z)\cdot\Kv(y'-z)\, (x-w)\cdot\Kv(x'-y')\Big]\notag\\
&\quad\quad -\Big[(x-z)\cdot\Kv(y'-z)\, \Kv(z-w)\cdot\Kv(x'-y') -\frac{1}{2}(x-z)\cdot\Kv(y'-z)\, \Kv(x-w)\cdot\Kv(x'-y')\notag\\
&\quad\quad\quad +(x-w)\cdot\Kv(y'-z)\, \Kv(z-w)\cdot\Kv(x'-y') +\frac{1}{2}\, \Kv(x-z)\cdot\Kv(x'-y')\notag\\
&\quad\quad\quad \times (x-w)\cdot\Kv(y'-z)\Big]\Bigg\} \times\frac{n(x',x)}{N_{c}^{2}-1}\Big[ N_{c}{\rm Tr}\Big\{\big(U^{\dagger}(y')-U^{\dagger}(x')\big)\big(U(x)-U(z)\big)\Big\}\notag\\
&\quad\quad -{\rm Tr}\Big\{T^{e}U^{\dagger}(z)\big(U(w)-U(z)\big)T^{e}\big(U^{\dagger}(y')-U^{\dagger}(x')\big)U(z)\Big\}\Big]\notag\\
&\quad -\frac{1}{(2\pi)^{3}}\, \frac{g^{4}}{4\pi^{4}}\, \frac{1}{k^{+}} \int_{0}^{\infty}\frac{d\xi}{2\pi}\Big[\frac{1}{\xi\,(1+\xi)}\Big]_{+} \int_{w,x',y',z,x}e^{-ik\cdot(w-y')}e^{-ik\cdot(z-y')\,\xi}\, \frac{1}{(x-w)^{2}}\notag\\
&\quad \times\Big[(x-w)\cdot\Kv(x'-y')\, \Kv(z-w)\cdot\Kv(y'-z) +\frac{1}{2}\, \Kv(x-z)\cdot\Kv(y'-z)\, (x-w)\cdot\Kv(x'-y')\Big]\notag\\
&\quad \times\frac{n(x',x)}{N_{c}^{2}-1}\Big[ N_{c}{\rm Tr}\Big\{\big(U^{\dagger}(y')-U^{\dagger}(x')\big)\big(U(x)-U(z)\big)\Big\}\notag\\
&\quad\quad -{\rm Tr}\Big\{T^{e}U^{\dagger}(z)\big(U(w)-U(z)\big)T^{e}\big(U^{\dagger}(y')-U^{\dagger}(x')\big)U(z)\Big\}\Big]\notag\\
&\quad -\frac{1}{(2\pi)^{3}}\, \frac{g^{4}}{4\pi^{4}}\, \frac{1}{k^{+}} \int_{0}^{\infty}\frac{d\xi}{2\pi} \int_{w,x',y',z,x}e^{-ik\cdot(w-y')}e^{-ik\cdot(z-y')\,\xi}\, \frac{\xi}{\xi\,(x-z)^{2}+(x-w)^{2}}\notag\\
&\quad \times\bigg\{\Big[(x-z)\cdot\Kv(x'-y')\, \Kv(z-w)\cdot\Kv(y'-z) -\frac{1}{2}(x-z)\cdot\Kv(x'-y')\, \Kv(x-w)\cdot\Kv(y'-z)\Big]\notag\\
&\quad\quad \times\frac{n(x',x)}{N_{c}^{2}-1}\Big[ N_{c}{\rm Tr}\Big\{\big(U^{\dagger}(y')-U^{\dagger}(x')\big)\big(U(x)-U(z)\big)\Big\}\notag\\
&\quad\quad\quad -{\rm Tr}\Big\{T^{e}U^{\dagger}(z)\big(U(w)-U(z)\big)T^{e}\big(U^{\dagger}(y')-U^{\dagger}(x')\big)U(z)\Big\}\Big]\notag\\
&\quad\quad -\Big[(x-z)\cdot\Kv(x'-z)\, \Kv(z-w)\cdot\Kv(y'-z) -\frac{1}{2}(x-z)\cdot\Kv(x'-z)\, \Kv(x-w)\cdot\Kv(y'-z)\Big]\notag\\
&\quad\quad \times\frac{n(x',x)}{N_{c}^{2}-1}\Big[ N_{c}{\rm Tr}\Big\{\big(U^{\dagger}(z)-U^{\dagger}(x')\big)\big(U(x)-U(z)\big)\Big\}\notag\\
&\quad\quad\quad -{\rm Tr}\Big\{T^{e}U^{\dagger}(z)\big(U(w)-U(z)\big)T^{e}\big(U^{\dagger}(z)-U^{\dagger}(x')\big)U(z)\Big\}\Big]\bigg\}\notag\\
&\quad +\frac{1}{(2\pi)^{3}}\, \frac{g^{4}}{8\pi^{5}}\, \frac{1}{k^{+}} \int_{w,x',w',z,x}e^{-ik\cdot(w-w')}\, \frac{(x-w)^{2}}{\big[(x-z)^{2}-(x-w)^{2}\big]}\notag\\
&\quad \times \log\left[\frac{(x-z)^{2}}{(x-w)^{2}}\right] \bigg[\Kv(x-z)\cdot\Kv(x'-z)\, \Kv(z-w)\cdot\Kv(w'-z)\notag\\
&\quad\quad -\frac{1}{2}\, \Kv(x-z)\cdot\Kv(x'-z)\, \Kv(x-w)\cdot\Kv(w'-z)\bigg]\notag\\
&\quad \times \frac{n(x',x)}{N_{c}^{2}-1}\Big[ N_{c}{\rm Tr}\Big\{\big(U^{\dagger}(z)-U^{\dagger}(x')\big)\big(U(x)-U(z)\big)\Big\}\notag\\
&\quad\quad -{\rm Tr}\Big\{T^{e}U^{\dagger}(z)\big(U(w)-U(z)\big)T^{e}\big(U^{\dagger}(z)-U^{\dagger}(x')\big)U(z)\Big\}\Big]\, +{\rm c.c.}\, ,
\end{align}
\endgroup


\paragraph{Row III.V.} The symmetrised three-$\rho$ part is
\begin{align}\label{eq:S-III.V-3}
&\Big\langle\frac{d^{3}N^{\rm III.V}_{{\rm sym}}}{d^{3}k}\Big\rangle=\frac{1}{(2\pi)^{3}}\, \frac{ig^{4}}{4\pi^{4}}\, \frac{1}{k^{+}} \int_{0}^{\infty}\frac{d\xi}{2\pi} \int_{w,y',z,y,x,x'}e^{-ik\cdot(w-y')}e^{-ik\cdot(z-y')\,\xi}\Bigg[ \mathfrak{K}^{(y')}_{\rm III}\notag\\
&\quad\quad \times \frac{N_{c}}{(N_{c}^{2}-4)(N_{c}^{2}-1)}\, \Big(+i\Big)\notag\\
&\quad\quad \times \sigma(x',x,y)\Big[{\rm Tr}\Big\{T^{h}U^{\dagger}(x)\big(U(y')-U(x')\big)D^{h}\big(U^{\dagger}(y)-U^{\dagger}(w)\big)U(x)\Big\}\Big] + \mathfrak{K}^{(z)}_{\rm III}\notag\\
&\quad\quad \times \frac{N_{c}}{(N_{c}^{2}-4)(N_{c}^{2}-1)}\, \Big(+i\Big)\notag\\
&\quad\quad \times \sigma(x',x,y)\Big[{\rm Tr}\Big\{T^{h}U^{\dagger}(x)\big(U(z)-U(x')\big)D^{h}\big(U^{\dagger}(y)-U^{\dagger}(w)\big)U(x)\Big\}\Big]\Bigg] +{\rm c.c.}
\end{align}

and the contact part is
\begingroup\small
\begin{align}\label{eq:S-III.V-2}
&\Big\langle\frac{d^{3}N^{\rm III.V}_{{\rm con}}}{d^{3}k}\Big\rangle=\frac{1}{(2\pi)^{3}}\, \frac{ig^{4}}{4\pi^{4}}\, \frac{1}{k^{+}} \int_{0}^{\infty}\frac{d\xi}{2\pi} \int_{w,y',z,y,x,x'}e^{-ik\cdot(w-y')}e^{-ik\cdot(z-y')\,\xi}\Bigg[\notag\\
&\quad\quad \mathfrak{K}^{(y')}_{\rm III} \times \frac{1}{N_{c}^{2}-1}\bigg[\Big(+\frac{1}{2}i\Big)\delta^{(2)}(x'-x)\, n(x,y)\Big[\notag\\
&\quad\quad\quad\quad {\rm Tr}\Big\{T^{e}U^{\dagger}(x)\big(U(y)-U(w)\big)T^{e}\big(U^{\dagger}(y')-U^{\dagger}(x)\big)U(x)\Big\}\Big]\notag\\
&\quad\quad\quad +\Big(+\frac{1}{2}i\Big)\delta^{(2)}(y-x)\, n(x,x')\Big[ {\rm Tr}\Big\{T^{e}U^{\dagger}(x)\big(U(x)-U(w)\big)T^{e}\big(U^{\dagger}(y')-U^{\dagger}(x')\big)U(x)\Big\}\Big]\notag\\
&\quad\quad\quad +\Big(-\frac{1}{2}i\Big)\delta^{(2)}(x'-y)\, n(y,x)\Big[ {\rm Tr}\Big\{T^{e}U^{\dagger}(x)\big(U(y)-U(w)\big)T^{e}\big(U^{\dagger}(y')-U^{\dagger}(y)\big)U(x)\Big\}\Big]\bigg]\notag\\
&\quad\quad + \mathfrak{K}^{(z)}_{\rm III} \times \frac{1}{N_{c}^{2}-1}\bigg[\Big(+\frac{1}{2}i\Big)\delta^{(2)}(x'-x)\, n(x,y)\Big[\notag\\
&\quad\quad\quad\quad {\rm Tr}\Big\{T^{e}U^{\dagger}(x)\big(U(y)-U(w)\big)T^{e}\big(U^{\dagger}(z)-U^{\dagger}(x)\big)U(x)\Big\}\Big]\notag\\
&\quad\quad\quad +\Big(+\frac{1}{2}i\Big)\delta^{(2)}(y-x)\, n(x,x')\Big[ {\rm Tr}\Big\{T^{e}U^{\dagger}(x)\big(U(x)-U(w)\big)T^{e}\big(U^{\dagger}(z)-U^{\dagger}(x')\big)U(x)\Big\}\Big]\notag\\
&\quad\quad\quad +\Big(-\frac{1}{2}i\Big)\delta^{(2)}(x'-y)\, n(y,x)\Big[ {\rm Tr}\Big\{T^{e}U^{\dagger}(x)\big(U(y)-U(w)\big)T^{e}\big(U^{\dagger}(z)-U^{\dagger}(y)\big)U(x)\Big\}\Big]\bigg]\Bigg]\notag\\
&\quad +{\rm c.c.}
\end{align}
\endgroup


\paragraph{Row III.VI.} The symmetrised three-$\rho$ part is
\begin{align}\label{eq:S-III.VI-3}
&\Big\langle\frac{d^{3}N^{\rm III.VI}_{{\rm sym}}}{d^{3}k}\Big\rangle=\frac{1}{(2\pi)^{3}}\, \frac{ig^{4}}{4\pi^{4}}\, \frac{1}{k^{+}} \int_{0}^{\infty}\frac{d\xi}{2\pi} \int_{w,y',z,y,x,x'}e^{-ik\cdot(w-y')}e^{-ik\cdot(z-y')\,\xi}\Bigg[ \mathfrak{K}^{(y')}_{\rm III}\notag\\
&\quad\quad \times \frac{N_{c}}{(N_{c}^{2}-4)(N_{c}^{2}-1)}\, \Big(+i\Big)\notag\\
&\quad\quad \times \sigma(x',y,x)\Big[{\rm Tr}\Big\{T^{h}U^{\dagger}(y)\big(U(x)-U(z)\big)D^{h}\big(U^{\dagger}(y')-U^{\dagger}(x')\big)U(y)\Big\}\Big] + \mathfrak{K}^{(z)}_{\rm III}\notag\\
&\quad\quad \times \frac{N_{c}}{(N_{c}^{2}-4)(N_{c}^{2}-1)}\, \Big(+i\Big)\notag\\
&\quad\quad \times \sigma(x',y,x)\Big[{\rm Tr}\Big\{T^{h}U^{\dagger}(y)\big(U(x)-U(z)\big)D^{h}\big(U^{\dagger}(z)-U^{\dagger}(x')\big)U(y)\Big\}\Big]\Bigg] +{\rm c.c.}
\end{align}

and the contact part is
\begin{align}\label{eq:S-III.VI-2}
&\Big\langle\frac{d^{3}N^{\rm III.VI}_{{\rm con}}}{d^{3}k}\Big\rangle=\frac{1}{(2\pi)^{3}}\, \frac{ig^{4}}{4\pi^{4}}\, \frac{1}{k^{+}} \int_{0}^{\infty}\frac{d\xi}{2\pi} \int_{w,y',z,y,x,x'}e^{-ik\cdot(w-y')}e^{-ik\cdot(z-y')\,\xi}\Bigg[ \mathfrak{K}^{(y')}_{\rm III}\notag\\
&\quad\quad \times \frac{1}{N_{c}^{2}-1}\bigg[\Big(+\frac{1}{2}i\Big)\delta^{(2)}(x'-y)\notag\\
&\quad\quad\quad \times n(y,x)\Big[-{\rm Tr}\Big\{T^{e}U^{\dagger}(y)\big(U(x)-U(z)\big)T^{e}\big(U^{\dagger}(y')-U^{\dagger}(y)\big)U(y)\Big\}\Big]\notag\\
&\quad\quad\quad +\Big(+\frac{1}{2}i\Big)\delta^{(2)}(y-x)\notag\\
&\quad\quad\quad \times n(x,x')\Big[-{\rm Tr}\Big\{T^{e}U^{\dagger}(x)\big(U(x)-U(z)\big)T^{e}\big(U^{\dagger}(y')-U^{\dagger}(x')\big)U(x)\Big\}\Big]\notag\\
&\quad\quad\quad +\Big(-\frac{1}{2}i\Big)\delta^{(2)}(x'-x)\notag\\
&\quad\quad\quad \times n(x,y)\Big[-{\rm Tr}\Big\{T^{e}U^{\dagger}(y)\big(U(x)-U(z)\big)T^{e}\big(U^{\dagger}(y')-U^{\dagger}(x)\big)U(y)\Big\}\Big]\bigg] + \mathfrak{K}^{(z)}_{\rm III}\notag\\
&\quad\quad \times \frac{1}{N_{c}^{2}-1}\bigg[\Big(+\frac{1}{2}i\Big)\delta^{(2)}(x'-y)\notag\\
&\quad\quad\quad \times n(y,x)\Big[-{\rm Tr}\Big\{T^{e}U^{\dagger}(y)\big(U(x)-U(z)\big)T^{e}\big(U^{\dagger}(z)-U^{\dagger}(y)\big)U(y)\Big\}\Big]\notag\\
&\quad\quad\quad +\Big(+\frac{1}{2}i\Big)\delta^{(2)}(y-x)\notag\\
&\quad\quad\quad \times n(x,x')\Big[-{\rm Tr}\Big\{T^{e}U^{\dagger}(x)\big(U(x)-U(z)\big)T^{e}\big(U^{\dagger}(z)-U^{\dagger}(x')\big)U(x)\Big\}\Big]\notag\\
&\quad\quad\quad +\Big(-\frac{1}{2}i\Big)\delta^{(2)}(x'-x)\notag\\
&\quad\quad\quad \times n(x,y)\Big[-{\rm Tr}\Big\{T^{e}U^{\dagger}(y)\big(U(x)-U(z)\big)T^{e}\big(U^{\dagger}(z)-U^{\dagger}(x)\big)U(y)\Big\}\Big]\bigg]\Bigg] +{\rm c.c.}
\end{align}


\subsection{Remarks on the finite terms}\label{sec:fin_summary}

\paragraph{Scales.} The finite terms contain no renormalisation or factorisation scale and no rapidity cut-off: the dependence on $\mu$ is in the running-coupling and DGLAP terms, and the dependence on $\vee$ and $\Lambda$ in the BFKL and JIMWLK terms. Under a dilatation of all transverse coordinates, $\v x\to\lambda\v x$, $\v k\to\v k/\lambda$, every finite term scales like the LO spectrum.

\paragraph{The symmetrised and the contact terms.} After the average over a colour-neutral projectile the rows that are cubic in the charges contribute in two ways. Their contact terms are two-point terms of the same order as the genuine two-$\rho$ terms and even in $\v k$; they are part of the result for every projectile, whatever its three-point function. They do not cancel among the rows at the level of the integrand, and they cannot cancel against the contact terms of the three-$\rho$ evolution terms, which multiply $\log(\vee/k^+)$ or $\log(k^+/\Lambda)$ and are part of the BFKL and JIMWLK terms. In the dilute limit the contact parts of the real rows are required to reproduce the two-gluon production cross section of $k_T$ factorisation (Section~\ref{sec:check_dilute}). The symmetrised terms are odd in $\v k$, proportional to the symmetrised ($d^{abc}$) three-point function $\sigma$ of the projectile and to C-odd structures of the target; only for a projectile with $\sigma=0$, for instance one with a Gaussian distribution of the charges, or for $SU(2)$, are the finite terms even in $\v k$.

\paragraph{Integrability.} With the groupings of Section~\ref{sec:fin_together} the finite terms are integrable at fixed positions of the observed gluon. The statement rests on exact identities, on power counting at all coincidences of the integration points with the colour factors as written, on the pointwise large-distance behaviour in all directions and on Fourier integrals at fixed values of the other coordinates (Appendix~\ref{app:tails} and Section~\ref{sec:lit_finite}); a complete numerical integration of the finite terms over all transverse positions with Wilson lines has not been carried out (Section~\ref{sec:openpoints}).
\par\endgroup

% ======================================================================
\section{Limits, consistency checks and comparison with the literature}\label{sec:checks}
% ======================================================================
% Sources: paper/notes/README.md, notes/lit_dglap_rc.tex, notes/lit_finite_dilute.tex,
% inputs/reviews/review_2.txt (secs. 3-6, 11, 13), limits_checked/*.
% Revised after sections dglap_rc and finite were written (notes/fix_soft.tex, fix_poles_dglap.tex,
% fix_singlet.tex). Revised for referee report 3 (notes/ref3_*.tex).

The calculation involves a large number of terms, with non-commuting charges and Wilson lines in a definite order. In this section we list the checks to which the result has been subjected and compare it with known results. Unless stated otherwise, the algebraic identities quoted below were verified both analytically and numerically: the expressions were evaluated on a transverse lattice with random $SU(2)$ and $SU(3)$ Wilson lines in the adjoint representation and with either a completely general ordered charge correlator $\langle\rho^b(\v u)\rho^c(\v v)\rangle$ or the colour-singlet one, Eq.~\eqref{eq:singlet}. Since the Wilson lines are random, an identity that holds on the lattice to round-off error cannot hold by accident. The scripts used for the numerical checks are available from the authors on request.

\subsection{Limits}\label{sec:limits}

\paragraph{No target, pure-gauge target.} For $U=1$ the operator $\bar\Omega^\dagger\Omega$ is the identity and every amplitude vanishes. In the result this is manifest: every term of every row, and every term of the BFKL, JIMWLK, DGLAP and running-coupling contributions, contains a factor that vanishes at $U=1$, such as $[U(\v x')-U(\v w')]^{ab'}\rho^{b'}(\v x')$, its analogue with three Wilson lines, or a combination of $f^{abc}$ and Wilson lines that reduces to $N_c-N_c-N_c+N_c=0$. A constant $U$ is a global colour rotation, which cannot produce gluons; all terms vanish for constant $U$ as well, by $f^{abc}U^{bd}U^{ce}=U^{ah}f^{hde}$ and $U^TU=1$.

\paragraph{Counting of charges.} The BFKL and JIMWLK terms are of order $g^4\rho^2$. No term of order $g^4\rho^4$ survives (Section~\ref{sec:fourrho}), and the three-$\rho$ terms that multiply a rapidity logarithm cancel in the symmetrised sense (Section~\ref{sec:threerho}). The finite terms are of order $g^4\rho^2$ and $g^4\rho^3$. After the average over the projectile the $g^4\rho^3$ terms do not drop out: their symmetrised part is proportional to the symmetrised three-point function, i.e.\ to the three-point function $\sigma$ of the projectile, Eq.~\eqref{eq:A2d}, and the commutators produced by the reordering of the charges leave contact terms of order $g^4\langle\rho\rho\rangle$, of the same order as the quadratic finite terms (Section~\ref{sec:fin_singlet}). The three-$\rho$ terms that multiply a rapidity logarithm cancel in the symmetrised sense for any three-point function, so that no logarithm multiplies $\sigma$.

\paragraph{Reality.} Groups I and III are written with their complex conjugates and are real. Group II is Hermitian by itself: its kernels are symmetric under $(\v x,\v w)\leftrightarrow(\v x',\v w')$ and so is its colour operator, which we verified numerically to $10^{-12}$.

\paragraph{Odd terms.} The three-$\rho$ terms carry an explicit factor $i$ (in Row~I.II, $N_{2b}$, it is in the Fourier transform of an odd kernel), and all kernels depend on $\v k$ through functions with $f(-s)=f(s)^*$. A three-$\rho$ term plus its complex conjugate is therefore odd under $\v k\to-\v k$ in its totally symmetrised part and even in the part produced by the commutators of the ordered charges (Section~\ref{sec:fin_singlet}). For a globally colour-invariant projectile with the symmetrised three-point function~\eqref{eq:A2d} the odd part is the sum of the symmetrised three-$\rho$ terms of Section~\ref{sec:fin_rows}: it is proportional to the three-point function $\sigma$ of the projectile and contains the target only through the C-odd traces with $D^h$, Eq.~\eqref{eq:traceid}. It is given row by row by Eqs.~\eqref{eq:S-I.II-3}, \eqref{eq:S-I.III-3}, \eqref{eq:S-I.V-3}, \eqref{eq:S-II.I-3}, \eqref{eq:S-III.II-3}, \eqref{eq:S-III.III-3}, \eqref{eq:S-III.V-3} and~\eqref{eq:S-III.VI-3}, for a general three-point function $\sigma$ (no model is assumed), and it is odd configuration by configuration of the target; for a weak target it starts with three gluon exchanges. (Special cases, not part of the result: the odd part is absent for a projectile with $\sigma=0$, such as a Gaussian, McLerran--Venugopalan-type projectile, and it vanishes after an average over a C-even target ensemble.) This agrees with the classical-field result that the leading-order spectrum is even configuration by configuration and that the odd part comes from the term cubic in the projectile charge, averages out for a Gaussian projectile and is ``defined by odderon exchanges''~\cite{McLerran:2016snu}, and with the counting of~\cite{Kovchegov:2018jun}, three scatterings in the projectile and three in the target. The even commutator part does \emph{not} vanish for any projectile: the three-$\rho$ rows contribute to the averaged spectrum through two-point contact terms $\propto\delta^{(2)}(\v x_i-\v x_j)\,n$, which have no counterpart for classical charges. All colour factors of the averaged terms, the symmetrised parts and the parity statements were verified at the operator level, with charges represented by generators on two or three sites and globally colour-invariant density matrices, including a C-odd state with $\sigma\neq0$ (Section~\ref{sec:fin_singlet}); the operator decomposition of each cubic row into its symmetrised and commutator parts, Eqs.~\eqref{eq:symord}--\eqref{eq:symr21}, was verified with explicit $SU(3)$ charges on three sites for all assignments of the charge points to the sites.

\paragraph{Comparison with Li and Skokov.} The symmetrised three-$\rho$ terms are the counterpart, for non-commuting charges, of the first saturation correction $M_{(1)}^*M_{(3)}+{\rm c.c.}\propto\rho_P^3$ of Li and Skokov~\cite{Li:2021ntm}, whose charges are classical; the contact terms have no classical counterpart. The colour structures agree: their two-Wilson-line term has a vanishing $d^{abc}$ projection, as the corresponding structures here, so that this sector does not contribute under (A2); the odd part comes from three Wilson lines and is C-odd in the target; and in our result the invariants with four Wilson lines cancel between Rows~II.I and~III.II in the sum of the symmetrised terms (their kernels are equal and opposite, $\mathfrak K^{(1)}_{\rm II.I}=-\mathfrak K_{\rm III.II}$), so that the odd part has at most three Wilson lines, as in~\cite{McLerran:2016snu,Li:2021ntm}. The parities in $\v k$ do not agree, and this is an open point. Li and Skokov find also a part of $M_{(1)}^*M_{(3)}$ that is even in $\v k$ configuration by configuration (their two-line term and the real part of their three-line kernel), whose $d^{abc}$ projection is not removed by any symmetry, whereas our symmetrised cubic terms are purely odd in the principal-value prescription. Conversely, the principal-value result does not reproduce only their odd part: in the two-Wilson-line sector, where no on-shell denominator occurs, their term (from the classical non-linear Weizs\"acker--Williams field) is even in $\v k$, while ours (Row~I.II, $N_{2b}$) is odd and has a different kernel. The difference can be traced to a relative phase $i$. The transverse kernel of the classical $O(g^3)$ Weizs\"acker--Williams field of~\cite{Li:2021ntm} occurs in our wave function, as the coefficient of $\log(\Lambda/k^+)$ in the two-$\rho$ one-gluon component, Eq.~\eqref{eq:Psigrhorho} (to $10^{-12}$), but with a factor $i$ relative to the LO coefficient~\eqref{eq:PsiLO}, whereas classically the $O(g^3)$ field is in phase with the $O(g)$ one. In the same way the two-charge structures $[\mf A^{(3)}]_{\rho\rho}$ (Row~I.II), the $f\rho$ part of $\mf B_2$ relative to its $\{\rho,\rho\}$ part (Rows~II.I, III.II) and the rescattering of a produced gluon on the rotated second charge after the collision ($\mf B_3$, Row~I.III) carry a phase that differs by $i$ from the classical field and from LCPT with retarded propagators~\cite{Kovchegov:2018jun}, while for the merging of two produced gluons after the collision ($\mf C_1$) the assignments agree. A possible origin is the prescription for $1/k^+$, i.e.\ the $x^-$ profile of the field (retarded, as in the classical solution, or symmetric, which corresponds to real LO coefficients in momentum space); we have not checked this. The $\v k$-even part of~\cite{Li:2021ntm} has not been confirmed by an independent calculation either, and a kernel-level comparison of the odd three-line parts, which requires all cubic rows in momentum space, has not been carried out. The two-line difference is invisible in the averaged spectrum under (A2), since both two-line terms vanish after the $d^{abc}$ projection, but it is present configuration by configuration and in correlations.

\paragraph{Scaling.} Under a dilatation of all transverse coordinates, $\v x\to\lambda\v x$, $\v k\to\v k/\lambda$, every finite term scales like the LO spectrum. The finite terms contain no renormalisation scale; the dependence on $\mu$ is confined to the collinear and running-coupling terms (Section~\ref{sec:dglap_rc}).

\paragraph{Dilute target.} For $U=1+O(\alpha)$ the LO spectrum reduces to the $k_T$-factorised form, Eq.~\eqref{eq:kTfact}, with the Lipatov-type vertex. In this limit the BFKL and JIMWLK terms reduce to the LO BFKL evolution of the projectile and of the target unintegrated gluon distributions, respectively. The comparison of the NLO terms with $k_T$ factorisation is the subject of Section~\ref{sec:check_dilute}.

\subsection{The dilute--dilute limit}\label{sec:check_dilute}

\paragraph{The limit.} We write $U^{ab}=\delta^{ab}+u^{ab}$, $u^{ab}=f^{acb}\alpha^c$, keep $O(u)$ in the amplitude and in the conjugate amplitude (two-gluon exchange with the target) and the two-point function of the projectile (one gluon exchanged with the projectile on each side); since the charges do not commute, products of two or three charges contribute to this order through $[\rho^a(\v x),\rho^b(\v y)]=if^{abc}\rho^c(\v x)\delta^{(2)}(\v x-\v y)$, while the symmetrised three-point function $\sigma$ is dropped. With $\langle\alpha^c(\v l)\alpha^d(-\v l')\rangle=\delta^{cd}(2\pi)^2\delta^{(2)}(\v l-\v l')T(\v l)$ and $\langle\rho^b(\v q')\rho^{b'}(-\v q)\rangle=\delta^{bb'}(2\pi)^2\delta^{(2)}(\v q-\v q')P(\v q)$ as in Eq.~\eqref{eq:kTfact}, and
\begin{align}\label{eq:dilute_not}
\mathcal N_{\rm d}\equiv\frac{\alpha_sN_c(N_c^2-1)}{\pi^2}\,\frac{S_\perp}{k^+}\,,\qquad\phi_P(\v q)\equiv\frac{P(\v q)}{\v q^2}\,,\qquad\phi_T(\v l)\equiv\v l^2\,T(\v l)\,,
\end{align}
the unintegrated gluon distributions of projectile and target, the LO spectrum is $\frac{d^3N_{\rm LO}}{d^3k}=\mathcal N_{\rm d}\int\frac{d^2q}{(2\pi)^2}\,\phi_P(\v q)\,\phi_T(\v k-\v q)/\v k^2$.

\paragraph{Real emission (Groups II and III).} At $O(u)$ the anticommutator parts of the two-gluon amplitude~\eqref{eq:A2gamp} cancel, $\mathcal A_B^{(2)}+\mathcal A^{(1)}_{AA}+\mathcal A^{(2)}_{AA}=\frac14[L_k,\mf A_p]+\frac14[L_p,\mf A_k]+O(u^2)$, with $L$ the LO amplitude before the Fourier transform, so that the two-gluon amplitude is linear in $\rho$, and with $\v q_a$ the momentum from the projectile, $\v q_b=\v k+\v p-\v q_a$ that from the target, $x=k^+/(k^++p^+)$,
\begin{align}\label{eq:dilute_real}
\frac{d^3N^{\rm (II+III)}}{d^3k}\bigg|_{\rm dilute}=\mathcal N_{\rm d}\int\frac{d^2q_a}{(2\pi)^2}\int_\Lambda^\vee\frac{dp^+}{p^+}\int d^2p\,\frac{\alpha_sN_c}{\pi^2}\,\phi_P(\v q_a)\,\phi_T(\v q_b)\,\frac{A_{gg}(\v q_a,\v q_b;\v k,\v p;x)}{\v q_a^2\,\v q_b^2}\,,
\end{align}
where $A_{gg}$ is the Born cross section for the production of two gluons in quasi-multi-Regge kinematics~\cite{Fadin:1998py}, in the form given in Ref.~\cite{Bartels:2006hg} (their $A_{\rm gluons}$). We verified this identity numerically at random kinematics for all $0<x<1$, for $SU(2)$ and $SU(3)$, and for $d_\perp$ transverse polarisations ($d_\perp=1.7,3,6$) against the $(D-2)$ terms of~\cite{Bartels:2006hg} (relative deviation $2\times10^{-15}$; there is no free normalisation). It requires the commutator (contact) parts generated by the independent-emission terms, $\mf B_2$ with $\{\rho,\rho\}$ and $\mathcal A^{(1,2)}_{AA}$: without them, or without $\mathcal A_{CA}$, the ratio deviates from one by factors $0.2$--$10$. In the limits $p^+\gg k^+$ and $p^+\ll k^+$, $A_{gg}\to\v q_a^2\v q_b^2/(\v k^2\v p^2)$: these are the real parts of one step of LO BFKL evolution of $\phi_P$ over $\log(\vee/k^+)$ and of $\phi_T$ over $\log(k^+/\Lambda)$, with $\bar\alpha_s=\alpha_sN_c/\pi$; the virtual parts are those of Sections~\ref{sec:check_BFKL} and~\ref{sec:check_JIMWLK}. The coefficient of $\log\mu^2$ and the scale of the coupling, $1-\frac{\alpha_s\beta_0}{4\pi}\log\frac{\v k^2}{\mu^2}$ in the local form~\eqref{eq:dg_local}, agree with~\cite{Fadin:1998py,Bartels:2006hg}. Since $dp^+/p^+=dx/(x(1-x))$ is, at fixed $k^+$, minus the rapidity measure of the second gluon, the finite part of Eq.~\eqref{eq:dilute_real} (the $x$ integral of $A_{gg}/(\v q_a^2\v q_b^2)-1/(\v k^2\v p^2)$) coincides identically with the two-gluon contribution minus its multi-Regge limit of~\cite{Bartels:2006hg} for $s_\Lambda\to\infty$ (numerically at three points, to $10^{-6}$ at $s_\Lambda=10^6$, the difference falling like $1/s_\Lambda$): for the real emission the plus-momentum cut-offs require no conversion, at fixed $\v p\neq0$. For $|\v p|\lesssim\Lambda|\v k|/k^+$ the cut $p^+>\Lambda$ removes part of the soft region at the rapidity of the observed gluon, and the remainder is not uniform there. The two rapidity logarithms refer to plus-momentum cut-offs; on the projectile side this is the asymmetric choice $s_0=\v k_J^2$ of~\cite{Bartels:2006hg}, on the target side they use the minus-momentum fraction. The two schemes differ by $\log(\v p^2/\v k^2)$-weighted LO kernels and trajectories (energy-scale terms, not constants times LO); for the real emission the shift cancels between the two ends, for the virtual trajectories (split of the loop-gluon window at $p^+=k^+$ instead of $y_p=y_k$) it has not been evaluated.

\paragraph{Virtual correction (Group I) for a point-like quark.} For $\rho^a(\v x)=t^a\delta^{(2)}(\v x)$, Row~I.I and the two- and three-$\rho$ parts of $\mf A^{(3)}$ vanish ($f^{abc}\{t^b,t^c\}=0$, scaleless integrals), and we evaluated the remaining Rows~I.II (one-$\rho$ part), I.III (principal value, cut-offs), I.IV and I.V in momentum space, in $d_\perp=2-\epsilon$, $\MS$, at $\mu=|\v k|$, at two kinematic points. To the accuracy of the numerical integration ($2\times10^{-4}$ in the coefficients of the logarithms, $2\times10^{-3}$ in the constant),
\begin{align}\label{eq:dilute_groupI}
\frac{\text{Group I}}{\rm LO}\bigg|_{\rm point\ quark}&=2\omega(\v q_a)\log\frac{\vee}{k^+}+2\omega(\v q_b)\log\frac{k^+}{\Lambda}\notag\\&\quad+\frac{\alpha_sN_c}{\pi}\Big[\frac2\epsilon\log\frac{k^+}{\Lambda}+\log^2\frac{k^+}{\Lambda}\Big]+C(\v k,\v l)\,,\notag\\
2\omega(\v q)&\equiv\frac{\alpha_sN_c}{\pi}\Big[\frac2\epsilon-\log\frac{\v q^2}{\mu^2}\Big],
\end{align}
with $\v q_a=\v k+\v l$, $\v q_b=-\v l$ ($-\v l$ the momentum from the target); $\omega$ has the form of the one-loop gluon trajectory of~\cite{Fadin:1998py} in our normalisation. (i) The coefficient of $\log(\vee/k^+)$ is $2\omega(\v q_a)$ exactly, pole and finite part; it comes from Rows~I.III and I.IV. (ii) The coefficient of $\log(k^+/\Lambda)$ is $2\omega(\v q_b)$, at a kinematic point where $2\omega(\v q_a)$ is very different, but only in the sum of Rows~I.II, I.III, I.IV and I.V (Row~I.II alone has $\omega(\v k)$ and $\omega(\v q_a)$ on the two terms of the vertex); on top of it there is the $\v l$-independent soft--collinear term $\frac{\alpha_sN_c}{\pi}[\frac2\epsilon L+L^2]$, $L=\log(k^+/\Lambda)$, whose dependence on $\mu$ and $\v k^2$ was not separated. (iii) The $\frac{11}{3}$ poles cancel between Rows~I.II and~I.V, and the unresolved part of Row~I.V cancels the whole self-energy part of Row~I.II, $\xi$ by $\xi$ to $10^{-13}$, including the $\log(\xi(1-\xi))$ that produces $\frac{67}{9}$: Group~I contains no $\frac{67}{9}$ (Eq.~\eqref{eq:dg_SEid}). (iv) The kinematics-dependent double logarithms cancel between Row~I.III and the resolved part of Row~I.V; the remaining $\frac{\alpha_sN_c}{\pi}L^2$ is part of the soft--collinear term, so that Group~I alone is not free of $\log^2(k^+/\Lambda)$. (v) The finite remainder $C(\v k,\v l)$ is a function, not a constant: $0.055$ and $0.305$ times LO at the two points. (vi) In this limit the result does not depend on the prescription for the vanishing energy denominators: the on-shell part of $\mf B_3$ is purely imaginary relative to LO and drops out.

\paragraph{What is and what is not tested.} The one-loop central-emission vertex~\cite{Fadin:1998py,DelDuca:1998cx,Bartels:2006hg} regulates the soft--collinear region by $\epsilon^{-2}(\boldsymbol\Delta^2/\mu^2)^\epsilon$, $\boldsymbol\Delta=\v q_a+\v q_b=\v k$, and has no rapidity cut-off; in Eq.~\eqref{eq:dilute_groupI} the same region is regulated by the plus-momentum cut-off. The two virtual parts therefore differ already in their pole terms by a function of $\boldsymbol\Delta$ that is not of the form $F_P(\v q_a)+F_T(\v q_b)$, the form that at fixed order could be absorbed into $\phi_P$ and $\phi_T$; the scheme-independent comparison requires the virtual and the real parts with matched regulators, including the real emission at $|\v p|\lesssim\Lambda|\v k|/k^+$ with the cut $p^+>\Lambda$, which in dimensional regularisation produces $\Lambda$-dependent soft--collinear terms of the type $\frac2\epsilon L$ and $L^2$. This comparison has not been carried out: whether the extra term of Eq.~\eqref{eq:dilute_groupI} cancels against that real contribution, and whether the finite remainder agrees with the one-loop vertex, is open. No disagreement with~\cite{Fadin:1998py,Bartels:2006hg} was found. The tests of this section are made at the level of the amplitudes and coefficients, Eqs.~\eqref{eq:A2gamp}, \eqref{eq:B2mom}, \eqref{eq:C1} and~\eqref{eq:A1}; the dilute limit of the coordinate-space finite terms as printed, and the equality of the printed colour-singlet contact terms with the commutator contributions required by Eq.~\eqref{eq:dilute_real}, have not been checked. For a general (not point-like) projectile Group~I has not been reduced in this limit.

\subsection{Leading order}

The LO spectrum, Eq.~\eqref{eq:LOsinglet}, coincides with the result of Kovchegov and Tuchin~\cite{Kovchegov:2001sc}, including its normalisation: for a quark projectile, $n\to C_F$, the prefactor is $\alpha_sC_F/(4\pi^4)$. The target operator $\mathcal O_{\rm LO}$ is the one of Kovner, Lublinsky and Weigert~\cite{Kovner:2006wr,Kovner:2006ge}.

\subsection{Projectile evolution: the BFKL term}\label{sec:check_BFKL}

\paragraph{Internal consistency.} The ordered form~\eqref{eq:BFKLrho} is an exact regrouping of the term-by-term list of all two-$\rho$ terms proportional to $\log(\vee/k^+)$, for an arbitrary ordered correlator (relative difference $\lesssim10^{-15}$ for $SU(2)$, $10^{-14}$ for $SU(3)$), and Eq.~\eqref{eq:BFKLsinglet} is its colour-singlet reduction to the same accuracy.

\paragraph{Exact one-gluon evolution of a colour-singlet projectile.} As an independent test we computed the change of the spectrum due to one additional gluon with $k^+<p^+<\vee$ in a colour-singlet projectile directly, without the expansion in $\rho$ used in the rest of the paper. Any set of colour sources $s$ at positions $\v x_s$, with charges $Q_s^e$, produces the observed gluon with the operator $\sum_{s,s'}Q^e_{s'}Q^e_sF(\v x_s,\v x_{s'})$, where
\begin{align}
F(\v a,\v b)=\int_{w,w'}e^{-i\v k\cdot(\v w-\v w')}\Kv(\v a-\v w)\cdot\Kv(\v b-\v w')\,{\rm Tr}\big[(U^\dagger(\v a)-U^\dagger(\v w))(U(\v b)-U(\v w'))\big]\,,
\end{align}
since the colour rotation of the sources cancels in the inclusive sum. For a quark--antiquark dipole at $\v x_1$, $\v x_2$ the three charges of the $q\bar qg$ colour singlet satisfy $Q_1^2=Q_2^2=C_F$, $Q_g^2=N_c$, $Q_1\cdot Q_2=1/(2N_c)$ and $Q_{1,2}\cdot Q_g=-N_c/2$, and the probability to emit the additional gluon at $\v z$ is $\frac{g^2C_F}{4\pi^3}\log\frac{\vee}{k^+}K_D(\v x_1,\v x_2,\v z)=\frac{\alpha_sC_F}{\pi^2}\log\frac{\vee}{k^+}K_D(\v x_1,\v x_2,\v z)$, i.e.\ $-2\mf N_2$ of Eq.~\eqref{eq:N2} for the dipole. With the LO spectrum of the dipole, $\frac{d^3N_{\rm LO}}{d^3k}=\frac{\alpha_s}{4\pi^4k^+}\mathcal D_{12}$ by Eq.~\eqref{eq:LOsinglet}, real minus virtual gives, exactly at finite $N_c$,
\begin{align}
\delta\Big[\frac{d^3N}{d^3k}\Big]&=\frac{\alpha_sN_c}{2\pi^2}\log\frac{\vee}{k^+}\int_zK_D(\v x_1,\v x_2,\v z)\big[\mathcal D_{1z}+\mathcal D_{z2}-\mathcal D_{12}\big]\times\frac{\alpha_s}{4\pi^4k^+}\,,\notag\\
\mathcal D_{ab}&=\frac{1}{2N_c}\big[F_{aa}+F_{bb}-F_{ab}-F_{ba}\big]\,,
\end{align}
the dipole BFKL (linearised BK) kernel acting on the projectile. Substituting $n=C_F[\delta_{11}+\delta_{22}-\delta_{12}-\delta_{21}]$ into Eq.~\eqref{eq:BFKLsinglet} reproduces this result pointwise in $\v z$. Since every symmetric $n(\v x,\v y)$ with vanishing row sums, i.e.\ every colour-singlet projectile, is a linear combination of such dipole correlators, and Eq.~\eqref{eq:BFKLsinglet} is linear in $n$, the BFKL term is the exact leading-logarithmic projectile evolution for every colour-singlet projectile. The calculation with fundamental charges is exact to all orders in the charges, which shows that the expansion in $\rho$ loses nothing for a dipole.

\paragraph{Literature.} This is the structure found by Kovner and Lublinsky for the single inclusive spectrum~\cite{Kovner:2006wr} and by Kovner, Lublinsky and Weigert~\cite{Kovner:2006ge}: the projectile correlator evolves with BFKL from the projectile rapidity down to the rapidity of the produced gluon. It is the first-order expansion of the leading-logarithmic factorisation theorem of Gelis, Lappi and Venugopalan~\cite{Gelis:2008rw}, and it satisfies the general structure of the NLO correction discussed by Gelis and Laidet~\cite{Gelis:2012ri}.

\subsection{Cubic terms on exact charge states: a quark and a dipole}\label{sec:check_cubic}

For a single quark, $\rho^a(\v x)=t^a\delta^{(2)}(\v x-\v x_0)$ in a colour-averaged state, the three-point function is exact: $\langle\rho^a\rho^b\rho^c\rangle=\frac1{N_c}{\rm tr}(t^at^bt^c)=\frac{1}{4N_c}(d^{abc}+if^{abc})$ at $\v x_0$, so that (A2) holds with $\sigma=\frac{(N_c^2-4)(N_c^2-1)}{4N_c^2}=C_F\frac{N_c^2-4}{2N_c}$ ($\frac{10}{9}$ for $SU(3)$; $d^{abc}t^at^bt^c=C_F\frac{N_c^2-4}{2N_c}$) and $n=C_F$. Each cubic colour structure then has the unique value $\frac{1}{4N_c}[\Delta+i\Phi]$ for an ordered product and $\frac{1}{4N_c}\Delta$ for a product with an anticommutator (${\rm tr}\,t^a\{t^b,t^c\}=\frac12d^{abc}$), with $\Delta=\sum C^{abc}d^{abc}$ and $\Phi=C^{abc}f^{abc}$ at $\v x_0$. Eq.~\eqref{eq:cubicavg} reproduces this term by term: the symmetrised part is the $d^{abc}$ half, the contact part the $f^{abc}$ half, and for an antiquark the first changes sign and the second does not. For a colour-singlet dipole, quark at $\v x_0$ and antiquark at $\v x_1$, $n=C_F\,s_is_j$ and $\sigma=\sigma_q\,s_is_js_k$, with $s=+1$ at the quark and $-1$ at the antiquark; the symmetrised part of a cubic term carries $s_is_js_k$ and is the part odd under the exchange of quark and antiquark, i.e.\ under the C-conjugation of the dipole, while the contact part carries products of two $s$'s and is even. The $\v k$-odd cubic term of a dipole is therefore the part of the spectrum that distinguishes the quark end from the antiquark end, $\frac12[\sigma(\v x_0,\v x_1)-\sigma(\v x_1,\v x_0)]$; it needs a C-odd target and vanishes for an orientation-averaged dipole, whereas the $\v k$-even cubic (contact) part does not. We evaluated all thirteen cubic colour structures with explicit $SU(2)$, $SU(3)$ and $SU(4)$ matrices for the quark, the antiquark, the dipole and the conjugate dipole, all assignments of the charge points, and random Wilson lines: the exact ordered expectation values agree with the split~\eqref{eq:cubicavg} to $5\times10^{-14}$, and the C-split $\frac12[\sigma(P)\mp\sigma(P^{\rm C})]$ to $8\times10^{-15}$. For a quark or a dipole neither part is enhanced: both are of the order of the quadratic terms (colour weights relative to the LO weight for a dense random $SU(3)$ target: on average $0.12$ (range $0.05$--$0.23$ over the structures) for the symmetrised and $0.8$ (range $0.4$--$1.1$) for the contact part; for a weak target the symmetrised part is suppressed by one power of the target field). For a projectile of $m$ independent dipoles $n$ and $\sigma$ are both proportional to $m$. The language of a ``first saturation correction'' therefore applies to the symmetrised part only for a coherent dense projectile with large $\sigma$; this is a property of the projectile model, not of the calculation. The kernels at the triple coincidence of the charge points, which a single quark requires, are integrable (Remark (iii) of Section~\ref{sec:fin_singlet}).

\subsection{Target evolution: the JIMWLK term}\label{sec:check_JIMWLK}

We computed $-H_{\rm JIMWLK}\,\mathcal O_{\rm LO}$ independently, with the Hamiltonian~\eqref{eq:HJIMWLK} in the normalisation of~\cite{Lublinsky:2016meo,Kovner:2013ona}, for a general ordered correlator. The term-by-term list of the two-$\rho$ terms proportional to $\log(k^+/\Lambda)$, the grouped form~\eqref{eq:JIMWLKrho} and $-H_{\rm JIMWLK}$ acting on the LO spectrum agree to $10^{-14}$ for $SU(2)$ and $SU(3)$, and the colour-singlet form~\eqref{eq:JIMWLKsinglet} agrees with the colour-singlet projection to $10^{-15}$. The JIMWLK term is therefore the JIMWLK evolution of the LO spectrum, with the correct sign and normalisation, for an arbitrary projectile and not only for a colour-singlet one. Together with Section~\ref{sec:check_BFKL}, this confirms at fixed order the factorisation of the rapidity logarithms at the rapidity of the produced gluon~\cite{Kovner:2006wr,Kovner:2006ge,Gelis:2008rw}.

\subsection{Wave function}\label{sec:check_wf}

The light-cone wave function from which the coefficients of $\Omega$ are obtained was computed to order $g^3$, with the same longitudinal cut-offs ($\vee=\Lambda e^{\delta Y}$) and $d_\perp=2-\epsilon$, by Lublinsky and Mulian~\cite{Lublinsky:2016meo}. Our components agree with theirs as follows:
\begin{center}
\renewcommand{\arraystretch}{1.2}
\small
\begin{tabular}{@{}p{4.6cm}p{2.3cm}p{6.6cm}@{}}
\toprule
component & Ref.~\cite{Lublinsky:2016meo} & comparison\\
\midrule
LO soft gluon, normalisation $\mf N_2$ & eqs.~2.50--2.52 & identical\\
two gluons, $O(g^2)$ & eqs.~3.95--3.96 & identical\\
one gluon, $O(g^3)$, linear in $\rho$ & eqs.~3.97--3.98 & identical for $N_f=0$, incl.\ $\frac{11}{3}$, $\frac{67}{9}-\frac{\pi^2}{3}$, $\log^2$ cut-off terms\\
one gluon, $O(g^3)$, $f^{abc}\{\rho^b,\rho^c\}$ & eqs.~3.100--3.101 & equal after the antisymmetrisation enforced by $f^{abc}\{\rho^b,\rho^c\}$\\
one gluon, $O(g^3)$, $\rho\rho\rho$ & eq.~3.102 & identical up to the operator ordering, see below\\
\bottomrule
\end{tabular}
\end{center}
In the two-$\rho$ component the two forms differ by terms symmetric under the exchange of the two charges, which do not contribute. The coordinate-space forms of the $\{\rho,\rho\}$ component, Eq.~\eqref{eq:A3tworho}, and of the coefficients go beyond the published material. The operator ordering of the three-$\rho$ component, $\rho^a(\v w)\rho^b(\v x)\rho^b(\v y)$, i.e.\ $[\mf A^{(3)}]_{\rho\rho\rho}=\mf A\,\mf N_2$, follows from third-order perturbation theory with operator-valued matrix elements; we verified it against a construction of $\Omega$ and an exact diagonalisation with finite-dimensional matrices representing the charges (agreement to $3\times10^{-16}$; the opposite ordering deviates at order one).

\subsection{The constant \texorpdfstring{$\frac{67}{9}-\frac{\pi^2}{3}$}{67/9-pi2/3}}\label{sec:lit_constant}

The term of Row I.II proportional to the LO structure is, with its complex conjugate,
\begin{align}\label{eq:constantLO}
\frac{\alpha_sN_c}{2\pi}\Big(\frac{67}{9}-\frac{\pi^2}{3}\Big)\frac{d^3N_{\rm LO}}{d^3k}\,.
\end{align}
It follows directly from the one-$\rho$ part of $\mf A^{(3)}$, Eq.~\eqref{eq:A3onerho}, and the normalisation of the LO amplitude. Before the cancellation of the ultraviolet $\beta_0$ logarithms against Row I.V, Row I.II with its complex conjugate is $\frac{\alpha_sN_c}{2\pi}\big[\frac{11}{3}\big(2\gamma_E+\frac12\log\frac{X^2\mu^2_{\MS}}{4}+\frac12\log\frac{X'^2\mu^2_{\MS}}{4}\big)+\frac{67}{9}-\frac{\pi^2}{3}\big]\times$LO, $X=x-w$, $X'=x'-w'$. This is, term by term at $N_f=0$, the bare one-gluon kernel $K_{JSJ}$ of the NLO JIMWLK Hamiltonian in the form of Kovner, Lublinsky, Skokov and Zhao~\cite{Kovner:2023vsy}, with the same sign, size and scale; since both are built from the wave function of~\cite{Lublinsky:2016meo}, this is a consistency check rather than an independent comparison. The inclusive quantities in which the same constant appears, the NLO BK kernel of Balitsky and Chirilli~\cite{Balitsky:2007feb}, the NLO BFKL kernel of Fadin and Lipatov~\cite{Fadin:1998py} and the NLO jet vertex of Bartels, Sabio Vera and Schwennsen~\cite{Bartels:2006hg}, have the coefficient $\frac{\alpha_sN_c}{4\pi}(\frac{67}{9}-\frac{\pi^2}{3})$, half of Eq.~\eqref{eq:constantLO}. These quantities are inclusive over the collinear gluon pair; as noted in~\cite{Kovner:2023vsy}, half of the constant sits there in the two-gluon terms, and in~\cite{Fadin:1998py} it appears only after the real two-gluon production is integrated over the pair at fixed produced momentum. The amplitude-level one-loop Lipatov vertex~\cite{DelDuca:1998cx} contains no $\frac{67}{9}$.

\paragraph{The constant is a property of the grouping.} Eq.~\eqref{eq:constantLO} is not a constant of the result. By the identity~\eqref{eq:dg_SEid}, the self-energy part of Row~I.II, which carries the whole constant $\frac{67}{9}$ together with the $\beta_0$ logarithm, is cancelled identically by the part of Row~I.V in which the target rotates the gluon pair at its centre of mass. The constant~\eqref{eq:constantLO} is therefore a property of the grouping of Rows~I.II and~I.V; in the grouping used in this paper (colour of the pair at the position of one of its gluons) the compensating terms are in the finite terms $N_{2a}$ and $N_{2b\,\rm rest}$ of Row~I.V. This is confirmed in momentum space in the dilute limit with a point-like projectile, Section~\ref{sec:check_dilute}, where Group~I contains no $\frac{67}{9}$ and its finite part is a function of the kinematics, not a multiple of LO. The value of the LO-proportional constant left in $N_{2b\,\rm rest}$ in our grouping, for a general projectile and target, has not been evaluated. We do not read Eq.~\eqref{eq:constantLO} as the LO-type constant of the single inclusive spectrum, and we make no statement about where ``the other half'' of the inclusive value resides.

\paragraph{Inclusive combinations.} We have not established that an inclusive combination of our collinear terms reproduces the value of the inclusive kernels, and there is no general reason why it should: in the one-hadron-inclusive gluon coefficient function of $H\to gX$ the constant $\frac{67}{9}$ appears only in the $\delta(1-z)$ part~\cite{Almasy:2011eq}, and the energy-weighted integral is the process-dependent total rate. For Row~II.IV the momentum-weighted ($\int d\xi\,\xi$) collinear terms give, exactly, $\frac{\alpha_sN_c}{4\pi}\big(\frac{65}{18}-\frac{4\pi^2}{3}\big)$ times LO in our scheme (zero in the scheme of CXY), plus a term of linearly polarised structure from the tensor term of Eq.~\eqref{eq:S-II.IVrem}, which is not proportional to LO; the corresponding regions of Rows~II.I, III.III and I.V are not of LO form and have not been integrated. In the dilute limit the real rows are exactly the two-gluon production cross section of quasi-multi-Regge kinematics, Eq.~\eqref{eq:dilute_real}; hence for a pair-inclusive quantity built from our amplitudes (virtual correction at fixed produced momentum plus two-gluon production integrated at fixed pair momentum) the real part contributes the value $\frac{\alpha_sN_c}{4\pi}(\frac{67}{9}-\frac{\pi^2}{3})$ of~\cite{Fadin:1998py}, provided the real rows are integrated with $d_\perp$ polarisations, as in our scheme.

\subsection{Collinear and running-coupling terms}\label{sec:lit_dglap}
% Sources: notes/lit_dglap_rc.tex, notes/fix_poles_dglap.tex; consistent with sec_dglap_rc.tex.

The collinear term comes from Row II.IV, in which the observed gluon is produced in the splitting $g\to gg$ of a gluon of momentum $k^+/\xi$ that was produced at LO; it is a final-state splitting. Compared with the forward hadron production results of Chirilli, Xiao and Yuan~\cite{Chirilli:2011km,Chirilli:2012jd} and of Altinoluk et al.~\cite{Altinoluk:2014eka}:
\begin{itemize}
\item The splitting function has the normalisation of the LO $P_{gg}$, $2N_c[\frac{\xi}{(1-\xi)_+}+\frac{1-\xi}{\xi}+\xi(1-\xi)]+\frac{\beta_0}{2}\delta(1-\xi)$ at $N_f=0$, and the coefficient $\beta_0/2$ of $\delta(1-\xi)$ is the one required by the momentum sum rule $\int_0^1d\xi\,\xi P_{gg}(\xi)=0$. The prefactor $\frac{\alpha_s}{2\pi}$, the factor $\xi^{-2}$ with the phase $e^{-i\v k\cdot(\v y-\v y')/\xi}$, the sign and the dependence on $\mu$ are those of a final-state $\MS$ fragmentation of the LO spectrum of the parent gluon.
\item The end point $\xi\to1$ (integrated gluon soft) is plus-subtracted into the JIMWLK term, as the rapidity divergence is subtracted into BK evolution in~\cite{Chirilli:2012jd}. The end point $\xi\to0$ (observed gluon soft with respect to its parent) is subtracted into the BFKL term; the combination of the two subtractions is exactly the ``double plus'' prescription of~\cite{Kovner:2023vsy}, for the same reason: both poles are part of the LL evolution. In the forward (hybrid) case $\xi$ is bounded from below and no subtraction at $\xi\to0$ arises.
\item The argument of the collinear logarithm is the transverse separation of the observed gluon in the amplitude and in the conjugate amplitude, $(\v y-\v y')/\xi$, whereas in the $\MS$ form of~\cite{Chirilli:2012jd,Altinoluk:2014eka} it is the separation of the parent, $\v y-\v y'$. The bare collinear term has no $\log\xi^2$; the two forms differ in the $d$-dimensional Jacobian of the $\MS$ counterterm, $\xi^{-d}$ for the fragmentation function used here, which is defined as an $\MS$ number density obeying the momentum sum rule, and $\xi^{-2}$ (the form of~\cite{Chirilli:2012jd}), by the finite, $\mu$-independent term $\Delta=\mathcal C\int d\xi\,\mathcal P_{gg}(\xi)\log\xi^2\,\mathcal G_0(\xi)$, Eq.~\eqref{eq:dg_Delta} (Section~\ref{sec:dg_msbar}). The convention $\xi^{-d}$ is the standard $\MS$ one~\cite{Catani:1996vz}; with it the gluon coefficient functions of $e^+e^-\to gX$~\cite{Altarelli:1979kv,Kretzer:2000yf} and $H\to gX$~\cite{Almasy:2011eq} are reproduced exactly, while with $\xi^{-2}$ they are off by $-\frac{\alpha_s}{2\pi}P(z)\log z^2$ times the Born term (Section~\ref{sec:dg_msbar}). The results of~\cite{Chirilli:2012jd,Altinoluk:2014eka} are meant for standard $\MS$ fragmentation functions, which obey the momentum sum rule; with them their collinear term misses $\Delta$, i.e.\ it is the coefficient for the functions $D^{\rm CXY}$, which do not. Our form coincides with the cut-off scheme of~\cite{Altinoluk:2014eka}, with $\mu^2_{\MS}=c_0^2\mu^2_{\rm cut}\approx1.26\,\mu^2_{\rm cut}$, $c_0=2e^{-\gamma_E}$. The subtraction at $\xi\to0$ is written explicitly; its left-over, the remnant~\eqref{eq:remnant}, is a finite term of Row~II.IV.
\item No counterpart arises to the initial-state DGLAP evolution of the projectile parton distribution of the hybrid formalism: the projectile here is described by an unintegrated charge correlator, whose small-$x$ evolution is the BFKL term.
\item The $\beta_0$ terms. The real emission produces no $\delta(1-\xi)$ term, and the virtual $\frac{11}{3}$ poles of the vertex correction (Row I.II) cancel against those of Row I.V, for any grouping of the finite terms (Section~\ref{sec:dg_poles}); which numbers $\frac{11}{6}$ are $\beta_0/(2N_c)$ is identified there by a BLM-type quark loop (the calculation itself is for $N_f=0$). The term $\frac{\beta_0}{2}\delta(1-\xi)$ of $P_{gg}$ is added and subtracted, the subtracted term being the running-coupling term; with it the LO coupling becomes $g^2(\mu)$ at the scale of the collinear subtraction, or $g^2(c_0^2/r^2)$ with $r$ the transverse separation conjugate to $\v k$. In the classification of~\cite{Kovner:2023vsy}, where of the $2\beta_0$ of the NLO JIMWLK kernel one is the ultraviolet running of the coupling and one a DGLAP-type resolution logarithm, our explicit $\beta_0$ terms are, in the grouping used here, of the second, collinear, type. This agrees in structure with Kovchegov and Weigert~\cite{Kovchegov:2007vf} and with Horowitz and Kovchegov~\cite{Horowitz:2010yg}, in whose calculation (quark--quark scattering, BLM resummation of quark bubbles) the observable includes collinear pairs whose invariant mass, i.e.\ the virtuality of the produced-gluon line, is below a fixed cut-off $\Lambda_{\rm coll}$, and the coupling of the emission vertex of the produced gluon is $\alpha_s(\Lambda^2_{\rm coll}e^{-5/3})$ (Eqs.~(3.10)--(3.16) of~\cite{Horowitz:2010yg}); $\Lambda_{\rm coll}$ does not depend on the exchanged momenta and plays the role of the factorisation scale of the fragmentation function~\cite{Kovchegov:2007vf}, while the remaining couplings run with $\v q^2$, $(\v k-\v q)^2$ and a complex scale $Q(\v q,\v k-\v q)$, Eqs.~(3.31), (3.32) of~\cite{Horowitz:2010yg}. The other couplings of~\cite{Horowitz:2010yg}, associated with the $t$-channel exchanges, are not renormalised in this calculation.
\item The poles. All poles $1/\epsilon$ are accounted for in Section~\ref{sec:dg_poles}: the ultraviolet poles of the vertex correction (Row~I.II) and of the recombination term (Row~I.V) cancel each other; the coupling counterterm removes the ultraviolet half of the normalisation of the dressed produced gluon ($Z_3$ in light-cone gauge); its infrared half completes the real splitting function to the full $\MS$ $\mathcal P_{gg}$, whose collinear pole is absorbed by the fragmentation counterterm. The poles proportional to $\log(k^+/\Lambda)$ cancel within the JIMWLK sector, and the only remainder, $\propto\log(1-k^+/\vee)=O(k^+/\vee)$, is a power correction outside the eikonal approximation, which we define as the leading term in $k^+/\vee$ and $\Lambda/k^+$ at fixed $\epsilon$ (Section~\ref{sec:dg_poles}); it is not a cancellation.
\end{itemize}

\subsection{Finite terms}\label{sec:lit_finite}
% Sources: notes/fix_soft.tex, notes/fix_singlet.tex, notes/lit_finite_dilute.tex, notes/fix_poles_dglap.tex.

The finite terms of Section~\ref{sec:finite} have the following properties. (i) All finite terms vanish for $U=1$ and for constant $U$. (ii) Group~II is real without its complex conjugate. (iii) In Group~II, Row~I, individual terms have denominators that vanish on whole surfaces in the transverse plane (e.g.\ $(\v x'-\v w')^2(\v x-\v z)^2=(\v x-\v w)^2(\v x'-\v z)^2$), but the sum is finite and continuous across them, so that no principal-value prescription is needed there; this is the mechanism known from the kernels $\log(X^2Y'^2/X'^2Y^2)/(X^2Y'^2-X'^2Y^2)$ of NLO JIMWLK~\cite{Balitsky:2013fea,Kovner:2013ona}. In Row~I.III, where the energy denominators of $\mf B_3$ do vanish, the principal value gives logarithms of absolute values, which in the combined kernel~\eqref{eq:K-z_rm_I_III} combine into logarithms of positive arguments, with no pole left (Section~\ref{sec:fin_conv}). (iv) The kernels are integrable in all coincidence limits, including those of the combined kernel of Row~I.III and those created by the contact terms of the averaged three-$\rho$ rows; the only contact kernel that is not integrable by itself, $K(\v x,\v x,\v z)=1/(\v x-\v z)^2$, comes with colour factors that vanish at $\v z=\v x$ (Section~\ref{sec:fin_integrable}). (v) The collinear limit of Row~II.IV agrees with collinear factorisation (Section~\ref{sec:lit_dglap}), and its limits $\xi\to0$ and $\xi\to1$ reproduce the coefficients of the projectile-side and target-side soft-gluon emission, as in the corresponding limits of the NLO jet vertex~\cite{Bartels:2006hg}. In the dilute limit the real emission agrees exactly with the quasi-multi-Regge two-gluon production of~\cite{Fadin:1998py,Bartels:2006hg}, Eq.~\eqref{eq:dilute_real}; for a point-like projectile the virtual part has the Regge trajectories, Eq.~\eqref{eq:dilute_groupI}, but its finite part has not been compared with the one-loop vertex (Section~\ref{sec:check_dilute}).

\paragraph{Soft remainders and double logarithms.} The rows in which the observed gluon recoils against a soft gluon contain, before the limit $\Lambda\to0$, double logarithms $\log^2(k^+/\Lambda)$ with coefficients given by the large-distance tails of the rows; their sum vanishes for an arbitrary projectile, in the symmetrised two- and three-$\rho$ sectors (Section~\ref{sec:fin_soft}). After the limit, the soft remainders of the plus prescriptions, which have $\log|z|/\v z^2$ tails row by row, cancel in four pairs, each made of a virtual (Group~I) and a real (Group~III) term with the same coefficient, colour structure and kernel, the recoil acting on the observed gluon in the amplitude and in the conjugate amplitude, respectively (Table~\ref{tab:pairs}). The real-minus-virtual combination is the analogue of the remainder $[A(y-z)-A(\bar y-z)]^2[s(y,z)s(z,\bar y)-s(y,\bar y)]$ of the kinematically constrained forward production of~\cite{Altinoluk:2014eka}, and the soft remainders can equivalently be absorbed into a JIMWLK term in which the soft gluon is restricted to rapidities below that of the observed gluon; Eq.~\eqref{eq:kinconstr}, as in the kinematically constrained formulations of~\cite{Altinoluk:2014eka,Ducloue:2016shw} (the constraint of~\cite{Ducloue:2016shw} has not been compared in detail). The analogous remainders of forward production were discussed in~\cite{Stasto:2013cha,Altinoluk:2014eka,Ducloue:2016shw}. The cancellation requires no condition on the Wilson lines at large distance. The term of Row~I.III that is proportional to $\log[(\v z-\v w)^2/(\v x-\v w)^2]$ and made the row non-integrable is cancelled by the soft remainder of Row~III.VI together with the angle-odd tail of the other terms of Row~I.III, after the Fourier integral over the position of the observed gluon (Section~\ref{sec:fin_GIIII}). The instantaneous terms of Group~III, which converge for $\vee\to\infty$ only as distributions, are written with $\pi\delta(s)-i\mathcal P\frac1s$ (Section~\ref{sec:fin_inst}).

\paragraph{Large-distance tail of Row II.IV.} The tail of Row~II.IV at large parent position, $\propto\frac{4\pi}{\v k^2}\int\frac{d^2Y}{\v Y^2}[-\frac{11}{6}+\frac16\cos2\varphi]G_\infty$, which is present also for a colour-neutral projectile, comes from a parent much softer than the observed gluon. At the level of the amplitudes the pole of the soft parent cancels between the two-gluon component of the wave function and the splitting of the LO gluon after the collision, which enter Group~II Row~I (second part), Group~II Row~IV and Group~III Row~III with the weights $1,1,-2$, so that the tail is cancelled by the same region of the other two rows, for a target of finite size or for a colour-neutral projectile (Section~\ref{sec:dg_tail}). We have also checked it term by term on the finite terms as written in Appendix~\ref{app:finite}: the second part of Row~II.I has the same tail as Row~II.IV, Row~III.III with its complex conjugate has minus twice that tail, and the sum vanishes, also in the BFKL and JIMWLK sectors (Section~\ref{sec:dg_tail}). The three rows have to be integrated together over the position of the pair, condition (C10) of Section~\ref{sec:fin_integrable}; the region in which only part of the pair is far does not contribute at this order.

% ======================================================================
\section{Conclusions}\label{sec:conclusions}
% ======================================================================
% Sources: all sections; sec_dglap_rc.tex (notes/fix_poles_dglap.tex, ref3_beta_scheme.tex, ref3_collinear.tex),
% sec_finite.tex (notes/fix_soft.tex, fix_singlet.tex, ref3_finiteness.tex, ref3_prescriptions.tex, ref3_cubic.tex),
% sec_checks.tex (ref3_dilute.tex, ref3_cubic.tex, ref3_liskokov.tex).

We have computed the single inclusive gluon spectrum at mid rapidity in the collision of a dilute projectile with a dense target to next-to-leading order in the Color Glass Condensate, in pure Yang--Mills theory. The calculation uses the evolution operator $\Omega$ of the projectile, written as a normal-ordered series in the soft gluon operators with coefficients that are functionals of the non-commuting valence charge density, and reduces the spectrum to products of these coefficients and of Wilson lines. We have organised the $O(g^4)$ cross section into three groups of terms and fifteen rows. The colour charges of the projectile have been treated as non-commuting operators throughout; their ordering, the contact terms left by their commutators and the symmetrised three-point function of the projectile are part of the result. The complete result, with pointers to all its terms, is summarised in Section~\ref{sec:result_summary}. Rapidity evolution of projectile and target, collinear evolution of the final state, the running of the coupling and the finite terms all emerge from this one calculation, which in this sense completes the NLO description of the single inclusive gluon spectrum within the eikonal, dilute-projectile approximation, up to the open points listed in Section~\ref{sec:openpoints}.

The rapidity-enhanced terms have the structure required by high-energy factorisation. The logarithm of the longitudinal cut-offs splits at the momentum of the produced gluon. The coefficient of $\log(\vee/k^+)$ is the BFKL evolution of the projectile charge correlator inserted into the LO spectrum, and the coefficient of $\log(k^+/\Lambda)$ is the JIMWLK Hamiltonian acting on the target operator of the LO spectrum, for an arbitrary ordered projectile correlator. Terms with four charges cancel, as required by the coherent-state structure of the LO wave function, and the terms with three charges cancel in the symmetrised sense; with non-commuting charges, their reordering produces part of the two-charge structures of both evolutions. That no rapidity logarithm can multiply the C-odd three-point function of the projectile follows from charge conjugation together with the factorised form of the logarithms; the explicit cancellation confirms it. The operator ordering of the charges, in particular of the three-charge term of the $O(g^3)$ wave function, is essential for these cancellations.

All poles $1/\epsilon$ of the calculation are accounted for, in conventional dimensional regularisation with the $\MS$ coupling and the $\MS$ fragmentation function. The ultraviolet poles of the vertex correction of the observed gluon and of the term in which the two gluons of the two-gluon component recombine cancel each other, together with their $\frac{11}{3}$ terms; this is a result, while the finite terms that accompany it depend on the grouping of the two rows: with the colour of the gluon pair evaluated at its centre of mass, the self-energy part of the vertex correction, including its constant $\frac{67}{9}-\frac{\pi^2}{3}$, is cancelled identically. The coupling is renormalised by the ultraviolet half of the normalisation of the dressed produced gluon, and the infrared half of that normalisation completes the real final-state splitting to the full $\MS$ splitting function $\mathcal P_{gg}$, whose collinear pole is absorbed into the fragmentation function. A left-over pole multiplied by $\log(1-k^+/\vee)=O(k^+/\vee)$ is a power correction outside the eikonal approximation, which we define as the leading term of the expansion in $k^+/\vee$ and $\Lambda/k^+$ at fixed $\epsilon$. The collinear term is the final-state DGLAP fragmentation of the produced gluon. Its end points $\xi\to1$ and $\xi\to0$ are subtracted into the JIMWLK and the BFKL term, respectively, which is the ``double plus'' prescription of~\cite{Kovner:2023vsy}. The fragmentation function is the standard $\MS$ number density; with this convention the one-hadron-inclusive gluon coefficient functions of $e^+e^-\to gX$ and of $H\to gX$ are reproduced exactly. The argument of the collinear logarithm is then the transverse separation of the observed gluon, conjugate to $\v k$, and the form of~\cite{Chirilli:2012jd} corresponds to fragmentation functions that differ from the standard ones by a finite, scale-independent convolution, which we give. The running-coupling term is the $\frac{\beta_0}{2}\delta(1-\xi)$ part of the collinear term, added and subtracted; with separate renormalisation and factorisation scales it leaves $\frac{\alpha_s\beta_0}{4\pi}\log(\mu_R^2/\mu_F^2)\times$LO. With it, the coupling of the LO spectrum runs with the transverse separation conjugate to the produced momentum or, in the local form, with the produced momentum itself. Which of the numbers $\frac{11}{6}$ of the calculation are $\beta_0/(2N_c)$ is identified by a BLM-type quark loop; the coefficient of the large-distance tail of the collinear row is not.

The finite terms are given Group by Group and Row by Row, with the exact cancellations carried out and every product of three charges written as its symmetrised product plus two-charge contact terms, for an arbitrary projectile and, after the average over a globally colour-invariant projectile, in colour-singlet form with traces of adjoint Wilson lines. All colour structures are written through differences of Wilson lines, as at leading order, and all longitudinal integrals with plus prescriptions in a momentum fraction. They contain no renormalisation scale and no rapidity cut-off. The soft remainders of the plus prescriptions in the rows in which the observed gluon recoils against a soft gluon cancel in pairs of a virtual and a real term with the same coefficient, colour structure and kernel, without any condition on the target at large distance; equivalently, they can be absorbed into a JIMWLK term in which the soft gluon is restricted to rapidities below that of the observed gluon. All kernels, including those of the contact terms, are integrable at all coincidences; the contact terms do not cancel among the rows, and they do not need to. The measured-momentum Fourier integrals are unique for any regulator of the measured gluon, and with them done first all remaining integrals converge absolutely in the grouping we give, except for principal values and one conditionally convergent collinear point, which is defined by the rotation-invariant regularisation in which its subtraction was obtained. The large-distance tail of the collinear row cancels against two other rows, including the region in which only part of the gluon pair is far. The only energy denominators that vanish inside the integration region occur in the final-state dressing of the observed gluon, in one row. They are defined as principal values. The instantaneous terms of the two-gluon interference, which converge for $\vee\to\infty$ only as distributions, come from $p^+\lesssim k^+\max\{1,1/(|\v k|\ell)\}$, $\ell$ the shortest transverse scale of the integrand, not from the cut-off, and they do not depend on the shape of the separation of soft and valence modes. In the averaged form the terms cubic in the charges contribute in two ways, both given row by row. Their symmetrised part is proportional to the symmetrised three-point function $\sigma$ of the projectile, for which no model is assumed; it is odd in $\v k$ configuration by configuration, and its colour factors are C-odd target structures with at most three Wilson lines, the four-line invariants cancelling between two rows. The commutators of the ordered charges leave contact terms of the same order as the quadratic terms, which are even in $\v k$ and present for any projectile. The split of each cubic row into these two parts is unique: under charge conjugation it is the split into the C-odd and C-even parts of the projectile, as we verified on a single quark and on a dipole computed exactly in charge space.

In the dilute--dilute limit the real emission reproduces exactly the two-gluon production cross section of quasi-multi-Regge kinematics~\cite{Fadin:1998py,Bartels:2006hg}, including its normalisation and its polarisation structure, and its finite part coincides with the corresponding remainder of~\cite{Bartels:2006hg}. For a point-like projectile the virtual correction has the gluon Regge trajectories as coefficients of the two rapidity logarithms and contains no $\frac{67}{9}$; in this limit the result does not depend on the prescription for the vanishing energy denominators.

\subsection{Open points}\label{sec:openpoints}

The following points are not settled by this paper.
\begin{enumerate}
\item[(i)] \emph{Dilute limit.} The comparison of the virtual correction with the one-loop central-emission vertex~\cite{Fadin:1998py,DelDuca:1998cx,Bartels:2006hg} has not been completed. It requires virtual and real parts with matched regulators; in particular, it is open whether the target-side term $\frac{\alpha_sN_c}{\pi}[\frac2\epsilon\log\frac{k^+}{\Lambda}+\log^2\frac{k^+}{\Lambda}]$ of the virtual correction for a point-like quark, Eq.~\eqref{eq:dilute_groupI}, cancels against the real emission from the region $|\v p|\lesssim(\Lambda/k^+)|\v k|$, which we have not computed, and whether the finite function $C(\v k,\v l)$ agrees with the one-loop vertex up to terms that can be absorbed into the unintegrated gluon distributions. Also not done are the conversion of the virtual trajectories between our plus-momentum cut-offs and the energy scales of~\cite{Bartels:2006hg}, the reduction of the virtual correction for a general projectile, the dilute limit of the coordinate-space finite terms as printed, and the comparison of the printed colour-singlet contact terms with the commutator contributions required in this limit.
\item[(ii)] \emph{Comparison with Li and Skokov.} Our symmetrised cubic terms and the first saturation correction of~\cite{Li:2021ntm} agree in their colour structures but not in their parity in $\v k$: their result has a $\v k$-even part, which our symmetrised terms do not have, and in the two-Wilson-line sector their even term has no counterpart here, while our odd two-line term has none there. We traced the difference to a relative phase $i$ in the two-charge structures that build the non-linear Weizs\"acker--Williams field and in the rescattering on the second charge after the collision; its origin, possibly the prescription for $1/k^+$ (retarded or symmetric $x^-$ profile), has not been determined, and the even part of~\cite{Li:2021ntm} has not been checked independently. A kernel-level comparison of the odd three-line terms, and the question whether our odd two-line terms add up to zero, are open.
\item[(iii)] \emph{Prescription.} The finite terms are given for the principal-value definition of the vanishing energy denominators of the final-state dressing of the observed gluon only (Section~\ref{sec:fin_conv}); the dependence of the result on this definition has not been studied here.
\item[(iv)] \emph{Numerical evaluation.} No numerical evaluation of the finite terms, and no plot, is given. Apart from the dilute limit, the finite terms are given in coordinate space only; a momentum-space form of all finite terms, and a complete numerical integration over all transverse positions with realistic Wilson lines, have not been carried out.
\item[(v)] \emph{Constant and $\beta_0$ terms.} The LO-proportional constant that is left in the finite terms of the recombination row in our grouping has not been evaluated for a general projectile and target, and the momentum-weighted collinear regions of three rows (II.I, III.III, I.V) have not been integrated. The logarithm of the emitter distance times the inverse saturation scale, with a $\beta_0$-type coefficient, that the resolving region of the recombination row contains by power counting, has not been evaluated with a target model. The $N_f$ terms are not part of this pure Yang--Mills calculation; the quark loop is used only for the identification of the $\beta_0$ terms.
\item[(vi)] \emph{Collinear limit.} The collinear convention is fixed by the general $\MS$ counterterm and by two angle-integrated coefficient functions; a comparison with the NLO coefficient of hadroproduction, $pp\to gX$, and a collinear limit of the hybrid formula have not been carried out. At the conditionally convergent collinear point of the recombination row, a continuation of the transverse indices to $d_\perp=2-\epsilon$ could give a finite local term; whether it is present has not been decided.
\item[(vii)] \emph{Status of the integrability statements.} The power counting of the contact kernels is numerical (exponent fits at generic configurations, two smooth target fields) together with analytic arguments, not a proof for every degenerate configuration; the kernel of the $f^{abc}\{\rho^b,\rho^c\}$ term of the vertex correction was estimated analytically only. The absence of a contribution from the region in which only part of the gluon pair is far was verified numerically for one charge and a target uniform at large distance; for a colour-neutral projectile and a non-uniform target it rests on an argument. The interchange of the $p^+$ integral with the joint large-distance limit was checked numerically for one of the two rows concerned. The statements on the instantaneous terms were checked with model functions of the correct short-distance behaviour, and the universality of the shift of the BFKL logarithm under a change of the soft--valence separation is argued from the form of the integrands, not checked row by row.
\end{enumerate}


\acknowledgments
% To be written by the authors.
ACKNOWLEDGMENTS.

\appendix
% ======================================================================
\section{States, Hamiltonian and light-cone matrix elements}\label{app:hamiltonian}
% ======================================================================

In this appendix we collect the definitions of the states, the parts of the eikonal light-cone Hamiltonian that are needed, and the matrix elements used in the LCPT calculations of Appendix~\ref{app:wavefunctions}. They are those of~\cite{Radhakrishnan:2026soft,Lublinsky:2016meo}, restricted to pure Yang--Mills theory.

\subsection{States and normalisation}\label{app:states}

The soft vacuum $\ket0$ is annihilated by all soft annihilation operators, $a^a_i(k^+,\v k)\ket0=0$ for $\Lambda<k^+<\vee$, and its energy is set to zero, $E_0=0$. The one- and two-gluon states are
\begin{align}
\ket{g^a_i(k)}&=\frac{a^{\dagger a}_i(k)}{(2\pi)^{3/2}}\ket0\,,&
\ket{g^a_i(k^+,\v z)}&=\frac{a^{\dagger a}_i(k^+,\v z)}{(2\pi)^{1/2}}\ket0\,,\notag\\
\ket{g^a_i(k)\,g^b_j(p)}&=\frac{a^{\dagger a}_i(k)\,a^{\dagger b}_j(p)}{(2\pi)^3}\ket0\,,&
\ket{g^a_i(k^+,\v z)\,g^b_j(p^+,\v z')}&=\frac{a^{\dagger a}_i(k^+,\v z)\,a^{\dagger b}_j(p^+,\v z')}{2\pi}\ket0\,,
\end{align}
with energies $E_g(k)=\v k^2/2k^+$ and $E_{gg}(k,p)=\v k^2/2k^++\v p^2/2p^+$, and normalisation
\begin{align}
\braket{g^b_j(p)|g^a_i(k)}&=\delta^{ab}\delta_{ij}\,\delta(k^+-p^+)\,\delta^{(2)}(\v k-\v p)\,,\notag\\
\braket{g^b_j(p^+,\v z')|g^a_i(k^+,\v z)}&=\delta^{ab}\delta_{ij}\,\delta(k^+-p^+)\,\delta^{(2)}(\v z-\v z')\,.
\end{align}
The coordinate-space state differs from the Fourier transform of the momentum-space state by a factor $2\pi$, $\ket{g(k^+,\v z)}=2\pi\int\frac{d^2k}{(2\pi)^2}e^{-i\v k\cdot\v z}\ket{g(k)}$. Consequently $a^a_i(k^+,\v k)\ket{g^b_j(q)}=(2\pi)^{3/2}\delta^{ab}\delta_{ij}\delta^{(3)}(k-q)\ket0$.

\subsection{Soft and valence modes}

The gauge field in the light-cone gauge is decomposed into soft and valence components,
\begin{align}
\underline{A}^a_i(x)&=\int_\Lambda^\vee\frac{dk^+}{2\pi}\int\frac{d^2k}{(2\pi)^2}\frac{1}{\sqrt{2k^+}}\Big(a^a_i(k^+,\v k)e^{-ik\cdot x}+a^{\dagger a}_i(k^+,\v k)e^{ik\cdot x}\Big)\,,\notag\\
\overline{A}^a_i(x)&=\int_\vee^\infty\frac{dk^+}{2\pi}\int\frac{d^2k}{(2\pi)^2}\frac{1}{\sqrt{2k^+}}\Big(a^a_i(k^+,\v k)e^{-ik\cdot x}+a^{\dagger a}_i(k^+,\v k)e^{ik\cdot x}\Big)\,.
\end{align}
The valence charge density is the colour charge density of the valence modes~\cite{Radhakrishnan:2026soft},
\begin{align}
\rho^a(-\v p)=-if^{abc}\int_\vee^\infty\frac{dk^+}{2\pi}\int\frac{d^2k}{(2\pi)^2}\,a^{\dagger b}_j(k^+,\v k)\,a^c_j(k^+,\v k+\v p)\,,
\end{align}
and satisfies the algebra~\eqref{eq:rho}. Inserting the decomposition into the light-cone Hamiltonian and dropping all terms suppressed by powers of the ratio of soft to valence longitudinal momenta (in particular all inverse longitudinal derivatives acting on valence fields), one obtains the interaction terms~\cite{Lublinsky:2016meo,Radhakrishnan:2026soft}
\begin{align}
H_g&=-g\int dx^-d^2x\,\big(\partial_i\underline A^a_i\big)\,f^{abc}\,\frac{1}{\partial^+}\big(\overline A^b_j\partial^+\overline A^c_j\big)\,,\notag\\
H_{ggg}&=-gf^{abc}\int dx^-d^2x\,\Big[\big(\partial_i\underline A^a_j\big)\underline A^b_i\underline A^c_j+\big(\partial_i\underline A^a_i\big)\frac{1}{\partial^+}\big(\underline A^b_j\partial^+\underline A^c_j\big)\Big]\,,\notag\\
H_{gg\text{-inst}}&=g^2f^{abc}f^{ade}\int dx^-d^2x\,\frac{1}{\partial^+}\big(\underline A^b_i\partial^+\underline A^c_i\big)\,\frac{1}{\partial^+}\big(\overline A^d_j\partial^+\overline A^e_j\big)\,,\notag\\
H_{gggg}&=\frac{g^2}{4}f^{abc}f^{ade}\int dx^-d^2x\,\Big[\underline A^b_i\underline A^c_j\underline A^d_i\underline A^e_j-2\,\frac{1}{\partial^+}\big(\underline A^b_i\partial^+\underline A^c_i\big)\frac{1}{\partial^+}\big(\underline A^d_j\partial^+\underline A^e_j\big)\Big]\,.
\end{align}
In terms of the soft creation and annihilation operators $H_g$ takes the form~\eqref{eq:Hg}, and
\begin{align}
H_{gg\text{-inst}}=\ &\int_\Lambda^\vee\frac{dp^+}{2\pi}\frac{dq^+}{2\pi}\int\frac{d^2p}{(2\pi)^2}\frac{d^2q}{(2\pi)^2}\frac{ig^2f^{abc}}{\sqrt{p^+q^+}}\notag\\&\times\Bigg(\frac{q^+\rho^a(-\v p-\v q)}{2(p^++q^+)^2}\,a^{\dagger b}_i(p^+,\v p)\,a^{\dagger c}_i(q^+,\v q)-{\rm h.c.}
\notag\\&
-\frac{(p^++q^+)\,\rho^a(-\v p+\v q)}{2(p^+-q^+)^2}\,a^{\dagger b}_i(p^+,\v p)\,a^c_i(q^+,\v q)\Bigg)\,.
\end{align}
The four-gluon interaction $H_{gggg}$ does not enter the coefficients used in this paper.
% CHOICE FOR THE AUTHORS: confirm that H_gggg (normal ordered) gives no O(g^3) one-gluon contribution (no H_gggg diagram appears in the calculation files).

\subsection{Matrix elements}\label{app:matrixelements}

The matrix elements of the interaction terms between the free states, which enter the LCPT expansion~\eqref{eq:lcpt}, are~\cite{Radhakrishnan:2026soft}
\begin{align}
\braket{g^a_i(k)|H_g|0}&=\frac{g\,k^i\rho^a(-\v k)}{4\pi^{3/2}|k^+|^{3/2}}\,,\label{eq:m1}\\
\braket{g^c_l(q)\,g^b_j(p)|H_g|g^a_i(k)}&=\delta^{ac}\delta_{il}\delta^{(3)}(k-q)\frac{g\,p^j\rho^b(-\v p)}{4\pi^{3/2}|p^+|^{3/2}}\notag\\&\quad+\delta^{ab}\delta_{ij}\delta^{(3)}(k-p)\frac{g\,q^l\rho^c(-\v q)}{4\pi^{3/2}|q^+|^{3/2}}\,,\label{eq:m2}\\
\braket{g^c_l(q)\,g^b_j(p)|H_{ggg}|g^a_i(k)}&=\frac{igf^{abc}\,\delta^{(3)}(k-p-q)}{8\pi^{3/2}\sqrt{k^+p^+q^+}}\Bigg[\Big(p^i-q^i+\frac{q^+-p^+}{k^+}k^i\Big)\delta_{jl}\notag\\&\hspace{1cm}+\Big(k^j+q^j-\frac{k^++q^+}{p^+}p^j\Big)\delta_{il}
\notag\\&\hspace{1cm}
+\Big(\frac{k^++p^+}{q^+}q^l-p^l-k^l\Big)\delta_{ij}\Bigg]\,,\label{eq:m3}\\
\braket{g^c_j(q)\,g^b_i(p)|H_{gg\text{-inst}}|0}&
=\frac{ig^2f^{abc}\,(q^+-p^+)\,\delta_{ij}\,\rho^a(-\v p-\v q)}{2(2\pi)^3\sqrt{p^+q^+}\,(p^++q^+)^2}\,,\label{eq:m4}\\
\braket{g^c_j(q)|H_{gg\text{-inst}}|g^b_i(p)}&
=\frac{ig^2f^{abc}\,(p^++q^+)\,\delta_{ij}\,\rho^a(\v p-\v q)}{2(2\pi)^3\sqrt{p^+q^+}\,(p^+-q^+)^2}\,.\label{eq:m5}
\end{align}
Eq.~\eqref{eq:m1} is the emission of a soft gluon by the valence charges, Eq.~\eqref{eq:m2} the emission of a second gluon in the presence of a first one, Eq.~\eqref{eq:m3} the triple-gluon vertex among soft gluons, Eq.~\eqref{eq:m4} the instantaneous emission of two soft gluons by the valence charges and Eq.~\eqref{eq:m5} the instantaneous interaction of a soft gluon with the valence charges, in which the longitudinal momentum $p^+-q^+$ is exchanged with the valence modes. All matrix elements are operators on the valence Hilbert space.

\subsection{Rules of light-cone perturbation theory}\label{app:lcptrules}

For the dressing of the soft vacuum the energy denominators are $E_0-E_n=-E_n$, Eq.~\eqref{eq:lcpt}. For the dressing of a one-gluon state $\ket{g^a_i(k)}$ the energy of the initial state is $E_g(k)$ and the denominators are $E_g(k)-E_n$; for example, the instantaneous transition of the gluon into a gluon of momentum $p$ is
\begin{align}
\ket{g^a_i(k)}\to\int dp^+\!\int d^2p\;\ket{g^b_j(p)}\,\frac{\braket{g^b_j(p)|H_{gg\text{-inst}}|g^a_i(k)}}{E_g(k)-E_g(p)}\,.
\end{align}
These denominators vanish inside the integration region, on $p^+\v k^2=k^+\v p^2$ in momentum space or on $p^+(\v x-\v u)^2=k^+(\v x-\v z)^2$ in coordinate space; they are understood as principal values (Appendix~\ref{app:distributions}, Section~\ref{sec:fin_conv}). The longitudinal momentum of every soft gluon, real or intermediate, is restricted to $\Lambda<p^+<\vee$, and the longitudinal momentum exchanged with the valence charges in Eqs.~\eqref{eq:m2} and~\eqref{eq:m5} satisfies $|p^+-q^+|>\Lambda$. The kinematic regions allowed for the transition $g(k^+)\to g(p^+)$ through the valence charges are shown in Figure~\ref{fig:kinem}.

\begin{figure}[t]
\centering
\tikzfig{kinem}
\caption{Allowed longitudinal momenta for the transition of a soft gluon of momentum $k^+$ into a soft gluon of momentum $p^+$ by the interaction with the valence charges, for $p^+>k^+$. The upper (purple) region corresponds to $\Lambda<p^+-k^+<\vee$, the green region to $\Lambda<k^+<\vee$ and $\Lambda<p^+<\vee$; the allowed region is their intersection. The case $p^+<k^+$, in which the gluon transfers longitudinal momentum to the valence charges, is the mirror image.}
\label{fig:kinem}
\end{figure}

% ======================================================================
\section{Light-cone wave functions and the coefficients of \texorpdfstring{$\Omega$}{Omega}}\label{app:wavefunctions}
% ======================================================================

The coefficients of $\Omega$ listed in Section~\ref{sec:coefficients} are obtained by computing the action of $\Omega$ on the soft vacuum and on a one-gluon state in two ways: from the expansion~\eqref{eq:OmegaExp}, and from LCPT, Eq.~\eqref{eq:lcpt}, with the matrix elements of Appendix~\ref{app:matrixelements}. The two are equated sector by sector in Fock space. In this appendix we give the Fock components, the diagrams from which they arise, and the steps that lead from the components to the coefficients. In the diagrams the horizontal line denotes the valence charges, curly lines soft gluons, and vertical dashed lines the intermediate states whose energies enter the denominators; a vertical bar on a line denotes an instantaneous interaction. Momenta are written in momentum space; the coordinate forms follow with the Fourier transforms of Appendix~\ref{app:fourier}.

\subsection{Matching conditions}

From Eq.~\eqref{eq:OmegaExp} and the normalisation of the states,
\begin{align}\label{eq:Omegavac}
\Omega\ket0=\ &\mf N\ket0+(2\pi)^{3/2}\int\frac{d^3p}{(2\pi)^3}\Big[\mf A^{(3)b}_{2j}(p)-\mf A^{(1)\dagger b}_{1j}(p)\Big]\ket{g^b_j(p)}\notag\\&+(2\pi)^3\int\frac{d^3p}{(2\pi)^3}\frac{d^3q}{(2\pi)^3}\,\mf B^{cb}_{2kj}(p,q)\ket{g^b_j(p)\,g^c_k(q)}+\dots
\end{align}
and
\begin{align}\label{eq:Omegag}
\Omega\ket{g^a_i(k)}=\ &\frac{1}{(2\pi)^{3/2}}\Big[\mf A^{(1)a}_{1i}(k)+\mf A^{(3)a}_{1i}(k)\Big]\ket0+\mf N\ket{g^a_i(k)}\notag\\&+\int\frac{d^3p}{(2\pi)^3}\,\mf B^{ba}_{3ji}(k,p)\ket{g^b_j(p)}
\notag\\&
+(2\pi)^{3/2}\int\frac{d^3p}{(2\pi)^3}\Big[\mf A^{(3)b}_{2j}(p)-\mf A^{(1)\dagger b}_{1j}(p)\Big]\ket{g^a_i(k)\,g^b_j(p)}
\notag\\&+(2\pi)^{3/2}\int\frac{d^3p}{(2\pi)^3}\,\mf C^{bca}_{1jki}(k,p)\ket{g^b_j(p)\,g^c_k(k-p)}+\dots,
\end{align}
where the dots denote states with more gluons. The vacuum sector fixes $\mf N$, $\mf A^{(1)}$, $\mf A^{(3)}_2$ and $\mf B_2$; the one-gluon sector fixes $\mf B_3$ and $\mf C_1$, and reproduces $\mf A^{(1)}_1$ and the product $\mf A\,$(one-gluon state) as a consistency check.

\subsection{The soft vacuum at order \texorpdfstring{$g$}{g} and the normalisation}

At order $g$ the valence charges emit one soft gluon (Figure~\ref{fig:vac_g}):
\begin{align}\label{eq:PsiLO}
\ket{\Psi^{\rm LO}}&=-\int_\Lambda^\vee dk^+\!\int d^2k\,\ket{g^a_i(k)}\frac{\braket{g^a_i(k)|H_g|0}}{E_g(k)}
\notag\\&=-\int_\Lambda^\vee dk^+\!\int d^2k\,\frac{g\,k^i\rho^a(-\v k)}{2\pi^{3/2}\sqrt{k^+}\,\v k^2}\ket{g^a_i(k)}
\notag\\&=\frac{ig}{2\pi^{3/2}}\int_\Lambda^\vee\frac{dk^+}{\sqrt{k^+}}\int_{x,z}\Kv^i(\v x-\v z)\rho^a(\v x)\ket{g^a_i(k^+,\v z)}\,.
\end{align}
Comparison with Eq.~\eqref{eq:Omegavac} gives $-\mf A^{(1)a}_{2i}(k^+,\v k)=\mf A^{(1)\dagger a}_{1i}(k^+,\v k)=\sqrt2\,g\,k^i\rho^a(-\v k)/(\sqrt{k^+}\v k^2)$, i.e.\ Eq.~\eqref{eq:A1}. The norm of the state to order $g^2$ fixes
\begin{align}
\mf N=1-\frac12\braket{\Psi^{\rm LO}|\Psi^{\rm LO}}=1-\frac{g^2}{8\pi^3}\int_\Lambda^\vee\frac{dp^+}{p^+}\int d^2p\,\frac{\rho^a(-\v p)\rho^a(\v p)}{\v p^2}\,,
\end{align}
which is Eq.~\eqref{eq:N2}.

\begin{figure}[t]
\centering
\tikzfig{vacuumtog}
\caption{Emission of a soft gluon by the valence charges at order $g$.}
\label{fig:vac_g}
\end{figure}

\subsection{Two soft gluons at order \texorpdfstring{$g^2$}{g2}}

Three diagrams contribute (Figure~\ref{fig:vac_gg}): the emission of two gluons by the valence charges, with the matrix elements~\eqref{eq:m1} and~\eqref{eq:m2},
\begin{align}
\ket{\Psi^{1}_{gg}}=\ &\frac12\int dq^+dk^+dp^+\!\int d^2q\,d^2k\,d^2p\,\ket{g^b_j(k-p)\,g^c_l(p)}\notag\\&\times\frac{\braket{g^b_j(k-p)\,g^c_l(p)|H_g|g^a_i(q)}\braket{g^a_i(q)|H_g|0}}{E_{gg}(p,k-p)\,E_g(q)}\,,
\end{align}
the emission of one gluon which splits into two through the triple-gluon vertex~\eqref{eq:m3},
\begin{align}
\ket{\Psi^{2}_{gg}}=\frac12\int dk^+dp^+dq^+\!\int d^2k\,d^2p\,d^2q\,\ket{g^b_j(q)\,g^c_l(p)}\frac{\braket{g^b_j(q)\,g^c_l(p)|H_{ggg}|g^a_i(k)}\braket{g^a_i(k)|H_g|0}}{E_{gg}(p,q)\,E_g(k)}\,,
\end{align}
and the instantaneous emission of the pair~\eqref{eq:m4},
\begin{align}
\ket{\Psi^{3}_{gg}}=-\frac12\int dk^+dp^+\!\int d^2k\,d^2p\,\ket{g^b_i(k-p)\,g^c_j(p)}\frac{\braket{g^b_i(k-p)\,g^c_j(p)|H_{gg\text{-inst}}|0}}{E_{gg}(k-p,p)}\,.
\end{align}
In the first diagram the product of the two charges is decomposed as
\begin{align}\label{eq:decomp}
\rho^b(-\v k+\v p)\,\rho^a(-\v p)=\frac{i}{2}f^{bad}\rho^d(-\v k)+\frac12\big\{\rho^b(-\v k+\v p),\rho^a(-\v p)\big\}\,.
\end{align}
The commutator term is linear in $\rho$ and is combined with the other two diagrams. The anticommutator term describes two independent emissions. Comparing with Eq.~\eqref{eq:Omegavac} gives the coefficient
\begin{align}\label{eq:B2mom}
\mf B^{ba}_{2ji}(k^+,\v k;p^+,\v p)=\ &\frac{ig^2f^{abc}\,p^jk^i\rho^c(-\v p-\v k)\sqrt{p^+k^+}}{2\,(p^+\v k^2+k^+\v p^2)}\Big[\frac{1}{k^+\v p^2}-\frac{1}{\v k^2p^+}\Big]\,\Theta(\vee-k^+-p^+)
\notag\\&
+\frac{g^2p^jk^i\,\{\rho^b(-\v p),\rho^a(-\v k)\}}{2\,\v p^2\v k^2\sqrt{p^+k^+}}\,\Theta(\vee-k^+)\,\Theta(\vee-p^+)
\notag\\&
-\frac{ig^2f^{abc}\rho^c(-\v p-\v k)\,(p+k)^l\,(p^mk^+-p^+k^m)}{(\v p+\v k)^2\sqrt{p^+k^+}\,(p^++k^+)\big(\frac{\v p^2}{p^+}+\frac{\v k^2}{k^+}\big)}\notag\\&\quad\times\Big[\frac{\delta_{lm}\delta_{ji}}{p^++k^+}-\frac{\delta_{im}\delta_{jl}}{k^+}-\frac{\delta_{jm}\delta_{il}}{p^+}\Big]\Theta(\vee-k^+-p^+)
\notag\\&
-\frac{ig^2f^{abc}\rho^c(-\v p-\v k)\,(p^+-k^+)\,\delta_{ji}}{2\sqrt{p^+k^+}\,(p^++k^+)^2\big(\frac{\v k^2}{k^+}+\frac{\v p^2}{p^+}\big)}\,\Theta(\vee-k^+-p^+)\,,
\end{align} 
where the first, third and fourth terms come from the commutator part of the first diagram, from the triple-gluon vertex and from the instantaneous emission, respectively, and the second term from the anticommutator part. The supports are discussed in Section~\ref{sec:coefficients}. With the Fourier transforms~\eqref{eq:FT_i2} and~\eqref{eq:FT_i7} this becomes Eq.~\eqref{eq:B2}.

\begin{figure}[t]
\centering
\tikzfig{vacuumto2g2rho}\qquad\tikzfig{vacuumto2g}\qquad\tikzfig{vacuumto2ginst}
\caption{Two soft gluons at order $g^2$: two emissions by the valence charges, emission followed by splitting, and instantaneous emission.}
\label{fig:vac_gg}
\end{figure}

\subsection{One soft gluon at order \texorpdfstring{$g^3$}{g3}}

\paragraph{Diagrams.} The one-gluon component at order $g^3$ receives contributions linear, quadratic and cubic in the charges.
\begin{itemize}
\item The self-energy of the emitted gluon through two triple-gluon vertices (Figure~\ref{fig:vac_g3_onerho}),
\begin{align}
\ket{\Psi^{1}_{g\rho}}=\ &-\frac12\int dk^+dp^+dq^+dr^+\!\int d^2k\,d^2p\,d^2q\,d^2r\,\ket{g^d_l(r)}\notag\\&\times\frac{\braket{g^d_l(r)|H_{ggg}|g^c_k(q)g^b_j(p)}\braket{g^c_k(q)g^b_j(p)|H_{ggg}|g^a_i(k)}\braket{g^a_i(k)|H_g|0}}{E_g(r)\,E_{gg}(p,q)\,E_g(k)}\,.
\end{align}
It is linear in $\rho$ and is the only diagram with an ultraviolet divergence of the transverse loop integral.
\item The diagrams with two valence vertices (Figure~\ref{fig:vac_g3_tworho}): the emitted gluon interacts instantaneously with the valence charges, gaining or losing longitudinal momentum; two gluons emitted by the valence charges recombine through the triple-gluon vertex; two gluons are emitted, one instantaneously, and one of them is absorbed by the valence charges; a gluon splits into two, one of which is absorbed by the valence charges. With the decomposition~\eqref{eq:decomp} each of them gives a term linear in $\rho$ (commutator) and a term $\propto f^{abc}\{\rho^b,\rho^c\}$.
\item The diagrams with three valence vertices (Figure~\ref{fig:vac_g3_threerho}): a gluon is emitted and reabsorbed by the valence charges before or after the emission of the observed gluon, in the three time orderings, and the corresponding normalisation term of Eq.~\eqref{eq:lcpt}.
\end{itemize}

\begin{figure}[t]
\centering
\tikzfig{loop}
\caption{Self-energy of the emitted gluon at order $g^3$.}
\label{fig:vac_g3_onerho}
\end{figure}

\begin{figure}[t]
\centering
\tikzfig{Misha_2rho_single_g}\qquad\tikzfig{Misha_Out_2rho_single_g}\qquad\tikzfig{Misha_2rho_3_single_g}\\[4mm]
\tikzfig{Misha_2rho_insta_single_g}\qquad\tikzfig{Misha_2rho_loop_single_g}
\caption{Diagrams with two valence vertices contributing to the one-gluon component at order $g^3$.}
\label{fig:vac_g3_tworho}
\end{figure}

\begin{figure}[t]
\centering
\tikzfig{Misha_2rho_loopVal1_single_g}\qquad\tikzfig{Misha_2rho_loopVal2_single_g}\qquad\tikzfig{Misha_2rho_loopVal3_single_g}
\caption{Diagrams with three valence vertices contributing to the one-gluon component at order $g^3$.}
\label{fig:vac_g3_threerho}
\end{figure}

\paragraph{Self-energy.} After the transverse loop integral in $d=2-\epsilon$ dimensions, with $p^+=\xi k^+$,
\begin{align}
\ket{\Psi^{1}_{g\rho}}=\ &-\int_\Lambda^\vee\frac{dk^+}{\sqrt{k^+}}\int d^2k\int_{\Lambda/k^+}^{1-\Lambda/k^+}d\xi\,\frac{g^3N_c\,\rho^a(-\v k)\,k^i}{16\pi^{7/2}\,\v k^2}\notag\\&\times\Big[-2+\xi(1-\xi)+\frac{1}{\xi(1-\xi)}\Big]\Big[-\frac2\epsilon+\log\frac{\v k^2\xi(1-\xi)}{\mu^2_{\MS}}\Big]\ket{g^a_i(k^+,\v k)}\,.
\end{align}
With $\delta=\Lambda/k^+\to0$,
\begin{align}\label{eq:xiintegrals}
\int_\delta^{1-\delta}d\xi\Big[-2+\xi(1-\xi)+\frac{1}{\xi(1-\xi)}\Big]&=-\frac{11}{6}-2\log\delta\,,\notag\\
\int_0^1d\xi\,\big[-2+\xi(1-\xi)\big]\log\big(\xi(1-\xi)\big)&=\frac{67}{18}\,,\notag\\
\int_\delta^{1-\delta}d\xi\,\frac{\log(\xi(1-\xi))}{\xi(1-\xi)}&=-\log^2\delta-\frac{\pi^2}{3}\,,
\end{align}
so that
\begin{align}
\ket{\Psi^{1}_{g\rho}}=\ &\int_\Lambda^\vee\frac{dk^+}{\sqrt{k^+}}\int d^2k\,\frac{g^3N_c\,\rho^a(-\v k)\,k^i}{32\pi^{7/2}\,\v k^2}\Big[\Big(\frac{11}{3}+4\log\frac{\Lambda}{k^+}\Big)\Big(-\frac2\epsilon+\log\frac{\v k^2}{\mu^2_{\MS}}\Big)\notag\\&\hspace{4cm}+2\log^2\frac{\Lambda}{k^+}-\frac{67}{9}+\frac{2\pi^2}{3}\Big]\ket{g^a_i(k^+,\v k)}\,.
\end{align}
This is the origin of the coefficient $\frac{11}{3}=\beta_0/N_c$ and of the constant $\frac{67}{9}$.

\paragraph{Sum of the one-$\rho$ terms.} Adding the commutator parts of the diagrams with two and three valence vertices, the one-gluon component linear in $\rho$ is
\begin{align}\label{eq:Psigrho}
\ket{\Psi_{g\rho}}=\ &\int_\Lambda^\vee\frac{dk^+}{\sqrt{k^+}}\int d^2k\,\frac{g^3N_c\,\rho^a(-\v k)\,k^i}{32\pi^{7/2}\,\v k^2}\Big[\Big(-\frac2\epsilon+\log\frac{\v k^2}{\mu^2_{\MS}}\Big)\Big(\frac{11}{3}+\log\frac{\Lambda}{k^+}-2\log\frac{\vee}{\Lambda}\Big)\notag\\&\hspace{3.5cm}+\frac{\pi^2}{3}-\frac{67}{9}+\log^2\frac{\Lambda}{k^+}-\log^2\frac{\vee}{k^+}\Big]\ket{g^a_i(k^+,\v k)}
\notag\\&
+\int_\Lambda^\vee\frac{dk^+}{\sqrt{k^+}}\int d^2k\int_{\Lambda/k^+}^{1-\Lambda/k^+}d\xi\,\frac{g^3N_c\,\rho^a(-\v k)\,k^i}{32\pi^{7/2}\sqrt\xi(1-\xi)\,\v k^2}\notag\\&\times\Big[(\xi-2)\log\xi+(\xi+1)\log(1-\xi)-(\xi-2)\Big(-\frac2\epsilon+\log\frac{\v k^2}{\mu^2_{\MS}}\Big)\Big]\ket{g^a_i(\xi k^+,\v k)}\,.
\end{align}
The second line, in which the observed gluon carries the momentum fraction $\xi$ of the parent, comes from the diagrams in which a gluon of momentum $k^+$ transfers part of its longitudinal momentum to the valence charges. To read off the coefficient of a gluon of momentum $k^+$ we relabel $\xi k^+\to k^+$ in this term; the parent momentum is then $k^+/\xi$, and the conditions $\Lambda<k^+/\xi<\vee$, $\xi<1-\Lambda\xi/k^+$ give $k^+/\vee<\xi<k^+/(k^++\Lambda)$. The $\xi$ integrals,
\begin{align}
&\int_{k^+/\vee}^{k^+/(k^++\Lambda)}d\xi\,\frac{2-\xi}{\xi(1-\xi)}=2\log\frac{\vee}{k^+}+\log\frac{k^+}{\Lambda}+\log\Big(1-\frac{k^+}{\vee}\Big)+O\Big(\frac{\Lambda}{k^+}\Big)\,,\notag\\
&\int_{k^+/\vee}^{k^+/(k^++\Lambda)}d\xi\,\frac{(\xi-2)\log\xi+(\xi+1)\log(1-\xi)}{\xi(1-\xi)}\notag\\&\quad=-\log^2\frac{\Lambda}{k^+}+\log^2\frac{\vee}{k^+}+\log^2\Big(1-\frac{k^+}{\vee}\Big)\notag\\&\quad+\log\Big(1-\frac{k^+}{\vee}\Big)\log\frac{\vee}{k^+}+O\Big(\frac{\Lambda}{k^+}\log\frac{\Lambda}{k^+}\Big)\,,
\end{align}
add $\log(\vee/k^+-1)+\log(\vee/\Lambda)$ to the coefficient of the pole and convert the constant of the first line into $\frac{\pi^2}{3}-\frac{67}{9}+\log^2(1-k^+/\vee)+\log(1-k^+/\vee)\log(\vee/k^+)$. The double logarithms of the cut-offs cancel, and the result is Eq.~\eqref{eq:A3onerho}.

\paragraph{Two-$\rho$ terms.} The anticommutator parts of the diagrams with two and three valence vertices give
\begin{align}\label{eq:Psigrhorho}
\ket{\Psi_{g\rho\rho}}=\ &\int_\Lambda^\vee\frac{dk^+}{\sqrt{k^+}}\int d^2k\,d^2p\,\frac{ig^3f^{abc}\{\rho^b(-\v p),\rho^c(-\v k+\v p)\}}{32\pi^{9/2}\,\v p^2\,\v k^2}\notag\\&\times\Bigg[\log\frac{\Lambda}{k^+}\,\frac{\v k^2p^i-(\v k\cdot\v p)k^i}{(\v k-\v p)^2}-2k^i\log\frac{\vee}{\Lambda}\notag\\&\quad-2k^i\log\frac{\Lambda\v k^2+k^+\v p^2}{\vee\v k^2+k^+\v p^2}
\notag\\&\quad
+\frac{(2\v p^2-2\v k\cdot\v p+\v k^2)k^i+\v k^2p^i}{(\v k-\v p)^2}\log\Big[1+\Big(\frac{\vee}{k^+}-1\Big)\frac{\v k^2}{(\v k-\v p)^2}\Big]
\notag\\&\quad
-\frac{\v p^2k^i+\v k^2p^i}{(\v k-\v p)^2}\log\frac{\vee-k^+}{\Lambda}-k^i\log\frac{\vee}{k^+}\Bigg]\ket{g^a_i(k)}\,,
\end{align}
after the $\xi$ integration. The terms proportional to $k^+/\Lambda$ of the individual instantaneous diagrams cancel in the sum. Setting $\Lambda\to0$ where possible and transforming to coordinate space gives Eq.~\eqref{eq:A3tworho}.

\paragraph{Three-$\rho$ terms.} The diagrams with three valence vertices give, in addition to commutator terms,
\begin{align}
\ket{\Psi^{6c}_{g\rho\rho\rho}}&=\int_\Lambda^\vee dk^+\!\int d^2k\,d^2p\,\frac{g^3\,k^i\,\rho^b(\v p)\rho^b(-\v p)\rho^a(-\v k)\sqrt{k^+}}{8\pi^{9/2}\,\v k^4}\notag\\&\quad\times\Big[\frac1\vee-\frac1\Lambda\Big]\ket{g^a_i(k^+,\v k)}\,,\notag\\
\ket{\Psi^{7}_{g\rho\rho\rho}}&=\int_\Lambda^\vee\frac{dk^+}{\sqrt{k^+}}\int d^2k\,d^2p\,\frac{g^3\,k^i\,\rho^a(-\v k)\rho^b(\v p)\rho^b(-\v p)}{16\pi^{9/2}\,\v k^2}\notag\\&\quad\times\Big[\frac{1}{\v p^2}\log\frac{\vee}{\Lambda}+\frac{2k^+}{\v k^2}\Big(\frac1\Lambda-\frac1\vee\Big)\Big]\ket{g^a_i(k^+,\v k)}\,.
\end{align}
The second line is the normalisation term of Eq.~\eqref{eq:lcpt}, with the single-emission matrix element to the left of the norm factor. The terms proportional to $1/\Lambda$ and $1/\vee$ cancel exactly between the two lines, although the charges stand in different orders: $\int d^2p\,\rho^b(\v p)\rho^b(-\v p)$ without the weight $1/\v p^2$ is local in coordinate space and commutes with $\rho^a$. What remains is
\begin{align}\label{eq:Psigrhorhorho}
\ket{\Psi_{g\rho\rho\rho}}=\log\frac{\vee}{\Lambda}\int_\Lambda^\vee\frac{dk^+}{\sqrt{k^+}}\int d^2k\,d^2p\,\frac{g^3\,\rho^a(-\v k)\,\rho^b(\v p)\,\rho^b(-\v p)\,k^i}{16\pi^{9/2}\,\v k^2\,\v p^2}\ket{g^a_i(k^+,\v k)}\,,
\end{align}
which gives $[\mf A^{(3)}_2]_{\rho\rho\rho}=\mf A^{(1)}_2\mf N_2$, Eq.~\eqref{eq:A3threerho}.

\subsection{One soft gluon in the initial state}

The coefficients $\mf B_3$ and $\mf C_1$ are fixed by the dressing of a one-gluon state $\ket{g^a_i(k)}$ at order $g^2$ and $g$, respectively. The diagrams for the transition into a one-gluon state are shown in Figure~\ref{fig:g_g}: the instantaneous interaction with the valence charges, for $p^+>k^+$ and $p^+<k^+$; the emission of a second gluon by the valence charges and the absorption of the first by the triple-gluon vertex, or vice versa; and the diagrams with two valence vertices, in which the first gluon is absorbed and a second one emitted, or the second one emitted before the first is absorbed. The splitting of the gluon into two soft gluons is shown in Figure~\ref{fig:g_gg}. The energy denominators are $E_g(k)-E_n$ (Appendix~\ref{app:lcptrules}). Comparison with Eq.~\eqref{eq:Omegag} gives
\begin{align}\label{eq:B3mom}
\mf B^{ba}_{3ji}(k^+,\v k;p^+,\v p)=\ &\frac{ig^2f^{abc}\sqrt{p^+k^+}\,(k^++p^+)}{(p^+-k^+)^2}\,\frac{\rho^c(-\v p+\v k)\,\delta_{ij}}{p^+\v k^2-k^+\v p^2}\notag\\&\quad\times\big[\Theta(p^+-k^+-\Lambda)+\Theta(k^+-\Lambda-p^+)\big]
\notag\\&
-\frac{g^2k^ip^j}{\v k^2\v p^2\sqrt{k^+p^+}}\Big[\{\rho^a(\v k),\rho^b(-\v p)\}-if^{abc}\rho^c(\v k-\v p)\,\frac{\v p^2k^++\v k^2p^+}{\v p^2k^+-\v k^2p^+}\Big]
\notag\\&
+\frac{2ig^2f^{abc}\rho^c(\v k-\v p)\,(k-p)^l\sqrt{k^+p^+}}{(k^+-p^+)^2}\,\frac{k^mp^+-k^+p^m}{\v k^2p^+-\v p^2k^+}\notag\\&\quad\times\Big[\frac{\delta_{jm}\delta_{il}}{p^+}-\frac{\delta_{im}\delta_{jl}}{k^+}-\frac{\delta_{lm}\delta_{ij}}{p^+-k^+}\Big]
\notag\\&\quad\times
\Bigg[\frac{p^+-k^+}{(\v k-\v p)^2}\,\Theta(p^+-k^+-\Lambda)+\frac{\Theta(k^+-\Lambda-p^+)}{\frac{\v k^2}{k^+}-\frac{(\v k-\v p)^2}{k^+-p^+}-\frac{\v p^2}{p^+}}\Bigg]
\end{align} 
and
\begin{align}\label{eq:C1mom}
\mf C^{bca}_{1jki}(k^+,\v k;p^+,\v p)=\ &-\frac{igf^{abc}\,(p^mk^+-p^+k^m)\sqrt{k^+p^+(k^+-p^+)}}{\sqrt2\,(k^+\v p-p^+\v k)^2}\notag\\&\times\Big[\frac{\delta_{im}\delta_{jk}}{k^+}-\frac{\delta_{ik}\delta_{mj}}{p^+}-\frac{\delta_{km}\delta_{ij}}{k^+-p^+}\Big]\Theta(k^+-\Lambda-p^+)\,.
\end{align}
The double poles $1/(k^+-p^+)^2$ of the instantaneous term and of the term with the triple-gluon vertex cancel each other at $|p^+-k^+|\to\Lambda$, so that no term proportional to $k^+/\Lambda$ arises. The coordinate forms, obtained with the Fourier transforms~\eqref{eq:FT_i3}--\eqref{eq:FT_i5}, are Eqs.~\eqref{eq:B3} and~\eqref{eq:C1}.

\begin{figure}[t]
\centering
\tikzfig{gtog_inst}\qquad\tikzfig{gandgtog}\qquad\tikzfig{gto3g_torho}\\[4mm]
\tikzfig{gtovactog_tworho}\qquad\tikzfig{gg_tworho}
\caption{Transition of a soft gluon into a soft gluon through the interaction with the valence charges at order $g^2$.}
\label{fig:g_g}
\end{figure}

\begin{figure}[t]
\centering
\tikzfig{g_togg}\qquad\tikzfig{gtog_plusg}
\caption{Splitting of a soft gluon into two soft gluons ($\rho$-independent, order $g$), and emission of a second soft gluon by the valence charges in the presence of the first one.}
\label{fig:g_gg}
\end{figure}

% ======================================================================
\section{Unitarity relations among the coefficients of \texorpdfstring{$\Omega$}{Omega}}\label{app:unitarity}
% ======================================================================

The coefficients of $\Omega$ that multiply annihilation operators are fixed by the unitarity condition $\Omega^\dagger\Omega=1$, imposed order by order in $g$. With the expansion~\eqref{eq:OmegaExp} and its conjugate, we normal order $\Omega^\dagger\Omega$ with the commutator~\eqref{eq:aadagger}, keeping only the terms needed for the single inclusive cross section at $O(g^4)$ and using $\mf N=\mf N^\dagger$. Since the coefficients are operators on the valence space, their order in each product is kept. The coefficients of the operator structures $1$, $a^\dagger$, $a^\dagger a^\dagger$ and $a^\dagger a^\dagger a$ must vanish (except for the identity), which gives
\begin{align}
\mf N&=1-\frac12\int\frac{d^3p}{(2\pi)^3}\,\mf A^{(1)b}_{1j}(p)\,\mf A^{(1)\dagger b}_{1j}(p)\,,\label{eq:Nunit}\\
\mf A^{(1)\dagger b}_{1j}(p)&=-\mf A^{(1)b}_{2j}(p)\,,\label{eq:A11unit}\\
\mf A^{(3)\dagger b}_{1j}(p)&=-\mf A^{(3)b}_{2j}(p)-\big[\mf N_2,\mf A^{(1)b}_{2j}(p)\big]
+\int\frac{d^3q}{(2\pi)^3}\Big[2\,\mf A^{(1)c}_{1k}(q)\,\mf B^{cb}_{2kj}(p,q)-\mf B^{\dagger cb}_{3kj}(p,q)\,\mf A^{(1)c}_{2k}(q)
\notag\\&\hspace{5.1cm}
-2\,\mf C^{\dagger dcb}_{1lkj}(p,q)\,\mf B^{cd}_{2kl}(q,p-q)\Big]\,,\label{eq:A31unit}\\
\mf B^{\dagger cb}_{1kj}(p,q)&=-\mf B^{cb}_{2kj}(p,q)+\mf C^{cbd}_{1kjl}(q+p,q)\,\mf A^{(1)d}_{2l}(q+p)\notag\\&\quad+\frac12\Big[\mf A^{(1)b}_{2j}(p)\,\mf A^{(1)c}_{2k}(q)+\mf A^{(1)c}_{2k}(q)\,\mf A^{(1)b}_{2j}(p)\Big]\,,\label{eq:B1unit}\\
\mf C^{\dagger dcb}_{2lkj}(p,q)&=-\mf C^{dcb}_{1lkj}(p,q)\,.\label{eq:C2unit}
\end{align}
In addition, the coefficients obey the symmetry relations $\mf B^{cb}_{1kj}(p,q)=\mf B^{bc}_{1jk}(q,p)$, $\mf B^{cb}_{2kj}(p,q)=\mf B^{bc}_{2jk}(q,p)$, $\mf C^{dcb}_{1lkj}(p,q)=\mf C^{cdb}_{1klj}(p,p-q)$, and similarly for $\mf C_2$, which follow from the symmetry of the products of creation and annihilation operators. The coefficient of $a^\dagger a$ gives a relation between $\mf B_3$, $\mf B_3^\dagger$, $\mf A^{(1)}$ and $\mf C_1$, which is not needed for the cross section but serves as a check of $\mf B_3$; in particular it fixes the sign of the non-instantaneous term of $\mf B_3$ for $p^+>k^+$.

In transverse coordinate space the relations read, with the Fourier conventions~\eqref{eq:FTconventions},
\begin{align}
\mf A^{(1)\dagger b}_{1j}(p^+,\v x)&=-\mf A^{(1)b}_{2j}(p^+,-\v x)\,,\notag\\
\mf A^{(3)\dagger b}_{1j}(p^+,\v x)&=-\mf A^{(3)b}_{2j}(p^+,-\v x)-\big[\mf N_2,\mf A^{(1)b}_{2j}(p^+,-\v x)\big]
\notag\\&\quad+\int\frac{dq^+}{2\pi}\int_y\Big[2\,\mf A^{(1)c}_{1k}(q^+,\v y)\,\mf B^{cb}_{2kj}(p^+,-\v x;q^+,-\v y)
\notag\\&\qquad
-\mf B^{\dagger cb}_{3kj}(p^+,\v x;q^+,-\v y)\,\mf A^{(1)c}_{2k}(q^+,-\v y)
\notag\\&\qquad-2\int_z\mf C^{\dagger dcb}_{1lkj}(p^+,\v x-\v z;q^+,\v z-\v y)\,\mf B^{dc}_{2lk}(p^+-q^+,-\v z;q^+,-\v y)\Big]\,,\notag\\
\mf B^{\dagger cb}_{1kj}(p^+,\v x;q^+,\v y)&=-\mf B^{cb}_{2kj}(p^+,-\v x;q^+,-\v y)\notag\\&\quad+\frac12\Big[\mf A^{(1)b}_{2j}(p^+,-\v x)\mf A^{(1)c}_{2k}(q^+,-\v y)\notag\\&\qquad+\mf A^{(1)c}_{2k}(q^+,-\v y)\mf A^{(1)b}_{2j}(p^+,-\v x)\Big]
\notag\\&\quad
+\int_z\mf C^{cbd}_{1kjl}(q^++p^+,\v z-\v x;q^+,\v x-\v y)\,\mf A^{(1)d}_{2l}(q^++p^+,-\v z)\,,\notag\\
\mf C^{\dagger dcb}_{2lkj}(p^+,\v z;q^+,\v x)&=-\mf C^{dcb}_{1lkj}(p^+,-\v z;q^+,-\v x)\,.
\end{align}

\paragraph{The commutator in $\mf A^{(3)}_1$.} The term $-[\mf N_2,\mf A^{(1)}_2]$ in Eq.~\eqref{eq:A31unit} is absent for commuting charges. It does not vanish here: $\mf N_2\propto\int\rho^a(\v x)\rho^a(\v y)K(\v x,\v y,\v z)$, and $[\rho^a(\v x)\rho^a(\v y),\rho^b(\v w)]\neq0$ for $\v x\neq\v y$, so that $[\mf N_2,\mf A^{(1)}_2]\propto g^3f^{abc}\{\rho^a,\rho^c\}\neq0$. With the ordering $[\mf A^{(3)}_2]_{\rho\rho\rho}=\mf A^{(1)}_2\mf N_2$ derived in Appendix~\ref{app:wavefunctions}, the first two terms of Eq.~\eqref{eq:A31unit} combine to
\begin{align}
-\big[\mf A^{(3)}_2\big]_{\rho\rho\rho}-\big[\mf N_2,\mf A^{(1)}_2\big]=-\mf N_2\,\mf A^{(1)}_2\,.
\end{align}
In the one-gluon amplitude~\eqref{eq:A3amp} this produces the combination~\eqref{eq:Nterms} of normalisation terms, which vanishes for $U=1$, as required by $\bar\Omega^\dagger\Omega=1$ at $U=1$. With the opposite ordering of the three-$\rho$ term, $\mf N_2\mf A^{(1)}_2$, the coefficient $\mf A^{(3)}_2$ would differ by $[\mf N_2,\mf A^{(1)}_2]$, a two-$\rho$ term proportional to $\log(\vee/\Lambda)$ equal and opposite to the term $-2k^i\log(\vee/\Lambda)$ of Eq.~\eqref{eq:A3tworho}, and the cross section would acquire a spurious two-$\rho$ term proportional to $\log(\vee/\Lambda)$ (which vanishes at $U=1$ but not otherwise).

% ======================================================================
\section{Fourier transforms}\label{app:fourier}
% ======================================================================

With the conventions~\eqref{eq:FTconventions}, $f(\v k)=\int_xe^{i\v k\cdot\v x}f(\v x)$ and $f(\v x)=\int\frac{d^2k}{(2\pi)^2}e^{-i\v k\cdot\v x}f(\v k)$, the transverse integrals needed to transform the wave functions and coefficients from momentum to coordinate space are~\cite{Radhakrishnan:2026soft}
\begin{align}
&\int\frac{d^2k}{2\pi}\,\frac{k^i}{\v k^2}\,e^{i\v k\cdot\v X}=\frac{iX^i}{\v X^2}\,,\label{eq:FT_K}\\
&\int\frac{d^2k}{2\pi}\frac{d^2p}{2\pi}\,\frac{k^ip^j}{\v k^2\,(a\v k^2+b\v p^2)}\,e^{i\v k\cdot\v X+i\v p\cdot\v Z}=-\frac{X^iZ^j}{\v Z^2\,(b\v X^2+a\v Z^2)}\,,\label{eq:FT_i2}\\
&\int\frac{d^2k}{2\pi}\frac{d^2p}{2\pi}\,\frac{1}{\xi\v k^2-\v p^2}\,e^{i\v k\cdot\v X+i\v p\cdot\v Z}=\frac{1}{\xi\v Z^2-\v X^2}\,,\label{eq:FT_i3}\\
&\int\frac{d^2k}{2\pi}\frac{d^2p}{2\pi}\,\frac{k^m-p^m}{(\v k-\v p)^2}\,\frac{p^l-\xi k^l}{\v p^2-\xi\v k^2}\,e^{i\v k\cdot\v X+i\v p\cdot\v Y}=-\frac{X^m+\xi Y^m}{\v X^2-\xi\v Y^2}\,\frac{X^l+Y^l}{(\v X+\v Y)^2}\,,\label{eq:FT_i4}\\
&\int\frac{d^2k}{2\pi}\frac{d^2p}{2\pi}\,\frac{k^m-p^m}{\v p^2-\xi\v k^2}\,\frac{p^l-\xi k^l}{(\v p-\xi\v k)^2}\,e^{i\v k\cdot\v X+i\v p\cdot\v Y}=-\frac{X^l+Y^l}{\v X^2-\xi\v Y^2}\,\frac{X^m+\xi Y^m}{(\v X+\xi\v Y)^2}\,,\label{eq:FT_i5}\\
&\int\frac{d^2k}{2\pi}\frac{d^2p}{2\pi}\,\frac{1}{\xi(1-\xi)\v k^2+\v p^2}\,e^{i\v k\cdot\v X+i\v p\cdot\v Z}=\frac{1}{\v X^2+\xi(1-\xi)\v Z^2}\,,\label{eq:FT_i7}\\
&\int\frac{d^2k}{2\pi}\,\frac{k^i}{\v k^2}\,\log\frac{\v k^2}{\mu^2}\,e^{i\v k\cdot\v X}=\frac{iX^i}{\v X^2}\Big[-2\gamma_E-\log\frac{\v X^2\mu^2}{4}\Big]\,.\label{eq:FT_log}
\end{align}
The integrals~\eqref{eq:FT_i3}--\eqref{eq:FT_i5}, whose denominators vanish inside the integration region, are understood as principal values. Eq.~\eqref{eq:FT_K} gives the Weizs\"acker--Williams field: in the conventions used here $\int_Xe^{i\v k\cdot\v X}\,X^i/\v X^2=2\pi i\,k^i/\v k^2$. Eq.~\eqref{eq:FT_log} follows from the angular integral $\int_0^{2\pi}d\theta\cos\theta\,e^{ikX\cos\theta}=2\pi iJ_1(kX)$ and $\int_0^\infty dt\,J_1(t)\log t=\log2-\gamma_E$; it relates the $\MS$ logarithm in momentum space to the combination
\begin{align}\label{eq:c0}
2\gamma_E+\log\frac{\v X^2\mu^2_{\MS}}{4}=\log\frac{\v X^2\mu^2_{\MS}}{c_0^2}\,,\qquad c_0=2e^{-\gamma_E}\,,
\end{align}
that appears in coordinate space.

% ======================================================================
\section{Prescriptions and distributions}\label{app:distributions}
% ======================================================================

We collect here the conventions used for the singular longitudinal and transverse integrals.

\paragraph{Principal values.} Energy denominators that vanish inside the region of integration, such as $p^+\v k^2-k^+\v p^2$ in momentum space or $p^+(\v x-\v u)^2-k^+(\v x-\v z)^2$ in coordinate space, arise from transitions between states of equal free energy and are understood as Cauchy principal values. Correspondingly, logarithms that arise from integrals over such denominators are logarithms of absolute values, $\log|\cdot|$. All finite terms of this paper refer to this prescription (Section~\ref{sec:fin_conv}).

\paragraph{Dimensional regularisation.} Transverse loop integrals are computed in $d=2-\epsilon$ dimensions, $\int\frac{d^2p}{(2\pi)^2}\to\mu^{\epsilon}\int\frac{d^dp}{(2\pi)^d}$, in the $\MS$ scheme with scale $\mu_{\MS}$. Scaleless integrals vanish, $\int d^dp/\v p^2=0$. Coordinate-space logarithms are written with the combination~\eqref{eq:c0}.

\paragraph{Plus prescriptions in $p^+$.} For a function $f(p^+)$ regular at $p^+=0$,
\begin{align}\label{eq:plusp}
\int_0^{P}dp^+\Big[\frac{1}{p^+}\Big]_+f(p^+)\equiv\int_0^{P}dp^+\,\frac{f(p^+)-f(0)}{p^+}\,.
\end{align}
The subtracted term, integrated from $\Lambda$, is $f(0)\log(P/\Lambda)$; when $f(0)$ is the JIMWLK (target-side) structure it is part of the JIMWLK term.

\paragraph{Plus prescriptions in the momentum fraction.} With $\xi=k^+/(p^++k^+)$ the momentum fraction of the observed gluon in its parent, and for a function $F(\xi)$,
\begin{align}\label{eq:plusxi}
\int_0^1d\xi\,\frac{\xi}{(1-\xi)_+}F(\xi)&\equiv\int_0^1d\xi\,\frac{\xi F(\xi)-F(1)}{1-\xi}\,,\notag\\
\int_0^1d\xi\,\Big[\frac{1-\xi}{\xi}\Big]_+F(\xi)&\equiv\int_0^1d\xi\,\frac{1-\xi}{\xi}\big[F(\xi)-F(0)\big]\,.
\end{align}
The first is the standard prescription at $\xi\to1$ (integrated gluon soft); the second subtracts the pole at $\xi\to0$ (observed gluon soft relative to its parent), which belongs to the projectile (BFKL) evolution. The combination of the two subtractions coincides with the ``double plus'' prescription of~\cite{Kovner:2023vsy}: with $N(\xi)=\xi^2+(1-\xi)^2+\xi^2(1-\xi)^2$, $N(0)=N(1)=1$,
\begin{align}
\int_0^1\frac{d\xi}{\xi_+(1-\xi)_+}\,N(\xi)\,F(\xi)&\equiv\int_0^1d\xi\Big[\frac{N(\xi)F(\xi)}{\xi(1-\xi)}-\frac{F(0)}{\xi}-\frac{F(1)}{1-\xi}\Big]\,,\notag\\
\frac{N(\xi)}{\xi(1-\xi)}&=\xi(1-\xi)+\frac{\xi}{1-\xi}+\frac{1-\xi}{\xi}\,,
\end{align}
and $\int_0^1\frac{d\xi}{\xi_+(1-\xi)_+}N(\xi)=-\frac{11}{6}$.

\paragraph{Integrals over $p^+$ up to $\vee\to\infty$.} Some terms have no dependence on $p^+$ other than an oscillating phase. Their integral converges for $\vee\to\infty$ in the sense of distributions,
\begin{align}\label{eq:distrib}
\int_0^\infty dp^+\,e^{-iap^+}=\pi\,\delta(a)-i\,\mathcal P\frac1a\,.
\end{align}

\paragraph{Recoil function.} When only the recoil phase of the observed gluon depends on the longitudinal momentum, the plus-prescription integral of the rows with $\mf B_3$ and $\mf C_1$ is the function (Appendix~\ref{app:closedforms})
\begin{align}\label{eq:Erecoil}
\mathcal E(s)=\int_0^1\frac{dt}{t}\big(e^{-ist}-1\big)\,,
\end{align}
which vanishes at $s=0$ and grows logarithmically for $|s|\to\infty$, $\mathcal E(s)=-\log|s|-\gamma_E-\frac{i\pi}{2}{\rm sign}(s)+O(1/s)$.

\paragraph{Longitudinal integrals.} The following integrals, with $\delta=\Lambda/k^+\to0$, produce the constants of the one-gluon wave function at order $g^3$, Eq.~\eqref{eq:xiintegrals}:
\begin{align}
\int_\delta^{1-\delta}d\xi\Big[-2+\xi(1-\xi)+\frac{1}{\xi(1-\xi)}\Big]&=-\frac{11}{6}-2\log\delta\,,\notag\\
\int_0^1d\xi\,\big[-2+\xi(1-\xi)\big]\log\big(\xi(1-\xi)\big)&=\frac{67}{18}\,,\notag\\
\int_\delta^{1-\delta}d\xi\,\frac{\log(\xi(1-\xi))}{\xi(1-\xi)}&=-\log^2\delta-\frac{\pi^2}{3}\,,
\end{align}
up to terms that vanish for $\delta\to0$.

% ======================================================================
\begingroup\sloppy
\section{Finite terms for an arbitrary projectile}\label{app:finite}
% ======================================================================
% The equations (files rows/H-*.tex) are generated mechanically by scripts/gen_rows.py from Sec. 8.1 of
% nlo_review/final_explicit/Final_terms_explicit_2_fixed.tex; do not edit them by hand.

This appendix lists the finite terms before the average over the projectile, Group by Group and Row by Row, after the exact cancellations of Section~\ref{sec:fin_together} and with every product of three charges written, by means of Eqs.~\eqref{eq:symord}--\eqref{eq:symr21}, as its symmetrised product (``sym'') plus two-$\rho$ contact terms (``con''); the genuine two-$\rho$ terms are marked ``quad''. The conventions are those of Section~\ref{sec:fin_conv}, and the kernels $\mathfrak K$ are those defined in Section~\ref{sec:fin_rows}. The complex conjugate is implied in Groups~I and~III and written as ``$+{\rm c.c.}$''. The colour tensors are written with differences of Wilson lines; the statement of what can and what cannot be factorised in this way is given in Section~\ref{sec:fin_singlet}. The one-$\rho$ (double commutator) terms of Eqs.~\eqref{eq:symord}--\eqref{eq:symr21} are not written; they have zero average for a projectile invariant under global colour rotations. The averages of the equations of this appendix according to Eqs.~\eqref{eq:A1n}--\eqref{eq:cubicavg} are the equations of Section~\ref{sec:fin_rows} with the same row and part.

\subsection{Group I}

\paragraph{Row I.II.} Symmetrised three-$\rho$ term, contact terms and genuine two-$\rho$ term, with the kernel~\eqref{eq:K-rm_I_II}:
\begin{align}\label{eq:H-I.II-3}
&\frac{d^{3}N^{\rm I.II}_{{\rm sym}}}{d^{3}k}=-\frac{1}{(2\pi)^{3}}\, \frac{g^{4}}{64\pi^{6}}\, \frac{1}{k^{+}} \int_{w,w'}e^{-ik\cdot(w-w')} \int_{x,y,x'} \times \mathfrak{K}_{\rm I.II}\notag\\
&\quad \times 2\, \Bigg\{-\Big[U^{ga}(x')-U^{ga}(w')\Big]f^{gde}\bigg[\Big[U^{db}(x)-U^{db}(w)\Big]U^{ec}(y)\notag\\
&\quad\quad\quad +U^{db}(w)\Big[U^{ec}(y)-U^{ec}(w)\Big]\bigg]\Bigg\}\big\{\rho^{a}(x')\rho^{b}(x)\rho^{c}(y)\big\}_{\rm sym} +{\rm c.c.}
\end{align}

\begin{align}\label{eq:H-I.II-2}
&\frac{d^{3}N^{\rm I.II}_{{\rm con}}}{d^{3}k}=-\frac{1}{(2\pi)^{3}}\, \frac{g^{4}}{64\pi^{6}}\, \frac{1}{k^{+}} \int_{w,w'}e^{-ik\cdot(w-w')} \int_{x,y,x'} \times \mathfrak{K}_{\rm I.II}\notag\\
&\quad \times i\, \Bigg\{-\Big[U^{ga}(x')-U^{ga}(w')\Big]f^{gde}\bigg[\Big[U^{db}(x)-U^{db}(w)\Big]U^{ec}(y)\notag\\
&\quad\quad\quad +U^{db}(w)\Big[U^{ec}(y)-U^{ec}(w)\Big]\bigg]\Bigg\} \times \Big[f^{abr}\, \delta^{(2)}(x'-x)\big\{\rho^{r}(x')\rho^{c}(y)\big\}_{\rm sym}\notag\\
&\quad\quad +f^{acr}\, \delta^{(2)}(x'-y)\big\{\rho^{r}(x')\rho^{b}(x)\big\}_{\rm sym}\Big] +{\rm c.c.}
\end{align}

\begin{align}\label{eq:H-I.II-2o}
&\frac{d^{3}N^{\rm I.II}_{{\rm quad}}}{d^{3}k}=+\frac{1}{(2\pi)^{3}}\, \frac{g^{4}\,N_{c}}{32\pi^{4}}\, \frac{1}{k^{+}}\Big(-\frac{67}{9}+\frac{\pi^{2}}{3}\Big) \int_{w,w',x,x'} e^{-ik\cdot(w-w')}\;\Kv(x'-w')\cdot\Kv(x-w)\notag\\
&\quad \times\Big[U^{ab'}(x')-U^{ab'}(w')\Big]\Big[U^{ab}(w)-U^{ab}(x)\Big]\rho^{b'}(x')\, \rho^{b}(x) + {\rm c.c.}
\end{align}


\paragraph{Row I.III.} Symmetrised three-$\rho$ term and contact terms, with the kernels~\eqref{eq:K-z_rm_I_III} and~\eqref{eq:K-w_rm_I_III}:
\begin{align}\label{eq:H-I.III-3}
&\frac{d^{3}N^{\rm I.III}_{{\rm sym}}}{d^{3}k}=\frac{1}{(2\pi)^{3}}\, \frac{ig^{4}}{8\pi^{5}}\, \frac{1}{k^{+}} \int_{w,w'}e^{-ik\cdot(w-w')} \int_{x,y,z,x'}\Bigg[ \mathfrak{K}^{(z)}_{\rm I.III}\notag\\
&\quad\quad \times \Bigg\{\Big[U^{ga}(x')-U^{ga}(w')\Big]f^{ghd}\, U^{db}(x)\Big[U^{hc}(z)-U^{hc}(y)\Big]\Bigg\}\, \big\{\rho^{a}(x')\rho^{b}(x)\rho^{c}(y)\big\}_{\rm sym}\notag\\
&\quad\quad + \mathfrak{K}^{(w)}_{\rm I.III} \times \Bigg\{\Big[U^{ga}(x')-U^{ga}(w')\Big]f^{ghd}\, U^{db}(x)\Big[U^{hc}(w)-U^{hc}(y)\Big]\Bigg\}\notag\\
&\quad\quad \times \big\{\rho^{a}(x')\rho^{b}(x)\rho^{c}(y)\big\}_{\rm sym}\Bigg] +{\rm c.c.}
\end{align}

\begin{align}\label{eq:H-I.III-2}
&\frac{d^{3}N^{\rm I.III}_{{\rm con}}}{d^{3}k}=\frac{1}{(2\pi)^{3}}\, \frac{ig^{4}}{8\pi^{5}}\, \frac{1}{k^{+}} \int_{w,w'}e^{-ik\cdot(w-w')} \int_{x,y,z,x'}\Bigg[ \mathfrak{K}^{(z)}_{\rm I.III}\notag\\
&\quad\quad \times \frac{i}{2}\, \Bigg\{\Big[U^{ga}(x')-U^{ga}(w')\Big]f^{ghd}\, U^{db}(x)\Big[U^{hc}(z)-U^{hc}(y)\Big]\Bigg\}\notag\\
&\quad\quad \times\Big[f^{abr}\, \delta^{(2)}(x'-x)\big\{\rho^{r}(x')\rho^{c}(y)\big\}_{\rm sym} +f^{bcr}\, \delta^{(2)}(x-y)\big\{\rho^{a}(x')\rho^{r}(x)\big\}_{\rm sym}\notag\\
&\quad\quad\quad +f^{acr}\, \delta^{(2)}(x'-y)\big\{\rho^{r}(x')\rho^{b}(x)\big\}_{\rm sym}\Big] + \mathfrak{K}^{(w)}_{\rm I.III}\notag\\
&\quad\quad \times \frac{i}{2}\, \Bigg\{\Big[U^{ga}(x')-U^{ga}(w')\Big]f^{ghd}\, U^{db}(x)\Big[U^{hc}(w)-U^{hc}(y)\Big]\Bigg\}\notag\\
&\quad\quad \times\Big[f^{abr}\, \delta^{(2)}(x'-x)\big\{\rho^{r}(x')\rho^{c}(y)\big\}_{\rm sym} +f^{bcr}\, \delta^{(2)}(x-y)\big\{\rho^{a}(x')\rho^{r}(x)\big\}_{\rm sym}\notag\\
&\quad\quad\quad +f^{acr}\, \delta^{(2)}(x'-y)\big\{\rho^{r}(x')\rho^{b}(x)\big\}_{\rm sym}\Big]\Bigg] +{\rm c.c.}
\end{align}


\paragraph{Row I.V.} Symmetrised three-$\rho$ term and contact terms, with the kernel~\eqref{eq:K-rm_I_V}:
\begin{align}\label{eq:H-I.V-3}
&\frac{d^{3}N^{\rm I.V}_{{\rm sym}}}{d^{3}k}=\frac{1}{(2\pi)^{3}}\, \frac{ig^{4}}{16\pi^{4}}\, \frac{1}{k^{+}} \int_{0}^{1}\frac{d\xi}{2\pi} \int_{x,z,y,y',x',w'}e^{-ik\cdot(z-w')} \times \mathfrak{K}_{\rm I.V}\notag\\
&\quad \times 2\, \Bigg\{\Big[U^{ga}(x')-U^{ga}(w')\Big]f^{gde}\bigg[\Big[U^{dc}(x)-U^{dc}(y')\Big]U^{eb}(z)\notag\\
&\quad\quad\quad +U^{dc}(y')\Big[U^{eb}(z)-U^{eb}(y)\Big]\bigg]\Bigg\} \times \big\{\rho^{a}(x')\rho^{b}(y)\rho^{c}(y')\big\}_{\rm sym} +{\rm c.c.}
\end{align}

\begin{align}\label{eq:H-I.V-2}
&\frac{d^{3}N^{\rm I.V}_{{\rm con}}}{d^{3}k}=\frac{1}{(2\pi)^{3}}\, \frac{ig^{4}}{16\pi^{4}}\, \frac{1}{k^{+}} \int_{0}^{1}\frac{d\xi}{2\pi} \int_{x,z,y,y',x',w'}e^{-ik\cdot(z-w')} \times \mathfrak{K}_{\rm I.V}\notag\\
&\quad \times i\, \Bigg\{\Big[U^{ga}(x')-U^{ga}(w')\Big]f^{gde}\bigg[\Big[U^{dc}(x)-U^{dc}(y')\Big]U^{eb}(z)\notag\\
&\quad\quad\quad +U^{dc}(y')\Big[U^{eb}(z)-U^{eb}(y)\Big]\bigg]\Bigg\} \times\Big[f^{abr}\, \delta^{(2)}(x'-y)\big\{\rho^{r}(x')\rho^{c}(y')\big\}_{\rm sym}\notag\\
&\quad\quad +f^{acr}\, \delta^{(2)}(x'-y')\big\{\rho^{r}(x')\rho^{b}(y)\big\}_{\rm sym}\Big] +{\rm c.c.}
\end{align}

Genuine two-$\rho$ term, with the kernel~\eqref{eq:K-rm_quad_rm_I_V}:
\begin{align}\label{eq:H-I.V-2o}
&\frac{d^{3}N^{\rm I.V}_{{\rm quad}}}{d^{3}k}=\frac{1}{(2\pi)^{3}}\, \frac{g^{4}}{8\pi^{4}}\, \frac{1}{k^{+}} \int_{0}^{1}\frac{d\xi}{2\pi} \int_{x,y,z,x',w'}e^{-ik\cdot(z-w')} \times \mathfrak{K}^{\rm quad}_{\rm I.V}\notag\\
&\quad \times\Bigg\{-\Big[U^{ab'}(x')-U^{ab'}(w')\Big]f^{adc}f^{fbe}\, U^{db}(x)\Big[U^{ce}(z)-U^{ce}(x)\Big]\notag\\
&\quad\quad +N_{c}\Big[U^{ab'}(x')-U^{ab'}(w')\Big]\Big[U^{af}(y)-U^{af}(x)\Big]\Bigg\}\rho^{b'}(x')\, \rho^{f}(y)\notag\\
&\quad +\frac{1}{(2\pi)^{3}}\, \frac{g^{4}N_{c}}{16\pi^{5}}\, \frac{1}{k^{+}} \int_{x,y,z,x',w'}e^{-ik\cdot(z-w')}\, e^{-ik\cdot(x-z)}\, \frac{11}{6}\, \Kv(x'-w')\cdot\Kv(y-x)\notag\\
&\quad \times \frac{(y-x)\cdot\Kv(y-z)}{(x-z)^{2}} \times\Big[U^{ab'}(x')-U^{ab'}(w')\Big]\Big[U^{af}(y)-U^{af}(x)\Big]\rho^{b'}(x')\, \rho^{f}(y)\, + {\rm c.c.}
\end{align}


\subsection{Group II}

\paragraph{Row II.I.} Symmetrised three-$\rho$ term and contact terms, with the kernels~\eqref{eq:K-1_rm_II_I} and~\eqref{eq:K-2_rm_II_I}:
\begin{align}\label{eq:H-II.I-3}
&\frac{d^{3}N^{\rm II.I}_{{\rm sym}}}{d^{3}k}=\frac{1}{(2\pi)^{3}}\, \frac{ig^{4}}{16\pi^{5}}\, \frac{1}{k^{+}} \int_{x,y,z,w,x',w'}e^{-ik\cdot(w-w')}\Bigg[ \mathfrak{K}^{(1)}_{\rm II.I}\notag\\
&\quad\quad \times 2\, \Bigg\{\Big[U^{ga}(w')-U^{ga}(x')\Big]\bigg[\Big[U^{gk}(w)-U^{gk}(z)\Big]f^{kbc} -f^{gkl}\, U^{kb}(z)\Big[U^{lc}(x)-U^{lc}(z)\Big]\bigg]\notag\\
&\quad\quad\quad +U^{ga}(x')\bigg[U^{mh}(z)\Big[U^{mb}(z)-U^{mb}(y)\Big]\Big[U^{gk}(w)-U^{gk}(z)\Big]f^{khc}\notag\\
&\quad\quad\quad\quad -f^{gml}\Big[U^{mb}(z)-U^{mb}(y)\Big]\Big[U^{lc}(x)-U^{lc}(z)\Big]\bigg]\Bigg\}\big\{\rho^{a}(x')\rho^{b}(y)\rho^{c}(x)\big\}_{\rm sym} + \mathfrak{K}^{(2)}_{\rm II.I}\notag\\
&\quad\quad \times 2\, \Bigg\{\Big[U^{gb}(w)-U^{gb}(x')\Big]\bigg[\Big[U^{gk}(w')-U^{gk}(z)\Big]f^{kca} -f^{gkl}\, U^{kc}(z)\Big[U^{la}(x)-U^{la}(z)\Big]\bigg]\notag\\
&\quad\quad\quad +U^{gb}(x')\bigg[U^{mh}(z)\Big[U^{mc}(z)-U^{mc}(y)\Big]\Big[U^{gk}(w')-U^{gk}(z)\Big]f^{kha}\notag\\
&\quad\quad\quad\quad -f^{gml}\Big[U^{mc}(z)-U^{mc}(y)\Big]\Big[U^{la}(x)-U^{la}(z)\Big]\bigg]\Bigg\}\big\{\rho^{a}(x)\rho^{b}(x')\rho^{c}(y)\big\}_{\rm sym}\Bigg]
\end{align}

\begin{align}\label{eq:H-II.I-2}
&\frac{d^{3}N^{\rm II.I}_{{\rm con}}}{d^{3}k}=\frac{1}{(2\pi)^{3}}\, \frac{ig^{4}}{16\pi^{5}}\, \frac{1}{k^{+}} \int_{x,y,z,w,x',w'}e^{-ik\cdot(w-w')}\Bigg[ \mathfrak{K}^{(1)}_{\rm II.I}\notag\\
&\quad\quad \times i\, \Bigg\{\Big[U^{ga}(w')-U^{ga}(x')\Big]\bigg[\Big[U^{gk}(w)-U^{gk}(z)\Big]f^{kbc} -f^{gkl}\, U^{kb}(z)\Big[U^{lc}(x)-U^{lc}(z)\Big]\bigg]\notag\\
&\quad\quad\quad +U^{ga}(x')\bigg[U^{mh}(z)\Big[U^{mb}(z)-U^{mb}(y)\Big]\Big[U^{gk}(w)-U^{gk}(z)\Big]f^{khc}\notag\\
&\quad\quad\quad\quad -f^{gml}\Big[U^{mb}(z)-U^{mb}(y)\Big]\Big[U^{lc}(x)-U^{lc}(z)\Big]\bigg]\Bigg\} \Big[f^{bcr}\notag\\
&\quad\quad\quad \times \delta^{(2)}(y-x)\big\{\rho^{a}(x')\rho^{r}(y)\big\}_{\rm sym} +f^{acr}\, \delta^{(2)}(x'-x)\big\{\rho^{r}(x')\rho^{b}(y)\big\}_{\rm sym}\Big] + \mathfrak{K}^{(2)}_{\rm II.I}\notag\\
&\quad\quad \times i\, \Bigg\{\Big[U^{gb}(w)-U^{gb}(x')\Big]\bigg[\Big[U^{gk}(w')-U^{gk}(z)\Big]f^{kca} -f^{gkl}\, U^{kc}(z)\Big[U^{la}(x)-U^{la}(z)\Big]\bigg]\notag\\
&\quad\quad\quad +U^{gb}(x')\bigg[U^{mh}(z)\Big[U^{mc}(z)-U^{mc}(y)\Big]\Big[U^{gk}(w')-U^{gk}(z)\Big]f^{kha}\notag\\
&\quad\quad\quad\quad -f^{gml}\Big[U^{mc}(z)-U^{mc}(y)\Big]\Big[U^{la}(x)-U^{la}(z)\Big]\bigg]\Bigg\} \Big[f^{abr}\notag\\
&\quad\quad\quad \times \delta^{(2)}(x-x')\big\{\rho^{r}(x)\rho^{c}(y)\big\}_{\rm sym} +f^{acr}\, \delta^{(2)}(x-y)\big\{\rho^{r}(x)\rho^{b}(x')\big\}_{\rm sym}\Big]\Bigg]
\end{align}

Genuine two-$\rho$ terms, with the kernels~\eqref{eq:K-rm_quad_a_rm_II_I}--\eqref{eq:K-rm_quad_d_rm_II_I}:
\begin{align}\label{eq:H-II.I-2a}
&\frac{d^{3}N^{\rm II.I}_{{\rm quad},a}}{d^{3}k}=\frac{1}{(2\pi)^{3}}\, \frac{g^{4}}{8\pi^{5}}\, \frac{1}{k^{+}} \int_{z,w,w',x,x'}e^{-ik\cdot(w-w')}\, \mathfrak{K}^{{\rm quad},a}_{\rm II.I}\notag\\
&\quad \times\bigg[\Big[U^{ac'}(w')-U^{ac'}(z)\Big]f^{c'bd'} -f^{ac'e'}\, U^{c'b}(z)\Big[U^{e'd'}(x')-U^{e'd'}(z)\Big]\bigg]\notag\\
&\quad \times\bigg[\Big[U^{ac}(w)-U^{ac}(z)\Big]f^{cbd} -f^{ace}\, U^{cb}(z)\Big[U^{ed}(x)-U^{ed}(z)\Big]\bigg]\rho^{d'}(x')\, \rho^{d}(x)
\end{align}

\begingroup\small
\begin{align}\label{eq:H-II.I-2b}
&\frac{d^{3}N^{\rm II.I}_{{\rm quad},b}}{d^{3}k}=-\frac{1}{(2\pi)^{3}}\, \frac{g^{4}}{8\pi^{5}}\, \frac{1}{k^{+}} \int_{z,w,w',x,x'}e^{-ik\cdot(w-w')}\, \mathfrak{K}^{{\rm quad},b}_{\rm II.I} \times\bigg[\Big[U^{ac'}(w')-U^{ac'}(z)\Big]f^{c'bd'}\notag\\
&\quad\quad -f^{ac'e'}\, U^{c'b}(z)\Big[U^{e'd'}(x')-U^{e'd'}(z)\Big]\bigg] \times\bigg[\Big[U^{ac}(w)-U^{ac}(z)\Big]f^{cbd}\notag\\
&\quad\quad -f^{ace}\, U^{cb}(z)\Big[U^{ed}(x)-U^{ed}(z)\Big]\bigg]\rho^{d'}(x')\, \rho^{d}(x)
\end{align}
\endgroup

\begin{align}\label{eq:H-II.I-2c}
&\frac{d^{3}N^{\rm II.I}_{{\rm quad},c}}{d^{3}k}=\frac{1}{(2\pi)^{3}}\, \frac{g^{4}}{8\pi^{5}}\, \frac{1}{k^{+}} \int_{z,w,w',x,x'}e^{-ik\cdot(w-w')}\, \mathfrak{K}^{{\rm quad},c}_{\rm II.I}\notag\\
&\quad \times\bigg[\Big[U^{ac'}(w')-U^{ac'}(z)\Big]f^{c'bd'} -f^{ac'e'}\, U^{c'b}(z)\Big[U^{e'd'}(x')-U^{e'd'}(z)\Big]\bigg]\notag\\
&\quad \times\bigg[\Big[U^{ac}(w)-U^{ac}(z)\Big]f^{cbd} -f^{ace}\, U^{cb}(z)\Big[U^{ed}(x)-U^{ed}(z)\Big]\bigg]\rho^{d'}(x')\, \rho^{d}(x)
\end{align}

\begin{align}\label{eq:H-II.I-2d}
&\frac{d^{3}N^{\rm II.I}_{{\rm quad},d}}{d^{3}k}=\frac{1}{(2\pi)^{3}}\, \frac{g^{4}}{8\pi^{5}}\, \frac{1}{k^{+}} \int_{z,w,w',x,x'}e^{-ik\cdot(w-w')}\, \mathfrak{K}^{{\rm quad},d}_{\rm II.I}\notag\\
&\quad \times\bigg[\Big[U^{ac'}(w')-U^{ac'}(z)\Big]f^{c'bd'} -f^{ac'e'}\, U^{c'b}(z)\Big[U^{e'd'}(x')-U^{e'd'}(z)\Big]\bigg]\notag\\
&\quad \times\bigg[\Big[U^{ac}(w)-U^{ac}(z)\Big]f^{cbd} -f^{ace}\, U^{cb}(z)\Big[U^{ed}(x)-U^{ed}(z)\Big]\bigg]\rho^{d'}(x')\, \rho^{d}(x)
\end{align}

\begin{align}\label{eq:H-II.I-2e}
&\frac{d^{3}N^{\rm II.I}_{{\rm quad},e}}{d^{3}k}=-\frac{1}{(2\pi)^{3}}\, \frac{g^{4}}{16\pi^{5}}\, \frac{1}{k^{+}} \int_{z,w,w',x,x'}e^{-ik\cdot(w-w')}\, \frac{1}{(z-w)^{2}\,(z-w')^{2}}\notag\\
&\quad \times\bigg[\Big[U^{ac'}(w')-U^{ac'}(z)\Big]f^{c'bd'} -f^{ac'e'}\, U^{c'b}(z)\Big[U^{e'd'}(x')-U^{e'd'}(z)\Big]\bigg]\notag\\
&\quad \times\bigg[\Big[U^{ac}(w)-U^{ac}(z)\Big]f^{cbd} -f^{ace}\, U^{cb}(z)\Big[U^{ed}(x)-U^{ed}(z)\Big]\bigg]\rho^{d'}(x')\, \rho^{d}(x)
\end{align}


\paragraph{Row II.IV.} With $\xi$ of Eq.~\eqref{eq:xi_defIV}:
\begin{align}\label{eq:H-II.IVrem}
&\frac{d^{3}N^{\rm II.IV}_{{\rm fin}}}{d^{3}k}=+\frac{1}{(2\pi)^{3}}\, \frac{g^{4}\,N_{c}}{4\pi^{3}}\, \frac{1}{k^{+}} \int_{k^{+}/\vee}^{1}\frac{d\xi}{2\pi}\, \frac{1-\xi}{\xi} \int_{x,y,x',y'}e^{-ik\cdot(y-y')\frac{1}{\xi}}\Big[\Kv(x-y)\cdot\Kv(x'-y')\notag\\
&\quad\quad -2\, \Kv(x-y)\cdot(y'-y)\, \Kv(x'-y')\cdot\Kv(y'-y)\Big]\notag\\
&\quad \times\Big[U^{ce'}(y') - U^{ce'}(x')\Big]\Big[U^{ce}(y) - U^{ce}(x)\Big]\rho^{e'}(x')\rho^{e}(x)\notag\\
&\quad -\frac{1}{(2\pi)^{3}}\, \frac{g^{4}\,N_{c}}{2\pi^{3}}\, \frac{1}{k^{2}}\, \frac{1}{k^{+}}\notag\\
&\quad \times \int_{x,x',z}\Kv(x-z)\cdot\Kv(x'-z)\Big[U^{ce'}(z)-U^{ce'}(x')\Big]\Big[U^{ce}(z)-U^{ce}(x)\Big]\rho^{e'}(x')\rho^{e}(x)
\end{align}


\subsection{Group III}

\paragraph{Row III.II.} Symmetrised three-$\rho$ term and contact terms, with the kernel~\eqref{eq:K-rm_III_II}:
\begin{align}\label{eq:H-III.II-3}
&\frac{d^{3}N^{\rm III.II}_{{\rm sym}}}{d^{3}k}=-\frac{1}{(2\pi)^{3}}\, \frac{ig^{4}}{8\pi^{5}}\, \frac{1}{k^{+}} \int_{x,y,x',z,w,w'}e^{-ik\cdot(w-w')}\, \mathfrak{K}_{\rm III.II}\notag\\
&\quad \times \Bigg\{-U^{gb}(x')\, U^{mh}(z)\Big[U^{ma}(z)-U^{ma}(y)\Big]\Big[U^{gk}(w)-U^{gk}(z)\Big]f^{khc}\notag\\
&\quad\quad +U^{gb}(x')\, f^{gml}\Big[U^{ma}(z)-U^{ma}(y)\Big]\Big[U^{lc}(x)-U^{lc}(z)\Big]\Bigg\}\, \big\{\rho^{a}(y)\rho^{b}(x')\rho^{c}(x)\big\}_{\rm sym} +{\rm c.c.}
\end{align}

\begin{align}\label{eq:H-III.II-2}
&\frac{d^{3}N^{\rm III.II}_{{\rm con}}}{d^{3}k}=-\frac{1}{(2\pi)^{3}}\, \frac{ig^{4}}{8\pi^{5}}\, \frac{1}{k^{+}} \int_{x,y,x',z,w,w'}e^{-ik\cdot(w-w')}\, \mathfrak{K}_{\rm III.II}\notag\\
&\quad \times \frac{i}{2}\, \Bigg\{-U^{gb}(x')\, U^{mh}(z)\Big[U^{ma}(z)-U^{ma}(y)\Big]\Big[U^{gk}(w)-U^{gk}(z)\Big]f^{khc}\notag\\
&\quad\quad +U^{gb}(x')\, f^{gml}\Big[U^{ma}(z)-U^{ma}(y)\Big]\Big[U^{lc}(x)-U^{lc}(z)\Big]\Bigg\} \Big[f^{abr}\notag\\
&\quad\quad \times \delta^{(2)}(y-x')\big\{\rho^{r}(y)\rho^{c}(x)\big\}_{\rm sym} +f^{bcr}\, \delta^{(2)}(x'-x)\big\{\rho^{a}(y)\rho^{r}(x')\big\}_{\rm sym}\notag\\
&\quad\quad +f^{acr}\, \delta^{(2)}(y-x)\big\{\rho^{r}(y)\rho^{b}(x')\big\}_{\rm sym}\Big] +{\rm c.c.}
\end{align}


\paragraph{Row III.III.} Symmetrised three-$\rho$ term and contact terms, with the kernels~\eqref{eq:K-y_rm_III} and~\eqref{eq:K-z_rm_III}:
\begin{align}\label{eq:H-III.III-3}
&\frac{d^{3}N^{\rm III.III}_{{\rm sym}}}{d^{3}k}=\frac{1}{(2\pi)^{3}}\, \frac{ig^{4}}{8\pi^{4}}\, \frac{1}{k^{+}} \int_{0}^{\infty}\frac{d\xi}{2\pi} \int_{w,x',y',z,y,x}e^{-ik\cdot(w-y')}e^{-ik\cdot(z-y')\,\xi}\Bigg[ \mathfrak{K}^{(y')}_{\rm III}\notag\\
&\quad\quad \times 2\, \Bigg\{f^{mnk}\Big[U^{ma}(y')-U^{ma}(x')\Big]\bigg[\Big[U^{nc}(z)-U^{nc}(x)\Big]U^{kb}(w)\notag\\
&\quad\quad\quad\quad +U^{nc}(x)\Big[U^{kb}(w)-U^{kb}(y)\Big]\bigg] \Bigg\}\big\{\rho^{a}(x')\rho^{b}(y)\rho^{c}(x)\big\}_{\rm sym} + \mathfrak{K}^{(z)}_{\rm III}\notag\\
&\quad\quad \times 2\, \Bigg\{f^{mnk}\Big[U^{ma}(z)-U^{ma}(x')\Big]\bigg[\Big[U^{nc}(z)-U^{nc}(x)\Big]U^{kb}(w)\notag\\
&\quad\quad\quad\quad +U^{nc}(x)\Big[U^{kb}(w)-U^{kb}(y)\Big]\bigg] \Bigg\}\big\{\rho^{a}(x')\rho^{b}(y)\rho^{c}(x)\big\}_{\rm sym}\Bigg] +{\rm c.c.}
\end{align}

\begin{align}\label{eq:H-III.III-2}
&\frac{d^{3}N^{\rm III.III}_{{\rm con}}}{d^{3}k}=\frac{1}{(2\pi)^{3}}\, \frac{ig^{4}}{8\pi^{4}}\, \frac{1}{k^{+}} \int_{0}^{\infty}\frac{d\xi}{2\pi} \int_{w,x',y',z,y,x}e^{-ik\cdot(w-y')}e^{-ik\cdot(z-y')\,\xi}\Bigg[ \mathfrak{K}^{(y')}_{\rm III}\notag\\
&\quad\quad \times i\, \Bigg\{f^{mnk}\Big[U^{ma}(y')-U^{ma}(x')\Big]\bigg[\Big[U^{nc}(z)-U^{nc}(x)\Big]U^{kb}(w)\notag\\
&\quad\quad\quad\quad +U^{nc}(x)\Big[U^{kb}(w)-U^{kb}(y)\Big]\bigg] \Bigg\} \Big[f^{abr}\, \delta^{(2)}(x'-y)\big\{\rho^{r}(x')\rho^{c}(x)\big\}_{\rm sym}\notag\\
&\quad\quad\quad +f^{acr}\, \delta^{(2)}(x'-x)\big\{\rho^{r}(x')\rho^{b}(y)\big\}_{\rm sym}\Big] + \mathfrak{K}^{(z)}_{\rm III}\notag\\
&\quad\quad \times i\, \Bigg\{f^{mnk}\Big[U^{ma}(z)-U^{ma}(x')\Big]\bigg[\Big[U^{nc}(z)-U^{nc}(x)\Big]U^{kb}(w)\notag\\
&\quad\quad\quad\quad +U^{nc}(x)\Big[U^{kb}(w)-U^{kb}(y)\Big]\bigg] \Bigg\} \Big[f^{abr}\, \delta^{(2)}(x'-y)\big\{\rho^{r}(x')\rho^{c}(x)\big\}_{\rm sym}\notag\\
&\quad\quad\quad +f^{acr}\, \delta^{(2)}(x'-x)\big\{\rho^{r}(x')\rho^{b}(y)\big\}_{\rm sym}\Big]\Bigg] +{\rm c.c.}
\end{align}

Genuine two-$\rho$ term:
\begin{align}\label{eq:H-III.III-2o}
&\frac{d^{3}N^{\rm III.III}_{{\rm quad}}}{d^{3}k}=\frac{1}{(2\pi)^{3}}\, \frac{g^{4}}{4\pi^{4}}\, \frac{1}{k^{+}} \int_{0}^{\infty}\frac{d\xi}{2\pi} \int_{w,x',y',z,x}e^{-ik\cdot(w-y')}e^{-ik\cdot(z-y')\,\xi}\notag\\
&\quad \times \frac{1}{\xi\,(x-z)^{2}+(x-w)^{2}} \times\Bigg\{\Kv(y'-z)\cdot\Kv(x'-y')\bigg[\frac{\xi\big[(x-z)^{2}-(x-w)^{2}\big]}{(z-w)^{2}\,(1+\xi)^{2}}\notag\\
&\quad\quad\quad +\frac{\xi}{1+\xi}\Big((x-z)\cdot\Kv(z-w)-\frac{1}{2}(x-z)\cdot\Kv(x-w)-\frac{(x-z)^{2}-(x-w)^{2}}{2\,(z-w)^{2}}\Big)\notag\\
&\quad\quad\quad +\frac{1}{1+\xi}\Big((x-w)\cdot\Kv(z-w)+\frac{1}{2}(x-w)\cdot\Kv(x-z)-\frac{(x-z)^{2}-(x-w)^{2}}{2\,(z-w)^{2}}\Big)\bigg]\notag\\
&\quad\quad +\frac{(x-z)^{2}-(x-w)^{2}}{(1+\xi)\,(x-w)^{2}}\Big[(x-w)\cdot\Kv(x'-y')\, \Kv(z-w)\cdot\Kv(y'-z)\notag\\
&\quad\quad\quad +\frac{1}{2}\, \Kv(x-z)\cdot\Kv(y'-z)\, (x-w)\cdot\Kv(x'-y')\Big]\notag\\
&\quad\quad -\Big[(x-z)\cdot\Kv(y'-z)\, \Kv(z-w)\cdot\Kv(x'-y')\notag\\
&\quad\quad\quad -\frac{1}{2}(x-z)\cdot\Kv(y'-z)\, \Kv(x-w)\cdot\Kv(x'-y')\notag\\
&\quad\quad\quad +(x-w)\cdot\Kv(y'-z)\, \Kv(z-w)\cdot\Kv(x'-y')\notag\\
&\quad\quad\quad +\frac{1}{2}\, \Kv(x-z)\cdot\Kv(x'-y')\, (x-w)\cdot\Kv(y'-z)\Big]\Bigg\}\notag\\
&\quad \times f^{c'd'a}\, U^{d'b}(z)\Big[U^{c'e'}(y')-U^{c'e'}(x')\Big]\rho^{e'}(x')\notag\\
&\quad \times\bigg[\Big[U^{ac}(w)-U^{ac}(z)\Big]f^{cbd} -f^{ace}\, U^{cb}(z)\Big[U^{ed}(x)-U^{ed}(z)\Big]\bigg]\rho^{d}(x)\notag\\
&\quad -\frac{1}{(2\pi)^{3}}\, \frac{g^{4}}{4\pi^{4}}\, \frac{1}{k^{+}} \int_{0}^{\infty}\frac{d\xi}{2\pi}\Big[\frac{1}{\xi\,(1+\xi)}\Big]_{+} \int_{w,x',y',z,x}e^{-ik\cdot(w-y')}e^{-ik\cdot(z-y')\,\xi}\, \frac{1}{(x-w)^{2}}\notag\\
&\quad \times\Big[(x-w)\cdot\Kv(x'-y')\, \Kv(z-w)\cdot\Kv(y'-z)\notag\\
&\quad\quad +\frac{1}{2}\, \Kv(x-z)\cdot\Kv(y'-z)\, (x-w)\cdot\Kv(x'-y')\Big]\notag\\
&\quad \times f^{c'd'a}\, U^{d'b}(z)\Big[U^{c'e'}(y')-U^{c'e'}(x')\Big]\rho^{e'}(x')\notag\\
&\quad \times\bigg[\Big[U^{ac}(w)-U^{ac}(z)\Big]f^{cbd} -f^{ace}\, U^{cb}(z)\Big[U^{ed}(x)-U^{ed}(z)\Big]\bigg]\rho^{d}(x)\notag\\
&\quad -\frac{1}{(2\pi)^{3}}\, \frac{g^{4}}{4\pi^{4}}\, \frac{1}{k^{+}} \int_{0}^{\infty}\frac{d\xi}{2\pi} \int_{w,x',y',z,x}e^{-ik\cdot(w-y')}e^{-ik\cdot(z-y')\,\xi}\, \frac{\xi}{\xi\,(x-z)^{2}+(x-w)^{2}}\notag\\
&\quad \times\bigg\{\Big[(x-z)\cdot\Kv(x'-y')\, \Kv(z-w)\cdot\Kv(y'-z)\notag\\
&\quad\quad\quad -\frac{1}{2}(x-z)\cdot\Kv(x'-y')\, \Kv(x-w)\cdot\Kv(y'-z)\Big]\notag\\
&\quad\quad \times f^{c'd'a}\, U^{d'b}(z)\Big[U^{c'e'}(y')-U^{c'e'}(x')\Big]\rho^{e'}(x')\notag\\
&\quad\quad -\Big[(x-z)\cdot\Kv(x'-z)\, \Kv(z-w)\cdot\Kv(y'-z)\notag\\
&\quad\quad\quad -\frac{1}{2}(x-z)\cdot\Kv(x'-z)\, \Kv(x-w)\cdot\Kv(y'-z)\Big]\notag\\
&\quad\quad \times f^{c'd'a}\, U^{d'b}(z)\Big[U^{c'e'}(z)-U^{c'e'}(x')\Big]\rho^{e'}(x')\bigg\}\notag\\
&\quad \times\bigg[\Big[U^{ac}(w)-U^{ac}(z)\Big]f^{cbd} -f^{ace}\, U^{cb}(z)\Big[U^{ed}(x)-U^{ed}(z)\Big]\bigg]\rho^{d}(x)\notag\\
&\quad +\frac{1}{(2\pi)^{3}}\, \frac{g^{4}}{8\pi^{5}}\, \frac{1}{k^{+}} \int_{w,x',w',z,x}e^{-ik\cdot(w-w')}\, \frac{(x-w)^{2}}{\big[(x-z)^{2}-(x-w)^{2}\big]}\notag\\
&\quad \times \log\left[\frac{(x-z)^{2}}{(x-w)^{2}}\right] \bigg[\Kv(x-z)\cdot\Kv(x'-z)\, \Kv(z-w)\cdot\Kv(w'-z)\notag\\
&\quad\quad -\frac{1}{2}\, \Kv(x-z)\cdot\Kv(x'-z)\, \Kv(x-w)\cdot\Kv(w'-z)\bigg]\notag\\
&\quad \times f^{c'd'a}\, U^{d'b}(z)\Big[U^{c'e'}(z)-U^{c'e'}(x')\Big]\rho^{e'}(x') \times\bigg[\Big[U^{ac}(w)-U^{ac}(z)\Big]f^{cbd}\notag\\
&\quad\quad -f^{ace}\, U^{cb}(z)\Big[U^{ed}(x)-U^{ed}(z)\Big]\bigg]\rho^{d}(x)\, +{\rm c.c.}
\end{align}


\paragraph{Row III.V.} Symmetrised three-$\rho$ term and contact terms, with the kernels~\eqref{eq:K-y_rm_III} and~\eqref{eq:K-z_rm_III}:
\begin{align}\label{eq:H-III.V-3}
&\frac{d^{3}N^{\rm III.V}_{{\rm sym}}}{d^{3}k}=\frac{1}{(2\pi)^{3}}\, \frac{ig^{4}}{4\pi^{4}}\, \frac{1}{k^{+}} \int_{0}^{\infty}\frac{d\xi}{2\pi} \int_{w,y',z,y,x,x'}e^{-ik\cdot(w-y')}e^{-ik\cdot(z-y')\,\xi}\Bigg[ \mathfrak{K}^{(y')}_{\rm III}\notag\\
&\quad\quad \times \Bigg\{f^{mkl}\Big[U^{ma}(y')-U^{ma}(x')\Big]U^{kb}(x)\Big[U^{lc}(y)-U^{lc}(w)\Big]\Bigg\}\big\{\rho^{a}(x')\rho^{b}(x)\rho^{c}(y)\big\}_{\rm sym}\notag\\
&\quad\quad + \mathfrak{K}^{(z)}_{\rm III}\notag\\
&\quad\quad \times \Bigg\{f^{mkl}\Big[U^{ma}(z)-U^{ma}(x')\Big]U^{kb}(x)\Big[U^{lc}(y)-U^{lc}(w)\Big]\Bigg\}\big\{\rho^{a}(x')\rho^{b}(x)\rho^{c}(y)\big\}_{\rm sym}\Bigg]\notag\\
&\quad +{\rm c.c.}
\end{align}

\begin{align}\label{eq:H-III.V-2}
&\frac{d^{3}N^{\rm III.V}_{{\rm con}}}{d^{3}k}=\frac{1}{(2\pi)^{3}}\, \frac{ig^{4}}{4\pi^{4}}\, \frac{1}{k^{+}} \int_{0}^{\infty}\frac{d\xi}{2\pi} \int_{w,y',z,y,x,x'}e^{-ik\cdot(w-y')}e^{-ik\cdot(z-y')\,\xi}\Bigg[ \mathfrak{K}^{(y')}_{\rm III}\notag\\
&\quad\quad \times \frac{i}{2}\, \Bigg\{f^{mkl}\Big[U^{ma}(y')-U^{ma}(x')\Big]U^{kb}(x)\Big[U^{lc}(y)-U^{lc}(w)\Big]\Bigg\}\Big[f^{abr}\notag\\
&\quad\quad\quad \times \delta^{(2)}(x'-x)\big\{\rho^{r}(x')\rho^{c}(y)\big\}_{\rm sym} +f^{bcr}\, \delta^{(2)}(x-y)\big\{\rho^{a}(x')\rho^{r}(x)\big\}_{\rm sym}\notag\\
&\quad\quad\quad +f^{acr}\, \delta^{(2)}(x'-y)\big\{\rho^{r}(x')\rho^{b}(x)\big\}_{\rm sym}\Big] + \mathfrak{K}^{(z)}_{\rm III}\notag\\
&\quad\quad \times \frac{i}{2}\, \Bigg\{f^{mkl}\Big[U^{ma}(z)-U^{ma}(x')\Big]U^{kb}(x)\Big[U^{lc}(y)-U^{lc}(w)\Big]\Bigg\}\notag\\
&\quad\quad \times \Big[f^{abr}\, \delta^{(2)}(x'-x)\big\{\rho^{r}(x')\rho^{c}(y)\big\}_{\rm sym} +f^{bcr}\, \delta^{(2)}(x-y)\big\{\rho^{a}(x')\rho^{r}(x)\big\}_{\rm sym}\notag\\
&\quad\quad\quad +f^{acr}\, \delta^{(2)}(x'-y)\big\{\rho^{r}(x')\rho^{b}(x)\big\}_{\rm sym}\Big]\Bigg] +{\rm c.c.}
\end{align}


\paragraph{Row III.VI.} Symmetrised three-$\rho$ term and contact terms, with the kernels~\eqref{eq:K-y_rm_III} and~\eqref{eq:K-z_rm_III}:
\begin{align}\label{eq:H-III.VI-3}
&\frac{d^{3}N^{\rm III.VI}_{{\rm sym}}}{d^{3}k}=\frac{1}{(2\pi)^{3}}\, \frac{ig^{4}}{4\pi^{4}}\, \frac{1}{k^{+}} \int_{0}^{\infty}\frac{d\xi}{2\pi} \int_{w,y',z,y,x,x'}e^{-ik\cdot(w-y')}e^{-ik\cdot(z-y')\,\xi}\Bigg[ \mathfrak{K}^{(y')}_{\rm III}\notag\\
&\quad\quad \times \Bigg\{f^{mnk}\Big[U^{ma}(y')-U^{ma}(x')\Big]U^{kb}(y)\Big[U^{nc}(x)-U^{nc}(z)\Big]\Bigg\}\, \big\{\rho^{a}(x')\rho^{b}(y)\rho^{c}(x)\big\}_{\rm sym}\notag\\
&\quad\quad + \mathfrak{K}^{(z)}_{\rm III} \times \Bigg\{f^{mnk}\Big[U^{ma}(z)-U^{ma}(x')\Big]U^{kb}(y)\Big[U^{nc}(x)-U^{nc}(z)\Big]\Bigg\}\notag\\
&\quad\quad \times \big\{\rho^{a}(x')\rho^{b}(y)\rho^{c}(x)\big\}_{\rm sym}\Bigg] +{\rm c.c.}
\end{align}

\begin{align}\label{eq:H-III.VI-2}
&\frac{d^{3}N^{\rm III.VI}_{{\rm con}}}{d^{3}k}=\frac{1}{(2\pi)^{3}}\, \frac{ig^{4}}{4\pi^{4}}\, \frac{1}{k^{+}} \int_{0}^{\infty}\frac{d\xi}{2\pi} \int_{w,y',z,y,x,x'}e^{-ik\cdot(w-y')}e^{-ik\cdot(z-y')\,\xi}\Bigg[ \mathfrak{K}^{(y')}_{\rm III}\notag\\
&\quad\quad \times \frac{i}{2}\, \Bigg\{f^{mnk}\Big[U^{ma}(y')-U^{ma}(x')\Big]U^{kb}(y)\Big[U^{nc}(x)-U^{nc}(z)\Big]\Bigg\}\notag\\
&\quad\quad \times\Big[f^{abr}\, \delta^{(2)}(x'-y)\big\{\rho^{r}(x')\rho^{c}(x)\big\}_{\rm sym} +f^{bcr}\, \delta^{(2)}(y-x)\big\{\rho^{a}(x')\rho^{r}(y)\big\}_{\rm sym}\notag\\
&\quad\quad\quad +f^{acr}\, \delta^{(2)}(x'-x)\big\{\rho^{r}(x')\rho^{b}(y)\big\}_{\rm sym}\Big] + \mathfrak{K}^{(z)}_{\rm III}\notag\\
&\quad\quad \times \frac{i}{2}\, \Bigg\{f^{mnk}\Big[U^{ma}(z)-U^{ma}(x')\Big]U^{kb}(y)\Big[U^{nc}(x)-U^{nc}(z)\Big]\Bigg\}\notag\\
&\quad\quad \times\Big[f^{abr}\, \delta^{(2)}(x'-y)\big\{\rho^{r}(x')\rho^{c}(x)\big\}_{\rm sym} +f^{bcr}\, \delta^{(2)}(y-x)\big\{\rho^{a}(x')\rho^{r}(y)\big\}_{\rm sym}\notag\\
&\quad\quad\quad +f^{acr}\, \delta^{(2)}(x'-x)\big\{\rho^{r}(x')\rho^{b}(y)\big\}_{\rm sym}\Big]\Bigg] +{\rm c.c.}
\end{align}

\par\endgroup

% ======================================================================
\begingroup\sloppy
\section{Large-distance behaviour and integrability of the finite terms}\label{app:tails}
% ======================================================================
% Adapted from the verified Sections 6.3-6.7 of the previous draft (paper/draft/sec_finite.tex;
% notes paper/notes/fix_soft.tex, fix_tail.tex, ref3_finiteness.tex, ref3_prescriptions.tex).

This appendix gives the reasons for the prescriptions of Section~\ref{sec:fin_together}: which terms of the finite part have non-integrable tails at large transverse distance, how these tails cancel, and in which sense the remaining integrals converge. The special functions introduced here are not needed to state the result; they are used only to exhibit its large-distance behaviour. We use the abbreviations
\begin{align}\label{eq:Cf}
\mathcal C_{\rm f}\equiv\frac{1}{(2\pi)^3}\,\frac{g^4}{16\pi^5}\,\frac{1}{k^+}=\frac{\mathcal C}{\pi}\,,\qquad P^i\equiv\Kv^i(x'-w')\,,
\end{align}
with $\mathcal C$ of Eq.~\eqref{eq:dg_not}. A term is called integrable if its singularities at all coincidence points are integrable and if it decays faster than $1/\v z^2$ at large transverse position $\v z$ of the integrated gluon (in practice like $\log|\v z|/|\v z|^3$). A term $\propto e^{-i\v k\cdot\v z}/\v z^2$, in which the measured momentum flows through $\v z$, converges as a Fourier integral, as $\int_we^{-i\v k\cdot\v w}\Kv(x-w)$ does at LO, and is not counted as a tail.

\subsection{Closed forms of the longitudinal integrals}\label{app:closedforms}

When only the recoil phase depends on $\xi$, the integrals of Eqs.~\eqref{eq:plus_xi} and~\eqref{eq:inst_xi} are, with a real variable $s$ (${\rm Ci}$, ${\rm Si}$, $E_1$ the cosine, sine and exponential integrals),
\begin{align}\label{eq:recoilfns}
\int_0^1d\xi\,\Big[\frac1\xi\Big]_+e^{-is\xi}&=\int_0^1\frac{d\xi}{\xi}\big(e^{-is\xi}-1\big)\equiv\mathcal E(s)={\rm Ci}(|s|)-\gamma_E-\log|s|-i\,{\rm Si}(s)\,,\notag\\
\int_0^\infty d\xi\,\Big[\frac{1}{\xi(1+\xi)}\Big]_+e^{-is\xi}&=\int_0^\infty\frac{d\xi}{\xi(1+\xi)}\big(e^{-is\xi}-1\big)\equiv\Psi(s)=\mathcal E_\infty(s)-\widehat{\mathcal E}(s)\,,\notag\\
\int_0^\infty\frac{d\xi}{1+\xi}\,e^{-is\xi}&\equiv\widehat{\mathcal E}(s)=e^{is}E_1(is)\,,\qquad
\mathcal E_\infty(s)\equiv-\gamma_E-\log|s|-\frac{i\pi}{2}\,{\rm sgn}(s)\,,\notag\\
\int_0^1d\xi\,e^{-is\xi}&=\frac{1-e^{-is}}{is}\,.
\end{align}
For $|s|\to\infty$, $\mathcal E(s)=\mathcal E_\infty(s)-\frac{e^{-is}}{is}+O(s^{-2})$, $\Psi(s)=\mathcal E_\infty(s)-\frac{1}{is}+O(s^{-2})$ and $\widehat{\mathcal E}(s)\simeq\frac{1}{is}$; for $s\to0$, $\mathcal E(s)\simeq-is$, $\Psi(s)=O(s\log|s|)$ and $\widehat{\mathcal E}(s)\simeq\mathcal E_\infty(s)$. All these functions satisfy $f(-s)=f(s)^*$. $\mathcal E$ belongs to Group~I, where the $\xi$ integral runs up to one, and $\Psi$ to Group~III, where it runs to infinity; at large $|s|$ both tend to the same function $\mathcal E_\infty$, including its odd part.

\subsection{Soft remainders of the plus prescriptions}\label{sec:fin_soft}

\paragraph{Origin.} In Rows~I.III and~I.V the observed gluon in the amplitude, and in Rows~III.III, III.V and~III.VI in the conjugate amplitude, is part of a two-gluon state: a hard gluon at $\v w$ (resp.\ $\v w'$) and a gluon with $p^+=\xi k^+$ at $\v z$. The observed gluon is at the $+$-momentum weighted centre, so that the phase is $e^{-i\v k\cdot\v w}e^{-i\xi\v k\cdot(\v z-\v w)}$ (Group~I), resp.\ $e^{i\v k\cdot\v w'}e^{-i\xi\v k\cdot(\v z-\v w')}$ (Group~III); in both cases the exponent is $-i\xi s$ with $s=\v k\cdot(\v z-\text{hard position})$. The longitudinal integral of a term $\frac{d\xi}{\xi}g(\xi)$ of a Group~I row is split as
\begin{align}\label{eq:softsplit}
\int_{\Lambda/k^+}^{1}\frac{d\xi}{\xi}\,e^{-i\xi s}g(\xi)&=\log\frac{k^+}{\Lambda}\,g(0)+g(0)\,\mathcal E(s)+\int_0^{1}\frac{d\xi}{\xi}\,e^{-i\xi s}\big[g(\xi)-g(0)\big]+O\Big(\frac{\Lambda}{k^+}\Big)\,,
\end{align}
and a term $\frac{d\xi}{\xi(1+\xi)}g(\xi)$ of a Group~III row in the same way, with $\Psi(s)$ in place of $\mathcal E(s)$. Here $g(0)$ is the coefficient of the row in the JIMWLK list, which is kept exactly as in Section~\ref{sec:JIMWLK}; the last two terms together are the plus-prescription term of Eq.~\eqref{eq:plus_xi} with $f(\xi)=e^{-i\xi s}g(\xi)$. We call $g(0)\,\mathcal E(s)$, resp.\ $g(0)\,\Psi(s)$, the soft remainder of the term; the last term of Eq.~\eqref{eq:softsplit} is integrable by itself. Since $|\mathcal E(s)|,|\Psi(s)|\simeq\log|s|$ for $|s|\gg1$, the recoil cuts the soft integral off at $\xi\simeq1/|\v k\cdot(\v z-\v w)|$, i.e.\ at $p^+/p_\perp\simeq k^+/|\v k|$ for $p_\perp\sim|\v z-\v w|^{-1}$: the soft gluon is restricted to rapidities below that of the observed gluon.

\paragraph{The double logarithms.} Each soft remainder behaves at large $\v z$ like $\frac{T}{\v z^2}\mathcal E_\infty(\v k\cdot\v z)$, where $T$ is the $1/\v z^2$ tail of $g(0)$, and is not integrable by itself. Before the split~\eqref{eq:softsplit}, i.e.\ at finite $\Lambda$, every such row contains $\pi\log^2(k^+/\Lambda)\,T_{\rm row}+O(\log(k^+/\Lambda))$, from the region in which the soft gluon is far from the observed one and is cut off only by the recoil; in the rows without recoil the tails multiply single logarithms only. The coefficient of the double logarithm in the sum of all rows is $\sum_{\rm rows}T_{\rm row}$. This sum vanishes for an arbitrary projectile, in the symmetrised two-$\rho$ and three-$\rho$ sectors (numerically $|\sum T|<10^{-7}$ against single tails of order one to six, for $SU(2)$ and $SU(3)$; changing any one weight gives a sum of order one), and hence the double logarithms $\log^2(k^+/\Lambda)$ cancel in the cross section. What is left at large distance in the logarithmic terms is only the infrared behaviour of the JIMWLK term for a non-singlet pair of charges; it vanishes for a colour-diagonal correlator, Eq.~\eqref{eq:singlet}, and is not affected by the finite terms. The vanishing of $\sum T$ is explained by the pair structure.

\paragraph{The pairs.} The soft remainders cancel in pairs. Each Group~I term (soft gluon absorbed by the observed gluon in the amplitude: virtual, recoil on $\v w$) has one Group~III partner (soft gluon emitted by the observed gluon in the conjugate amplitude: real, recoil on $\v w'$) with the same coefficient, the same colour structure and the same kernel, except that $\v w$ is replaced by $\v w'$. With the colour operators of the Group~I members in their ordered form,
\begin{align}\label{eq:CABCD}
\mathcal L^a&\equiv\big[U^{ab'}(x')-U^{ab'}(w')\big]\rho^{b'}(x')\,,\notag\\
C_A&\equiv f^{abd}\,\mathcal L^a\,U^{de}(x)\big[U^{bc}(w)-U^{bc}(y)\big]\rho^e(x)\rho^c(y)\,,\notag\\
C_D&\equiv f^{abd}\,\mathcal L^a\,U^{de}(x)\big[U^{bc}(z)-U^{bc}(y)\big]\rho^e(x)\rho^c(y)\,,\notag\\
C_B&\equiv\mathcal L^a\Big[f^{adc}\,U^{db}(z)\,U^{ce}(w)\,\{\rho^e(y),\rho^b(x)\}-f^{abc}\,U^{cd}(y)\,U^{be}(x)\,\{\rho^d(y),\rho^e(x)\}\Big]\,,\notag\\
C_C&\equiv\big[U^{ab'}(x')-U^{ab'}(w')\big]\Big[N_c\,U^{ah}(x)-f^{adc}f^{hbe}\,U^{db}(z)\,U^{ce}(w)\Big]\rho^{b'}(x')\,\rho^h(x)\,,
\end{align}
and with the variables of Row~I.III (soft gluon at $\v z$, hard gluon at $\v w$, charges at $\v x'$, $\v x$, $\v y$; the terms of Row~I.V relabelled accordingly, and in Group~III $\v y'$ renamed $\v w'$), each pair has the form
\begin{align}\label{eq:pair}
\text{pair }X:\qquad\mathcal C_{\rm f}\,\alpha_X\,V_X\,C_X\Big[\kappa_X(w)\,\mathcal E\big(\v k\cdot(\v z-\v w)\big)-\kappa_X(w')\,\Psi\big(\v k\cdot(\v z-\v w')\big)\Big]
\end{align}
under $\int_{w,w'}e^{-ik\cdot(w-w')}\int_{x,y,x',z}$, with the entries of Table~\ref{tab:pairs}.
\begin{table}[t]
\centering
\renewcommand{\arraystretch}{1.3}
\footnotesize\setlength{\tabcolsep}{3.5pt}
\begin{tabular}{@{}c@{\hspace{5pt}}l@{\hspace{5pt}}l@{\hspace{5pt}}c@{\hspace{5pt}}c@{\hspace{5pt}}c@{\hspace{5pt}}l@{}}
\toprule
pair & Group I & Group III & $\alpha_X$ & $V_X$ & $\kappa_X(u)$ & tail $T_X$\\
\midrule
A & I.III, $\mathfrak K^{(w)}_{\rm I.III}$ & III.V & $-2i$ & $P\cdot\Kv(y-w)$ & $K(x,u,z)$ & $-2i\,P\cdot\Kv(y-w)\,C_A$\\
B & I.V, cubic & III.III, cubic & $-i$ & $P\cdot\Kv(y-w)$ & $K(x,u,z)$ & $-i\,P\cdot\Kv(y-w)\,C_B$\\
C & I.V, quad & III.III, quad & $-1$ & $P\cdot\Kv(x-w)$ & $K(x,u,z)-2K(w,u,z)$ & $+P\cdot\Kv(x-w)\,C_C$\\
D & I.III, $\mathfrak K^{(z)}_{\rm I.III}$ (last) & III.VI & $+2i$ & $P\cdot\Kv(x-w)$ & $K(y,u,z)$ & $+2i\,P\cdot\Kv(x-w)\,C_D$\\
\bottomrule
\end{tabular}
\caption{The four pairs of soft remainders, Eq.~\eqref{eq:pair}. $T_X\equiv\lim_{|z|\to\infty}\v z^2\alpha_XV_X\kappa_X(w)C_X$ is the $1/\v z^2$ tail of the Group~I member; that of the Group~III member is $-T_X$. ``Cubic'' and ``quad'' denote the cubic and the genuine two-$\rho$ part of a row; for pair~D the Group~I member is the last term of $\mathfrak K^{(z)}_{\rm I.III}$. In pair C, $C_C$ contains $U(\v w)$ also in its second term, and $K(w,u,z)$ is $1/(\v z-\v w)^2$ at $u=w$ and $K(w,w',z)$ at $u=w'$. In pair B the two recoil terms of Row~I.V are combined after the relabelling $\v y\leftrightarrow\v y'$ of the integrated charge positions, and in pair D Row~III.VI is written with $\v x\leftrightarrow\v y$.}
\label{tab:pairs}
\end{table}
That the Group~III member has the colour structure of its Group~I partner follows from $U^{d'b}(z)U^{cb}(z)=\delta^{d'c}$, $f^{c'd'a}f^{dd'a}=N_c\delta^{c'd}$ and a relabelling of colour indices only; no charges are reordered, so these are operator identities, valid for any $U$ at any $\v z$, and they hold separately for the symmetrised and for the contact terms, since the two members of a pair have the same ordered colour operator and hence the same decomposition~\eqref{eq:symord}--\eqref{eq:symr21}. They were also verified numerically, together with the statement that each entry of Table~\ref{tab:pairs} equals the row's own coefficient, in the two- and three-$\rho$ sectors for $SU(2)$ and $SU(3)$ (to $2\times10^{-15}$).

\paragraph{The two recoil sides combine without a condition on the target.} Each pair can be written as
\begin{align}\label{eq:rmv}
&\kappa(w)\,\mathcal E\big(\v k\cdot(\v z-\v w)\big)-\kappa(w')\,\Psi\big(\v k\cdot(\v z-\v w')\big)\notag\\
&\quad=\big[\kappa(w)-\kappa(w')\big]\mathcal E\big(\v k\cdot(\v z-\v w)\big)
+\kappa(w')\Big[\mathcal E\big(\v k\cdot(\v z-\v w)\big)-\Psi\big(\v k\cdot(\v z-\v w')\big)\Big]\,.
\end{align}
The first term is a real-minus-virtual combination times a kinematic logarithm: $\kappa(w)-\kappa(w')=O(|\v w-\v w'|/|\v z|^3)$ while $|\mathcal E|\le\log|\v z|+{\rm const}$; it has the structure of the remainder found in forward production~\cite{Altinoluk:2014eka} and is integrable. The second term is the mismatch of the recoil on $\v w$ and on $\v w'$: $\mathcal E(\v k\cdot(\v z-\v w))-\Psi(\v k\cdot(\v z-\v w'))=\log\frac{|\v k\cdot(\v z-\v w')|}{|\v k\cdot(\v z-\v w)|}+O(1/|\v k\cdot\v z|)$; the odd parts $-\frac{i\pi}{2}{\rm sgn}$ are the same on the two sides and cancel, and the remainder is $O(1/|\v z|)$ outside the strip $|\v k\cdot\v z|\lesssim|\v k||\v w-\v w'|$ and bounded by logarithms inside it, which gives a convergent integral. The colour structure $C_X$, with its whole dependence on $U(\v z)$, is a common factor of both terms. Therefore the tails of the two members of a pair are $\pm T_X/\v z^2$ exactly, as operators and for every value of $U(\v z)$: the parts linear in $U(\v z)$, the $\v z$-independent parts and the parts with $U(\v z)U^T(\v z)$ cancel together, and nothing has to be assumed about the Wilson line at large distance beyond $|U^{ab}|\le1$; for a target of finite size, $U(\v z)\to1$, the cancellation is a special case. The complex-conjugate half is the complex conjugate of each pair and cancels pair by pair in the same way. Numerically, $\v z^2$ times the sum of the two members decreases like $1/|\v z|$ in every direction, i.e.\ the sum is $O(|\v z|^{-3})$, for all four pairs, $SU(2)$ and $SU(3)$, in the two- and three-$\rho$ sectors, in the colour-singlet projection, with $U(\v z)=1$ and with a random $U(\v z)$ in every direction; the angular mean of $[\kappa_X(w)-\kappa_X(w')]V_XC_X$ falls like $|\v z|^{-4}$.

\paragraph{Kinematic constraint.} The pairing leaves the JIMWLK term unchanged. Equivalently, the soft remainders can be combined with the $\log(k^+/\Lambda)$ terms of the same pairs: with $L=\log(k^+/\Lambda)$,
\begin{align}\label{eq:kinconstr}
&\kappa(w)\int_{\Lambda/k^+}^{1}\frac{d\xi}{\xi}\,e^{-i\xi\,\v k\cdot(\v z-\v w)}-\kappa(w')\int_{\Lambda/k^+}^\infty\frac{d\xi}{\xi(1+\xi)}\,e^{-i\xi\,\v k\cdot(\v z-\v w')}\notag\\
&\quad=L\big[\kappa(w)-\kappa(w')\big]+\kappa(w)\,\mathcal E\big(\v k\cdot(\v z-\v w)\big)-\kappa(w')\,\Psi\big(\v k\cdot(\v z-\v w')\big)+O\Big(\frac{\Lambda}{k^+}\Big)\,.
\end{align}
In this form each kernel of the JIMWLK term that carries a recoil phase comes with its own rapidity interval, $\simeq\log\frac{k^+}{\Lambda|\v k\cdot(\v z-\v w)|}-\gamma_E-\frac{i\pi}{2}{\rm sgn}(\v k\cdot(\v z-\v w))$ for $|\v k\cdot(\v z-\v w)|\gg1$ (resp.\ with $\v w'$) and $\log\frac{k^+}{\Lambda}$ for $|\v k\cdot(\v z-\v w)|\ll1$: the soft gluon has a rapidity below that of the observed gluon, as in the kinematically constrained formulations of forward production~\cite{Altinoluk:2014eka,Ducloue:2016shw} (the constraint of~\cite{Ducloue:2016shw} has not been compared in detail). For pairs~A--C the terms with $L$ are the contributions of these rows to the JIMWLK term; for pair~D only the real half is in the JIMWLK term (see the next subsection).

\subsection{The kernel of Row I.III}\label{sec:fin_GIIII}

The term of $\mathfrak K^{(z)}_{\rm I.III}$, Eq.~\eqref{eq:K-z_rm_I_III}, that multiplies $K(y,w,z)\,P\cdot\Kv(x-w)$ is of a different kind. With $a_t=(x-w)-t(x-z)$,
\begin{align}\label{eq:ell}
\ell(z,w,x)\equiv-\frac12\log\frac{(z-w)^2}{(x-w)^2}=-\int_0^1\frac{dt}{t}\Big[1-\frac{(x-w)\cdot a_t}{a_t^2}\Big]\,,
\end{align}
so that $\ell\,\Kv^i(x-w)$ is the LO kernel minus the LO kernel with the source shifted by the recoil; the soft integral is cut off at $t\simeq|\v x-\v w|/|\v z-\v w|$, without a phase, and nothing of this term is in the JIMWLK list. The term is integrable at $\v z\to\v w$ ($\ell K(y,w,z)\sim\log|\v z-\v w|/|\v z-\v w|$), but at large $\v z$ it has a tail $\propto\ell/\v z^2$ with $\ell\to-\log\frac{|\v z|}{|\v x-\v w|}$, an isotropic tail without the odd part of $\mathcal E_\infty$. In the kernel~\eqref{eq:K-z_rm_I_III} it is therefore written as
\begin{align}
\ell=\Big[\ell-\mathcal E\big(\v k\cdot(\v z-\v w)\big)\Big]+\mathcal E\big(\v k\cdot(\v z-\v w)\big)\,,\qquad
\mathcal E\big(\v k\cdot(\v z-\v w)\big)=\int_0^1d\xi\Big[\frac1\xi\Big]_+e^{-ik\cdot(z-w)\,\xi}\,.
\end{align}
The second part is the Group~I member of pair~D. The first part is combined with the other terms of $\mathfrak K^{(z)}_{\rm I.III}$, the first three groups of terms of Eq.~\eqref{eq:K-z_rm_I_III} (24 terms when written out). Denoting by $\widehat{\mathfrak K}$ the kernel $\mathfrak K^{(z)}_{\rm I.III}$ without its last term, at fixed $\v x,\v y,\v x',\v w'$ and $|\v z|\to\infty$,
\begin{align}\label{eq:Mhatclaim}
\int d^2w\,e^{-i\v k\cdot\v w}\,\widehat{\mathfrak K}=O\Big(\frac{\log|\v z|}{|\v z|^3}\Big)+\big(\text{terms}\propto e^{-i\v k\cdot\v z}\big)\,,
\end{align}
i.e.\ after the $\v w$ integration there is no non-oscillating $1/\v z^2$ term in any direction, for any $U(\v z)$: the colour factor does not depend on $\v w$ and is a spectator. The argument is as follows.
\begin{enumerate}
\item The first three groups of terms of the kernel have the large-$\v z$ form
\begin{align}\label{eq:saw}
\frac{\varphi}{\v z^2}\,\frac{P\cdot e_\perp}{|\v x-\v w|}+O\Big(\frac{\log|\v z|}{|\v z|^3}\Big)\,,\qquad-\pi<\varphi<\pi\,,
\end{align}
where $\varphi$ is the angle from $\v x-\v w$ to $\v z$ and $e_\perp$ the unit vector $(\v x-\v w)/|\v x-\v w|$ rotated by $+\frac\pi2$. It comes from the terms with the arccosine, has zero angular mean, but does not vanish pointwise; Eq.~\eqref{eq:saw} was established numerically.
\item By Eq.~\eqref{eq:recoilfns} and $K(y,w,z)\to1/\v z^2$, the term with $\ell-\mathcal E$ tends to $\frac{P\cdot\Kv(x-w)}{\v z^2}\big[\gamma_E+\log(|\v k||\v x-\v w|)+c(\theta)\big]$, up to terms $\propto e^{-i\v k\cdot(\v z-\v w)}$, with $c(\theta)=\log|\cos\theta|+\frac{i\pi}{2}{\rm sgn}(\cos\theta)$ and $\theta$ the angle between $\v k$ and $\v z$.
\item With $\v r=\v x-\v w$ and $e^{-i\v k\cdot\v w}=e^{-i\v k\cdot\v x}e^{i\v k\cdot\v r}$,
\begin{align}\label{eq:sawid}
\int d^2r\,e^{i\v k\cdot\v r}\,\Kv^i(r)\Big[\gamma_E+\log(|\v k||\v r|)+c(\theta)\Big]&=+\frac{2\pi i\,\hat k^i}{|\v k|}\,\chi(\theta)\,,\notag\\
\int d^2r\,e^{i\v k\cdot\v r}\,\varphi\,\frac{e_\perp^i}{|\v r|}&=-\frac{2\pi i\,\hat k^i}{|\v k|}\,\chi(\theta)\,,\notag\\
\chi(\theta)&\equiv\log|2\cos\theta|+\frac{i\pi}{2}\,{\rm sgn}(\cos\theta)\,,
\end{align}
which follow from the radial integrals $\int_0^\infty dr\,e^{i\kappa r}=i\mathcal P\frac1\kappa+\pi\delta(\kappa)$ and $\int_0^\infty dr\,e^{i\kappa r}[\gamma_E+\log(|\v k|r)]=\frac{i}{\kappa}\big[-\log\frac{|\kappa|}{|\v k|}+\frac{i\pi}{2}{\rm sgn}\,\kappa\big]$, $\kappa=|\v k|\cos\alpha$ ($\alpha$ the angle of $\v r$ with $\v k$; with a convergence factor $e^{-\epsilon r}$), followed by the angular integrals. The two integrals cancel, for every $\theta$. Each of the two Fourier integrals is unique, i.e.\ independent of the regulator and of the order of the radial and angular integrals (Section~\ref{sec:fin_integrable}); what would be only conditionally convergent is the subsequent $\v z$ integral of each term alone, since $\chi$ has zero angular mean. In their sum nothing conditional is left.
\end{enumerate}
Eq.~\eqref{eq:sawid} was checked numerically with the convergence factor $e^{-\epsilon r}$ and with two anisotropic factors (deviations $\le1.2\times10^{-3}$ at $\epsilon=10^{-4}$, falling like $\epsilon$), and Eq.~\eqref{eq:Mhatclaim} with the full kernels: at $|\v z|=2560$ the two parts of $\v z^2\int_we^{-ik\cdot w}\widehat{\mathfrak K}$ are $6.933-2.690i$ and $-6.906+2.690i$, against the prediction $\pm(6.909-2.687i)$, and $\v z^2$ times their sum decreases like $1/|\v z|$ in all directions tested.

The order of the integrations is the definition of the measured momentum (the unobserved gluon is summed over at fixed $\v k$), not a choice: at fixed $\v w$ the kernel has a $1/\v z^2$ tail whose angular mean does not vanish, so that the other order does not exist; after the $\v w$ integration, which is unique for any regulator of the measured gluon, the tail is absent in every direction and the $\v z$ integral converges absolutely. In summary, the term of Row~I.III with $\ell\,K(y,w,z)\,P\cdot\Kv(x-w)$ is cancelled at large distance (1) in its $\log|\v z|$ part by Row~III.VI, which has the same coefficient and colour structure and the kernel with $\v w\to\v w'$ (pair~D), and (2) in its remaining $1/\v z^2$ part by the tail~\eqref{eq:saw} of the other terms of the same row, after the integration over $\v w$. For pair~D only the real half, from Row~III.VI, is in the JIMWLK term; the virtual half is not, since Row~I.III carries $\ell$ in place of $\log(k^+/\Lambda)$ in this term. The term with the kernel $\mathfrak K^{(w)}_{\rm I.III}$ is cancelled at large distance by Row~III.V (pair~A). At $\v z\to\v w$ both terms are integrable.

\subsection{Instantaneous terms of Group III}\label{sec:fin_inst}

The instantaneous terms of Rows~III.III, III.V and~III.VI, the last term of $\mathfrak K^{(y')}_{\rm III}$ and the kernel $\mathfrak K^{(z)}_{\rm III}$, depend on $\xi$ only through the phase, and have no pointwise limit for $\vee\to\infty$. With $s=\v k\cdot(\v z-\v y')$ and the $\v y'$ integral done first, they are given by the distribution~\eqref{eq:inst_xi}. They are finite. With $\v r=\v y'-\v z$, the combination of the term with $\Kv(x'-y')$ and the Wilson lines at $\v y'$ and of the term with $\Kv(x'-z)$ and the Wilson lines at $\v z$ vanishes at $\v r\to0$, so that $\Kv(y'-z)$ times this combination is bounded there. Integrating first over the component of $\v r$ perpendicular to $\v k$, symmetrically at large distance, gives a function of the component $r_\parallel$ along $\v k$ on which $\pi\delta(s)$ and $\mathcal P\frac1s$ act as on a continuous function, and the phase makes the $r_\parallel$ integral converge at large distance. The two terms must stay together: each alone grows like $\log\vee$, which is the BFKL term. Numerically, with a smooth test function, the longitudinal integral cut at $\vee$ approaches the distributional value like $k^+/\vee$, and the distributional form reproduces the limit to $7\times10^{-8}$. The symmetric cut-off at large distance is sufficient but not necessary: the only part of the perpendicular integral that is not absolutely convergent is independent of $s$ and is annihilated by $\int ds\,e^{-is}[\pi\delta(s)-i\mathcal P\frac1s]=0$, so that any cut-off at large distance gives the same result.

Although the integrand does not decay in $\xi$ at fixed transverse positions, the instantaneous terms do not come from $p^+\sim\vee$, where the eikonal approximation fails. After the $\v y'$ integral the distribution in $\xi$ falls like $(1+\xi)^{-2}$ for Wilson lines that are differentiable, so that $p^+>P$ contributes $O(k^+/P)$, or $O((k^+/P)^\alpha)$ if the Wilson lines are only H\"older continuous with exponent $\alpha$. The bulk comes from $p^+\lesssim k^+\max\{1,1/(|\v k|\ell)\}$, with $\ell$ the shortest transverse scale of the bracket; the eikonal approximation requires $\vee$ to be large compared to this. With a smooth separation $f(p^+/\vee)$ of soft and valence modes ($f(0)=1$) the limit is the same for every $f$, and the whole result depends on $f$ only through $\log\frac{\vee}{k^+}\to\log\frac{\vee}{k^+}+c_f$, $c_f=\int_0^\infty\frac{du}{u}[f(u)-\theta(1-u)]$, in the BFKL term, a change of the rapidity-factorisation scale compensated by the evolution of the projectile. These statements were checked with model brackets that have the short-distance behaviour of the actual one, not with the actual bracket with Wilson lines; that the same $c_f$ appears in every row is argued from the $1/p^+$ form of the integrands, not checked row by row.

\subsection{Conditions and convergence}\label{sec:fin_integrable}

With the pairs of Section~\ref{sec:fin_soft} and the order of integration of Section~\ref{sec:fin_GIIII}, every finite term is integrable over all transverse positions at fixed positions $\v w$, $\v w'$ of the observed gluon, under the following conditions.
\begin{itemize}
\item[(C1)] $\Lambda\to0$ and $\vee\to\infty$ are taken after the subtractions~\eqref{eq:fin_def}, at fixed $\epsilon$ (leading power of the eikonal approximation, Section~\ref{sec:dg_poles}); terms $O(\Lambda/k^+)$ and $O(k^+/\vee)$ are dropped, and the terms $\log(1-k^+/\vee)$ go with the JIMWLK structures of their rows.
\item[(C2)] The identities that relate the colour structures of the two members of a pair use only a relabelling of colour indices and $U(\v z)U^T(\v z)=1$; they are operator identities, valid for any $U$. The statements about the $\v z$ integrand hold at fixed values of all other coordinates after the relabellings of integration variables listed in Table~\ref{tab:pairs}.
\item[(C3)] In Row~I.III the $\v w$ integral, i.e.\ the Fourier transform at the measured momentum, is done before the $\v z$ integral (Section~\ref{sec:fin_GIIII}).
\item[(C4)] Terms $\propto e^{-i\v k\cdot\v z}/\v z^2$ (from the end point $\xi\to1$ in Group~I, and from the regions in which the observed gluon is at $\v z$) are treated as convergent Fourier integrals.
\item[(C5)] No assumption on $U(\v z)$ at large distance is made; the orthogonality of the adjoint Wilson lines is used.
\item[(C6)] The vanishing energy denominators of $\mf B_3$ are principal values.
\item[(C7)] In the instantaneous terms of Group~III the $\v y'$ integral is done before the $\xi$ integral (Section~\ref{sec:fin_inst}).
\item[(C8)] The large-$\v z$ form~\eqref{eq:saw} is established numerically; its Fourier identity~\eqref{eq:sawid} is analytic.
\item[(C9)] The cancellations were verified on a projectile of charges at two sites (all assignments of the charge positions to the sites, including the contact terms of the three-$\rho$ products), with random classical $\rho$ for the three-$\rho$ part, for a general two-$\rho$ coefficient matrix and for its colour-singlet projection.
\item[(C10)] In the region in which both final-state gluons are far from the projectile together ($\v w,\v w',\v z\to\infty$ at fixed $\v w-\v z$), the finite terms of Row~II.IV have the large-distance tail of Section~\ref{sec:dg_tail}. It cancels at the level of the amplitudes against the same region of the genuine two-$\rho$ parts of Row~II.I and of Row~III.III, which therefore have to be integrated together with Row~II.IV under one integral over the position of the pair; the three rows carry the tail with the weights $1\!:\!1\!:\!-2$ (Section~\ref{sec:dg_tail}, Table~\ref{tab:tails}). The region in which only part of the pair is far does not contribute at this order: the three rows are, at fixed $\xi$, $|\mathcal A_B-\mathcal A_{CA}|^2$ with the measured gluon Fourier transformed, and in $\mathcal A_B-\mathcal A_{CA}$ the joint-limit coefficient cancels. This was verified numerically on the full $\v w$ plane for one charge and a target uniform at large distance. For a colour-neutral projectile and a target that is not uniform at large distance the argument is that the non-uniform remainder multiplies $\int_x\Kv(\v x-\v Y)\rho(\v x)=O(\v Y^{-2})$, which was not checked numerically. For a coloured projectile and a target that is not uniform at large distance a remainder proportional to the total colour charge of the projectile is left; it is the non-singlet term of the JIMWLK formula and is absent for a colour-neutral projectile.
\end{itemize}

\paragraph{Convergence.} All integrals over a coordinate that carries the measured phase ($\v w$, $\v w'$, $\v y'$, $\v y-\v y'$, and $\v z$ in the terms $\propto e^{-i\v k\cdot\v z}$) are Fourier integrals of locally integrable, tempered functions whose transforms are continuous at the measured momentum. For such a function $f$ and any regulator $\phi(\epsilon\v r)$ with $\phi(0)=1$ and integrable Fourier transform (an exponential or Gaussian factor, of any direction dependence, or a smooth cut-off of any shape),
\begin{align}\label{eq:FTlemma}
\lim_{\epsilon\to0}\int d^2r\,e^{-i\v k\cdot\v r}\,\phi(\epsilon\v r)\,f(\v r)=\hat f(\v k)\,,
\end{align}
since the left-hand side is the convolution of $\hat f$ with an approximate identity. These integrals, like $\int_we^{-i\v k\cdot\v w}\Kv(x-w)$ at LO, are therefore unique and independent of the order of their radial and angular parts. With them done first, every remaining integral, over the positions of the unobserved gluon and of the charges and over the longitudinal fractions, converges absolutely in the grouping of Section~\ref{sec:fin_together}, with two exceptions. The first are the principal values $\mathcal P\frac1s$ of the instantaneous terms, which act on H\"older-continuous coefficient functions. The second is the collinear point $\v z\to\v x$ of Row~I.V, where the kernel $\mathfrak K^{\rm quad}_{\rm I.V}$ together with the $\frac{11}{6}$ term behaves as a pure second harmonic in the angle, $a_2\cos2(\varphi-\varphi_0)/|\v z-\v x|^2$ (mean and first harmonic vanish; checked numerically down to $|\v z-\v x|=10^{-4}$): it is conditionally convergent, gives zero with any rotation-invariant cut-off around $\v z=\v x$ and a finite, different value for a non-circular one. The rotation-invariant choice is the one consistent with the $\frac{11}{6}$ subtraction, which comes from the dimensionally regularised (rotation-invariant) integral; an absolutely convergent form is obtained by subtracting $\theta(R_0-|\v z-\v x|)\,a(\varphi)/|\v z-\v x|^2$, which integrates to zero with a rotation-invariant measure. If the transverse indices are continued to $d_\perp=2-\epsilon$, the traceless tensor $\hat r^i\hat r^m-\frac12\delta^{im}$ acquires a trace $O(\epsilon)$, which multiplies the $1/\epsilon$ of the collinear subtraction and could give a finite local term; whether it is present has not been decided (Section~\ref{sec:openpoints}). The two integrals of Eq.~\eqref{eq:sawid} and the two orders of the $\xi$ and $\v y-\v y'$ integrals in Row~II.IV were evaluated with different regulators and orders, with the same result. Hence none of (C3), (C7) and the orders of integration in Eq.~\eqref{eq:sawid} and in Row~II.IV is a prescription that could shift the result by a finite amount. The checks are exact identities, power counting at all coincidences, pointwise tails in all directions and Fourier integrals at fixed values of the other coordinates; a complete numerical integration of the finite terms over all transverse positions with Wilson lines has not been carried out.

\paragraph{Contact terms.} All contact kernels are integrable. Every contact line was power-counted at every coincidence of its integration points (pairs and all larger clusters, in particular the coincidences of the merged charge point with $\v z$, $\v w$, $\v w'$, $\v y'$ and the other charge), with the colour factors as written: each line is absolutely integrable at each coincidence. Only one kernel is not integrable by itself, the instantaneous kernel $\mathfrak K^{(z)}_{\rm III}$ with the charge at $\v x'$ put at $\v x$, which contains $1/(\v x-\v z)^2$; at $\v z\to\v x$ its colour factors vanish identically, so that the line is integrable for any target with H\"older continuous Wilson lines. The kernel $\mathfrak K_{\rm I.II}$ is homogeneous of degree $-1$ in $(\v x-\v w,\v y-\v w)$, bounded at $\v x\to\v y$ and $\propto\log r/r$ at $\v x\to\v w$, and integrable (analytic estimate; its coordinate form is not available). The large-distance tails of the contact lines cancel pair by pair as at separate points. These statements were checked numerically, by fits of the exponents at scales $10^{-2}$ to $10^{-5}$ for three base configurations per line and two smooth, non-uniform $SU(3)$ Wilson-line fields; this is power counting, not a proof for every degenerate configuration.
\par\endgroup


\bibliographystyle{JHEP}
\bibliography{refs}

\end{document}
